% This chapter is a faithful blueprint adaptation of the indicated part of
% Cao--Zhu's revised TeX source. Source-line sentinels preserve auditability.

\chapter{Ricci Flow on Three-manifolds}
\label{chap:cz-three-manifold-ricci-flow}

\bigskip
We will use the Ricci flow to study the topology of compact
orientable three-manifolds. Let $M$ be a compact three-dimensional
orientable manifold. Arbitrarily given a Riemannian metric on the
manifold, we evolve it by the Ricci flow. The basic idea is to
understand the topology of the underlying manifold by studying
long-time behavior of the solution of the Ricci flow. We have seen
in Chapter 5 that for a compact three-manifold with positive Ricci
curvature as initial data, the solution to the Ricci flow tends,
up to scalings, to a metric of positive constant curvature.
Consequently, a compact three-manifold with positive Ricci
curvature is diffeomorphic to the round three-sphere or a metric
quotient of it.

However, for general initial metrics, the Ricci flow may develop
singularities in some parts while it keeps smooth in other parts.
Naturally one would like to cut off the singularities and continue
to run the Ricci flow. If the Ricci flow still develops
singularities after a while, one can do the surgeries and run the
Ricci flow again. By repeating this procedure, one will get a kind
of ``weak" solution to the Ricci flow. Furthermore, if the ``weak"
solution has only a finite number of surgeries at any finite time
interval and one can remember what had been cut during the
surgeries, and if the ``weak" solution has a well-understood
long-time behavior, then one will also get the topology structure
of the initial manifold. This theory of surgically modified Ricci
flow was first developed by Hamilton \cite{Ha97} for compact
four-manifolds and further developed more recently by Perelman
\cite{P2} for compact orientable three-manifolds.

The main purpose of this chapter is to give a complete and
detailed discussion of Perelman's work on the Ricci flow with
surgery on three-manifolds. More specifically, Sections 7.1-7.2
give a detailed exposition of section 12 of Perelman's first paper
\cite{P1}; Sections 7.3-7.6 give a detailed exposition of sections
2-7 of Perelman's second paper \cite{P2}, except the general
collapsing result, Theorem 7.4 of \cite{P2}, claimed by Perelman
(a special case of which has been proved by Shioya-Yamaguchi
\cite{ShY}). In Section 7.7, we combine the long time behavior
result in Section 7.6 and the collapsing result of
Shioya-Yamaguchi \cite{ShY} to present a proof of Thurston's
geometrization conjecture.



\section{Canonical Neighborhood Structures}

Let us call a Riemannian metric on a compact orientable
three-dimen\-sional manifold \textbf{normalized}\index{normalized}
if the eigenvalues of its curvature operator at every point are
bounded by $\frac{1}{10}\geq \lambda \geq \mu \geq \nu \geq
-\frac{1}{10}$, and every geodesic ball of radius one has volume
at least one. By the evolution equation of the curvature and the
maximum principle, it is easy to see that any solution to the
Ricci flow with (compact and three-dimensional) normalized initial
metric exists on a maximal time interval $[0,t_{\max})$ with
$t_{\max}>1$.

Consider a smooth solution $g_{ij}(x,t)$ to the Ricci flow on
$M\times [0,T)$, where $M$ is a compact orientable three-manifold
and $T<+\infty$. After rescaling, we may always assume the initial
metric $g_{ij}(\cdot,0)$ is normalized. By Theorem 5.3.2, the
solution $g_{ij}(\cdot,t)$ then satisfies the pinching estimate
\begin{equation}
R \geq (-\nu)[\log(-\nu) + \log(1+t) -3] %\eqno (7.1.1)
\end{equation} whenever $\nu < 0$ on $M\times [0,T)$. Recall the function
$$
y= f(x) = x(\log x - 3), \ \mbox{ for } e^2 \leq x < +\infty,
$$
is increasing and convex with range $-e^2 \leq y < +\infty$, and
its inverse function is also increasing and satisfies
$$
\lim_{y\rightarrow +\infty} f^{-1}(y)/y = 0.
$$
We can rewrite the pinching estimate (7.1.1) as \begin{equation}
Rm(x,t) \geq-[f^{-1}(R(x,t)(1+t))/(R(x,t)(1+t))]R(x,t) %\eqno(7.1.2)
\end{equation} on $M\times [0,T)$.

Suppose that the solution $g_{ij}(\cdot,t)$ becomes singular as
$t\rightarrow T$. Let us take a sequence of times $t_k \rightarrow
T$, and a sequence of points $p_k \in M$ such that for some
positive constant $C$, $|Rm|(x,t) \leq CQ_k$ with
$Q_k=|Rm(p_k,t_k)|$ for all $x \in M$ and $t \in [0,t_k]$. Thus,
$(p_k,t_k)$ is a sequence of (almost) maximum points. By applying
Hamilton's compactness theorem and Perelman's no local collapsing
theorem I as well as the pinching estimate (7.1.2), a sequence of
the scalings of the solution $g_{ij}(x,t)$ around the points $p_k$
with factors $Q_k$ converges to a nonflat complete
three-dimensional orientable ancient $\kappa$-solution (for some
$\kappa > 0$). For an arbitrarily given $\varepsilon>0$, the
canonical neighborhood theorem (Theorem 6.4.6) in the previous
chapter implies that each point in the ancient $\kappa$-solution
has a neighborhood which is either an evolving $\varepsilon$-neck,
or an evolving $\varepsilon$-cap, or a compact (without boundary)
positively curved manifold. This gives the structure of
singularities coming from a sequence of (almost) maximum points.

However the above argument does not work for singularities coming
from a sequence of points $(y_k,\tau_k)$ with $\tau_k \rightarrow T$
and $|Rm(y_k,\tau_k)| \rightarrow +\infty$ when $|Rm(y_k,\tau_k)|$
is not comparable with the maximum of the curvature at the time
$\tau_k$, since we cannot take a limit directly. In \cite{P1},
Perelman developed a refined rescaling argument to obtain the
following singularity structure theorem. We remark that our
statement of the singularity structure theorem below is slightly
different from Perelman's original statement (cf. Theorem 12.1 of
\cite{P1}). While Perelman assumed the condition of
$\kappa$-noncollapsing on scales less than $r_0$, we assume that the
initial metric is normalized so that from the rescaling argument one
can get the $\kappa$-noncollapsing \emph{on all scales} for the
limit solutions.

%\vskip 0.2cm \noindent \noindent{\bf Theorem 7.1.1}(\textbf
%CZ-SOURCE-LINE-16438
\begin{theorem}[singularity structure of three-dimensional Ricci flow]
  \label{thm:cz-three-dimensional-singularity-structure}
  \dcref{Cao--Zhu Theorem 7.1.1}
  \group{Canonical Neighborhood Structures}
  \source{cao-zhu}
 Given $\varepsilon>0$ and $T_0 >1$, one can
find $r_0>0$ with the following property. If $g_{ij}(x,t), x\in M$
and $t\in [0,T)$ with $1 < T \leq T_0$, is a solution to the Ricci
flow on a compact orientable three-manifold $M$ with normalized
initial metric, then for any point $(x_0,t_0)$ with $t_0\geq 1$
and $Q=R(x_0,t_0)\geq r_0^{-2}$, the solution in $\{(x,t)\ |\
d_{t_0}^2(x,x_0)<\varepsilon^{-2} Q^{-1},\ \ t_0-\varepsilon^{-2}
Q^{-1}\leq t\leq t_0\}$ is, after scaling by the factor $Q$,
$\varepsilon$-close $($in the $C^{[\varepsilon^{-1}]}$-topology$)$
to the corresponding subset of some orientable ancient
$\kappa$-solution $($for some $\kappa >0)$.
\end{theorem}

%\vskip 0.1cm \noindent {\bf Proof.} \
\begin{proof} The proof is basically along the line sketched by Perelman
(cf. section 12 of \cite{P1}). However, the proof of Step 2
follows, with some modifications, the argument given in the
initial notes of Kleiner-Lott \cite{KL} on Perelman's first paper
\cite{P1}. Also, we give a proof of Step 4 which is different from
both Perelman \cite{P1} and Kleiner-Lott \cite{KL}.

Since the initial metric is normalized, it follows from the no
local collapsing theorem I or I' (and their proofs) that there is
a positive constant $\kappa$, depending only on $T_0$, such that
the solution in Theorem 7.1.1 is $\kappa$-noncollapsed on all
scales less than $\sqrt{T_0}$. Let $C(\varepsilon)$ be a positive
constant larger than or equal to $\varepsilon^{-2}$. It suffices
to prove that there exists $r_0
>0$ such that for any point $(x_0,t_0)$ with $t_0 \geq 1$ and
$Q=R(x_0,t_0) \geq r_0^{-2}$, the solution in the parabolic region
$\{ (x,t) \in M \times [0,T)\ | \ d^2_{t_0}(x,x_0)<
C(\varepsilon)Q^{-1}, t_0-C(\varepsilon)Q^{-1}\leq t \leq t_0\}$
is, after scaling by the factor $Q$, $\varepsilon$-close to the
corresponding subset of some orientable ancient $\kappa$-solution.
The constant $C(\varepsilon)$ will be determined later.

We argue by contradiction. Suppose for some $\varepsilon>0$, there
exist a sequence of solutions $(M_k,g_k(\cdot,t))$ to the Ricci
flow on compact orientable three-manifolds with normalized initial
metrics, defined on the time intervals [0,$T_k)$ with $1<T_k \leq
T_0$, a sequence of positive numbers $r_k\rightarrow 0$, and a
sequence of points $x_k\in M_k$ and times $t_k\geq 1$ with
$Q_k=R_k(x_k,t_k)\geq r_k^{-2}$ such that each solution
$(M_k,g_k(\cdot,t))$ in the parabolic region $\{(x,t) \in
M_k\times [0,T_k)\ |\ d^2_{t_k}(x,x_k)<C(\varepsilon)Q_k^{-1},
t_k-C(\varepsilon)Q_k^{-1}\leq t\leq t_k\}$ is not, after scaling
by the factor $Q_k$, $\varepsilon$-close to the corresponding
subset of any orientable ancient $\kappa$-solution, where $R_k$
denotes the scalar curvature of $(M_k,g_k)$.

For each solution $(M_k,g_k(\cdot,t))$, we may adjust the point
$(x_k,t_k)$ with $t_k\geq \frac{1}{2}$ and with $Q_k=R_k(x_k,t_k)$
to be as large as possible so that the conclusion of the theorem
fails at $(x_k,t_k)$, but holds for any $(x,t)\in M_k\times
[t_k-H_kQ_k^{-1},t_k]$ satisfying $R_k(x,t)\geq 2Q_k$, where
$H_k=\frac{1}{4}r_k^{-2}\rightarrow +\infty$ as $k\rightarrow
+\infty.$  Indeed, suppose not, by setting
$(x_{k_1},t_{k_1})=(x_k,t_k)$, we can choose a sequence of points
$(x_{k_l},t_{k_l})\in M_k\times
[t_{k_{(l-1)}}-H_kR_k(x_{k_{(l-1)}},t_{k_{(l-1)}})^{-1},t_{k_{(l-1)}}]$
such that $R_k(x_{k_l},t_{k_l})\geq
2R_k(x_{k_{(l-1)}},t_{k_{(l-1)}})$ and the conclusion of the
theorem fails at $(x_{k_l},t_{k_l})$ for each $l =2, 3, \ldots.$
 Since the solution is smooth, but
$$
R_k(x_{k_l},t_{k_l})\geq 2R_k(x_{k_{(l-1)}},t_{k_{(l-1)}})\geq
\cdots \geq 2^{l-1}R_k(x_{k},t_{k}),
$$
and
\begin{align*}
t_{k_l}&  \geq
t_{k_{(l-1)}}-H_kR_k(x_{k_{(l-1)}},t_{k_{(l-1)}})^{-1}\\
&  \geq
t_{k}-H_k\sum\limits_{i=1}^{l-1}\frac{1}{2^{i-1}}R_k(x_{k},t_{k})^{-1}\\
&  \geq  \frac{1}{2},
\end{align*}
this process must terminate after a finite number of steps and the
last element fits.

Let\; $(M_k,\,\tilde{g}_k(\cdot,t),\,x_k)$\, be\, the\, rescaled\,
solutions\, obtained\, by\, rescaling $(M_k,g_k(\cdot,t))$ around
$x_k$ with the factors $Q_k=R_k(x_k,t_k)$ and shifting the time
$t_k$ to the new time zero. Denote by $\tilde{R}_k$ the rescaled
scalar curvature.  We will show that a subsequence of the
orientable rescaled solutions $(M_k,\tilde{g}_k(\cdot,t),x_k)$
converges in the $C^{\infty}_{loc}$ topology to an orientable
ancient $\kappa$-solution, which is a contradiction. In the
following we divide the argument into four steps.

\medskip
{\it Step} 1. \ First of all, we need a local bound on curvatures.
The following lemma is the Claim 1 of Perelman in his proof Theorem
12.1 of \cite{P1}.

%\vskip 0.1cm \noindent {\bf Lemma 7.1.2} \ \ \emph{
%CZ-SOURCE-LINE-16534
\begin{lemma}[backward parabolic curvature control near blow-up points]
  \label{lem:cz-blowup-backward-parabolic-curvature-control}
  \dcref{Cao--Zhu Lemma 7.1.2}
  \group{Canonical Neighborhood Structures}
  \source{cao-zhu}

For each $(\bar{x},\bar{t})$ with $t_k-\frac{1}{2}H_kQ_k^{-1}\leq
\bar{t}\leq t_k$, we have $R_k(x,t)\leq 4\bar{Q}_k$ whenever
$\bar{t}-c\bar{Q}_k^{-1}\leq t\leq \bar{t}$ and
$d^2_{\bar{t}}(x,\bar{x})\leq c\bar{Q}_k^{-1}$, where
$\bar{Q}_k=Q_k+R_k(\bar{x},\bar{t})$ and $c>0$ is a small
universal constant.
\end{lemma}

%\vskip 0.1cm \noindent {\bf Proof.}
\begin{proof}
  \uses{thm:cz-ancient-kappa-elliptic-estimates}
Consider any point $(x,t)\in B_{\bar{t}}(\bar{x},
(c\bar{Q}_k^{-1})^{\frac{1}{2}})\times
[\bar{t}-c\bar{Q}_k^{-1},\bar{t}]$ with $c>0$ to be determined. If
$R_k(x,t)\leq 2Q_k$, there is nothing to show. If $R_k(x,t)>
2Q_k$, consider a space-time curve $\gamma$ from $(x,t)$ to
$(\bar{x},\bar{t})$ that goes straight from $(x,t)$ to
$(x,\bar{t})$ and goes from $(x,\bar{t})$ to $(\bar{x},\bar{t})$
along a minimizing geodesic (with respect to the metric
$g_k(\cdot,\bar{t}))$. If there is a point on $\gamma$ with the
scalar curvature $2Q_k$, let $y_0$ be the nearest such point to
$(x,t)$. If not, put $y_0=(\bar{x},\bar{t}).$ On the segment of
$\gamma$ from $(x,t)$ to $y_0$, the scalar curvature is at least
$2Q_k$. According to the choice of the point $(x_k,t_k)$, the
solution along the segment is $\varepsilon$-close to that of some
ancient $\kappa$-solution. It follows from Theorem 6.4.3 (ii) that
$$
|\nabla(R_k^{-\frac{1}{2}})|\leq 2\eta \; \mbox{ and } \;
\left|\frac{\partial}{\partial t}(R_k^{-1})\right|\leq 2\eta
$$
on the segment. (Here, without loss of generality, we may assume
$\varepsilon$ is suitably small). Then by choosing $c>0$
(depending only on $\eta$) small enough we get the desired
curvature bound by integrating the above derivative estimates
along the segment. This proves the lemma.
\end{proof}

\smallskip
{\it Step} 2. \  Next we want to show that for each $A<+\infty$,
there exist a positive constant $C(A)$ (independent of $k$) such
that the curvatures of the rescaled solutions
$\tilde{g}_k(\cdot,t)$ at the new time $t =0$ (corresponding to
the original times $t_k$) satisfy the estimate
$$
|\widetilde{Rm}_{k}|(y,0) \leq C(A)
$$
whenever $d_{\tilde{g}_k(\cdot,0)}(y,x_k) \leq A$ and $k \geq 1$.
(This is just a weak version of the Claim 2 of Perelman in his
proof of Theorem 12.1 of \cite{P1}. The first detailed exposition
on this weak version was given by Kleiner-Lott in their June 2003
notes \cite{KL}. We now follow their argument as in \cite{KL} here
with some modifications.)

For each $\rho \geq 0$, set
$$
M(\rho)=\sup\{\tilde{R}_k(x,0)\ |\ k\geq 1, x\in M_k \; \mbox{
with } \;d_{0}(x,x_k)\leq \rho\}
$$
and
$$
\rho_0=\sup\{\rho \geq 0\  |\ \ M(\rho)<+\infty \}.
$$
By the pinching estimate (7.1.1), it suffices to show
$\rho_0=+\infty.$

Note that $\rho_0>0$ by applying Lemma 7.1.2 with
$(\bar{x},\bar{t})=(x_k,t_k).$ We now argue by contradiction to
show $\rho_0=+\infty.$ Suppose not, we may find (after passing to
a subsequence if necessary) a sequence of points $y_k\in M_k$ with
$d_{0}(x_k,y_k)\rightarrow \rho_0<+\infty$ and
$\tilde{R}_k(y_k,0)\rightarrow +\infty$. Let $\gamma_k(\subset
M_k)$ be a minimizing geodesic segment from $x_k$ to $y_k$. Let
$z_k\in \gamma_k$ be the point on $\gamma_k$ closest to $y_k$ with
$\tilde{R}_k(z_k,0)=2$, and let $\beta_k$ be the subsegment of
$\gamma_k$ running from $z_k$ to $y_k$. By Lemma 7.1.2 the length
of $\beta_k$ is bounded away from zero independent of $k$. By the
pinching estimate (7.1.1), for each $\rho<\rho_0$, we have a
uniform bound on the curvatures on the open balls
$B_{0}(x_k,\rho)\subset (M_k,\tilde{g}_k)$. The injectivity radii
of the rescaled solutions $\tilde{g}_k$ at the points $x_k$ and
the time $t=0$ are also uniformly bounded from below by the
$\kappa$-noncollapsing property. Therefore by Lemma 7.1.2 and
Hamilton's compactness theorem (Theorem 4.1.5), after passing to a
subsequence, we can assume that the marked sequence
$(B_{0}(x_k,\rho_0),\tilde{g}_k(\cdot,0),x_k)$ converges in the
$C_{\rm loc}^{\infty}$ topology to a marked (noncomplete) manifold
$(B_{\infty},\tilde{g}_{\infty},x_{\infty})$, the segments
$\gamma_k$ converge to a geodesic segment (missing an endpoint)
$\gamma_{\infty}\subset B_{\infty}$ emanating from $x_{\infty}$,
and $\beta_k$ converges to a subsegment $\beta_{\infty}$ of
$\gamma_{\infty}$. Let $\bar{B}_{\infty}$ denote the completion of
$(B_{\infty},\tilde{g}_{\infty})$, and $y_{\infty}\in
\bar{B}_{\infty}$ the limit point of $\gamma_{\infty}$.

Denote by $\tilde{R}_{\infty}$ the scalar curvature of
$(B_{\infty},\tilde{g}_{\infty})$. Since the rescaled scalar
curvatures $\tilde{R}_k$ along $\beta_k$ are at least 2, it
follows from the choice of the points $(x_k,0)$ that for any
$q_0\in \beta_{\infty}$, the manifold
$(B_{\infty},\tilde{g}_{\infty})$ in $\{q\in B_{\infty}| \ \ {\rm
dist}^2_{\tilde{g}_{\infty}}(q,q_0)<
C(\varepsilon)(\tilde{R}_{\infty}(q_0))^{-1}\}$ is
$2\varepsilon$-close to the corresponding subset of (a time slice
of) some orientable ancient $\kappa$-solution. Then by Theorem
6.4.6, we know that the orientable ancient $\kappa$-solution at
each point $(x,t)$ has a radius $r$, $0 < r <
C_1(2\varepsilon)R(x,t)^{-\frac{1}{2}}$, such that its canonical
neighborhood $B$, with $B_t(x,r) \subset  B \subset B_t(x,2r)$, is
either an evolving $2\varepsilon$-neck, or an evolving
$2\varepsilon$-cap, or a compact manifold (without boundary)
diffeomorphic to a metric quotient of the round three-sphere
$\mathbb{S}^3$, and moreover the scalar curvature is between
$(C_2(2\varepsilon))^{-1}R(x,t)$ and $C_2(2\varepsilon)R(x,t)$,
where $C_1(2\varepsilon)$ and $C_2(2\varepsilon)$ are the positive
constants in Theorem 6.4.6.

We now choose $C(\varepsilon) = \max \{
2C_1^2(2\varepsilon),\varepsilon^{-2}\}$. By the local curvature
estimate in Lemma 7.1.2, we see that the scalar curvature
$\tilde{R}_{\infty}$ becomes unbounded when approaching
$y_{\infty}$ along $\gamma_{\infty}$. This implies that the
canonical neighborhood around $q_0$ cannot be a compact manifold
(without boundary) diffeomorphic to a metric quotient of the round
three-sphere $\mathbb{S}^3$. Note that $\gamma_{\infty}$ is
shortest since it is the limit of a sequence of shortest
geodesics. Without loss of generality, we may assume $\varepsilon$
is suitably small (say, $\varepsilon \leq \frac{1}{100}$). These
imply that as $q_0$ gets sufficiently close to $y_{\infty}$, the
canonical neighborhood around $q_0$ cannot be an evolving
$2\varepsilon$-cap. Thus we conclude that each $q_0 \in
\gamma_{\infty}$ sufficiently close to $y_{\infty}$ is the center
of an evolving $2\varepsilon$-neck.

Let
$$
U=\bigcup\limits_{q_0\in
\gamma_{\infty}}B(q_0,4\pi(\tilde{R}_{\infty}(q_0))^{-\frac{1}{2}})\
\ (\subset (B_{\infty}, \tilde{g}_{\infty})),
$$
where $B(q_0,4\pi(\tilde{R}_{\infty}(q_0))^{-\frac{1}{2}})$ is the
ball centered at $q_0\in B_{\infty}$ with the radius
$4\pi(\tilde{R}_{\infty}(q_0))^{-\frac{1}{2}}$. Clearly $U$ has
nonnegative sectional curvature by the pinching estimate (7.1.1).
Since the metric $\tilde{g}_{\infty}$ is cylindrical at any point
$q_0\in \gamma_{\infty}$ which is sufficiently close to
$y_{\infty}$, we see that the metric space $\overline{U}=U\cup
\{y_{\infty}\}$ by adding in the point $y_{\infty}$, is locally
complete and intrinsic near $y_{\infty}$. Furthermore $y_{\infty}$
cannot be an interior point of any geodesic segment in
$\overline{U}$. This implies the curvature of $\overline{U}$ at
$y_{\infty}$ is nonnegative in the Alexandrov sense. It is a basic
result in Alexandrov space theory (see for example Theorem 10.9.3
and Corollary 10.9.5 of \cite{BBI}) that there exists a
three-dimensional tangent cone $C_{y_{\infty}}\overline{U}$ at
$y_{\infty}$ which is a metric cone. It is clear that its aperture
is $\leq 10\varepsilon$, thus the tangent cone is nonflat.

Pick a point $p\in C_{y_{\infty}}\overline{U}$ such that the
distance from the vertex $y_{\infty}$ to $p$ is one and it is
nonflat around $p$. Then the ball $B(p,\frac{1}{2})\subset
C_{y_{\infty}}\overline{U}$ is the Gromov-Hausdorff limit of the
scalings of a sequence of balls
$B_0(p_k,s_k)\subset(M_k,\tilde{g}_{k}(\cdot,0))$ by some factors
$a_k$, where $s_k\rightarrow 0^+$. Since the tangent cone is
three-dimensional and nonflat around $p$, the factors $a_k$ must
be comparable with $\tilde{R}_k(p_k,0)$. By using the local
curvature estimate in Lemma 7.1.2, we actually have the
convergence in the $C^{\infty}_{\rm loc}$ topology for the
solutions $\tilde{g}_{k}(\cdot,t)$ on the balls $B_0(p_k,s_k)$ and
over some time interval $t\in [-\delta, 0]$ for some sufficiently
small $\delta>0$. The limiting ball $B(p,\frac{1}{2})\subset
C_{y_{\infty}}\overline{U}$ is a piece of the nonnegative curved
and nonflat metric cone whose radial directions are all Ricci
flat. On the other hand, by applying Hamilton's strong maximum
principle to the evolution equation of the Ricci curvature tensor
as in the proof of Lemma 6.3.1, the limiting ball
$B(p,\frac{1}{2})$ would split off all radial directions
isometrically (and locally). Since the limit is nonflat around
$p$, this is impossible. Therefore we have proved that the
curvatures of the rescaled solutions $\tilde{g}_k(\cdot,t)$ at the
new times $t=0$ (corresponding to the original times $t_k$) stay
uniformly bounded at bounded distances from $x_k$ for all $k$.

We have proved that for each $A<+\infty$, the curvature of the
marked manifold $(M_k,\tilde{g}_k(\cdot,0),x_k)$ at each point
$y\in M_k$ with distance from $x_k$ at most $A$ is bounded by
$C(A)$. Lemma 7.1.2 extends this curvature control to a backward
parabolic neighborhood centered at $y$ whose radius depends only
on the distance from $y$ to $x_k$. Thus by Shi's local derivative
estimates (Theorem 1.4.2) we can control all derivatives of the
curvature in such backward parabolic neighborhoods. Then by using
the $\kappa$-noncollapsing and Hamilton's compactness theorem
(Theorem 4.1.5), we can take a $C_{\rm loc}^{\infty}$ subsequent
limit to obtain $(M_{\infty},
\tilde{g}_{\infty}(\cdot,t),x_{\infty})$, which is
$\kappa$-noncollapsed on all scales and is defined on a space-time
open subset of $M_{\infty}\times(-\infty,0]$ containing the time
slice $M_{\infty}\times\{0\}$. Clearly it follows from the
pinching estimate (7.1.1) that the limit $(M_{\infty},
\tilde{g}_{\infty}(\cdot,0),x_{\infty})$ has nonnegative curvature
operator (and hence nonnegative sectional curvature).

\medskip
{\it Step} 3. \ We further claim that the limit $(M_{\infty},
\tilde{g}_{\infty}(\cdot,0),x_{\infty})$ at the time slice $\{ t=0
\}$ has bounded curvature.

We know that the sectional curvature of the limit $(M_{\infty},
\tilde{g}_{\infty}(\cdot,0),x_{\infty})$ is nonnegative
everywhere. Argue by contradiction. Suppose the curvature of
$(M_{\infty}, \tilde{g}_{\infty}(\cdot,0),x_{\infty})$ is not
bounded, then by Lemma 6.1.4, there exists a sequence of points
$q_j \in M_{\infty}$ diverging to infinity such that their scalar
curvatures $\tilde{R}_{\infty}(q_j,0) \rightarrow +\infty$ as $j
\rightarrow +\infty$ and
$$
\tilde{R}_{\infty}(x,0) \leq 4\tilde{R}_{\infty}(q_j,0)
$$
for $x \in B(q_j, j/\sqrt{\tilde{R}_{\infty}(q_j,0)})\subset
(M_{\infty},\tilde{g}_{\infty}(\cdot,0))$. By combining with Lemma
7.1.2 and the $\kappa$-noncollapsing, a subsequence of the
rescaled and marked manifolds
$(M_{\infty},\tilde{R}_{\infty}(q_j,0)\tilde{g}_{\infty}(\cdot,0),q_j)$
converges in the $C^{\infty}_{\rm loc}$ topology to a smooth
nonflat limit $Y$. By Proposition 6.1.2, the new limit $Y$ is
isometric to a metric product $N \times \mathbb{R}$ for some
two-dimensional manifold $N$. On the other hand, in view of the
choice of the points $(x_k,t_k)$, the original limit $(M_{\infty},
\tilde{g}_{\infty}(\cdot,0),x_{\infty})$ at the point $q_j$ has a
canonical neighborhood which is either a $2\varepsilon$-neck, a
$2\varepsilon$-cap, or a compact manifold (without boundary)
diffeomorphic to a metric quotient of the round $\mathbb{S}^3$. It
follows that for $j$ large enough, $q_j$ is the center of a
$2\varepsilon$-neck of radius
$(\tilde{R}_{\infty}(q_j,0))^{-\frac{1}{2}}$. Without loss of
generality, we may further assume that $2\varepsilon <
\varepsilon_0$, where $\varepsilon_0$ is the positive constant
given in Proposition 6.1.1. Since
$(\tilde{R}_{\infty}(q_j,0))^{-\frac{1}{2}}\rightarrow 0$ as
$j\rightarrow +\infty$, this contradicts Proposition 6.1.1. So the
curvature of $(M_{\infty}, \tilde{g}_{\infty}(\cdot,0))$ is
bounded.

\medskip
{\it Step} 4. \  Finally we want to extend the limit backwards in
time to $-\infty$.

By Lemma 7.1.2 again, we now know that the limiting solution
$(M_{\infty},$ $\tilde{g}_{\infty}(\cdot,t))$ is defined on a
backward time interval $[-a,0]$ for some $a>0$.

Denote by
\begin{align*}
t' = \inf \{ \tilde{t}&  | \mbox{ we can take a
smooth limit on }\; (\tilde{t},0]  \mbox { (with bounded} \\
& \mbox {curvature at each time slice) from a subsequence}\\
&  \mbox {of the convergent rescaled solutions }\; \tilde{g}_k\}.
\end{align*}
We first claim that there is a subsequence of the rescaled
solutions $\tilde{g}_k$ which converges in the $C^{\infty}_{\rm
loc}$ topology to a smooth limit
$(M_{\infty},\tilde{g}_{\infty}(\cdot,t))$ on the maximal time
interval $(t',0]$.

Indeed, let $t'_k$ be a sequence of negative numbers such that
$t'_k\rightarrow t'$ and there exist smooth limits
$(M_{\infty},\tilde{g}_{\infty}^{k}(\cdot,t))$ defined on
$(t'_k,0]$. For each $k$, the limit has nonnegative sectional
curvature and has bounded curvature at each time slice. Moreover
by Lemma 7.1.2, the limit has bounded curvature on each
subinterval $[-b,0]\subset(t'_k,0]$. Denote by $\tilde{Q}$ the
scalar curvature upper bound of the limit at time zero (where
$\tilde{Q}$ is the same for all $k$). Then we can apply the
Li-Yau-Hamilton estimate (Corollary 2.5.7) to get
$$
\tilde{R}_{\infty}^{k}(x,t)\leq\tilde{Q} \left(\frac{-t'_k}{t-t'_k}\right),
$$
where $\tilde{R}_{\infty}^{k}(x,t)$ are the scalar curvatures of
the limits $(M_{\infty},\tilde{g}_{\infty}^{k}(\cdot,t))$.  Hence
by the definition of convergence and the above curvature
estimates, we can find a subsequence of the above convergent
rescaled solutions $\tilde{g}_k$ which converges in the
$C^{\infty}_{\rm loc}$ topology to a smooth limit
$(M_{\infty},\tilde{g}_{\infty}(\cdot,t))$ on the maximal time
interval $(t',0]$.

We next claim that $t'=-\infty$.

Suppose not, then by Lemma 7.1.2, the curvature of the limit
$(M_{\infty},$ $\tilde{g}_{\infty}(\cdot,t))$ becomes unbounded as
$t\rightarrow t'>-\infty$. By applying the maximum principle to
the evolution equation of the scalar curvature, we see that the
infimum of the scalar curvature is nondecreasing in time. Note
that $\tilde{R}_{\infty}(x_{\infty},0) =1$. Thus there exists some
point $y_{\infty}\in M_{\infty}$ such that
$$
\tilde{R}_{\infty}\left(y_{\infty},t'+\frac{c}{10}\right)<\frac{3}{2}
$$
where $c>0$ is the universal constant in Lemma 7.1.2. By using
Lemma 7.1.2 again we see that the limit
$(M_{\infty},\tilde{g}_{\infty}(\cdot,t))$ in a small neighborhood
of the point $(y_{\infty},t'+\frac{c}{10})$ extends backwards to
the time interval $[t'-\frac{c}{10},t'+\frac{c}{10}]$. We remark
that the distances at time $t$ and time $0$ are roughly equivalent
in the following sense \begin{equation}
d_t(x,y) \geq d_0(x,y) \geq d_t(x,y)-{\rm const.}  %\eqno(7.1.3)
\end{equation} for any $x,y \in M_{\infty}$ and $t \in (t',0]$. Indeed from
the Li-Yau-Hamilton inequality (Corollary 2.5.7) we have the
estimate
$$
\tilde{R}_\infty(x,t) \le \tilde{Q}\left(\frac{-t'}{t-t'}\right),\quad
\mbox{on }\; M_\infty\times(t',0].
$$
By applying Lemma 3.4.1 (ii), we have
$$
d_t(x,y)\le d_0(x,y)+30(-t')\sqrt{\tilde{Q}}
$$
for any $x,y\in M_\infty$ and $t\in(t',0]$. On the other hand,
since the curvature of the limit metric
$\tilde{g}_\infty(\cdot,t)$ is nonnegative, we have
$$
d_t(x,y)\ge d_0(x,y)
$$
for any $x,y\in M_\infty$ and $t\in(t',0]$. Thus we obtain the
estimate (7.1.3).

Let us still denote by $(M_k,\tilde{g}_k(\cdot,t))$ the
subsequence which converges on the maximal time interval
$(t'\!,0]$. Consider the rescaled sequence
$(\!M_k,\tilde{g}_k(\cdot,t))$ with the marked points $x_k$
replaced by the associated sequence of points $y_k \rightarrow
y_{\infty}$ and the (original unshifted) times $t_k$ replaced by
any $s_k \in
[t_k+(t'-\frac{c}{20})Q_k^{-1},t_k+(t'+\frac{c}{20})Q_k^{-1}]$. It
follows from Lemma 7.1.2 that for $k$ large enough, the rescaled
solutions $(M_k,\tilde{g}_k(\cdot,t))$ at $y_k$ satisfy
$$
\tilde{R}_k(y_k,t) \leq 10
$$
for all $t \in [t'-\frac{c}{10},t'+\frac{c}{10}]$. By applying the
same arguments as in the above Step 2, we conclude that for any $A
> 0$, there is a positive constant $C(A)<+\infty$ such that
$$
\tilde{R}_k(x,t) \leq C(A)
$$
for all $(x,t)$ with ${d}_t(x,y_k) \leq A$ and $t \in
[t'-\frac{c}{20},t'+\frac{c}{20}]$. The estimate (7.1.3) implies
that there is a positive constant $A_0$ such that for arbitrarily
given small $\epsilon' \in(0,\frac{c}{100})$, for $k$ large
enough, there hold
$$
{d}_t(x_k,y_k) \leq A_0
$$
for all $t \in [t'+ \epsilon',0]$. By combining with Lemma 7.1.2,
we then conclude that for any $A>0$, there is a positive constant
$\tilde{C}(A)$ such that for $k$ large enough, the rescaled
solutions $(M_k,\tilde{g}_k(\cdot,t))$ satisfy
$$
\tilde{R}_k(x,t) \leq \tilde{C}(A)
$$
for all $x \in \tilde{B}_0(x_k,A)$ and $t \in [t' -
\frac{c}{100}(C(A))^{-1},0]$.

Now, by taking convergent subsequences from the (original)
rescaled solutions $(M_k,\tilde{g}_k(\cdot,t),x_k)$, we see that
the limiting solution $(M_{\infty},\tilde{g}_{\infty}(\cdot,t))$
is defined on a space-time open subset of $M_{\infty} \times
(-\infty,0]$ containing $M_{\infty} \times [t',0]$. By repeating
the argument of Step 3 and using Lemma 7.1.2, we further conclude
the limit $(M_{\infty},\tilde{g}_{\infty}(\cdot,t))$ has uniformly
bounded curvature on $M_{\infty} \times [t',0]$. This is a
contradiction.

Therefore we have proved a subsequence of the rescaled solutions
$(M_k,$ $\tilde{g}_k(\cdot,t),x_k)$ converges to an orientable
ancient $\kappa$-solution, which gives the desired contradiction.
This completes the proof of the theorem.
%$$\eqno \#$$ \vskip 0.3cm
\end{proof}

We remark that this singularity structure theorem has been
extended by Chen and the second author in \cite{CZ05F} to the
Ricci flow on compact four-manifolds with positive isotropic
curvature.


\section{Curvature Estimates for Smooth Solutions}

Let us consider solutions to the Ricci flow on compact orientable
three-manifolds with normalized initial metrics. The above
singularity structure theorem of Perelman (Theorem 7.1.1) tells us
that the solutions around high curvature points are sufficiently
close to ancient $\kappa$-solutions. It is thus reasonable to
expect that the elliptic type estimate (Theorem 6.4.3) and the
curvature estimate via volume growth (Theorem 6.3.3) for ancient
$\kappa$-solutions are heritable to general solutions of the Ricci
flow on three-manifolds. The main purpose of this section is to
establish such curvature estimates. In the fifth section of this
chapter, we will further extend these estimates to surgically
modified solutions.

The first result of this section is an extension of the elliptic
type estimate (Theorem 6.4.3) by Perelman (cf. Theorem 12.2 of
\cite{P1}). This result is reminiscent of the second step in the
proof of Theorem 7.1.1.

%\vskip 0.2cm\noindent{\bf Theorem 7.2.1} (Perelman \cite{P1}) \ \emph{
%CZ-SOURCE-LINE-16941
\begin{theorem}[curvature control from almost-Euclidean volume]
  \label{thm:cz-almost-euclidean-volume-curvature-control}
  \dcref{Cao--Zhu Theorem 7.2.1}
  \group{Curvature Estimates for Smooth Solutions}
  \source{cao-zhu}

For any $A<+\infty$, there exist $K=K(A)<+\infty$ and
$\alpha=\alpha(A)>0$ with the following property. Suppose we have
a solution to the Ricci flow on a three-dimensional, compact and
orientable manifold $M$ with normalized initial metric. Suppose
that for some $x_0\in M$ and some $r_0>0$ with $r_0<\alpha$, the
solution is defined for $0\le t\le r_0^2$ and satisfies
$$
|Rm|(x,t)\le r_0^{-2},\quad \text{for }\; 0\le t\le
r_0^2,\;d_0(x,x_0)\le r_0,
$$
and
$$
\Vol_0(B_0(x_0,r_0))\ge A^{-1}r_0^3.
$$
Then $R(x,r_0^2)\le Kr_0^{-2}$ whenever $d_{r_0^2}(x,x_0)<Ar_0$.
\end{theorem}

% \vskip 0.1cm\noindent{\bf Proof.}
\begin{proof}
  \uses{lem:cz-local-ancient-model-alternative}
  \uses{thm:cz-three-dimensional-singularity-structure, thm:cz-ancient-kappa-elliptic-estimates, lem:cz-blowup-backward-parabolic-curvature-control}
Given any large $A>0$ and letting $\alpha>0$ be chosen later,  by
Perelman's no local collapsing theorem II (Theorem 3.4.2), there
exists a positive constant $\kappa=\kappa(A)$ (independent of
$\alpha$) such that any complete solution satisfying the
assumptions of the theorem is $\kappa$-noncollapsed on scales $\le
r_0$ over the region $\{(x,t)\ |\ \frac{1}{5}r_0^2\le t\le
r_0^2,\;d_t(x,x_0)\le 5Ar_0\}.$ Set
$$
\varepsilon = \min \left\{ \frac{1}{4}\varepsilon_0,
\frac{1}{100}\right\},
$$
where $\varepsilon_0$ is the positive constant in Proposition
6.1.1. We first prove the following assertion.

\medskip
{\bf Claim.} \ For the above fixed $\varepsilon>0$, one can find
$K=K(A,\varepsilon)<+\infty$ such that if we have a
three-dimensional complete orientable solution with normalized
initial metric and satisfying
$$
|Rm|(x,t)\le r_0^{-2}\quad \mbox{for }\; 0\le t\le
r_0^2,\;d_0(x,x_0)\le r_0,
$$
and
$$
\Vol_0(B_0(x_0,r_0))\ge A^{-1}r_0^3
$$
for some $x_0 \in M$ and some $r_0 > 0$, then for any point $x\in
M$ with $d_{r_0^2}(x,x_0)<3Ar_0$, either $$R(x,r_0^2)< Kr_0^{-2}$$
or the subset $\{(y,t)\ |\ d_{r_0^2}^2(y,x)\leq \varepsilon^{-2}
R(x,r_0^2)^{-1},\;r_0^2-\varepsilon^{-2} R(x,r_0^2)^{-1}\le t\le
r_0^2\}$ around the point $(x,r_0^2)$ is $\varepsilon$-close to
the corresponding subset of an orientable ancient
$\kappa$-solution.
\medskip

Notice that in this assertion we don't impose the restriction of
$r_0<\alpha$, so we can consider for the moment $r_0>0$ to be
arbitrary in proving the above claim. Note that the assumption on
the normalization of the initial metric is just to ensure the
pinching estimate. By scaling, we may assume $r_0=1$. The proof of
the claim is essentially adapted from that of Theorem 7.1.1. But
we will meet the difficulties of adjusting points and verifying a
local curvature estimate.

Suppose that the claim is not true. Then there exist a sequence of
solutions $(M_k,g_k(\cdot,t))$ to the Ricci flow satisfying the
assumptions of the claim with the origins $x_{0_k}$, and a
sequence of positive numbers $K_k\rightarrow\infty$, times $t_k=1$
and points $x_k\in M_k$ with $d_{t_k}(x_k,x_{0_k})<3A$ such that
$Q_k=R_k(x_k,t_k)\ge K_k$ and the solution in $\{(x,t)\ |\
t_k-C(\varepsilon)Q_k^{-1}\le t\le t_k,\;d_{t_k}^2(x,x_k)\le
C(\varepsilon)Q_k^{-1}\}$ is not, after scaling by the factor
$Q_k$, $\varepsilon$-close to the corresponding subset of any
orientable ancient $\kappa$-solution, where $R_k$ denotes the
scalar curvature of $(M_k,g_k(\cdot,t))$ and $C(\varepsilon)(\geq
\varepsilon^{-2})$ is the constant defined in the proof of Theorem
7.1.1. As before we need to first adjust the point $(x_k,t_k)$
with $t_k\ge\frac{1}{2}$ and $d_{t_k}(x_k,x_{0_k})<4A$ so that
$Q_k=R_k(x_k,t_k)\ge K_k$ and the conclusion of the claim fails at
$(x_k,t_k)$, but holds for any $(x,t)$ satisfying $R_k(x,t)\ge
2Q_k,$ $t_k-H_kQ_k^{-1}\le t\le t_k$ and
$d_t(x,x_{0_k})<d_{t_k}(x_k,x_{0_k})
+H_k^{\frac{1}{2}}Q_k^{-\frac{1}{2}}$, where
$H_k=\frac{1}{4}K_k\rightarrow\infty,$ as $k\rightarrow+\infty$.

Indeed, by starting with $(x_{k_1},t_{k_1})=(x_k,1)$ we can choose
$(x_{k_2},t_{k_2})\in M_k\times (0,1]$ with
$t_{k_1}-H_kR_k(x_{k_1},t_{k_1})^{-1}\le t_{k_2}\le t_{k_1}$, and
$d_{t_{k_2}}(x_{k_2},x_{0_k})<d_{t_{k_1}}(x_{k_1},x_{0_k})
+H_k^\frac{1}{2}R_k(x_{k_1},t_{k_1})^{-\frac{1}{2}}$ such that
$R_k(x_{k_2},t_{k_2})\ge 2R_k(x_{k_1},t_{k_1})$ and the conclusion
of the claim fails at $(x_{k_2},t_{k_2})$; otherwise we have the
desired point. Repeating this process, we can choose points
$(x_{k_i},t_{k_i})$, $i=2, \ldots, j$, such that
$$
R_k(x_{k_i},t_{k_i})\ge 2R_k(x_{k_{i-1}},t_{k_{i-1}}),
$$
$$
t_{k_{i-1}}-H_kR_k(x_{k_{i-1}},t_{k_{i-1}})^{-1} \le t_{k_{i}}\le
t_{k_{i-1}},
$$
$$
d_{t_{k_{i}}}(x_{k_{i}},x_{0_{k}})
<d_{t_{k_{i-1}}}(x_{k_{i-1}},x_{0_k})+H_k^\frac{1}{2}
R_k(x_{k_{i-1}},t_{k_{i-1}})^{-\frac{1}{2}},
$$
and the conclusion of the claim fails at the points
$(x_{k_i},t_{k_i})$, $i=2, \ldots, j$. These inequalities imply
$$
R_k(x_{k_{j}},t_{k_{j}})\ge 2^{j-1}R_k(x_{k_1},t_{k_1})\ge
2^{j-1}K_k,
$$
$$
1\ge t_{k_{j}}\ge t_{k_{1}}
-H_k\sum_{i=0}^{j-2}\frac{1}{2^i}R_k(x_{k_1},t_{k_1})^{-1}
\ge\frac{1}{2},
$$
and
$$
d_{t_{k_{j}}}(x_{k_j},x_{0_k}) <d_{t_{k_1}}(x_{k_1},x_{0_k})
+H_k^\frac{1}{2}\sum_{i=0}^{j-2}\frac{1}{(\sqrt{2})^i}
R_k(x_{k_1},t_{k_1})^{-\frac{1}{2}}<4A.
$$
Since the solutions are smooth, this process must terminate after
a finite number of steps to give the desired point, still denoted
by $(x_k,t_k).$

For each adjusted $(x_k,t_k)$, let $[t',t_k]$ be the maximal
subinterval of $[t_k-\frac{1}{2}\varepsilon^{-2}Q_k^{-1},t_k]$ so
that the conclusion of the claim with $K=2Q_k$ holds on
\begin{align*}
& P\left(x_k,t_k,\frac{1}{10}H_k^{\frac{1}{2}}Q_k^{-\frac{1}{2}},t'-t_k\right)\\
&=\left\{ ({x},{t}) \ |\ {x} \in
B_{{t}}\left(x_k,\frac{1}{10}H_k^{\frac{1}{2}}Q_k^{-\frac{1}{2}}\right),
{t} \in [t',t_k]\right\}
\end{align*}
for all sufficiently large $k$. We now want to show $t' =
t_k-\frac{1}{2}\varepsilon^{-2}Q_k^{-1}$.

Consider the scalar curvature $R_k$ at the point $x_k$ over the
time interval $[t',t_k]$. If there is a time $\tilde{t} \in
[t',t_k]$ satisfying $R_k(x_k,\tilde{t}) \geq 2Q_k$, we let
$\tilde{t}$ be the first of such time from $t_k$. Then the
solution $(M_k,g_k(\cdot,t))$ around the point $x_k$ over the time
interval $[\tilde{t}-
\frac{1}{2}\varepsilon^{-2}Q_k^{-1},{\tilde{t}}]$ is
$\varepsilon$-close to some orientable ancient $\kappa$-solution.
Note from the Li-Yau-Hamilton inequality that the scalar curvature
of any ancient $\kappa$-solution is pointwise nondecreasing in
time. Consequently, we have the following curvature estimate
$$
R_k(x_k,t) \leq 2(1+\varepsilon)Q_k
$$
for $ t \in [\tilde{t}- \frac{1}{2}\varepsilon^{-2}Q_k^{-1},t_k]$
(or $t \in [t',t_k]$ if there is no such time $\tilde{t}$).  By
combining with the elliptic type estimate for ancient
$\kappa$-solutions (Theorem 6.4.3) and the Hamilton-Ivey pinching
estimate, we further have \begin{equation}
|Rm(x,t)| \leq 5\omega(1)Q_k  %\eqno (7.2.1)
\end{equation} for all $x \in B_t(x_k,(3Q_k)^{-\frac{1}{2}})$ and $t \in
[\tilde{t}- \frac{1}{2}\varepsilon^{-2}Q_k^{-1},t_k]$ (or $t \in
[t',t_k]$) and all sufficiently large $k$, where $\omega$ is the
positive function in Theorem 6.4.3.

For any point $(x,t)$ with $ \tilde{t}-
\frac{1}{2}\varepsilon^{-2}Q_k^{-1} \leq t \leq t_k$ (or $t \in
[t',t_k]$) and $d_t(x,x_k) \leq
\frac{1}{10}H_k^{\frac{1}{2}}Q_k^{-\frac{1}{2}}$, we divide the
discussion into two cases.

\medskip
{\it Case} (1): $d_t(x_k,x_{0_k}) \leq
\frac{3}{10}H_k^{\frac{1}{2}}Q_k^{-\frac{1}{2}}.$
\begin{align}
d_t(x,x_{0_k})&  \leq
d_t(x,x_k) + d_t(x_k,x_{0_k})\\
&  \leq\frac{1}{10}H_k^{\frac{1}{2}}Q_k^{-\frac{1}{2}}
+ \frac{3}{10}H_k^{\frac{1}{2}}Q_k^{-\frac{1}{2}}\nonumber\\
&  \leq \frac{1}{2}H_k^{\frac{1}{2}}Q_k^{-\frac{1}{2}}.\nonumber
\end{align} %\eqno (7.2.2)

\medskip
{\it Case} (2): $d_t(x_k,x_{0_k}) >
\frac{3}{10}H_k^{\frac{1}{2}}Q_k^{-\frac{1}{2}}.$

\smallskip
{}From the curvature bound $(7.2.1)$ and the assumption, we apply
Lemma 3.4.1(ii) with $r_0 = Q_k^{-\frac{1}{2}}$ to get
$$
\frac{d}{dt}(d_t(x_k,x_{0_k})) \geq -20(\omega(1) +
1)Q_k^{\frac{1}{2}},
$$
and then for $k$ large enough,
\begin{align*}
d_t(x_k,x_{0_k})&  \leq d_{\hat{t}}(x_k,x_{0_k})
+ 20(\omega(1)+1)\varepsilon^{-2}Q_k^{-\frac{1}{2}}\\
&  \leq d_{\hat{t}}(x_k,x_{0_k}) +
\frac{1}{10}H_k^{\frac{1}{2}}Q_k^{-\frac{1}{2}},
\end{align*}
where $\hat{t} \in (t,t_k]$ satisfies the property that
$d_s(x_k,x_{0_k}) \geq
\frac{3}{10}H_k^{\frac{1}{2}}Q_k^{-\frac{1}{2}}$ whenever $s \in
[t,\hat{t}]$. So we have
\begin{align}
d_t(x,x_{0_k})&  \leq d_t(x,x_k) + d_t(x_k,x_{0_k})\\
&  \leq\frac{1}{10}H_k^{\frac{1}{2}}Q_k^{-\frac{1}{2}} +
d_{\hat{t}}(x_k,x_{0_k})
+ \frac{1}{10}H_k^{\frac{1}{2}}Q_k^{-\frac{1}{2}} \nonumber\\
&  \leq d_{t_k}(x_k,x_{0_k}) +
\frac{1}{2}H_k^{\frac{1}{2}}Q_k^{-\frac{1}{2}},\nonumber
\end{align}   %\eqno (7.2.3)
for all sufficiently large $k$. Then the combination of (7.2.2),
(7.2.3) and the choice of the points $(x_k,t_k)$ implies $t'=
t_k-\frac{1}{2}\varepsilon^{-2}Q_k^{-1}$ for all sufficiently
large $k$. (Here we also used the maximality of the subinterval
$[t',t_k]$ in the case that there is no time in $ [t',t_k]$ with
$R_k(x_k,\cdot) \geq 2Q_k$.)

Now we rescale the solutions $(M_k,g_k(\cdot,t))$ into
$(M_k,\tilde{g}_k(\cdot,t))$ around the points $x_k$ by the
factors $Q_k=R_k(x_k,t_k)$ and shift the times $t_k$ to the new
times zero. Then the same arguments from Step 1 to Step 3 in the
proof of Theorem 7.1.1 prove that a subsequence of the rescaled
solutions $(M_k,\tilde{g}_k(\cdot,t))$ converges in the
$C^\infty_{\rm loc}$ topology to a limiting (complete) solution
$(M_{\infty},\tilde{g}_{\infty}(\cdot,t))$, which is defined on a
backward time interval $[-a,0]$ for some $a>0$. (The only
modification is in Lemma 7.1.2 of Step 1 by further requiring
$t_k- \frac{1}{4}\varepsilon^{-2}Q_k^{-1} \leq \bar{t} \leq t_k$).

We next study how to adapt the argument of Step 4 in the proof of
Theorem 7.1.1. As before, we have a maximal time interval
$(t_{\infty},0]$ for which we can take a smooth limit
$(M_{\infty},\tilde{g}_{\infty}(\cdot,t),x_{\infty})$ from a
subsequence of the rescaled solutions
$(M_k,\tilde{g}_k(\cdot,t),x_k)$. We want to show
$t_{\infty}=-\infty$.

Suppose not; then $t_{\infty} > -\infty.$ Let $c>0$ be a positive
constant much smaller than $\frac{1}{10}\varepsilon^{-2}$. Note
that the infimum of the scalar curvature is nondecreasing in time.
Then we can find some point $y_{\infty} \in M_{\infty}$ and some
time $t=t_{\infty} + \theta$ with $0<\theta<\frac{c}{3}$ such that
$\tilde{R}_{\infty}(y_{\infty},t_{\infty}+\theta) \leq
\frac{3}{2}$.

Consider the (unrescaled) scalar curvature $R_k$ of
$(M_k,{g}_k(\cdot,t))$ at the point $x_k$ over the time interval
$[t_k+(t_{\infty}+\frac{\theta}{2})Q_k^{-1},t_k]$. Since the
scalar curvature $\tilde{R}_{\infty}$ of the limit on
$M_{\infty}\times [t_{\infty}+\frac{\theta}{3},0]$ is uniformly
bounded by some positive constant $C$, we have the curvature
estimate $$R_k(x_k,t) \leq 2CQ_k$$ for all $t \in
[t_k+(t_{\infty}+\frac{\theta}{2})Q_k^{-1},t_k]$ and all
sufficiently large $k$. Then by repeating the same arguments as in
deriving (7.2.1), (7.2.2) and (7.2.3), we deduce that the
conclusion of the claim with $K=2Q_k$ holds on the parabolic
neighborhood
$P(x_k,t_k,\frac{1}{10}H_k^{\frac{1}{2}}Q_k^{-\frac{1}{2}},
(t_{\infty}+\frac{\theta}{2})Q_k^{-1})$ for all sufficiently large
$k$.

Let $(y_k,t_k+(t_{\infty}+\theta_k)Q_k^{-1})$ be a sequence of
associated points and times in the (unrescaled) solutions
$(M_k,g_k(\cdot,t))$ so that after rescaling, the sequence
converges to the $(y_{\infty},t_{\infty}+\theta)$ in the limit.
Clearly $\frac{\theta}{2}\leq \theta_k \leq 2\theta$ for all
sufficiently large $k$. Then, by considering the scalar curvature
$R_k$ at the point $y_k$ over the time interval
$[t_k+(t_{\infty}-\frac{c}{3})Q_k^{-1},t_k+(t_{\infty}+\theta_k)Q_k^{-1}]$,
the above argument (as in deriving the similar estimates
(7.2.1)-(7.2.3)) implies that the conclusion of the claim with
$K=2Q_k$ holds on the parabolic neighborhood
$P(y_k,t_k,\frac{1}{10}H_k^{\frac{1}{2}}Q_k^{-\frac{1}{2}},
(t_{\infty}-\frac{c}{3})Q_k^{-1})$ for all sufficiently large $k$.
In particular, we have the curvature estimate
$$
R_k(y_k,t) \leq 4(1+\varepsilon)Q_k
$$
for $t \in [t_k+(t_{\infty}-\frac{c}{3})Q_k^{-1},
t_k+(t_{\infty}+\theta_k)Q_k^{-1}]$ for all sufficiently large
$k$.

We now consider the rescaled sequence $(M_k,\tilde{g}_k(\cdot,t))$
with the marked points replaced by $y_k$ and the times replaced by
$s_k \in [t_k+(t_{\infty}-\frac{c}{4})Q_k^{-1},
t_k+(t_{\infty}+\frac{c}{4})Q_k^{-1}]$. By applying the same
arguments from Step 1 to Step 3 in the proof of Theorem 7.1.1 and
the Li-Yau-Hamilton inequality as in Step 4 of Theorem 7.1.1, we
conclude that there is some small constant $a'>0$ such that the
original limit $(M_{\infty},\tilde{g}_{\infty}(\cdot,t))$ is
actually well defined on $M_{\infty} \times [t_{\infty}-a',0]$
with uniformly bounded curvature. This is a contradiction.
Therefore we have checked the claim.

To finish the proof, we next argue by contradiction. Suppose there
exist sequences of positive numbers $K_k\rightarrow+\infty$,
$\alpha_k\rightarrow0$, as $k\rightarrow+\infty$, and a sequence
of solutions $(M_k,g_k(\cdot,t))$ to the Ricci flow satisfying the
assumptions of the theorem with origins $x_{0_k}$ and with radii
$r_{0_k}$ satisfying $r_{0_k}<\alpha_k$ such that for some points
$x_k\in M_k$ with $d_{r^2_{0_k}}(x_k,x_{0_k})<Ar_{0_k}$ we have
\begin{equation}
R(x_k,r^2_{0_k})>K_kr^{-2}_{0_k} %\eqno(7.2.4)
\end{equation} for all $k$. Let $(M_k,\hat{g}_k(\cdot,t),x_{0_k})$ be the
rescaled solutions of $(M_k,g_k(\cdot,t))$ around the origins
$x_{0_k}$ by the factors $r^{-2}_{0_k}$ and shifting the times
$r^2_{0_k}$ to the new times zero. The above claim tells us that
for $k$ large, any point $(y,0) \in
(M_k,\hat{g}_k(\cdot,0),x_{0_k})$ with
$d_{\hat{g}_k(\cdot,0)}(y,x_{0_k}) < 3A$ and with the rescaled
scalar curvature $\hat{R}_k(y,0) > K_k$ has a canonical
neighborhood which is either a $2\varepsilon$-neck, or a
$2\varepsilon$-cap, or a compact manifold (without boundary)
diffeomorphic to a metric quotient of the round three-sphere. Note
that the pinching estimate (7.1.1) and the condition
$\alpha_k\rightarrow0$ imply any subsequential limit of the
rescaled solutions $(M_k,\hat{g}_k(\cdot,t),x_{0_k})$ must have
nonnegative sectional curvature. Thus the same argument as in Step
2 of the proof of Theorem 7.1.1 shows that for all sufficiently
large $k$, the curvatures of the rescaled solutions at the time
zero stay uniformly bounded at those points whose distances from
the origins $x_{0_k}$ do not exceed $2A$. This contradicts (7.2.4)
for $k$ large enough.

Therefore we have completed the proof of the theorem.
%$$\eqno \#$$ \vskip 0.3cm
\end{proof}

The next result is a generalization of the curvature estimate via
volume growth in Theorem 6.3.3 (ii) where the condition on the
curvature lower bound over a time interval is replaced by that at
a time slice only.

%\vskip 0.2cm\noindent{\bf Theorem 7.2.2 }(Perelman \cite{P1}) \ \emph{
%CZ-SOURCE-LINE-17277
\begin{theorem}[curvature control from lower volume ratio]
  \label{thm:cz-lower-volume-ratio-curvature-control}
  \dcref{Cao--Zhu Theorem 7.2.2}
  \group{Curvature Estimates for Smooth Solutions}
  \source{cao-zhu}

For any $w>0$ there exist $\tau=\tau(w)>0$, $K=K(w)<+\infty$,
$\alpha=\alpha(w)>0$ with the following property. Suppose we have
a three-dimensional, compact and orientable solution to the Ricci
flow defined on $M\times[0,T)$ with normalized initial metric.
Suppose that for some radius $r_0>0$ with $r_0<\alpha$ and a point
$(x_0,t_0)\in M\times[0,T)$ with $T
> t_0\geq4\tau r^2_0$, the solution on the ball $B_{t_0}(x_0,r_0)$
satisfies
$$
Rm(x,t_0)\geq-r_0^{-2}\ \ \ on\ B_{t_0}(x_0,r_0),
$$
$$
and\ \ \ \Vol_{t_0}(B_{t_0}(x_0,r_0))\geq wr_0^3.
$$
Then $R(x,t)\leq Kr_0^{-2}$ whenever $t\in[t_0-\tau r^2_0,t_0]$
and $d_t(x,x_0)\leq\frac{1}{4}r_0.$
\end{theorem}

%\vskip 0.1cm \noindent{\textbf{Proof.}} \
\begin{proof}
  \uses{thm:cz-local-volume-growth-curvature-estimate, thm:cz-almost-euclidean-volume-curvature-control} The following argument basically follows
the proof of Theorem 12.3 of Perelman \cite{P1}.

 If we knew that
$$
Rm(x,t)\geq-r_0^{-2}
$$
for all $t\in[0,t_0]$ and $d_t(x,x_0)\leq r_0$, then we could just
apply Theorem 6.3.3 (ii) and take $\tau(w)=\tau_0(w)/2$,
$K(w)=C(w)+2B(w)/\tau_0(w)$. Now fix these values of $\tau$ and
$K$.

We argue by contradiction. Consider a three-dimensional, compact
and orientable solution $g_{ij}(t)$ to the Ricci flow with
normalized initial metric, a point $(x_0,t_0)$ and some radius
$r_0>0$ with $r_0<\alpha$, for $\alpha>0$ a sufficiently small
constant to be determined later, such that the assumptions of the
theorem do hold whereas the conclusion does not. We first claim
that we may assume that any other point $(x',t')$ and radius
$r'>0$ with the same property has either $t'>t_0$ or $t'<t_0-2\tau
r_0^2$, or $2r'>r_0$. Indeed, suppose otherwise. Then there exist
$(x_0',t_0')$ and $r_0'$ with $t_0'\in[t_0-2\tau r_0^2,t_0]$ and
$r_0'\leq\frac{1}{2}r_0$, for which the assumptions of the theorem
hold but the conclusion does not. Thus, there is a point $(x,t)$
such that
$$
t\in[t_0'-\tau(r_0')^2,t_0']\subset\left[t_0-2\tau
r^2_0-\frac{\tau}{4}r_0^2,t_0\right]
$$
$$
\mbox{and}\ \ \ R(x,t)> K(r_0')^{-2}\geq4Kr_0^{-2}.
$$
If the point $(x_0',t_0')$ and the radius $r_0'$ satisfy the claim
then we stop, and otherwise we iterate the procedure. Since
$t_0\geq4\tau r^2_0$ and the solution is smooth, the iteration
must terminate in a finite number of steps, which provides the
desired point and the desired radius.

Let $\tau'\geq0$ be the largest number such that \begin{equation}
Rm(x,t)\geq-r_0^{-2}  %\eqno(7.2.5)
\end{equation} whenever $t\in[t_0-\tau'r_0^2,t_0]$ and $d_t(x,x_0)\leq r_0$.
If $\tau'\geq2\tau$, we are done by Theorem 6.3.3 (ii). Thus we
may assume $\tau'<2\tau$.  By applying Theorem 6.3.3(ii), we know
that at time $t'=t_0-\tau'r_0^2$, the ball $B_{t'}(x_0,r_0)$ has
\begin{equation}
\Vol_{t'}(B_{t'}(x_0,r_0))\geq \xi(w) r_0^3 %\eqno(7.2.6)
\end{equation} for some positive constant $\xi(w)$ depending only on $w$. We
next claim that there exists a ball (at time $t'=t_0-\tau'r_0^2$)
$B_{t'}(x',r')\subset B_{t'}(x_0,r_0)$ with \begin{equation}
\Vol_{t'}(B_{t'}(x',r'))\geq\frac{1}{2}\alpha_3(r')^3 %\eqno(7.2.7)
\end{equation} and with \begin{equation}
\frac{r_0}{2} > r'\geq c(w)r_0 %\eqno(7.2.8)
\end{equation} for some small positive constant $c(w)$ depending only on $w$,
where $\alpha_3$ is  the volume of the unit ball $\mathbb{B}^3$ in
the Euclidean space $\mathbb{R}^3$. (The following argument in
deriving (7.2.7) and (7.2.8) is a standard one in the Alexandrov
space theory and has nothing to do with the Ricci flow. Our
presentation here is inspired from Lemma 53.1 of Kleiner-Lott notes
\cite{KL}.)

Indeed, suppose that it is not true. Then after rescaling, there
is a sequence of Riemannian manifolds $M_i,\ i=1,2,\ldots,$ with
balls $B(x_i,1)\subset M_i$ so that
$$
Rm\geq-1\ \ \ \mbox{on}\ B(x_i,1) \leqno(7.2.5)'
$$
and
$$
\Vol(B(x_i,1))\geq \xi(w) \leqno(7.2.6)'
$$
for all $i$, but all balls $B(x_i',r_i')\subset B(x_i,1)$ with
$\frac{1}{2} > r_i'\geq\frac{1}{i}$ satisfy \begin{equation}
\Vol(B(x_i',r_i'))<\frac{1}{2}\alpha_3(r_i')^3. %\eqno(7.2.9)
\end{equation} It follows from basic results in Alexandrov space theory (see
for example Theorem 10.7.2 and Theorem 10.10.10 of \cite{BBI})
that, after taking a subsequence, the marked balls
$(B(x_i,1),x_i)$ converge in the Gromov-Hausdorff topology to a
marked length space $(B_\infty,x_\infty)$ with curvature bounded
from below by $-1$ in the Alexandrov space sense, and the
associated Riemannian volume forms $d\Vol_{M_i}$ over
$(B(x_i,1),x_i)$ converge weakly to the Hausdorff measure $\mu$ of
$B_\infty$. It is well-known that the Hausdorff dimension of any
Alexandrov space is either an integer or infinity (see for example
Theorem 10.8.2 of \cite{BBI}). Then by (7.2.6)$'$, we know the
limit $(B_\infty,x_\infty)$ is a three-dimensional Alexandrov
space of curvature $\geq -1$. In the Alexandrov space theory, a
point $p \in B_\infty$ is said to be
\textbf{regular}\index{regular} if the tangent cone of $B_\infty$
at $p$ is isometric to $\mathbb{R}^3$. It is also a basic result
in Alexandrov space theory (see for example Corollary 10.9.13 of
\cite{BBI}) that the set of regular points in $B_\infty$ is dense
and for each regular point there is a small neighborhood which is
almost isometric to an open set of the Euclidean space
$\mathbb{R}^3$. Thus for any $\varepsilon>0$, there are balls
$B(x'_\infty,r'_\infty)\subset B_\infty$ with $0 < r'_\infty <
\frac{1}{3}$ and satisfying
$$
\mu(B(x'_\infty,r'_\infty))\geq(1-\varepsilon)\alpha_3(r'_\infty)^3.
$$
This is a contradiction with (7.2.9).

Without loss of generality, we may assume
$w\leq\frac{1}{4}\alpha_3$. Since $\tau'<2\tau$, it follows from
the choice of the point $(x_0,t_0)$ and the radius $r_0$ and
(7.2.5), (7.2.7), (7.2.8) that the conclusion of the theorem holds
for $(x',t')$ and $r'$. Thus we have the estimate
$$
R(x,t)\leq K(r')^{-2}
$$
whenever $t\in[t'-\tau(r')^2,t']$ and
$d_t(x,x')\leq\frac{1}{4}r'$. For $\alpha > 0$ small, by combining
with the pinching estimate (7.1.1), we have
$$
|Rm(x,t)|\leq K'(r')^{-2}
$$
whenever $t\in[t'-\tau(r')^2,t']$ and
$d_t(x,x')\leq\frac{1}{4}r'$, where $K'$ is some positive constant
depending only on $K$. Note that this curvature estimate implies
the evolving metrics are equivalent over a suitable subregion of
$\{(x,t)\ |\ t\in[t'-\tau(r')^2,t'] \ \ \mbox{and} \ \
d_t(x,x')\leq\frac{1}{4}r'\}$. Now we can apply Theorem 7.2.1 to
choose $\alpha=\alpha(w)>0$ so small that \begin{equation}
R(x,t)\leq\tilde{K}(w)(r')^{-2}
\leq\tilde{K}(w)c(w)^{-2}r_0^{-2} %\eqno(7.2.10)
\end{equation} whenever $t\in[t'-\frac{\tau}{2}(r')^2,t']$ and
$d_t(x,x')\leq10r_0$. Then the combination of (7.2.10) with the
pinching estimate (7.1.2) would imply
\begin{displaymath}
\begin{split}
Rm(x,t)&  \geq-[f^{-1}(R(x,t)(1+t))/(R(x,t)(1+t))]R(x,t)\\[1mm]
    &  \geq-\frac{1}{2}r_0^{-2}
\end{split}
\end{displaymath}
on the region $\{(x,t)\ |\ t\in[t'-\frac{\tau}{2}(r')^2,t']$ and
$d_t(x,x_0)\leq r_0\}$ when $\alpha=\alpha(w)>r_0$ small enough.
This contradicts the choice of $\tau'$. Therefore we have proved
the theorem.
%$$\eqno \#$$\vskip 0.3cm
\end{proof}

The combination of the above two theorems immediately gives the
following consequence.

%\vskip 0.2cm \noindent{\bf Corollary 7.2.3 } \emph{
%CZ-SOURCE-LINE-17441
\begin{corollary}[localized curvature control under volume and sectional bounds]
  \label{cor:cz-localized-volume-sectional-curvature-control}
  \dcref{Cao--Zhu Corollary 7.2.3}
  \group{Curvature Estimates for Smooth Solutions}
  \source{cao-zhu}

For any $w>0$ and $A<+\infty$, there exist $\tau=\tau(w,A)>0$,
$K=K(w,A)<+\infty$, and $\alpha=\alpha(w,A)>0$ with the following
property. Suppose we have a three-dimensional, compact and
orientable solution to the Ricci flow defined on $M\times[0,T)$
with normalized initial metric. Suppose that for some radius
$r_0>0$ with $r_0 <\alpha$ and a point $(x_0,t_0)\in M\times[0,T)$
with $T > t_0\geq 4\tau r^2_0$, the solution on the ball
$B_{t_0}(x_0,r_0)$ satisfies
$$
Rm(x,t_0)\geq-r_0^{-2}\ \ \ on\ B_{t_0}(x_0,r_0),
$$
$$
and\ \ \ \Vol_{t_0}(B_{t_0}(x_0,r_0))\geq wr_0^3.
$$
Then $R(x,t)\leq Kr_0^{-2}$ whenever $t\in[t_0-\tau r^2_0,t_0]$
and $d_t(x,x_0)\leq Ar_0.$
\end{corollary}
%$$\eqno \#$$ \vskip 0.3cm

We can also state the previous corollary in the following version.

%\vskip 0.2cm \noindent{\bf Corollary 7.2.4 }(Perelman \cite{P1})
%CZ-SOURCE-LINE-17464
\begin{corollary}[pseudolocal curvature control on small balls]
  \label{cor:cz-small-ball-pseudolocal-curvature-control}
  \dcref{Cao--Zhu Corollary 7.2.4}
  \group{Curvature Estimates for Smooth Solutions}
  \source{cao-zhu}

For any $w>0$ one can find $\rho = \rho(w)>0$ such that if
$g_{ij}(t)$ is a complete solution to the Ricci flow defined on
$M\times[0,T)$ with $T>1$ and with normalized initial metric,
where $M$ is a three-dimensional, compact and orientable manifold,
and if $B_{t_0}(x_0,r_0)$ is a metric ball at time $t_0\geq1$,
with $r_0<\rho$, such that$$\min\{ Rm(x,t_0)\ |\ x\in
B_{t_0}(x_0,r_0)\}=-r_0^{-2},$$ then
$$
\Vol_{t_0}(B_{t_0}(x_0,r_0))\leq wr_0^3.
$$
\end{corollary}

%\vskip 0.1cm\noindent{\textbf{Proof.}} \
\begin{proof}
  \uses{cor:cz-localized-volume-sectional-curvature-control}
We argue by contradiction. Suppose for any $\rho>0$, there is a
solution and a ball $B_{t_0}(x_0,r_0)$ satisfying the assumption
of the corollary with $r_0<\rho$, $t_0\geq1$, and with
$$
\min\{ Rm(x,t_0)\ |\ x\in B_{t_0}(x_0,r_0)\}=-r_0^{-2},
$$
but
$$
\Vol_{t_0}(B_{t_0}(x_0,r_0))>wr_0^3.
$$
We can apply Corollary 7.2.3 to get
$$
R(x,t)\leq Kr_0^{-2}
$$
whenever $t\in[t_0-\tau r_0^2,t_0]$ and $d_t(x,x_0)\leq2r_0$,
provided $\rho>0$ is so small that $4\tau\rho^2\leq1$ and
$\rho<\alpha$, where $\tau,\ \alpha$ and $K$ are the positive
constants obtained in Corollary 7.2.3. Then for $r_0<\rho$ and
$\rho>0$ sufficiently small, it follows from the pinching estimate
(7.1.2) that
\begin{displaymath}
\begin{split}
Rm(x,t)&  \geq -[f^{-1}(R(x,t)(1+t))/(R(x,t)(1+t))]R(x,t)\\[1mm]
           &  \geq-\frac{1}{2}r_0^{-2}
\end{split}
\end{displaymath}
in the region $\{(x,t)\ |\ t\in[t_0-\tau(r_0)^2,t_0]$ and
$d_t(x,x_0)\leq 2r_0\}$. In particular, this would imply
$$
\min\{ Rm(x,t_0)|x\in B_{t_0}(x_0,r_0)\}>-r_0^{-2}.
$$
This contradicts the assumption.
%$$\eqno \#$$ \vskip 1cm
\end{proof}


\section{Ricci Flow with Surgery}

One of the central themes of the Ricci flow theory is to give a
classification of all compact orientable three-manifolds. As we
mentioned before, the basic idea is to obtain long-time behavior
of solutions to the Ricci flow. However the solutions will in
general become singular in finite time. Fortunately, we now
understand the precise structures of the solutions around the
singularities, thanks to Theorem 7.1.1. When a solution develops
singularities, one can perform geometric surgeries by cutting off
the canonical neighborhoods around the singularities and gluing
back some known pieces, and then continue running the Ricci flow.
By repeating this procedure, one hopes to get a kind of ``weak"
solution. In this section we will give a detailed description of
this surgery procedure (cf. \cite{Ha97, P2}) and define a global
``weak" solution to the Ricci flow. This section is a detailed
exposition of sections 3 and 4 of Perelman \cite{P2}.

Given any $\varepsilon>0$, based on the singularity structure
theorem (Theorem 7.1.1), we can get a clear picture of the
solution near the singular time as follows.

Let $(M,g_{ij}(\cdot,t))$ be a maximal solution to the Ricci flow
on the maximal time interval $[0,T)$ with $T<+\infty$, where $M$
is a connected compact orientable three-manifold and the initial
metric is normalized. For the given $\varepsilon>0$ and the
solution $(M,g_{ij}(\cdot,t))$, we can find $r_0>0$ such that each
point $(x,t)$ with $R(x,t)\ge r_0^{-2}$ satisfies the derivative
estimates \begin{equation} |\nabla R(x,t)|<\eta R^{\frac{3}{2}}(x,t)\quad
\mbox{and }\; \left|\frac{\partial}{\partial t}R(x,t)\right|<\eta
R^2(x,t),
%\eqno(7.3.1)
\end{equation} where $\eta>0$ is a universal constant, and has a canonical
neighborhood which is either an evolving $\varepsilon$-neck, or an
evolving $\varepsilon$-cap, or a compact positively curved
manifold (without boundary). In the last case the solution becomes
extinct at the maximal time $T$ and the manifold $M$ is
diffeomorphic to the round three-sphere $\mathbb{S}^3$ or a metric
quotient of $\mathbb{S}^3$ by Theorem 5.2.1.

Let $\Omega$ denote the set of all points in $M$ where the
curvature stays bounded as $t\rightarrow T$. The gradient
estimates in (7.3.1) imply that $\Omega$ is open and that
$R(x,t)\rightarrow\infty$ as $t\rightarrow T$ for each $x\in
M\setminus \Omega$.

If $\Omega$ is empty, then the solution becomes extinct at time
$T$. In this case, either the manifold $M$ is compact and
positively curved, or it is entirely covered by evolving
$\varepsilon$-necks and evolving $\varepsilon$-caps shortly before
the maximal time $T$. So the manifold $M$ is diffeomorphic to
either $\mathbb{S}^3$, or a metric quotient of the round
$\mathbb{S}^3$, or $\mathbb{S}^2\times \mathbb{S}^1$, or
$\mathbb{RP}^3\#\mathbb{RP}^3$. The reason is as follows. Clearly,
we only need to consider the situation that the manifold $M$ is
entirely covered by evolving $\varepsilon$-necks and evolving
$\varepsilon$-caps shortly before the maximal time $T$. If $M$
contains a cap $C$, then there is a cap or a neck adjacent to the
neck-like end of $C$. The former case implies that $M$ is
diffeomorphic to $\mathbb{S}^3$, $\mathbb{RP}^3$, or
$\mathbb{RP}^3\#\mathbb{RP}^3$. In the latter case, we get a new
longer cap and continue. Finally, we must end up with a cap,
producing a $\mathbb{S}^3$, $\mathbb{RP}^3$, or
$\mathbb{RP}^3\#\mathbb{RP}^3$. If $M$ contains no caps, we start
with a neck $N$. By connecting with the necks that are adjacent to
the boundary of $N$, we get a longer neck and continue. After a
finite number of steps, the resulting neck must repeat itself.
Since $M$ is orientable, we conclude that $M$ is diffeomorphic to
$\mathbb{S}^2\times \mathbb{S}^1$.

We can now assume that $\Omega$ is nonempty. By using the local
derivative estimates of Shi (Theorem 1.4.2), we see that as
$t\rightarrow T$ the solution $g_{ij}(\cdot,t)$ has a smooth limit
$\bar{g}_{ij}(\cdot)$ on $\Omega$. Let $\bar{R}(x)$ denote the
scalar curvature of $\bar{g}_{ij}$. For any $\rho<r_0$, let us
consider the set
$$
\Omega_\rho=\{x\in\Omega\ |\ \bar{R}(x)\le\rho^{-2}\}.
$$
By the evolution equation of the Ricci flow, we see that the
initial metric $g_{ij}(\cdot,0)$ and the limit metric
$\overline{g}_{ij}(\cdot)$ are equivalent over any fixed region
where the curvature remains uniformly bounded. Note that for any
fixed $x\in
\partial \Omega$, and any sequence of points $x_j\in \Omega$ with
$x_j\rightarrow x$ with respect to the initial metric
$g_{ij}(\cdot,0)$, we have $\overline{R}(x_j)\rightarrow +\infty$.
In fact, if there were a subsequence $x_{j_k}$ so that
$\lim_{k\rightarrow \infty}\overline{R}(x_{j_k})$ exists and is
finite, then it would follow from the gradient estimates (7.3.1)
that $\overline{R}$ is uniformly bounded in some small
neighborhood of $x\in \partial \Omega$ (with respect to the
induced topology of the initial metric $g_{ij}(\cdot,0)$); this is
a contradiction. From this observation and the compactness of the
initial manifold, we see that $\Omega_{\rho}$ is compact (with
respect to the metric $\overline{g}_{ij}(\cdot)$).

For further discussions, let us introduce the following
terminologies. Denote by $\mathbb{I}$ an interval.

Recall that an {\bf $\varepsilon$-neck}\index{$\varepsilon$-neck}
(of radius $r$) is an open set with a Riemannian metric, which is,
after scaling the metric with factor $r^{-2}$, $\varepsilon$-close
to the standard neck $\mathbb{S}^2\times \mathbb{I}$ with the
product metric, where $\mathbb{S}^2$ has constant scalar curvature
one and $\mathbb{I}$ has length $2\varepsilon^{-1}$ and the
$\varepsilon$-closeness refers to the $C^{[\varepsilon^{-1}]}$
topology.

A metric on $\mathbb{S}^2\times \mathbb{I}$, such that each point
is contained in some $\varepsilon$-neck, is called an {\bf
$\varepsilon$-tube}\index{$\varepsilon$-tube}, or an {\bf
$\varepsilon$-horn}\index{$\varepsilon$-horn}, or a {\bf double
$\varepsilon$-horn}\index{double $\varepsilon$-horn}, if the
scalar curvature stays bounded on both ends, or stays bounded on
one end and tends to infinity on the other, or tends to infinity
on both ends, respectively.

A metric on $\mathbb{B}^3$ or $\mathbb{RP}^3\setminus
\bar{\mathbb{B}}^3$ is called a {\bf
$\varepsilon$-cap}\index{$\varepsilon$-cap} if the region outside
some suitable compact subset is an $\varepsilon$-neck. A metric on
$\mathbb{B}^3$ or $\mathbb{RP}^3\setminus \bar{\mathbb{B}}^3$ is
called an {\bf capped $\varepsilon$-horn}\index{capped
$\varepsilon$-horn} if each point outside some compact subset is
contained in an $\varepsilon$-neck and the scalar curvature tends
to infinity on the end.

Now take any $\varepsilon$-neck in $(\Omega,\bar{g}_{ij})$ and
consider a point $x$ on one of its boundary components. If $x\in
\Omega\setminus\Omega_\rho$, then there is either an
$\varepsilon$-cap or an $\varepsilon$-neck, adjacent to the
initial $\varepsilon$-neck. In the latter case we can take a point
on the boundary of the second $\varepsilon$-neck and continue.
This procedure can either terminate when we get into $\Omega_\rho$
or an $\varepsilon$-cap, or go on indefinitely, producing an
$\varepsilon$-horn. The same procedure can be repeated for the
other boundary component of the initial $\varepsilon$-neck.
Therefore, taking into account that $\Omega$ has no compact
components, we conclude that each $\varepsilon$-neck of
$(\Omega,\bar{g}_{ij})$ is contained in a subset of $\Omega$ of
one of the following types:
\begin{align}
\text{(a)}&\quad \text{an $\varepsilon$-tube with boundary
components in $\Omega_\rho$, or} \nonumber\\[1mm]
\text{(b)}&\quad \text{an $\varepsilon$-cap with boundary
in $\Omega_\rho$, or} \nonumber\\[1mm]
\text{(c)}&\quad \text{an $\varepsilon$-horn with boundary
in $\Omega_\rho$, or}\\[1mm]  % (7.3.2)
\text{(d)}&\quad \text{a capped $\varepsilon$-horn, or} \nonumber\\[1mm]
\text{(e)}&\quad \text{a double $\varepsilon$-horn.} \nonumber
\end{align}

Similarly, each $\varepsilon$-cap of $(\Omega,\bar{g}_{ij})$ is
contained in a subset of $\Omega$ of either type (b) or type (d).

It is clear that there is a definite lower bound (depending on
$\rho$) for the volume of subsets of types (a), (b) and (c), so
there can be only a finite number of them. Thus we conclude that
there is only a finite number of components of $\Omega$ containing
points of $\Omega_\rho$, and every such component has a finite
number of ends, each being an $\varepsilon$-horn. On the other
hand, every component of $\Omega$, containing no points of
$\Omega_\rho$, is either a capped $\varepsilon$-horn, or a double
$\varepsilon$-horn. If we look at the solution at a slightly
earlier time, the above argument shows each $\varepsilon$-neck or
$\varepsilon$-cap of $(M,g_{ij}(\cdot,t))$ is contained in a
subset of types (a) or (b), while the $\varepsilon$-horns, capped
$\varepsilon$-horns and double $\varepsilon$-horns (at the maximal
time T) are connected together to form $\varepsilon$-tubes and
$\varepsilon$-caps at the times $t$ shortly before $T$.


\begin{figure}[h]
\begin{picture}(320, 240)
\qbezier(15,90)(0,50)(20,20)\qbezier(20,20)(50,10)(100,50)\qbezier(100,50)(130,60)(170,70)
\qbezier(15,90)(50,140)(100,130)\qbezier(100,130)(130,140)(170,150)\qbezier(130,130)(110,90)(130,70)
\qbezier(130,140)(126,135)(130,130)\qbezier(130,70)(126,65)(130,60)
\qbezier[8](130,140)(134,135)(130,130)\qbezier[8](130,70)(134,65)(130,60)
\qbezier(130,130)(150,135)(170,150)\qbezier(130,70)(150,70)(170,70)
\qbezier(170,70)(190,78)(200,80)\qbezier(170,70)(190,68)(200,70)
\qbezier(200,80)(210,78)(230,70)\qbezier(200,70)(210,68)(230,70)
\qbezier(200,80)(196,75)(200,70)\qbezier[8](200,80)(204,75)(200,70)
\qbezier(230,70)(260,70)(280,70)\qbezier(230,70)(260,60)(280,50)\qbezier(280,70)(310,60)(280,50)
\qbezier(260,70)(256,65)(260,60)\qbezier[8](260,70)(264,65)(260,60)
\qbezier(170,150)(200,160)(220,170)\qbezier(170,150)(200,150)(220,150)
\qbezier[8](200,160)(204,155)(200,150)\qbezier(200,160)(196,155)(200,150)
\qbezier(220,170)(260,230)(290,200)\qbezier(220,150)(260,150)(290,170)\qbezier(290,200)(310,200)(290,170)
\qbezier(230,160)(260,160)(290,190)\qbezier(250,165)(260,185)(280,180)
\qbezier(30,60)(50,60)(80,100)\qbezier(40,65)(50,85)(75,90)
\qbezier(238,160)(230,155)(235,150)\qbezier[8](238,160)(241,155)(235,150)
\qbezier(280,180)(290,175)(290,170)\qbezier[10](280,180)(278,175)(290,170)

\put(30,140){\makebox(0,0)[bl]{$\Omega_{\rho}$}}
\put(40,140){\vector(1,-1){10}}
\put(200,190){\makebox(0,0)[bl]{$\Omega_{\rho}$}}
\put(210,190){\vector(1,-1){10}}
\put(160,115){\makebox(0,0)[bl]{$\varepsilon$-horn}}
\put(160,125){\vector(-1,1){10}}
\put(270,130){\makebox(0,0)[bl]{$\varepsilon$-tube}}
\put(270,140){\vector(-1,1){10}}
\put(160,48){\makebox(0,0)[bl]{double $\varepsilon$-horn}}
\put(190,58){\vector(0,1){10}}
 \put(250,30){\makebox(0,0)[bl]{capped $\varepsilon$-horn}}
 \put(270,40){\vector(0,1){10}}

\end{picture}
\end{figure}



Hence, by looking at the solution at times shortly before $T$, we
see that the topology of $M$ can be reconstructed as follows: take
the components $\Omega_j$, $1\le j\le k$, of $\Omega$ which
contain points of $\Omega_\rho$, truncate their
$\varepsilon$-horns, and glue to the boundary components of
truncated $\Omega_j$ a finite collection of tubes
$\mathbb{S}^2\times\mathbb{I}$ and caps $\mathbb{B}^3$ or
$\mathbb{RP}^3 \setminus \bar{\mathbb{B}}^3$. Thus, $M$ is
diffeomorphic to a connected sum of $\bar{\Omega}_j$, $1\le j\le
k$, with a finite number of copies of $\mathbb{S}^2\times
\mathbb{S}^1$ (which correspond to gluing a tube to two boundary
components of the same $\Omega_j$), and a finite number of copies
of $\mathbb{RP}^3$. Here $\bar{\Omega}_j$ denotes $\Omega_j$ with
each $\varepsilon$-horn one point compactified. More
geometrically, one can get $\bar{\Omega}_j$ in the following way:
in every $\varepsilon$-horn of $\Omega_j$ one can find an
$\varepsilon$-neck, cut it along the middle two-sphere, remove the
horn-shaped end, and glue back a cap (or more precisely, a
differentiable three-ball). Thus to understand the topology of
$M$, one only needs to understand the topologies of the compact
orientable three-manifolds $\bar{\Omega}_j$, $1\le j\le k$.

Naturally one can evolve each $\bar{\Omega}_j$ by the Ricci flow
again and, when singularities develop again, perform the above
surgery for each $\varepsilon$-horn to get new compact orientable
three-manifolds. By repeating this procedure indefinitely, it will
likely give a long-time ``weak" solution to the Ricci flow. The
following abstract definition for this kind of ``weak" solution
was introduced by Perelman in \cite{P2} .

%\vskip 0.2cm \noindent{\bf Definition 7.3.1}\ \
%CZ-SOURCE-LINE-17757
\begin{definition}[Ricci flow with surgery]
  \label{def:cz-ricci-flow-with-surgery}
  \dcref{Cao--Zhu Definition 7.3.1}
  \group{Ricci Flow with Surgery}
  \source{cao-zhu}

Suppose we are given a (finite or countably infinite) collection
of three-dimensional smooth solutions $g^k_{ij}(t)$ to the Ricci
flow defined on $M_k\times [t^-_k,t^+_k)$ and go singular as
$t\rightarrow t^+_k$, where each manifold $M_k$ is compact and
orientable, possibly disconnected with only a finite number of
connected components. Let $(\Omega_k,\bar{g}^k_{ij})$ be the
limits of the corresponding solutions $g^k_{ij}(t)$ as
$t\rightarrow t^+_k$, as above. Suppose also that for each $k$ we
have $t^-_k=t^+_{k-1}$, and that
$(\Omega_{k-1},\bar{g}^{k-1}_{ij})$ and $(M_k,g^k_{ij}(t^-_k))$
contain compact (possibly disconnected) three-dimensional
submanifolds with smooth boundary which are isometric. Then by
identifying these isometric submanifolds, we say the collection of
solutions $g^k_{ij}(t)$ is a solution to the \textbf{Ricci flow
with surgery}\index{Ricci flow with surgery} (or a
\textbf{surgically modified solution} to the Ricci flow)
\index{surgically modified solution} on the time interval which is
the union of all $[t^-_k,t^+_k)$, and say the times $t^+_k$ are
\textbf{surgery
times}\index{surgery times}. %\vskip 0.3cm
\end{definition}

To get the topology of the initial manifold from the solution to
the Ricci flow with surgery, one has to overcome the following two
difficulties:
\begin{itemize}
\item[(i)] how to prevent the surgery times from accumulating?
\item[(ii)] how to obtain the long time behavior of the solution
to the Ricci flow with surgery?
\end{itemize}

Thus it is natural to consider those solutions having ``good"
properties. For any arbitrarily fixed positive number
$\varepsilon$, we will only consider those solutions to the Ricci
flow with surgery which satisfy the following \textbf{a priori
assumptions (with accuracy $\varepsilon$)}\index{a priori
assumptions (with accuracy $\varepsilon$)}.

\medskip
{\bf Pinching assumption.}\index{pinching assumption} \ The
eigenvalues $\lambda \geq \mu \geq \nu$ of the curvature operator
of the solution to the Ricci flow with surgery at each point and
each time satisfy \begin{equation}
R \geq (-\nu)[\log(-\nu) + \log(1+t) -3] %\eqno (7.3.3)
\end{equation} whenever $\nu < 0$.

\medskip
{\bf Canonical neighborhood assumption (with accuracy $\boldsymbol
\varepsilon$).}\index{canonical neighborhood assumption (with
accuracy $\varepsilon$)} \ For any given $\varepsilon >0$, there
exist positive constants $C_1$ and $C_2$ depending only on
$\varepsilon$, and a nonincreasing positive function $r:
[0,+\infty) \rightarrow (0,+\infty)$ such that at each time $t>0$,
every point $x$ where scalar curvature $R(x,t)$ is at least
$r^{-2}(t)$ has a neighborhood $B$, with $B_t(x,\sigma)\subset B
\subset B_t(x,2\sigma)$ for some $ 0<\sigma <
C_1R^{-\frac{1}{2}}(x,t)$, which falls into one of the following
three categories:
\begin{itemize}
\item[(a)] $B$ is a \textbf{strong
$\varepsilon$-neck}\index{$\varepsilon$-neck!strong} (in the sense
$B$ is the slice at time $t$ of the parabolic neighborhood
$\{(x',t')\ |\ x'\in B, t'\in[t-R(x,t)^{-1},t]\}$, where the
solution is well defined on the whole parabolic neighborhood and
is, after scaling with factor $R(x,t)$ and shifting the time to
zero, $\varepsilon$-close (in the $C^{[\varepsilon^{-1}]}$
topology) to the subset $(\mathbb{S}^2 \times \mathbb{I}) \times
[-1,0]$ of the evolving standard round cylinder with scalar
curvature 1 to $\mathbb{S}^2$ and length $2\varepsilon^{-1}$ to
$\mathbb{I}$ at time zero), or \item[(b)] $B$ is an
$\varepsilon$-cap, or \item[(c)] $B$ is a compact manifold
(without boundary) of positive sectional curvature.
\end{itemize}
Furthermore, the scalar curvature in $B$ at time $t$ is between
$C^{-1}_2R(x,t)$ and $C_2R(x,t)$, satisfies the gradient estimates
\begin{equation} |\nabla R|<\eta R^{\frac{3}{2}}\;  \mbox{ and }\;
\left|\frac{\partial R}{\partial t}\right|<\eta R^2, %\eqno (7.3.4)
\end{equation} and the volume of $B$ in case (a) and case (b) satisfies
$$
(C_2R(x,t))^{-\frac{3}{2}} \leq \Vol_t(B) \leq \varepsilon
\sigma^3.
$$
Here $\eta$ is a universal positive constant.

\medskip
Without loss of generality, we always assume the above constants
$C_1$ and $C_2$ are twice bigger than the corresponding constants
$C_1(\frac{\varepsilon}{2})$ and $C_2(\frac{\varepsilon}{2})$ in
Theorem 6.4.6 with the accuracy $\frac{\varepsilon}{2}$.

We remark that the above definition of the canonical neighborhood
assumption is slightly different from that of Perelman in
\cite{P2} in two aspects: (1) it allows the parameter $r$ to
depend on time; (2) it also includes an volume upper bound for the
canonical neighborhoods of types (a) and (b).

Arbitrarily given a compact orientable three-manifold with a
Riemannian metric, by scaling, we may assume the Riemannian metric
is normalized. In the rest of this section and the next section,
we will show the Ricci flow with surgery, with the normalized
metric as initial data, has a long-time solution which satisfies
the above a priori assumptions and has only a finite number of
surgery times at each finite time interval. The construction of
the long-time solution will be given by an induction argument.

First, for the arbitrarily given compact orientable normalized
three-dimensional Riemannian manifold $(M,g_{ij}(x))$, the Ricci
flow with it as initial data has a maximal solution $g_{ij}(x,t)$
on a maximal time interval $[0,T)$ with $T>1$. It follows from
Theorem 5.3.2 and Theorem 7.1.1 that the a priori assumptions
(with accuracy $\varepsilon$) hold for the smooth solution on
$[0,T)$. If $T=+\infty$, we have the desired long time solution.
Thus, without loss of generality, we may assume the maximal time
$T<+\infty$ so that the solution goes singular at time $T$.

Suppose that we have a solution to the Ricci flow with surgery,
with the normalized metric as initial data, satisfying the a
priori assumptions (with accuracy $\varepsilon$), defined on
$[0,T)$ with $T<+\infty$, going singular at time $T$ and having
only a finite number of surgery times on $[0,T)$. Let $\Omega$
denote the set of all points where the curvature stays bounded as
$t\rightarrow T$. As we have seen before, the canonical
neighborhood assumption implies that $\Omega$ is open and that
$R(x,t)\rightarrow\infty$ as $t\rightarrow T$ for all $x$ lying
outside $\Omega$. Moreover, as $t\rightarrow T$, the solution
$g_{ij}(x,t)$ has a smooth limit $\bar{g}_{ij}(x)$ on $\Omega$.

For some $\delta>0$ to be chosen much smaller than $\varepsilon$,
we let $\rho=\delta r(T)$ where $r(t)$ is the positive
nonincreasing function in the definition of the canonical
neighborhood assumption. We consider the corresponding compact set
$$
\Omega_\rho=\{x\in\Omega\ |\ \bar{R}(x)\leq\rho^{-2}\},
$$
where $\bar{R}(x)$ is the scalar curvature of $\bar{g}_{ij}$. If
$\Omega_\rho$ is empty, the manifold (near the maximal time $T$)
is entirely covered by $\varepsilon$-tubes, $\varepsilon$-caps and
compact components with positive curvature. Clearly, the number of
the compact components is finite. Then in this case the manifold
(near the maximal time $T$) is diffeomorphic to the union of a
finite number of copies of $\mathbb{S}^3$, or metric quotients of
the round $\mathbb{S}^3$, or $\mathbb{S}^2 \times \mathbb{S}^1$,
or a connected sum of them. Thus when $\Omega_\rho$ is empty, the
procedure stops here, and we say the \textbf{solution becomes
extinct}\index{solution becomes extinct}. We now assume
$\Omega_\rho$ is not empty. Then we know that every point
$x\in\Omega\setminus\Omega_\rho$ lies in one of the subsets of
$\Omega$ listed in (7.3.2), or in a compact component with
positive curvature, or in a compact component which is contained
in $\Omega\setminus\Omega_\rho$ and is diffeomorphic to either
$\mathbb{S}^3$, or $\mathbb{S}^2 \times \mathbb{S}^1$ or
$\mathbb{RP}^3\#\mathbb{RP}^3$. Note again that the number of the
compact components is finite. Let us throw away all the compact
components lying in $\Omega\setminus\Omega_\rho$ and all the
compact components with positive curvature, and then consider
those components $\Omega_j$, $1\leq j\leq k$, of $\Omega$ which
contain points of $\Omega_\rho$. (We will consider those
components of $\Omega\setminus\Omega_\rho$ consisting of capped
$\varepsilon$-horns and double $\varepsilon$-horns later). We will
perform surgical procedures, as we roughly described before, by
finding an $\varepsilon$-neck in every horn of $\Omega_j$, $1\leq
j\leq k$, and then cutting it along the middle two-sphere,
removing the horn-shaped end, and gluing back a cap.

In order to maintain the a priori assumptions \textbf{with the
same accuracy} after the surgery, we will need to find sufficient
``fine" necks in the $\varepsilon$-horns and to glue sufficient
``fine" caps. Note that $\delta>0$ will be chosen much smaller
than $\varepsilon>0$. The following lemma due to Perelman
\cite{P2} gives us the ``fine" necks in the $\varepsilon$-horns.
(At the first sight, we should also cut off all those
$\varepsilon$-tubes and $\varepsilon$-caps in the surgery
procedure. However, in general we are not able to find a ``fine"
neck in an $\varepsilon$-tube or in an $\varepsilon$-cap, and
surgeries at ``rough" $\varepsilon$-necks will certainly lose some
accuracy. If we perform surgeries at the necks with some fixed
accuracy $\varepsilon$ in the high curvature region at each
surgery time, then it is possible that the errors of surgeries may
accumulate to a certain amount so that at some later time we
cannot recognize the structure of very high curvature regions.
This prevents us from carrying out the whole process in finite
time with a finite number of steps. This is the reason why we will
only perform the surgeries at the $\varepsilon$-horns.)

%\vskip 0.2cm\noindent{\bf Lemma 7.3.2}
%CZ-SOURCE-LINE-17943
\begin{lemma}[existence of cutoff necks in horns]
  \label{lem:cz-cutoff-neck-existence-in-horns}
  \dcref{Cao--Zhu Lemma 7.3.2}
  \group{Ricci Flow with Surgery}
  \source{cao-zhu}

Given $0<\varepsilon \leq \frac{1}{100}, 0<\delta<\varepsilon$ and
$0<T<+\infty$, there exists a radius $0<h<\delta\rho$, depending
only on $\delta$ and $r(T)$, such that if we have a solution to
the Ricci flow with surgery, with a normalized metric as initial
data, satisfying the a priori assumptions $($with accuracy
$\varepsilon),$ defined on $[0,T)$, going singular at time $T$ and
having only a finite number of surgery times on $[0,T)$, then for
each point $x$ with $h(x)=\bar{R}^{-\frac{1}{2}}(x)\leq h$ in an
$\varepsilon$-horn of $(\Omega,\bar{g}_{ij})$ with boundary in
$\Omega_\rho$, the neighborhood
$B_T(x,\delta^{-1}h(x))\triangleq\{y\in\Omega\ |\ d_T(y,x)\leq
\delta^{-1}h(x)\}$ is a strong $\delta$-neck $($i.e.,
$B_T(x,\delta^{-1}h(x)) \times [T-h^2(x),T]$ is, after scaling
with factor $h^{-2}(x)$, $\delta$-close $($in the
$C^{[\delta^{-1}]}$ topology$)$ to the corresponding subset of the
evolving standard round cylinder $\mathbb{S}^2 \times \mathbb{R}$
over the time interval $[-1,0]$ with scalar curvature $1$ at time
zero$)$.
\end{lemma}

\vskip 0.5cm
\begin{center}
\setlength{\unitlength}{2mm}
\begin{picture}(60,8)
\linethickness{0.5pt}

\qbezier[10](50,0)(48,2.5)(50,5)

\qbezier(50,0)(52,2.5)(50,5)

\qbezier[5](25,1.7)(24,2.5)(25,3.3)
\qbezier(25,1.7)(26,2.5)(25,3.3)

\qbezier(26,1.9)(26.5,2.5)(27,3.1) \qbezier(27,2)(27.5,2.5)(28,3)
\qbezier(28,1.9)(28.5,2.5)(29,3.1)
\qbezier(29,1.9)(29.5,2.5)(30,3.1)
\qbezier(30,2)(30.5,2.5)(31,3)\qbezier(32,1.9)(32.5,2.5)(33,3.1)
\qbezier(31,2)(31.5,2.5)(32,3)

\qbezier(33,1.4)(34,2.5)(33,3.6)

\qbezier(10,2.5)(48,4)(50,5)

\qbezier(10,2.5)(48,1)(50,0)

\put(26,-2){{strong $\delta$-neck}}
\put(50,-3){\textbf{$\Omega_{\rho}$}}
\end{picture}
\end{center}

\vskip 1.0cm

%\noindent{\bf Proof.}
\begin{proof}
  \uses{thm:cz-three-dimensional-singularity-structure, thm:cz-two-dimensional-ancient-kappa-classification} The following proof is essentially given by Perelman
(cf. 4.3 of \cite{P2}).

We argue by contradiction. Suppose that there exists a sequence of
solutions $g^k_{ij}(\cdot,t)$, $k=1,2,\ldots$, to the Ricci flow
with surgery, satisfying the a priori assumptions (with accuracy
$\varepsilon$), defined on $[0,T)$ with limit metrics
$(\Omega^k,\bar{g}^k_{ij}), k=1, 2, \ldots$, and points $x^k$,
lying inside an $\varepsilon$-horn of $\Omega^k$ with boundary in
$\Omega_\rho^k$, and having $h(x^k)\rightarrow0$ such that the
neighborhoods $B_T(x^k,\delta^{-1}h(x^k))=\{y\in\Omega^k\ |\
d_T(y,x^k)\leq \delta^{-1}h(x^k)\}$ are not strong $\delta$-necks.

Let $\widetilde{g}^k_{ij}(\cdot,t)$ be the solutions obtained by
rescaling by the factor $\bar{R}(x^k)=h^{-2}(x^k)$ around $x^k$
and shifting the time $T$ to the new time zero. We now want to
show that a subsequence of
$\widetilde{g}^k_{ij}(\cdot,t),k=1,2,\ldots$, converges to the
evolving round cylinder, which will give a contradiction.

Note that $\widetilde{g}^k_{ij}(\cdot,t), k=1, 2, \ldots,$ are
solutions modified by surgery. So, we cannot apply Hamilton's
compactness theorem directly since it is stated only for smooth
solutions. For each (unrescaled) surgical solution
$\bar{g}^k_{ij}(\cdot,t)$, we pick a point $z^k$, with
$\bar{R}(z^k) = 2C_2^2(\varepsilon)\rho^{-2},$ in the
$\varepsilon$-horn of $(\Omega^k,\bar{g}^k_{ij})$ with boundary in
$\Omega_\rho^k$, where $C_2(\varepsilon)$ is the positive constant
in the canonical neighborhood assumption. From the definition of
$\varepsilon$-horn and the canonical neighborhood assumption, we
know that each point $x$ lying inside the $\varepsilon$-horn of
$(\Omega^k,\bar{g}^k_{ij})$ with
$d_{\bar{g}^k_{ij}}(x,\Omega_\rho^k) \geq
d_{\bar{g}^k_{ij}}(z^k,\Omega_\rho^k)$ has a strong
$\varepsilon$-neck as its canonical neighborhood. Since
$h(x^k)\rightarrow0$, each $x^k$ lies deeply inside an
$\varepsilon$-horn. Thus for each positive $A < +\infty$, the
rescaled (surgical) solutions $\widetilde{g}^k_{ij}(\cdot,t)$ with
the marked origins $x^k$ over the geodesic balls
$B_{\widetilde{g}^k_{ij}(\cdot,0)}(x^k,A)$, centered at $x^k$ of
radii $A$ (with respect to the metrics
$\widetilde{g}^k_{ij}(\cdot,0)$), will be smooth on some uniform
(size) small time intervals for all sufficiently large $k$
whenever the curvatures of the rescaled solutions
$\widetilde{g}^k_{ij}$ at $t=0$ in
$B_{\widetilde{g}^k_{ij}(\cdot,0)}(x^k,A)$ are uniformly bounded.
In such a situation, Hamilton's compactness theorem is applicable.
Then we can apply the same argument as in the second step of the
proof of Theorem 7.1.1 to conclude that for each $A<+\infty$,
there exists a positive constant $C(A)$ such that the curvatures
of the rescaled solutions $\widetilde{g}^k_{ij}(\cdot,t)$ at the
new time $0$ satisfy the estimate $$|\widetilde{R}m_k|(y,0) \leq
C(A)$$ whenever $d_{\widetilde{g}^k_{ij}(\cdot,0)}(y,x^k) \leq A$
and $k \geq 1$; otherwise we would get a piece of a non-flat
nonnegatively curved metric cone as a blow-up limit, which
contradicts Hamilton's strong maximum principle. Moreover, by
Hamilton's compactness theorem (Theorem 4.1.5), a subsequence of
the rescaled solutions $\widetilde{g}^k_{ij}(\cdot,t)$ converges
to a $C^\infty_{loc}$ limit $\widetilde{g}^\infty_{ij}(\cdot,t)$,
defined on a spacetime set which is relatively open in the half
spacetime $\{t\leq 0\}$ and contains the time slice $t=0$.

By the pinching assumption, the limit is a complete manifold with
nonnegative sectional curvature. Since $x^k$ was contained in an
$\varepsilon$-horn with boundary in $\Omega^k_\rho$ and
$h(x^k)/\rho\rightarrow0$, the limiting manifold has two ends.
Thus, by Toponogov's splitting theorem, the limiting manifold
admits a metric splitting $\Sigma^2\times \mathbb{R}$, where
$\Sigma^2$ is diffeomorphic to the two-sphere $\mathbb{S}^2$
because $x^k$ was the center of a strong $\varepsilon$-neck.

By combining with the canonical neighborhood assumption (with
accuracy $\varepsilon$), we see that the limit is defined on the
time interval $[-1,0]$ and is $\varepsilon$-close to the evolving
standard round cylinder. In particular, the scalar curvature of
the limit at time $t=-1$ is $\varepsilon$-close to $1/2$.

Since $h(x^k)/\rho\rightarrow0$, each point in the limiting
manifold at time $t=-1$ also has a strong $\varepsilon$-neck as
its canonical neighborhood. Thus the limit is defined at least on
the time interval $[-2,0]$ and the limiting manifold at time
$t=-2$ is, after rescaling, $\varepsilon$-close to the standard
round cylinder.

By using the canonical neighborhood assumption again, every point
in the limiting manifold at time $t=-2$ still has a strong
$\varepsilon$-neck as its canonical neighborhood. Also note that
the scalar curvature of the limit at $t=-2$ is not bigger than
$1/2+\varepsilon$. Thus the limit is defined at least on the time
interval $[-3,0]$ and the limiting manifold at time $t=-3$ is,
after rescaling, $\varepsilon$-close to the standard round
cylinder. By repeating this argument we prove that the limit
exists on the ancient time interval $(-\infty, 0]$.

The above argument also shows that at every time, each point of
the limit has a strong $\varepsilon$-neck as its canonical
neighborhood. This implies that the limit is $\kappa$-noncollaped
on all scales for some $\kappa>0$. Therefore, by Theorem 6.2.2,
the limit is the evolving round cylinder
$\mathbb{S}^2\times\mathbb{R}$, which gives the desired
contradiction.
%$$\eqno \#$$ \vskip 0.3cm
\end{proof}

In the above lemma, the property that the radius $h$ depends only
on $\delta$ and the time $T$ but is independent of the surgical
solution is crucial; otherwise we will not be able to cut off
enough volume at each surgery to guarantee the number of surgeries
being finite in each finite time interval. We also remark that the
above proof actually proves a stronger result: the parabolic
region $\{(y,t)\ |\ y\in
B_T(x,\delta^{-1}h(x)),t\in[T-\delta^{-2}h^2(x),T]\}$ is, after
scaling with factor $h^{-2}(x)$, $\delta$-close (in the
$C^{[\delta^{-1}]}$ topology) to the corresponding subset of the
evolving standard round cylinder $\mathbb{S}^2 \times \mathbb{R}$
over the time interval $[-\delta^{-2},0]$ with scalar curvature
$1$ at the time zero. This fact will be used later in the proof of
Proposition 7.4.1.

We next want to construct ``fine" caps. Take a rotationally
symmetric metric on $\mathbb{R}^3$ with nonnegative sectional
curvature and positive scalar curvature such that outside some
compact set it is a semi-infinite standard round cylinder (i.e.
the metric product of a ray with the round two-sphere of scalar
curvature 1). We call such a metric on $\mathbb{R}^3$  a
\textbf{standard capped infinite cylinder}\index{standard capped
infinite cylinder}. By the short-time existence theorem of Shi
(Theorem 1.2.3), the Ricci flow with a standard capped infinite
cylinder as initial data has a complete solution on a maximal time
interval $[0,T)$ such that the curvature of the solution is
bounded on $\mathbb{R}^3 \times [0,T']$ for each $0<T'<T$.  Such a
solution is called a \textbf{standard
solution}\index{solution!standard} by Perelman \cite{P2}.

The following result, proved by Chen and the second author in
\cite{CZ05F}, gives the curvature estimate for standard solutions.
This curvature estimate in the special case, when the dimension is
three and the initial metric is rotationally symmetric, was first
claimed by Perelman in \cite{P2}.

%\vskip 0.2cm \noindent \textbf{Proposition 7.3.3} \emph{
%CZ-SOURCE-LINE-18138
\begin{proposition}[standard solution from an asymptotically cylindrical cap]
  \label{prop:cz-asymptotically-cylindrical-standard-solution}
  \dcref{Cao--Zhu Proposition 7.3.3}
  \group{Ricci Flow with Surgery}
  \source{cao-zhu}

Let $g_{ij}$ be a complete Riemannian metric on $\mathbb{R}^{n}$
($n \geq 3$) with nonnegative curvature operator and positive
scalar curvature which is asymptotic to a round cylinder of scalar
curvature $1$ at infinity.  Then there is a complete solution
$g_{ij}({\cdot,t})$ to the Ricci flow, with $g_{ij}$ as initial
metric, which exists on the time interval $[0,\frac{n-1}{2})$, has
bounded curvature at each time $t \in [0,\frac{n-1}{2})$, and
satisfies the estimate
$$
  R(x,t)\geq \frac{C^{-1}}{\frac{n-1}{2}-t}
$$
for some $C$ depending only on the initial metric $g_{ij}$.
\end{proposition}

%\vskip 0.2cm \noindent{\textbf{Proof}}.
\begin{proof}
  \uses{lem:cz-standard-solution-spatial-curvature-comparison, lem:cz-standard-solution-maximal-time-curvature-blowup}
  \uses{thm:cz-ancient-kappa-elliptic-estimates, thm:cz-local-volume-growth-curvature-estimate}
Since the initial metric has bounded curvature operator and a
positive lower bound on its scalar curvature, the Ricci flow has a
solution $g_{ij}(\cdot,t)$ defined on a maximal time interval
$[0,T)$ with $T<\infty$ which has bounded curvature on
$\mathbb{R}^n \times [0,T']$ for each $0<T'<T$. By Proposition
2.1.4, the solution $g_{ij}(\cdot,t)$ has nonnegative curvature
operator for all $t \in [0, T)$.  Note that the injectivity radius
of the initial metric has a positive lower bound. As we remarked
at the beginning of Section 3.4, the same proof of Perelman's no
local collapsing theorem I concludes that $g_{ij}({\cdot,t})$ is
$\kappa$-noncollapsed on all scales less than $\sqrt{T}$ for some
$\kappa>0$ depending only on the initial metric.

We will first prove the following assertion.

\medskip
{\bf Claim 1.} \ There is a positive function
$\omega:[0,\infty)\rightarrow[0,\infty)$ depending only on the
initial metric and $\kappa$ such that
$$
  R(x,t)\leq R(y,t)\omega(R(y,t)d_{t}^{2}(x,y))
$$
for all $x,y\in \mathbb{R}^{n}$, $t\in[0,T)$.

\medskip
The proof is similar to that of Theorem 6.4.3. Notice that the
initial metric has nonnegative curvature operator and its scalar
curvature satisfies the bounds \begin{equation}
C^{-1}\leqslant R(x)\leqslant C %\eqno (7.3.5)
\end{equation} for some positive constant $C$. By the maximum principle, we
know $T\geq\frac{1}{2nC}$ and $R(x,t)\leq 2C$ for
$t\in[0,\frac{1}{4nC}]$. The assertion is clearly true for
$t\in[0,\frac{1}{4nC}]$.

Now fix $(y,t_{0})\in \mathbb{R}^{n}\times[0,T)$ with
$t_{0}\geq\frac{1}{4nC}$. Let $z$ be the closest point to $y$ with
the property $R(z,t_{0})d^{2}_{t_{0}}(z,y)=1$ (at time $t_{0}$).
Draw a shortest geodesic from $y$ to $z$ and choose a point
$\tilde{z}$ on the geodesic satisfying $d_{t_{0}}(z,\tilde{z})
=\frac{1}{4}R(z,t_{0})^{-\frac{1}{2}}$, then we have
$$
 R(x,t_{0})\leq\frac{1}{(\frac{1}{2}R(z,t_{0})^{-\frac{1}{2}})^{2}}
    \qquad \mbox{on}\ \
 B_{t_{0}}\left(\tilde{z},\frac{1}{4}R(z,t_{0}\right)^{-\frac{1}{2}}).
$$

Note that $R(x,t)\geqslant C^{-1}$ everywhere by the evolution
equation of the scalar curvature. Then by the Li-Yau-Hamilton
inequality (Corollary 2.5.5), for all $(x,t)\in
B_{t_{0}}(\tilde{z},\frac{1}{8nC}R(z,t_{0})^{-\frac{1}{2}})
\times[t_{0}-(\frac{1}{8nC}R(z,t_{0})^{-\frac{1}{2}})^{2},t_{0}]
$, we have
\begin{equation*}
\begin{split}
R(x,t)&  \leq \left(\frac{t_{0}}{t_{0}-\left(\frac{1}{8n\sqrt{C}}\right)^{2}}\right)
\frac{1}{\left(\frac{1}{2}R(z,t_{0})^{-\frac{1}{2}}\right)^{2}}\\
&  \leq \left[\frac{1}{8nC}R(z,t_{0})^{-\frac{1}{2}}\right]^{-2}.
\end{split}
\end{equation*}
Combining this with the $\kappa$-noncollapsing property, we have
$$
\Vol
\left(B_{t_{0}}\left(\tilde{z},\frac{1}{8nC}R(z,t_{0})^{-\frac{1}{2}}\right)\right)
\geq \kappa \left(\frac{1}{8nC}R(z,t_{0})^{-\frac{1}{2}}\right)^{n}
$$
and then
$$
\Vol \left(B_{t_{0}}\left(z,8R(z,t_{0})^{-\frac{1}{2}}\right)\right) \geq \kappa
\left(\frac{1}{64nC}\right)^{n}\left(8R(z,t_{0})^{-\frac{1}{2}}\right)^{n}.
$$
So by Theorem 6.3.3 (ii), we have
$$
R(x,t_{0})\leq C(\kappa)R(z,t_{0})\quad \text{for all }\; x\in
B_{t_{0}}\left(z,4R(z,t_{0})^{-\frac{1}{2}}\right).
$$
Here and in the following we denote by $C(\kappa)$ various
positive constants depending only on $\kappa,n$ and the initial
metric.

Now by the Li-Yau-Hamilton inequality (Corollary 2.5.5) and local
gradient estimate of Shi (Theorem 1.4.2), we obtain
\begin{align*}
R(x,t)\leq C(\kappa)R(z,t_{0})\; \text{ and }\;
\left|\frac{\partial}{\partial t}R\right|(x,t) \leq
C(\kappa)(R(z,t_{0}))^{2}
\end{align*}
for all $ (x,t)\in
B_{t_{0}}(z,2R(z,t_{0})^{-\frac{1}{2}}))\times[t_{0}-
(\frac{1}{8nC}R(z,t_{0})^{-\frac{1}{2}})^{2},t_{0}]$. Therefore by
combining with the Harnack estimate (Corollary 2.5.7), we obtain
\begin{align*}
R(y,t_{0})&  \geq C(\kappa)^{-1}R(z,t_{0}-C(\kappa)^{-1}R(z,t_{0})^{-1})\\
&  \geq C(\kappa)^{-2} R(z,t_{0})
\end{align*}

Consequently, we have showed that there is a constant $C(\kappa)$
such that
$$
\Vol \left(B_{t_{0}}\left(y,R(y,t_{0})^{-\frac{1}{2}}\right)\right) \geq
C(\kappa)^{-1}\left(R(y,t_{0})^{-\frac{1}{2}}\right)^{n}
$$
and
$$
R(x,t_{0})\leq C(\kappa)R(y,t_{0}) \quad \text{for all }\; x\in
B_{t_{0}}\left(y,R(y,t_{0})^{-\frac{1}{2}}\right).
$$
In general, for any $r\geq R(y,t_{0})^{-\frac{1}{2}}$, we have
$$
\Vol (B_{t_{0}}(y,r)) \geq
C(\kappa)^{-1}(r^{2}R(y,t_{0}))^{-\frac{n}{2}}r^{n}.
$$
By\; applying\; Theorem\; 6.3.3(ii)\; again,\; there\; exists\;
a\; positive\; constant $\omega(r^{2}R(y,t_{0}))$ depending only
on the constant $r^{2}R(y,t_{0})$ and $\kappa$ such that
$$
R(x,t_{0})\leq R(y,t_{0})\omega(r^{2}R(y,t_{0}))\quad \mbox{for
all }\;  x\in B_{t_{0}}\left(y,\frac{1}{4}r\right).
$$
This proves the desired Claim 1.

Now we study the asymptotic behavior of the solution at infinity.
For any $0<t_{0}<T$, we know that the metrics $g_{ij}(x,t)$ with
$t\in [0,t_{0}]$ has uniformly bounded curvature.  Let $x_{k}$ be
a sequence of points with $d_{0}(x_{0},x_{k})\rightarrow \infty$.
After passing to a subsequence, $g_{ij}(x,t)$ around $x_{k}$ will
converge to a solution to the Ricci flow on $\mathbb{R}\times
\mathbb{S}^{n-1}$ with round cylinder metric of scalar curvature 1
as initial data.  Denote the limit by $\tilde{g}_{ij}$. Then by
the uniqueness theorem (Theorem 1.2.4), we have
$$
\tilde{R}(x,t) =\frac{\frac{n-1}{2}}{\frac{n-1}{2}-t}\quad
\text{for all }\; t\in[0,t_{0}].
$$
It follows that $T\leq \frac{n-1}{2}$. In order to show
$T=\frac{n-1}{2}$, it suffices to prove the following assertion.

\medskip
{\bf Claim 2.} \ Suppose $T<\frac{n-1}{2}$. Fix a point $x_{0}\in
\mathbb{R}^{n} $, then there is a $\delta>0$, such that for any
$x\in M$ with $d_{0}(x,x_{0})\geq \delta^{-1}$, we have
$$
R(x,t)\leq 2C+\frac{n-1}{\frac{n-1}{2}-t}\quad \text{ for all }\;
t\in[0,T),
$$
where $C$ is the constant in (7.3.5).

\medskip
In view of Claim 1, if Claim 2 holds, then
\begin{align*}
\sup_{M^{n}\times[0,T)}R(y,t) &  \leq
\omega\left(\delta^{-2}\left(2C+\frac{n-1}{\frac{n-1}{2}-T}\right)\right)
\left(2C+\frac{n-1}{\frac{n-1}{2}-T}\right)\\
&  < \infty
\end{align*}
which will contradict the definition of $T$.

To show Claim 2, we argue by contradiction. Suppose for each
$\delta>0$, there is a point $(x_{\delta},t_{\delta})$ with
$0<t_{\delta}<T $ such that
$$
R(x_{\delta},t_{\delta})>2C+\frac{n-1}{\frac{n-1}{2}-t_{\delta}}\;
\mbox{ and }\; d_{0}(x_{\delta},x_{0})\geq \delta^{-1}.
$$
Let
$$
\bar{t}_{\delta}=\sup\left\{t\; \Big|\; \sup_{M^{n} \setminus
B_{0}(x_{0},\delta^{-1})}
R(y,t)<2C+\frac{n-1}{\frac{n-1}{2}-t}\right\}.
$$
Since $\lim\limits_{d_{0}(y,x_{0})\rightarrow\infty}R(y,t)
=\frac{\frac{n-1}{2}}{\frac{n-1}{2}-t}$ and
$\sup_{M\times[0,\frac{1}{4nC}]}R(y,t)\leq 2C$, we know
$\frac{1}{4nC}\leq \bar{t}_{\delta}\leq t_{\delta}$ and there is a
$\bar{x}_{\delta}$ such that
$d_{0}(x_{0},\bar{x}_{\delta})\geq\delta^{-1}$ and
$R(\bar{x}_{\delta},\bar{t}_{\delta})
=2C+\frac{n-1}{\frac{n-1}{2}-\bar{t}_{\delta}}.$ By Claim 1 and
Hamilton's compactness theorem (Theorem 4.1.5), for $\delta\
\rightarrow 0$ and after taking a subsequence, the metrics
$g_{ij}(x,t)$ on $B_{0}(\bar{x}_{\delta},\frac{\delta^{-1}}{2})$
over the time interval $[0,\bar{t}_{\delta}]$ will converge to a
solution $\tilde{g}_{ij}$ on $\mathbb{R}\times \mathbb{S}^{n-1}$
with the standard metric of scalar curvature 1 as initial data
over the time interval $[0,\bar{t}_{\infty}]$, and its scalar
curvature satisfies
\begin{align*}
\tilde{R}(\bar{x}_{\infty},\bar{t}_{\infty})
&  = 2C+\frac{n-1}{\frac{n-1}{2}-\bar{t}_{\infty}},\\
\tilde{R}(x,t)&  \leqslant
2C+\frac{n-1}{\frac{n-1}{2}-\bar{t}_{\infty}},
 \quad \mbox{for all }\;  t\in [0,\bar{t}_{\infty}],
\end{align*}
where $(\bar{x}_{\infty},\bar{t}_{\infty})$ is the limit of
$(\bar{x}_{\delta},\bar{t}_{\delta})$. On the other hand, by the
uniqueness theorem (Theorem 1.2.4) again, we know
$$
\tilde{R}(\bar{x}_{\infty},\bar{t}_{\infty})
=\frac{\frac{n-1}{2}}{\frac{n-1}{2}-\bar{t}_{\infty}}
$$
which is a contradiction.  Hence we have proved Claim 2 and then
have verified $T=\frac{n-1}{2}$.

Now we are ready to show \begin{equation} R(x,t)\geq
\frac{\tilde{C}^{-1}}{\frac{n-1}{2}-t}, \quad \text{for all }\;
(x,t)\in \mathbb{R}^{n}\times\Big[0,\frac{n-1}{2}\Big),
%\eqno (7.3.6)
\end{equation} for some positive constant $\tilde{C}$ depending only on the
initial metric.

For any $(x,t)\in \mathbb{R}^{n}\times[0,\frac{n-1}{2})$, by Claim
1 and $\kappa$-noncollapsing, there is a constant $C(\kappa)>0$
such that
$$
\Vol_{t}(B_{t}(x,{R(x,t)}^{-\frac{1}{2}}))\geq
C(\kappa)^{-1}({R(x,t)}^{-\frac{1}{2}})^{n}.
$$
Then by the well-known volume estimate of Calabi-Yau (see for
example \cite{Y76} or \cite{ScY}) for complete manifolds with
$\Ric\geq 0$, for any $a\geq 1$, we have
$$
\Vol_{t}(B_{t}(x,a{R(x,t)}^{-\frac{1}{2}}))\geq
C(\kappa)^{-1}\frac{a}{8n}({R(x,t)}^{-\frac{1}{2}})^{n}.
$$
On the other hand, since $(\mathbb{R}^{n},g_{ij}(\cdot,t))$ is
asymptotic to a cylinder of scalar curvature
${(\frac{n-1}{2})}/{(\frac{n-1}{2}-t)}$, for sufficiently large
$a>0$, we have
$$
\Vol_{t}\left(B_{t}\left(x,a\sqrt{\frac{n-1}{2}-t}\right)\right)\leq
C(n)a\left(\frac{n-1}{2}-t\right)^{\frac{n}{2}}.
$$
Combining the two inequalities, for  all sufficiently large $a$,
we have:
\begin{align*}
C(n)a\left(\frac{n-1}{2}-t\right)^{\frac{n}{2}}&  \geq
\Vol_{t}\left(B_{t}\left(x,a\left(\frac{\sqrt{\frac{n-1}{2}
-t}}{R(x,t)^{-\frac{1}{2}}}\right){R(x,t)}^{-\frac{1}{2}}\right)\right)\\
&  \geq C(\kappa)^{-1}\frac{a}{8n}\left(\frac{\sqrt{\frac{n-1}{2}
-t}}{{R(x,t)}^{-\frac{1}{2}}}\right) \left({R(x,t)}^{-\frac{1}{2}}\right)^{n},
\end{align*}
which gives the desired estimate (7.3.6). Therefore the proof of
the proposition is complete.
%$$\eqno \#$$ \vskip 0.3cm
\end{proof}

We now fix a standard capped infinite cylinder for dimension $n=3$
as follows.  Consider the semi-infinite standard round cylinder
$N_0 = \mathbb{S}^2 \times (-\infty,4)$ with the metric $g_0$ of
scalar curvature 1. Denote by $z$ the coordinate of the second
factor $(-\infty,4)$. Let $f$ be a smooth nondecreasing convex
function on $(-\infty,4)$ defined by \begin{equation}
   \left\{
   \begin{array}{lll}
    f(z) = 0, \ \ \ z\leq 0,
         \\[3mm]
    f(z) = ce^{-\frac{P}{z}}, \ \ \ z \in (0,3], \\[3mm]
    f(z) \mbox{ is strictly convex on } z \in [3,3.9], \\[3mm]
    f(z) = -\frac{1}{2}\log(16-z^2), \ \ \ z \in [3.9,4),
\end{array}
\right.         %%%  NEW (7.3.7)
\end{equation} where the (small) constant $c>0$ and (big) constant $P>0$ will
be determined later. Let us replace the standard metric $g_0$ on
the portion $\mathbb{S}^2 \times [0,4)$ of the semi-infinite
cylinder by $\hat{g} = e^{-2f}g_0$. Then the resulting metric
$\hat{g}$ will be smooth on $\mathbb{R}^3$ obtained by adding a
point to $\mathbb{S}^2 \times (-\infty,4)$ at $z=4$. We denote by
$C(c,P) = (\mathbb{R}^3,\hat{g})$. Clearly, $C(c,P)$ is a standard
capped infinite cylinder.

We next use a compact portion of the standard capped infinite
cylinder $C(c,P)$ and the $\delta$-neck obtained in Lemma 7.3.2 to
perform the following surgery procedure due to Hamilton
\cite{Ha97}.

Consider the metric $\bar{g}$ at the maximal time $T<+\infty$.
Take an $\varepsilon$-horn with boundary in $\Omega_\rho$. By
Lemma 7.3.2, there exists a $\delta$-neck $N$ of radius
$0<h<\delta \rho$ in the $\varepsilon$-horn. By definition,
$(N,h^{-2}\bar{g})$ is $\delta$-close (in the $C^{[\delta^{-1}]}$
topology) to the standard round neck $\mathbb{S}^2\times
\mathbb{I}$ of scalar curvature 1 with
$\mathbb{I}=(-\delta^{-1},\delta^{-1})$. Using the parameter $z
\in \mathbb{I}$, we see the above function $f$ is defined on the
$\delta$-neck $N$.

Let\; us\; cut\; the\; $\delta$-neck\; $N$\; along\; the\;
middle\; (topological)\; two-sphere $N\bigcap\{z=0\}$. Without
loss of generality, we may assume that the right hand half portion
$N\bigcap\{z\geq 0\}$ is contained in the horn-shaped end. Let
$\varphi$ be a smooth bump function with $\varphi = 1$ for
$z\leq2$, and $\varphi = 0$ for $z\geq3$. Construct a new metric
$\tilde{g}$ on a (topological) three-ball $\mathbb{B}^3$ as
follows \begin{equation}
   \tilde{g} =  \begin{cases}
    \bar{g}, &\quad z= 0, \\[1mm]
    e^{-2f}\bar{g},&\quad  z \in [0,2], \\[1mm]
    \varphi e^{-2f}\bar{g} + (1-\varphi)e^{-2f}h^2g_0,
&\quad z \in [2,3], \\[1mm]
    h^2e^{-2f}g_0, &\quad z\in [3,4].
\end{cases}  % \eqno (7.3.8)
\end{equation} The surgery is to replace the horn-shaped end by the cap
$(\mathbb{B}^3,\tilde{g})$. We call such surgery procedure a
\textbf{$\delta$-cutoff surgery}\index{$\delta$-cutoff surgery}.

The following lemma determines the constants $c$ and $P$ in the
$\delta$-cutoff surgery so that the pinching assumption is
preserved under the surgery (cf 4.4 of Perelman \cite{P2}).

%\vskip 0.2cm \noindent \textbf{Lemma 7.3.4} (\textbf
%CZ-SOURCE-LINE-18464
\begin{lemma}[preservation of Hamilton--Ivey pinching through surgery]
  \label{lem:cz-surgery-preserves-hamilton-ivey-pinching}
  \dcref{Cao--Zhu Lemma 7.3.4}
  \group{Ricci Flow with Surgery}
  \source{cao-zhu}
 There
are universal positive constants $\delta_0$, $c_0$ and $P_0$ such
that if we take a $\delta$-cutoff surgery at a $\delta$-neck of
radius $h$ at time $T$ with $\delta \leq \delta_0$ and $h^{-2}
\geq 2e^2\log(1+T)$, then we can choose $c=c_0$ and $P=P_0$ in the
definition of $f(z)$ such that after the surgery, there still
holds the pinching condition \begin{equation} \tilde{R} \geq
(-\tilde{\nu})[\log(-\tilde{\nu}) + \log(1+T) -3]
%\eqno (7.3.9)
\end{equation} whenever $\tilde{\nu} < 0$, where $\tilde{R}$ is the scalar
curvature of the metric $\tilde{g}$ and $\tilde{\nu}$ is the least
eigenvalue of the curvature operator of $\tilde{g}$. Moreover,
after the surgery, any metric ball of radius
$\delta^{-\frac{1}{2}}h$ with center near the tip $($i.e., the
origin of the attached cap$)$ is, after scaling with factor
$h^{-2}$, $\delta^{\frac{1}{2}}$-close $($in the
$C^{[\delta^{-\frac{1}{2}}]}$ topology$)$ to the corresponding
ball of the standard capped infinite cylinder $C(c_0,P_0)$.
\end{lemma}

%\vskip 0.1cm \noindent \textbf{Proof}.
\begin{proof}
First, we consider the metric $\tilde{g}$ on the portion $\{0 \leq
z \leq 2\}$. Under the conformal change $\tilde{g}=
e^{-2f}\bar{g}$, the curvature tensor $\tilde{R}_{ijkl}$ is given
by
\begin{align*}
\tilde{R}_{ijkl} &= e^{-2f}\Big[\bar{R}_{ijkl} + |\bar{\nabla}
f|^2(\bar{g}_{il}\bar{g}_{jk}-\bar{g}_{ik}\bar{g}_{jl}) + (f_{ik}
+ f_if_k)\bar{g}_{jl} \\
&\qquad + (f_{jl} + f_jf_l)\bar{g}_{ik} - (f_{il} +
f_if_l)\bar{g}_{jk} - (f_{jk} + f_jf_k)\bar{g}_{il}\Big].
\end{align*}
If $\{\bar{F}_a = \bar{F}^i_a\frac{\partial}{\partial x^i}\}$ is
an orthonormal frame for $\bar{g}_{ij}$, then $\{\tilde{F}_a =
e^f\bar{F}_a = \tilde{F}^i_a\frac{\partial}{\partial x^i}\}$ is an
orthonormal frame for $\tilde{g}_{ij}$. Let
\begin{align*}
\bar{R}_{abcd}
&= \bar{R}_{ijkl}\bar{F}^i_a\bar{F}^j_b\bar{F}^k_c\bar{F}^l_d,\\
\tilde{R}_{abcd}
&=\tilde{R}_{ijkl}\tilde{F}^i_a\tilde{F}^j_b\tilde{F}^k_c\tilde{F}^l_d,
\end{align*}
then
\begin{align}
\tilde{R}_{abcd} &= e^{2f}\Big[\bar{R}_{abcd} + |\bar{\nabla}
f|^2(\delta_{ad}\delta_{bc}-\delta_{ac}\delta_{bd}) + (f_{ac} +
f_af_c)\delta_{bd} \\
&\qquad + (f_{bd} + f_bf_d)\delta_{ac} - (f_{ad} +
f_af_d)\delta_{bc} - (f_{bc} + f_bf_c)\delta_{ad}\Big], \nonumber
%\eqno(7.3.10)
\end{align}
and \begin{equation} \tilde{R} = e^{2f}(\bar{R} + 4\bar{\triangle} f
-2|\bar{\nabla}f|^2). %\eqno (7.3.11)
\end{equation} Since
$$
\frac{df}{dz}= ce^{-\frac{P}{z}}\frac{P}{z^2}, \ \
\frac{d^2f}{dz^2} =
ce^{-\frac{P}{z}}\left(\frac{P^2}{z^4}-\frac{2P}{z^3}\right),
$$
then for any small $\theta >0$, we may choose $c>0$ small and
$P>0$ large such that for $z \in [0,3]$, we have \begin{equation} |e^{2f} - 1|
+ \left|\frac{df}{dz}\right| + \left|\left(\frac{df}{dz}\right)^2\right| <
\theta \frac{d^2f}{dz^2}, \ \ \frac{d^2f}{dz^2} < \theta.
%\eqno(7.3.12)
\end{equation}

On the other hand, by the definition of $\delta$-neck of radius
$h$, we have
$$
|\bar{g} - h^2g_0|_{g_0} < \delta h^2,
$$
$$
|\overset{o}{\nabla^j}\bar{g}|_{g_0} < \delta h^2, \; \mbox{ for
}\; 1 \leq j\leq [\delta^{-1}],
$$
where $g_0$ is the standard metric of the round cylinder
$\mathbb{S}^2 \times \mathbb{R}$. Note that in three dimensions,
we can choose the orthonormal frame
$\{\bar{F}_1,\bar{F}_2,\bar{F}_3 \}$ for the metric $\bar{g}$ so
that its curvature operator is diagonal in the orthonormal frame
$\{ \sqrt{2}\bar{F}_2\wedge \bar{F}_3, \sqrt{2}\bar{F}_3\wedge
\bar{F}_1, \sqrt{2}\bar{F}_1\wedge \bar{F}_2 \}$ with eigenvalues
$\bar{\nu} \leq \bar{\mu} \leq \bar{\lambda}$ and
$$
\bar{\nu}=2\bar{R}_{2323}, \ \ \bar{\mu}=2\bar{R}_{3131}, \ \
\bar{\lambda}=2\bar{R}_{1212}.
$$
Since $h^{-2}\bar{g}$ is $\delta$-close to the standard round
cylinder metric $g_0$ on the $\delta$-neck, we have \begin{equation} \left\{
\begin{array}{lll}
|\bar{R}_{3131}| + |\bar{R}_{2323}| < \delta^{\frac{7}{8}}h^{-2}, \\[1mm]
|\bar{R}_{1212} - \frac{1}{2}h^{-2}| < \delta^{\frac{7}{8}}h^{-2}, \\[1mm]
|\bar{F}_3 - h^{-1}\frac{\partial}{\partial z}|_{g_0} <
\delta^{\frac{7}{8}}h^{-1},\\[1mm]
\end{array}
\right. %\eqno (7.3.13)
\end{equation} for suitably small $\delta > 0$. Since $\bar{\nabla}_az =
\bar{\nabla} z(\bar{F}_a)$ and $\bar{\nabla}_a\bar{\nabla}_b z =
\bar{\nabla}^2 z(\bar{F}_a,\bar{F}_b)$, it follows that
\begin{align*}
|\bar{\nabla}_3z - h^{-1}| &< \delta^{\frac{7}{8}}h^{-1},\\
|\bar{\nabla}_1z| + |\bar{\nabla}_2z|
&<\delta^{\frac{7}{8}}h^{-1},
\end{align*}
and
$$
|\bar{\nabla}_a\bar{\nabla}_b z| < \delta^{\frac{7}{8}}h^{-2},\;
\mbox{ for }\; 1 \leq a, b \leq 3.
$$
By combining with
$$
\bar{\nabla}_af = \frac{df}{dz}\bar{\nabla}_az, \ \
\bar{\nabla}_a\bar{\nabla}_bf =
\frac{df}{dz}\bar{\nabla}_a\bar{\nabla}_bz +
\frac{d^2f}{dz^2}\bar{\nabla}_az\bar{\nabla}_bz
$$
and (7.3.12), we get \begin{equation}
\begin{cases}
|\bar{\nabla}_af| < 2\theta h^{-1}\frac{d^2f}{dz^2}, &
\mbox{for }\; 1 \leq a \leq 3, \\[1mm]
   |\bar{\nabla}_a\bar{\nabla}_b f| <
\delta^{\frac{3}{4}}h^{-2}\frac{d^2f}{dz^2},
& \mbox{unless }\;  a=b=3, \\[1mm]
|\bar{\nabla}_3\bar{\nabla}_3f - h^{-2}\frac{d^2f}{dz^2}| <
\delta^{\frac{3}{4}}h^{-2}\frac{d^2f}{dz^2}.
\end{cases} %\eqno (7.3.14)
\end{equation} By combining (7.3.10) and (7.3.14), we have \begin{equation} \left\{
\begin{array}{lll}
\tilde{R}_{1212} \geq \bar{R}_{1212} - (\theta^{\frac{1}{2}}
+ \delta^{\frac{5}{8}})h^{-2}\frac{d^2f}{dz^2}, \\[1mm]
\tilde{R}_{3131} \geq \bar{R}_{3131} + (1 - \theta^{\frac{1}{2}}
- \delta^{\frac{5}{8}})h^{-2}\frac{d^2f}{dz^2},  \\[1mm]
\tilde{R}_{2323} \geq \bar{R}_{2323} + (1 - \theta^{\frac{1}{2}}
- \delta^{\frac{5}{8}})h^{-2}\frac{d^2f}{dz^2}, \\[1mm]
|\tilde{R}_{abcd}| \leq (\theta^{\frac{1}{2}} +
\delta^{\frac{5}{8}})h^{-2}\frac{d^2f}{dz^2}, \qquad
\mbox{otherwise},
\end{array}
\right. %\eqno (7.3.15)
\end{equation} where $\theta$ and $\delta$ are suitably small. Then it
follows that
$$
\tilde{R} \geq \bar{R} + [4-6(\theta^{\frac{1}{3}} +
\delta^{\frac{1}{2}})]h^{-2}\frac{d^2f}{dz^2},
$$
$$
-\tilde{\nu} \leq -\bar{\nu} - [2-2(\theta^{\frac{1}{3}} +
\delta^{\frac{1}{2}})]h^{-2}\frac{d^2f}{dz^2},
$$
for suitably small $\theta$ and $\delta$.

If $0<-\tilde{\nu} \leq e^2$, then by the assumption that $h^{-2}
\geq 2e^2\log(1+T)$, we have
\begin{align*}
\tilde{R} &\geq \bar{R}\\
&  \geq \frac{1}{2}h^{-2} \\
&  \geq e^2\log(1+T)\\[2mm]
&  \geq (-\tilde{\nu})[\log(-\tilde{\nu}) + \log(1+T) -3].
\end{align*}
While if $-\tilde{\nu} > e^2$, then by the pinching estimate of
$\bar{g}$, we have
\begin{align*}
\tilde{R}&  \geq \bar{R}\\
&  \geq (-\bar{\nu})[\log(-\bar{\nu}) + \log(1+T) -3] \\
&  \geq (-\tilde{\nu})[\log(-\tilde{\nu}) + \log(1+T) -3].
\end{align*}
So we have verified the pinching condition on the portion $\{0
\leq z \leq 2\}$.

Next, we consider the metric $\tilde{g}$ on the portion $\{2 \leq
z \leq 4\}$. Let $\theta$ be a fixed suitably small positive
number. Then the constant $c=c_0$ and $P=P_0$ are fixed. So $\zeta
= \min_{z \in [1,4]} \frac{d^2f}{dz^2}
> 0$ is also fixed. By the same argument as in the derivation of (7.3.15)
from (7.3.10), we see that the curvature of the metric $\hat{g} =
e^{-2f}g_0$ of the standard capped infinite cylinder $C(c_0,P_0)$
on the portion $\{ 1\leq z \leq 4\}$ is bounded from below by
$\frac{2}{3}\zeta >0$.  Since $h^{-2}\bar{g}$ is $\delta$-close to
the standard round metric $g_0$, the metric $h^{-2}\tilde{g}$
defined by (7.3.8) is clearly $\delta^{\frac{3}{4}}$-close to the
metric $\hat{g} = e^{-2f}g_0$ of the standard capped infinite
cylinder on the portion $\{ 1\leq z \leq 4 \}$. Thus as $\delta$
is sufficiently small, the curvature operator of $\tilde{g}$ on
the portion $\{2 \leq z \leq 4\}$ is positive. Hence the pinching
condition (7.3.9) holds trivially on the portion $\{2 \leq z \leq
4\}$.

The last statement in Lemma 7.3.4 is obvious from the definition
(7.3.8).
%$$\eqno \#$$\vskip 0.3cm
\mbox{ \ \ }\end{proof}

Recall from Lemma 7.3.2 that the $\delta$-necks at a time $t>0$,
where we performed Hamilton's surgeries, have their radii $0 < h <
\delta\rho =\delta^2r(t)$. Without loss of generality, we may
assume the positive nonincreasing function $r(t)$ in the
definition of the canonical neighborhood assumption is less than
$1$ and the universal constant $\delta_0$ in Lemma 7.3.4 is also
less than $1$. We define a positive function $\bar{\delta}(t)$ by
\begin{equation} \bar{\delta}(t) = \min\left\{\frac{1}{{2e^2\log(1+t)}},
\delta_0\right\} \quad \mbox{for } \; t\in [0,+\infty). %\eqno (7.3.16)
\end{equation}

{}From now on, we always assume $0 < \delta < \bar{\delta}(t)$ for
any $\delta$-cutoff surgery at time $t>0$ and assume $c=c_0$ and
$P=P_0$. As a result, the standard capped infinite cylinder and
the standard solution are also fixed. The following lemma, which
was claimed by Perelman in \cite{P2}, gives the canonical
neighborhood structure for the fixed standard solution.

%\vskip 0.2cm \noindent \textbf{Lemma 7.3.5} \emph{
%CZ-SOURCE-LINE-18677
\begin{lemma}[canonical neighborhoods in the standard solution]
  \label{lem:cz-standard-solution-canonical-neighborhoods}
  \dcref{Cao--Zhu Lemma 7.3.5}
  \group{Ricci Flow with Surgery}
  \source{cao-zhu}

Let $g_{ij}(x,t)$ be the above fixed standard solution to the
Ricci flow on $\mathbb{R}^3\times [0,1)$.  Then for any
$\varepsilon>0$, there is a positive constant $C(\varepsilon)$
such that each point $(x,t) \in \mathbb{R}^3\times [0,1)$ has an
open neighborhood $B$, with $B_t(x,r)\subset B \subset B_t(x,2r)$
for some $0<r<C(\varepsilon)R(x,t)^{-\frac{1}{2}}$, which falls
into one of the following two categories: either
\begin{itemize}
\item[(a)] $B$ is  an $\varepsilon$-cap, or \item[(b)] $B$ is an
$\varepsilon$-neck and it is the slice at the time $t$ of the
parabolic neighborhood
$B_t(x,\varepsilon^{-1}R(x,t)^{-\frac{1}{2}}) \times [t-\min
\{R(x,t)^{-1},t\},t]$, on which the standard solution is, after
scaling with the factor $R(x,t)$ and shifting the time $t$ to
zero, $\varepsilon$-close $($in the $C^{[\varepsilon^{-1}]}$
topology$)$ to the corresponding subset of the evolving standard
cylinder $\mathbb{S}^2 \times \mathbb{R}$ over the time interval
$[-\min \{tR(x,t),1\},0]$ with scalar curvature $1$ at the time
zero.
\end{itemize}
\end{lemma}

%\vskip 0.2cm \noindent{\textbf{Proof}}.
\begin{proof}
The proof of the lemma is reduced to two assertions. We now state
and prove the first assertion which takes care of those points
with times close to $1$.

% \vskip 0.1cm\noindent\textbf{Assertion 1}  \emph{
%CZ-SOURCE-LINE-18707
\begin{proposition}[late-time standard solution modeled by ancient solutions]
  \label{prop:cz-standard-solution-late-time-ancient-model}
  \dcref{Cao--Zhu Assertion 1}
  \group{Ricci Flow with Surgery}
  \source{cao-zhu}

For any $\varepsilon>0$, there is a positive number
$\theta=\theta(\varepsilon)$ with $0<\theta<1$ such that for any
$(x_0,t_0)\in \mathbb{R}^3\times [\theta,1)$, the standard
solution on the parabolic neighborhood
$B_{t_0}(x,\varepsilon^{-1}R(x_0,t_0)^{-\frac{1}{2}}) \times
[t_0-\varepsilon^{-2}R(x_0,t_0)^{-1},t_0]$ is well-defined and is,
after scaling with the factor $R(x_0,t_0)$, $\varepsilon$-close
(in the $C^{[\varepsilon^{-1}]}$ topology) to the corresponding
subset of some orientable ancient $\kappa$-solution.
\end{proposition}

We argue by contradiction. Suppose Assertion 1 is not true, then
there exist $\bar\varepsilon>0$ and a sequence of points
$(x_{k},t_k)$ with $t_{k}\rightarrow 1$, such that for each $k$,
the standard solution on the parabolic neighborhood
$$
B_{t_k}(x_k,\bar\varepsilon^{-1}R(x_k,t_k)^{-\frac{1}{2}}) \times
[t_k-\bar\varepsilon^{-2}R(x_k,t_k)^{-1},t_k]
$$
is not, after scaling by the factor $R(x_k,t_k)$,
$\bar\varepsilon$-close to the corresponding subset of any ancient
$\kappa$-solution. Note that by Proposition 7.3.3, there is a
constant $C>0$ (depending only on the initial metric, hence it is
universal ) such that $R(x,t)\geq C^{-1}/(1-t)$. This implies
$$
\bar\varepsilon^{-2}R(x_k,t_k)^{-1} \leq
C\bar\varepsilon^{-2}(1-t_k)<t_k ,
$$
and\; then\; the\; standard\; solution\; on\; the\; parabolic\;
neighborhood\; $B_{t_k}(x_k,$
$\bar\varepsilon^{-1}R(x_k,t_k)^{-\frac{1}{2}}) \times
[t_k-\bar\varepsilon^{-2}R(x_k,t_k)^{-1},t_k]$ is well-defined for
$k$ large. By Claim 1 in the proof of Proposition 7.3.3, there is
a positive function $\omega:[0,\infty)\rightarrow[0,\infty)$ such
that
$$
  R(x,t_k)\leq R(x_k,t_k)\omega(R(x_k,t_k)d_{t_k}^{2}(x,x_k))
$$
for all $x\in \mathbb{R}^3$.  Now by scaling the standard solution
$g_{ij}(\cdot,t)$ around $x_k$ with the factor $R(x_k,t_k)$ and
shifting the time $t_k$ to zero, we get a sequence of the rescaled
solutions $
\tilde{g}^{k}_{ij}(x,\tilde{t})=R(x_k,t_k)g_{ij}(x,t_k+\tilde{t}/R(x_k,t_k))$
to the Ricci flow defined on $\mathbb{R}^3$ with $\tilde{t}\in
[-R(x_k,t_k)t_k,0]$. We denote the scalar curvature and the
distance of the rescaled metric $\tilde{g}^{k}_{ij}$ by
$\tilde{R}^{k}$ and $\tilde{d}$. By combining with Claim 1 in the
proof of Proposition 7.3.3 and the Li-Yau-Hamilton inequality, we
get
\begin{align*}
\tilde{R}^{k}(x,0)& \leq \omega(\tilde{d}_{0}^{2}(x,x_k))\\
\tilde{R}^{k}(x,\tilde{t})& \leq
\frac{R(x_k,t_k)t_k}{\tilde{t}+R(x_k,t_k)t_k}
\omega(\tilde{d}_{0}^{2}(x,x_k))
\end{align*}
for any $x\in \mathbb{R}^3$ and $\tilde{t}\in (-R(x_k,t_k)t_k,0]$.
Note that $R(x_k,t_k)t_k\rightarrow \infty$ by Proposition 7.3.3.
We have shown in the proof of Proposition 7.3.3 that the standard
solution is $\kappa$-noncollapsed on all scales less than $1$ for
some $\kappa>0$. Then from the $\kappa$-noncollapsing property,
the above curvature estimates and Hamilton's compactness theorem,
we know $\tilde{g}^{k}_{ij}(x,\tilde{t})$ has a convergent
subsequence (as $k\rightarrow\infty$) whose limit is an ancient,
$\kappa$-noncollapsed, complete and orientable solution with
nonnegative curvature operator. This limit must have bounded
curvature by the same proof of Step 3 in the proof of Theorem
7.1.1. This gives a contradiction. Hence Assertion 1 is proved.

We now fix the constant $\theta(\varepsilon)$ obtained in
Assertion 1. Let $O$ be the tip of the standard capped infinite
cylinder $\mathbb{R}^3$ (it is rotationally symmetric about $O$ at
time $0$, and it remains so as $t>0$ by the uniqueness Theorem
1.2.4).

%\vskip 0.1cm\noindent\textbf{Assertion 2} \emph{
%CZ-SOURCE-LINE-18783
\begin{proposition}[early-time standard solution cap-or-neck structure]
  \label{prop:cz-standard-solution-early-time-cap-neck}
  \dcref{Cao--Zhu Assertion 2}
  \group{Ricci Flow with Surgery}
  \source{cao-zhu}

There are constants $B_1(\varepsilon)$ and $B_2(\varepsilon)$
depending only on $\varepsilon$, such that if $(x_0,t_0)\in
\mathbb{R}^3\times [0,\theta)$ with $d_{t_0}(x_0,O)\leq
B_{1}(\varepsilon)$, then there is a $0<r<B_2(\varepsilon)$ such
that $B_{t_0}(x_0,r)$ is an $\varepsilon$-cap; if $(x_0,t_0)\in
\mathbb{R}^3\times [0,\theta)$ with $d_{t_0}(x_0,O)\geq
B_1(\varepsilon)$, then the parabolic neighborhood
$B_{t_0}(x_0,\varepsilon^{-1}R(x_0,t_0)^{-\frac{1}{2}})$ $\times
[t_0-\min \{R(x_0,t_0)^{-1},t_0\},t_0]$ is after scaling with the
factor $R(x_0,t_0)$ and shifting the time $t_0$ to zero,
$\varepsilon$-close (in the $C^{[\varepsilon^{-1}]}$ topology) to
the corresponding subset of the evolving standard cylinder
$\mathbb{S}^2 \times \mathbb{R}$ over the time interval $[-\min
\{t_0R(x_0,t_0),1\},0]$ with scalar curvature $1$ at time zero.
\end{proposition}

Since the standard solution exists on the time interval $[0,1)$,
there is a constant $B_{0}(\varepsilon)$ such that the curvatures
on $[0,\theta(\varepsilon)]$ are uniformly bounded by
$B_{0}(\varepsilon)$. This implies that the metrics in
$[0,\theta(\varepsilon)]$ are equivalent.  Note that the initial
metric is asymptotic to the standard capped infinite cylinder. For
any sequence of points $x_{k}$ with $d_{0}(O,x_{k})\rightarrow
\infty$, after passing to a subsequence, $g_{ij}(x,t)$ around
$x_{k}$ will converge to a solution to the Ricci flow on
$\mathbb{R}\times \mathbb{S}^{2}$ with round cylinder metric of
scalar curvature 1 as initial data. By the uniqueness theorem
(Theorem 1.2.4), the limit solution must be the standard evolving
round cylinder.  This implies that there is a constant
$B_{1}(\varepsilon)>0$ depending on $\varepsilon$ such that for
any $(x_0,t_0)$ with $t_0\leq \theta(\varepsilon)$ and
$d_{t_0}(x_0,O)\geq B_{1}(\varepsilon)$, the standard solution on
the parabolic neighborhood
$B_{t_0}(x_0,\varepsilon^{-1}R(x_0,t_0)^{-\frac{1}{2}}) \times
[t_0-\min \{R(x_0,t_0)^{-1},t_0\},t_0]$ is, after scaling with the
factor $R(x_0,t_0)$, $\varepsilon$-close to the corresponding
subset of the evolving round cylinder. Since the solution is
rotationally symmetric around $O$, the cap neighborhood structures
of those points $x_0$ with $d_{t_0}(x_0,O)\leq B_{1}(\varepsilon)$
follow directly. Hence Assertion 2 is proved.

The combination of these two assertions proves the lemma.
%$$\eqno \#$$ \vskip 0.3cm
\end{proof}

Since there are only a finite number of horns with the other end
connected to $\Omega_\rho$, \emph{we perform only a finite number
of such $\delta$-cutoff surgeries at time $T$}. Besides those
horns, there could be capped horns and double horns which lie in
$\Omega \setminus \Omega_\rho$. As explained before, they are
connected to form tubes or capped tubes at any time slightly
before $T$. So \emph{we can regard the capped horns and double
horns (of $\Omega \setminus \Omega_\rho$) to be extinct and throw
them away at time $T$}. We only need to remember that the
connected sums were broken there. Remember that \emph{we have
thrown away all compact components, either lying in $\Omega
\setminus\Omega_\rho$ or with positive sectional curvature}, each
of which is diffeomorphic to either $\mathbb{S}^3$, or a metric
quotient of $\mathbb{S}^3$, or $\mathbb{S}^2 \times \mathbb{S}^1$
or $\mathbb{RP}^3\#\mathbb{RP}^3$. So we have also removed a
finite number of copies of $\mathbb{S}^3$, or metric quotients of
$\mathbb{S}^3$, or $\mathbb{S}^2 \times \mathbb{S}^1$ or
$\mathbb{RP}^3\#\mathbb{RP}^3$ at the time $T$. Let us agree to
\textbf{declare extinct every compact component either with
positive sectional curvature or lying in $\Omega
\setminus\Omega_\rho$}; in particular, this allows us to exclude
the components with positive sectional curvature from the list of
canonical neighborhoods.

\medskip
\emph{In summary, our surgery at time $T$ consists of the
following four procedures:}

(1) \emph{perform $\delta$-cutoff surgeries for all
$\varepsilon$-horns, whose other ends are  connected to
$\Omega_\rho$;}

(2) \emph{declare extinct every compact component which has
positive sectional curvature;}

(3) \emph{throw away all capped horns and double horns lying in
$\Omega \setminus \Omega_\rho$;}

(4) \emph{declare extinct all compact components lying in $\Omega
\setminus\Omega_\rho$.}

\medskip \noindent
({\it In Sections $7.6$ and $7.7,$ we will add one more procedure
by declaring extinct every compact component which has nonnegative
scalar curvature.})

\medskip
 By Lemma 7.3.4, after performing
surgeries at time $T$, the pinching assumption (7.3.3) still holds
for the surgically modified manifold. With this surgically
modified manifold (possibly disconnected) as initial data, we now
continue our solution under the Ricci flow until it becomes
singular again at some time $T'(>T)$. Therefore, we have extended
the solution to the Ricci flow with surgery, originally defined on
$[0,T)$ with $T<+\infty$, to the new time interval $[0,T')$ with
$T'>T$. By the proof of Theorem 5.3.2, we see that the solution to
the Ricci flow with surgery also satisfies the pinching assumption
on $[0,T')$. It remains to verify the canonical neighborhood
assumption (with accuracy $\varepsilon$) for the solution on the
time interval $[T,T')$ and to prove that this extension procedure
works indefinitely (unless it becomes extinct at some finite time)
and that there exists at most a finite number of surgeries at
every finite time interval. We leave these arguments to the next
section.

Before we end this section, we check the following two results of
Perelman in \cite{P2} which will be used in the next section to
estimate the Li-Yau-Perelman distance of space-time curves which
stretch to surgery regions. The proofs are basically given by
Perelman (cf. 4.5 and 4.6 of \cite{P2}).

%\vskip 0.2cm \noindent{\textbf{Lemma 7.3.6}}
%CZ-SOURCE-LINE-18901
\begin{lemma}[surgery stability near an almost-standard cap]
  \label{lem:cz-surgery-stability-near-standard-cap}
  \dcref{Cao--Zhu Lemma 7.3.6}
  \group{Ricci Flow with Surgery}
  \source{cao-zhu}

For any $0<\varepsilon \leq 1/100$, $1<A<+\infty$ and
$0<\theta<1$, one can find
$\bar{\delta}=\bar{\delta}(A,\theta,\varepsilon)$ with the
following property. Suppose we have a solution to the Ricci flow
with surgery which satisfies the a priori assumptions $($with
accuracy $\varepsilon)$ on $[0,T]$ and is obtained from a compact
orientable three-manifold by a finite number of $\delta$-cutoff
surgeries with each $\delta<\bar{\delta}$. Suppose we have a
cutoff surgery at time $T_0\in(0,T)$, let $x_0$ be any fixed point
on the gluing caps $($i.e., the regions affected by the cutoff
surgeries at time $T_0),$ and let $T_1=min\{T,T_0+\theta h^2\}$,
where $h$ is the cutoff radius around $x_0$ obtained in Lemma
$7.3.2.$ Then either
\begin{itemize}
\item[(i)] the solution is defined on
$P(x_0,T_0,Ah,T_1-T_0)\triangleq\{(x,t)\ |\ x\in B_t(x_0,Ah),\\
 t\in[T_0,T_1]\}$ and is, after scaling with factor $h^{-2}$ and
shifting time $T_0$ to zero, $A^{-1}$-close to a corresponding
subset of the standard solution, or \item[(ii)] the assertion (i)
holds with $T_1$ replaced by some time $t^+\in(T_0,T_1)$, where
$t^+$ is a surgery time; moreover, for each point in
$B_{T_0}(x_0,Ah)$, the solution is defined for $t\in[T_0,t^+)$ but
is not defined past $t^+$ $($i.e., the whole ball
$B_{T_0}(x_0,Ah)$ is cut off at the time $t^+).$
\end{itemize}
\end{lemma}

%\vskip 0.1cm \noindent{\bf Proof.}
\begin{proof}
  \uses{lem:cz-surgery-preserves-hamilton-ivey-pinching}
Let $Q$ be the maximum of the scalar curvature of the standard
solution in the time interval $[0,\theta]$ and choose a large
positive integer $N$ so that $\Delta
t=\frac{(T_1-T_0)}{N}<\varepsilon\eta^{-1}Q^{-1}h^2$, where the
positive constant $\eta$ is given in the canonical neighborhood
assumption. Set $t_k=T_0+k\Delta t$, $k=0, 1, \ldots, N$.

{}From Lemma 7.3.4, the geodesic ball $B_{T_0}(x_0,A_0h)$ at time
$T_0$, with $A_0=\delta^{-\frac{1}{2}}$ is, after scaling with
factor $h^{-2}$, $\delta^{\frac{1}{2}}$-close to the corresponding
ball in the standard capped infinite cylinder with the center near
the tip. Assume first that for each point in $B_{T_0}(x_0,A_0h)$,
the solution is defined on $[T_0,t_1]$. By the gradient estimates
(7.3.4) in the canonical neighborhood assumption and the choice of
$\Delta t$ we have a uniform curvature bound on this set for
$h^{-2}$-scaled metric. Then by the uniqueness theorem (Theorem
1.2.4), if $\delta^{\frac{1}{2}}\rightarrow0$ (i.e.
$A_0=\delta^{-\frac{1}{2}} \rightarrow +\infty$), the solution
with $h^{-2}$-scaled metric will converge to the standard solution
in the $C^\infty_{\rm loc}$ topology. Therefore we can define
$A_1$, depending only on $A_0$ and tending to infinity with $A_0$,
such that the solution in the parabolic region
$P(x_0,T_0,A_1h,t_1-T_0)\triangleq\{(x,t)\ |\ x\in
B_t(x_0,A_1h),t\in[T_0,T_0+(t_1-T_0)]\}$ is, after scaling with
factor $h^{-2}$ and shifting time $T_0$ to zero, $A^{-1}_1$-close
to the corresponding subset in the standard solution. In
particular, the scalar curvature on this subset does not exceed
$2Qh^{-2}$. Now if for each point in $B_{T_0}(x_0,A_1h)$ the
solution is defined on $[T_0,t_2]$, then we can repeat the
procedure, defining $A_2$, such that the solution in the parabolic
region $P(x_0,T_0,A_2h,t_2-T_0)\triangleq\{(x,t)\ |\ x\in
B_t(x_0,A_2h),t\in[T_0,T_0+(t_2-T_0)]\}$ is, after scaling with
factor $h^{-2}$ and shifting time $T_0$ to zero, $A^{-1}_2$-close
to the corresponding subset in the standard solution. Again, the
scalar curvature on this subset still does not exceed $2Qh^{-2}$.
Continuing this way, we eventually define $A_N$. Note that $N$ is
depends only on $\theta$ and $\varepsilon$. Thus there exists a
positive $\bar{\delta}=\bar{\delta}(A,\theta,\varepsilon)$ such
that for $\delta<\bar{\delta}$, we have $A_0>A_1>\cdots>A_N>A$,
and assertion (i) holds when the solution is defined on
$B_{T_0}(x_0,A_{(N-1)}h)\times [T_0,T_1]$.

The above argument shows that either assertion (i) holds, or there
exists some $k$ ($0\leq k\leq N-1$) and a surgery time
$t^+\in(t_k,t_{k+1}]$ such that the solution on
$B_{T_0}(x_0,A_kh)$ is defined on $[T_0,t^+)$, but for some point
of this set it is not defined past $t^+$. Now we consider the
latter case. Clearly the above argument also shows that the
parabolic region $P(x_0,T_0,A_{k+1}h,t^+-T_0)\triangleq\{(x,t)\ |\
x\in B_t(x,A_{k+1}h),t\in[T_0,t^+)\}$ is, after scaling with
factor $h^{-2}$ and shifting time $T_0$ to zero,
$A^{-1}_{k+1}$-close to the corresponding subset in the standard
solution. In particular, as time tends to $t^+$, the ball
$B_{T_0}(x_0,A_{k+1}h)$ keeps on looking like a cap. Since the
scalar curvature on $B_{T_0}(x_0,A_kh) \times [T_0,t_k]$ does not
exceed $2Qh^{-2}$, it follows from the pinching assumption, the
gradient estimates in the canonical neighborhood assumption and
the evolution equation of the metric that the diameter of the set
$B_{T_0}(x_0,A_kh)$ at any time $t \in [T_0,t^+)$ is bounded from
above by $4 \delta^{-\frac{1}{2}}h$. These imply that no point of
the ball $B_{T_0}(x_0,A_kh)$ at any time near $t^+$ can be the
center of a $\delta$-neck for any $0 < \delta <
\bar{\delta}(A,\theta,\varepsilon)$ with
$\bar{\delta}(A,\theta,\varepsilon)
> 0$ small enough, since $4 \delta^{-\frac{1}{2}}h$ is much
smaller than $\delta^{-1}h$. However the solution disappears
somewhere in the set $B_{T_0}(x_0,A_kh)$ at time $t^+$ by a cutoff
surgery and the surgery is always done along the middle two-sphere
of a $\delta$-neck. So the set $B_{T_0}(x_0,A_kh)$  at  time $t^+$
is a part of a capped horn. (Recall that we have declared extinct
every compact component with positive curvature and every compact
component lying in $\Omega \setminus \Omega_{\rho}$). Hence for
each point of $B_{T_0}(x_0,A_kh)$ the solution terminates at
$t^+$. This proves assertion (ii).
%$$\eqno \#$$ \vskip 0.3cm
\end{proof}

%\noindent{\bf Corollary 7.3.7} (Perelman \cite{P2}) \emph{
%CZ-SOURCE-LINE-19009
\begin{corollary}[large reduced length for curves entering surgery caps]
  \label{cor:cz-surgery-cap-curves-large-reduced-length}
  \dcref{Cao--Zhu Corollary 7.3.7}
  \group{Ricci Flow with Surgery}
  \source{cao-zhu}

For any $l<\infty$ one can find $A=A(l)<\infty$ and
$\theta=\theta(l)$, $0<\theta<1$, with the following property.
Suppose we are in the situation of the lemma above, with
$\delta<\bar{\delta}(A,\theta,\varepsilon)$. Consider smooth
curves $\gamma$ in the set $B_{T_0}(x_0,Ah)$, parametrized by
$t\in[T_0,T_{\gamma}]$, such that $\gamma(T_0)\in
B_{T_0}(x_0,\frac{Ah}{2})$ and either $T_{\gamma}=T_1<T$, or
$T_{\gamma}<T_1$ and $\gamma(T_{\gamma})\in \partial
B_{T_0}(x_0,Ah)$, where $x_0$ is any fixed point on a gluing cap
at $T_0$ and $T_1=min\{T,T_0+\theta h^2\}$. Then
$$
\int^{T_{\gamma}}_{T_0}(R(\gamma(t),t)+|\dot{\gamma}(t)|^2)dt>l.
$$
\end{corollary}

%\vskip 0.1cm\noindent{\bf Proof.}
\begin{proof}
  \uses{prop:cz-asymptotically-cylindrical-standard-solution, lem:cz-surgery-stability-near-standard-cap}
We know from Proposition 7.3.3 that on the standard solution,
\begin{align*}
\int^{\theta}_{0}Rdt
&  \geq{\rm  const.}\,\int^{\theta}_0(1-t)^{-1}dt\\
&  = -{\rm  const.}\,\cdot\log(1-\theta).
\end{align*}
By choosing $\theta=\theta(l)$ sufficiently close to 1 we have the
desired estimate for the standard solution.

Let us consider the first case: $T_{\gamma}=T_1<T$. For
$\theta=\theta(l)$ fixed above, by Lemma 7.3.6, our solution in
the subset $B_{T_0}(x_0,Ah)$ and in the time interval
$[T_0,T_{\gamma}]$ is, after scaling with factor $h^{-2}$ and
shifting time $T_0$ to zero, $A^{-1}$-close to the corresponding
subset in the standard solution for any sufficiently large $A$. So
we have
\begin{align*}
\int^{T_{\gamma}}_{T_0}(R(\gamma(t),t) + |\dot{\gamma}(t)|^2)dt
&  \geq {\rm const.}\,\int^{\theta}_0(1-t)^{-1}dt\\
&  = -{\rm const.}\,\cdot\log(1-\theta).
\end{align*}
Hence we have obtained the desired estimate in the first case.

We now consider the second case: $T_{\gamma}<T_1$ and
$\gamma(T_{\gamma})\in \partial B_{T_0}(x_0,Ah)$. Let
$\theta=\theta(l)$ be chosen above and let $Q=Q(l)$ be the maximum
of the scalar curvature on the standard solution in the time
interval $[0,\theta]$.

On the standard solution, we can choose $A=A(l)$ so large that for
each $t\in[0,\theta]$,
\begin{align*}
{\rm dist}_t(x_0,\partial B_0(x_0,A))
&  \geq {\rm dist}_0(x_0,\partial B_0(x_0,A))-4(Q+1)t\\
&  \geq A-4(Q+1)\theta\\
&  \geq \frac{4}{5}A
\end{align*}
and
$$
{\rm dist}_t\left(x_0,\partial
B_0\left(x_0,\frac{A}{2}\right)\right)\leq\frac{A}{2},
$$
where we have used Lemma 3.4.1(ii) in the first inequality. Now
our solution in the subset $B_{T_0}(x_0,Ah)$ and in the time
interval $[T_0,T_{\gamma}]$ is, after scaling with factor $h^{-2}$
and shifting time $T_0$ to zero, $A^{-1}$-close to the
corresponding subset in the standard solution. This implies that
for $A=A(l)$ large enough
$$
\frac{1}{5}Ah\leq\int^{T_{\gamma}}_{T_0}|\dot{\gamma}(t)|dt
\leq\left(\int^{T_{\gamma}}_{T_0}|
\dot{\gamma}(t)|^2dt\right)^{\frac{1}{2}}\cdot(T_{\gamma}-T_0)^{\frac{1}{2}},
$$
Hence
$$
\int^{T_{\gamma}}_{T_0}(R(\gamma(t),t) +
|\dot{\gamma}(t)|^2)dt\geq \frac{A^2}{25\theta}>l.
$$
This proves the desired estimate.
%$$\eqno\#$$
\end{proof}


\section{Justification of the Canonical Neighborhood Assumptions}

We continue the induction argument for the construction of a
long-time solution to the Ricci flow with surgery. Let us recall
what we have done in the previous section. Let $\varepsilon$ be an
arbitrarily given positive constant satisfying $0<\varepsilon \leq
1/100$. For an arbitrarily given compact orientable normalized
three-manifold, we evolve it by the Ricci flow. We may assume that
the solution goes singular at some time $0<t^+_1<+\infty$ and know
that the solution satisfies the a priori assumptions (with
accuracy $\varepsilon$) on $[0,t^+_1)$ for a nonincreasing
positive function $r=r_1(t)$ (defined on $[0,+\infty)$). Suppose
that we have a solution to the Ricci flow with surgery, defined on
$[0,t^+_k)$ with $0<t^+_1<t^+_2<\cdots<t^+_k<+\infty$, satisfying
the a priori assumptions (with accuracy $\varepsilon$) for some
nonincreasing positive function $r=r_k(t)$ (defined on
$[0,+\infty)$), going singular at time $t^+_k$ and having
$\delta_i$-cutoff surgeries at each time $t^+_i$, $1\leq i\leq
k-1$, where $\delta_i < \bar{\delta}(t^+_i)$ for each $1\leq i\leq
k-1$. Then for any $0<\delta_k < \bar{\delta}(t^+_k)$, we can
perform $\delta_k$-cutoff surgeries at the time $t^+_k$ and extend
the solution to the interval $[0,t^+_{k+1})$ with
$t^+_{k+1}>t^+_k$. Here $\bar{\delta}(t)$ is the positive function
defined in (7.3.16). We have already shown in Lemma 7.3.4 that the
extended solution still satisfies the pinching assumption on
$[0,t^+_{k+1})$.

In view of Theorem 7.1.1, there always is a nonincreasing positive
function $r=r_{k+1}(t)$, defined on $[0,+\infty)$, such that the
canonical neighborhood assumption (with accuracy $\varepsilon$)
holds on the extended time interval $[0,t^+_{k+1})$ with the
positive function $r=r_{k+1}(t)$. Nevertheless, in order to
prevent the surgery times from accumulating, the key is to choose
the nonincreasing positive functions $r=r_i(t), i= 1, 2, \ldots$,
uniformly. That is, to justify the canonical neighborhood
assumption (with accuracy $\varepsilon$) for the indefinitely
extending solution, we need to show that there exists a
nonincreasing positive function $\widetilde{r}(t)$, defined on
$[0,+\infty)$, which is independent of $k$, such that the above
chosen nonincreasing positive functions satisfy
$$r_i(t) \geq \widetilde{r}(t), \ \
\mbox{ on } \ \ [0,+\infty),$$ for all $i = 1, 2, \ldots, k+1$.

By a further restriction on the positive function
$\bar{\delta}(t)$, we can verify this after proving the following
assertion which was stated by Perelman in \cite{P2}.

%\vskip 0.2cm \noindent{\bf Proposition 7.4.1}
%CZ-SOURCE-LINE-19138
\begin{proposition}[canonical-neighborhood induction across surgery times]
  \label{prop:cz-canonical-neighborhood-induction-surgery}
  \dcref{Cao--Zhu Proposition 7.4.1}
  \group{Canonical-Neighborhood Induction}
  \source{cao-zhu}
 Given any small $\varepsilon>0$, there exist
decreasing sequences $0<\widetilde{r}_j<\varepsilon$,
$\kappa_j>0$, and $0<\widetilde{\delta}_j<\varepsilon^2 $,
$j=1,2,\ldots$, with the following property. Define the positive
function $\widetilde{\delta}(t)$ on $[0,+\infty)$ by
$\widetilde{\delta}(t) = \widetilde{\delta}_j$ for $t \in
[(j-1)\varepsilon^2,j\varepsilon^2)$. Suppose there is a
surgically modified solution, defined on $[0,T)$ with $T
<+\infty$, to the Ricci flow which satisfies the following:
\begin{itemize}
\item[(1)] it starts on a compact orientable three-manifold with
normalized initial metric, and \item[(2)] it has only a finite
number of surgeries such that each surgery at a time $t \in (0,T)$
is a $\delta(t)$-cutoff surgery with
$$
0 < \delta(t) \leq \min \{\widetilde{\delta}(t),\bar{\delta}(t)\}.
$$
\end{itemize}
Then on each time interval
$[(j-1)\varepsilon^2,j\varepsilon^2]\bigcap [0,T)$, $ j = 1, 2,
\cdots $, the solution satisfies the $\kappa_j$-noncollapsing
condition on all scales less than $\varepsilon$ and the canonical
neighborhood assumption $($with accuracy $\varepsilon)$ with
$r=\widetilde{r}_j$.
\end{proposition}

Here and in the following, we call a (three-dimensional)
surgically modified solution $g_{ij}(t), 0 \leq t <T$,
\textbf{$\kappa$-noncollapsed} \index{$\kappa$-noncollapsed} at
$(x_0,t_0)$ on the scales less than $\rho$ (for some
$\kappa>0,\rho>0$) if it satisfies the following property:
whenever $r < \rho$ and
$$
|Rm(x,t)| \leq r^{-2}
$$
for all those $(x,t) \in P(x_0,t_0,r,-r^2)=\{ (x',t') \ |\ x'\in
B_{t'}(x_0,r), t' \in [t_0 -r^2,t_0] \}$, for which the solution
is defined, we have
$$
\Vol_{t_0}(B_{t_0}(x_0,r)) \geq \kappa r^3.
$$
Before we give the proof of the proposition, we need to verify a
$\kappa$-non\-collapsing estimate which was given by Perelman in
\cite{P2}.

%\vskip 0.2cm\noindent{\bf Lemma 7.4.2}  \emph{
%CZ-SOURCE-LINE-19186
\begin{lemma} [noncollapsing and canonical-neighborhood induction step]
  \label{lem:cz-surgery-noncollapsing-canonical-induction-step}
  \dcref{Cao--Zhu Lemma 7.4.2}
  \group{Canonical-Neighborhood Induction}
  \source{cao-zhu}
  \uses{prop:cz-canonical-neighborhood-induction-surgery}

Given any $0<\varepsilon \leq \bar{\varepsilon}_0 $ $($for some
sufficiently small universal constant $\bar{\varepsilon}_0),$
suppose we have constructed the sequences satisfying the
proposition for $1\leq j\leq m$ $($for some positive integer $m)$.
Then there exists $\kappa>0$, such that for any $r$,
$0<r<\varepsilon$, one can find
$\widetilde{\delta}=\widetilde{\delta}(r,\varepsilon)$,
$0<\widetilde{\delta}<\varepsilon^2$, which may also depend on the
already constructed sequences, with the following property.
Suppose we have a solution with a compact orientable normalized
three-manifold as initial data, to the Ricci flow with finite
number of surgeries on a time interval $[0,\bar{T}]$ with
$m\varepsilon^2\leq \bar{T}<(m+1)\varepsilon^2$, satisfying the
assumptions and the conclusions of Proposition $7.4.1$ on
$[0,m\varepsilon^2)$, and the canonical neighborhood assumption
$($with accuracy $\varepsilon)$ with $r$ on
$[m\varepsilon^2,\bar{T}]$, as well as $0<\delta(t) \leq \min
\{\widetilde{\delta},\bar{\delta}(t)\}$ for any $\delta$-cutoff
surgery with $\delta =\delta(t)$ at a time
$t\in[(m-1)\varepsilon^2,\bar{T}]$. Then the solution is
$\kappa$-noncollapsed on $[0,\bar{T}]$ for all scales less than
$\varepsilon$.
\end{lemma}

%\vskip 0.1cm\noindent{\bf Proof.}
\begin{proof}
  \uses{lem:cz-surgery-cap-reduced-length-barrier, lem:cz-rescaled-ball-backward-extension, lem:cz-low-curvature-point-backward-extension, lem:cz-backward-control-or-capwise-comparability, lem:cz-rescaled-balls-uniform-curvature-backward-existence, lem:cz-extended-limit-uniform-spacetime-curvature}
  \uses{cor:cz-surgery-cap-curves-large-reduced-length}
Consider a parabolic neighborhood
$$
P(x_0,t_0,r_0,-r^2_0)\triangleq\{(x,t)\ |\ x\in
B_t(x_0,r_0),t\in[t_0-r^2_0,t_0]\}
$$
with $m\varepsilon^2\leq t_0\leq \bar{T}$ and $0<r_0<\varepsilon$,
where the solution satisfies $|Rm|\leq r^{-2}_0$, whenever it is
defined. We will use an argument analogous to the proof of Theorem
3.3.2 (no local collapsing theorem I) to prove \begin{equation}
\Vol_{t_0}(B_{t_0}(x_0,r_0)) \geq \kappa
r^3_0. %\eqno(7.4.1)
\end{equation}

Let $\eta$ be the universal positive constant in the definition of
the canonical neighborhood assumption. Without loss of generality,
we always assume $\eta \geq 10$. Firstly, we want to show that one
may assume $r_0 \geq \frac{1}{2\eta}r$.

Obviously, the curvature satisfies the estimate
$$
|Rm(x,t)|\leq 20r^{-2}_0,
$$
for those $(x,t) \in
P(x_0,t_0,\frac{1}{2\eta}r_0,-\frac{1}{8\eta}r^2_0) = \{(x,t)\ |\
x\in
B_t(x_0,\frac{1}{2\eta}r_0),t\in[t_0-\frac{1}{8\eta}r^2_0,t_0]\}$,
for which the solution is defined. When $r_0< \frac{1}{2\eta}r$,
we can enlarge $r_0$ to some $r'_0\in[r_0,r]$ so that
$$
|Rm|\leq 20r'^{-2}_0
$$
on $P(x_0,t_0,\frac{1}{2\eta}r'_0,-\frac{1}{8\eta}r'^2_0)$
(whenever it is defined), and either the equality holds somewhere
in $P(x_0,t_0,\frac
{1}{2\eta}r'_0,-(\frac{1}{8\eta}{r'}^2_0+\epsilon'))$ for
arbitrarily small $\epsilon'>0$, or $r'_0=r$.

In the case that the equality holds somewhere, it follows from the
pinching assumption that we have
$$
R > 10r'^{-2}_0
$$
somewhere in $P(x_0,t_0,\frac{1}{2\eta}r'_0,-(\frac{1}{8\eta}{r'}
^2_0+\epsilon'))$ for arbitrarily small $\epsilon'>0$. Here,
without loss of generality, we have assumed $r$ is suitably small.
Then by the gradient estimates in the definition of the canonical
neighborhood assumption, we know $$R(x_0,t_0) > r'^{-2}_0 \geq
r^{-2}.$$ Hence the desired noncollapsing estimate (7.4.1) in this
case follows directly from the canonical neighborhood assumption.
(Recall that we have excluded every compact component which has
positive sectional curvature in the surgery procedure and then we
have excluded them from the list of canonical neighborhoods. Here
we also used the standard volume comparison when the canonical
neighborhood is an $\varepsilon$-cap.)

While in the case that $r'_0 = r$, we have the curvature bound
$$
|Rm(x,t)|\leq \left(\frac{1}{2\eta}r\right)^{-2},
$$
for those $(x,t) \in
P(x_0,t_0,\frac{1}{2\eta}r,-(\frac{1}{2\eta}r)^{2}) = \{(x,t)\ |\
x\in
B_t(x_0,\frac{1}{2\eta}r),t\in[t_0-(\frac{1}{2\eta}r)^2,t_0]\}$,
for which the solution is defined. It follows from the standard
volume comparison that we only need to verify the noncollapsing
estimate (7.4.1) for $r_0 = \frac{1}{2\eta}r$. Thus we have
reduced the proof to the case $r_0\geq \frac{1}{2\eta}r$.

Recall from Theorem 3.3.2 that if a solution is smooth everywhere,
we can get a lower bound for the volume of the ball
$B_{t_0}(x_0,r_0)$ as follows: define $\tau(t)=t_0-t$ and consider
Perelman's reduced volume function and the Li-Yau-Perelman
distance associated to the point $x_0$; take a point $\bar{x}$ at
the time $ t=\varepsilon^2$ so that the Li-Yau-Perelman distance
$l$ attains its minimum
$l_{\min}(\tau)=l(\bar{x},\tau)\leq\frac{3}{2}$ for $\tau
=t_0-\varepsilon^2$; use it to obtain an upper bound for the
Li-Yau-Perelman distance from $x_0$ to each point of
$B_0(\bar{x},1)$, thus getting a lower bound for Perelman's
reduced volume at $\tau=t_0$; apply the monotonicity of Perelman's
reduced volume to deduce a lower bound for Perelman's reduced
volume at $\tau$ near $0$, and then get the desired estimate for
the volume of the ball $B_{t_0}(x_0,r_0)$. Now since our solution
has undergone surgeries, we need to localize this argument to the
region which is unaffected by surgery.

We call a space-time curve in the solution track
\textbf{admissible}\index{admissible curve} if it stays in the
space-time region unaffected by surgery, and we call a space-time
curve in the solution track a \textbf{barely admissible
curve}\index{barely admissible curve} if it is on the boundary of
the set of admissible curves.

 First of all, we want to estimate the $\mathcal{L}$-length of a
barely admissible curve.

\medskip
{\bf Claim.} \ For any $L<\infty$ one can find
$\bar{\delta}=\bar{\delta}(L,r,\widetilde{r}_m,\varepsilon)>0$
with the following property. Suppose that we have a curve
$\gamma$, parametrized by $t\in[T_0,t_0]$, $(m-1)\varepsilon^2\leq
T_0<t_0$, such that $\gamma(t_0)=x_0$, $T_0$ is a surgery time,
and $\gamma(T_0)$ lies in the gluing cap. Suppose also each
$\delta$-cutoff surgery at a time in
$[(m-1)\varepsilon^2,\bar{T}]$ has $\delta \leq \bar{\delta}$.
Then we have an estimate \begin{equation}
\int^{t_0}_{T_0}\sqrt{t_0-t}(R_+(\gamma(t),t)
+|\dot{\gamma}(t)|^2)dt\geq L  %\eqno(7.4.2)
\end{equation} where $R_+=\max\{R,0\}$.

\medskip
Since $r_0\geq \frac{1}{2\eta}r$ and $|Rm|\leq r^{-2}_0$ on
$P(x_0,t_0,r_0,-r^2_0)$ (whenever it is defined), we can require
$\bar{\delta}>0$, depending on $r$ and $\widetilde{r}_m$, to be so
small that $\gamma(T_0)$ does not lie in the region
$P(x_0,t_0,r_0,-r^2_0)$. Let $\Delta t$ be the maximal number such
that $\gamma|_{[t_0-\Delta t,t_0]}\subset P(x_0,t_0,r_0,-\Delta
t)$ (i.e., $t_0-\Delta t$ is the first time when $\gamma$ escapes
the parabolic region $P(x_0,t_0,r_0,-r^2_0)).$ Obviously we only
need to consider the case:
$$
\int^{t_0}_{t_0 -\Delta t}\sqrt{t_0 -
t}(R_+(\gamma(t),t)+|\dot{\gamma}(t)|^2)dt<L.
$$

We observe that $\Delta t$ can be bounded from below in terms of
$L$ and $r_0$. Indeed, if $\Delta t \geq r_0^2$, there is nothing
to prove. Thus we may assume $\Delta t<r_0^2$. By the curvature
bound $|Rm|\leq r^{-2}_0$ on $P(x_0,t_0,r_0,-r^2_0)$ and the Ricci
flow equation we see
$$
\int^{t_0}_{t_0-\Delta t}|\dot{\gamma}(t)|dt\geq cr_0
$$
for some universal positive constant $c$. On the other hand, by
the Cauchy-Schwarz inequality, we have
\begin{align*}
\int^{t_0}_{t_0-\Delta t}|\dot{\gamma}(t)|dt & \leq
\left(\int^{t_0}_{t_0-\Delta t}\sqrt{t_0-t}(R_++|
\dot{\gamma}|^2)dt\right)^{\frac{1}{2}} \cdot\left(\int^{t_0}_{t_0-\Delta
t}\frac{1}{\sqrt{t_0-t}}dt\right)^{\frac{1}{2}}\\
&  \leq (2L)^{\frac{1}{2}}(\Delta t)^{\frac{1}{4}},
\end{align*}
which implies \begin{equation}
(\Delta t)^{\frac{1}{2}}\geq\frac{c^2r^2_0}{2L}. %\eqno (7.4.3)
\end{equation} Thus
\begin{align*}
\int^{t_0}_{T_0}\sqrt{t_0-t}(R_++|\dot{\gamma}|^2)dt& \geq
\int^{t_0-\Delta
t}_{T_0}\sqrt{t_0-t}(R_++|\dot{\gamma}|^2)dt\\
&  \geq (\Delta t)^{\frac{1}{2}}\int^{t_0-\Delta
t}_{T_0}(R_++|\dot{\gamma}|^2)dt\\
&  \geq \left(\min \left\{\frac{c^2r^2_0}{2L},r_0 \right\}\right)
\int^{t_0-\Delta t}_{T_0}(R_++|\dot{\gamma}|^2)dt,
\end{align*}
while by Corollary 7.3.7, we can find
$\bar{\delta}=\bar{\delta}(L,r,\widetilde{r}_m,\varepsilon)>0$ so
small that
$$
\int^{t_0-\Delta t}_{T_0}(R_++|\dot{\gamma}|^2)dt \geq
L\left(\min\left\{\frac{c^2r^2_0}{2L},r_0 \right\}\right)^{-1}.
$$
Then we have proved the desired assertion (7.4.2).

Recall that for a curve $\gamma$, parametrized by
$\tau=t_0-t\in[0,\bar{\tau}],$ with $\gamma(0)=x_0$ and
$\bar{\tau}\leq t_0-(m-1)\varepsilon^2$, we have
$L(\gamma)=\int^{\bar{\tau}}_0\sqrt{\tau}(R+|\dot{\gamma}|^2)d\tau$.
We can also define $L_+(\gamma)$ by replacing $R$ with $R_+$ in
the previous formula. Recall that $R\geq-1$ at the initial time
$t=0$ for the normalized initial manifold. Recall that the
surgeries occur at the parts where the scalar curvatures are very
large. Thus we can apply the maximum principle to conclude that
the solution with surgery still satisfies $R\geq -1$  everywhere
in space-time. This implies \begin{equation} L_+(\gamma)\leq
L(\gamma)+(2\varepsilon^2)^{\frac{3}{2}}.
%\eqno(7.4.4)
\end{equation} By applying the assertion (7.4.2), we now choose
$\tilde{\delta}>0$ (depending on $r$, $\varepsilon$ and
$\widetilde{r}_m$) such that as each $\delta$-cutoff surgery at
the time interval $t\in[(m-1)\varepsilon^2,T]$ has $\delta \leq
\tilde{\delta}$, every barely admissible curve $\gamma$ from
$(x_0,t_0)$ to a point $(x,t)$ (with
$t\in[(m-1)\varepsilon^2,t_0)$) has
$$
L_+(\gamma)\geq22\sqrt{2}.
$$
Thus if the Li-Yau-Perelman distance from $(x_0,t_0)$ to a point
$(x,t)$ (with $t\in[(m-1)\varepsilon^2,t_0)$) is achieved by a
space-time curve which is not admissible, then its Li-Yau-Perelman
distance has \begin{equation}
l\geq\frac{L_+-(2\varepsilon^2)^\frac{3}{2}}{2\sqrt{2}
\varepsilon}>10\varepsilon^{-1}. %\eqno (7.4.5)
\end{equation} We also observe that the absolute value of $l(x_0,\tau)$ is
very small as $\tau$ closes to zere. Thus the maximum principle
argument in Corollary 3.2.6 still works for our solutions with
surgery because barely admissible curves do not attain the
minimum. So we conclude that
$$
l_{\min}(\bar{\tau})=\min\{l(x,\bar{\tau})\ |\ x \mbox{ lies in
the solution manifold at }t_0-\bar{\tau}\} \leq\frac{3}{2}
$$
for $\bar{\tau}\in(0,t_0-(m-1)\varepsilon^2]$. In particular,
there exists a minimizing curve $\gamma$ of
$l_{\min}(t_0-(m-1)\varepsilon^2)$, defined on
$\tau\in[0,t_0-(m-1)\varepsilon^2]$ with $\gamma(0)=x_0$, such
that
\begin{align}
L_+(\gamma)
& \leq \frac{3}{2}\cdot2\sqrt{2}\varepsilon+2\sqrt{2}\varepsilon^3\\
&  \leq 5\varepsilon, \nonumber
\end{align}  %\eqno (7.4.6)
since $0<\varepsilon \leq \bar{\varepsilon}_0$ with
$\bar{\varepsilon}_0$ sufficiently small (to be further
determined). Consequently, there exists a point
$(\bar{x},\bar{t})$ on the minimizing curve $\gamma$ with
$\bar{t}\in[(m-1)\varepsilon^2+\frac{1}{4}\varepsilon^2,
(m-1)\varepsilon^2+\frac{3}{4}\varepsilon^2]$ (i.e.,
$\tau\in[t_0-(m-1)\varepsilon^2-\frac{3}{4}\varepsilon^2,
t_0-(m-1)\varepsilon^2-\frac{1}{4}\varepsilon^2])$ such that \begin{equation}
R(\bar{x},\bar{t})\leq25\widetilde{r}^{-2}_m. %\eqno(7.4.7)
\end{equation} Otherwise, we have
\begin{align*}
L_+(\gamma)& \geq
\int^{t_0-(m-1)\varepsilon^2-\frac{1}{4}\varepsilon^2}_{t_0-(m-1)
\varepsilon^2-\frac{3}{4}\varepsilon^2}
\sqrt{\tau}R(\gamma(\tau),t_0-\tau)d\tau\\
&  > 25\widetilde{r}^{-2}_m\sqrt{\frac{1}{4}\varepsilon^2}
\left(\frac{1}{2}\varepsilon^2\right)\\
&  > 5\varepsilon,
\end{align*}
since $0<\widetilde{r}_m<\varepsilon$. This contradicts (7.4.6).

Next we want to get a lower bound for Perelman's reduced volume of
a ball around $\bar{x}$ of radius about $\widetilde{r}_m$ at some
time slightly before $\bar{t}$.

Denote by $\theta_1 = \frac{1}{16}\eta^{-1}$ and $\theta_2 =
\frac{1}{64}\eta^{-1}$, where $\eta$ is the universal positive
constant in the gradient estimates (7.3.4). Since the solution
satisfies the canonical neighborhood assumption on the time
interval $[(m-1)\varepsilon^2,m\varepsilon^2)$, it follows from
the gradient estimates (7.3.4) that \begin{equation}
R(x,t)\leq 400\widetilde{r}^{-2}_m  %\eqno (7.4.8)
\end{equation} for those $(x,t)\in P(\bar{x},\bar{t},\theta_1\widetilde{r}_m,
-\theta_2\widetilde{r}^2_m)\triangleq\{(x',t')\ |\ x'\in
B_{t'}(\bar{x},\theta_1\widetilde{r}_m),t'\in[\bar{t}
-\theta_2\widetilde{r}^2_m,\bar{t}]\}$, for which the solution is
defined. And since the scalar curvature at the points where the
$\delta$-cutoff surgeries occur in the time interval
$[(m-1)\varepsilon^2,m\varepsilon^2)$ is at least
$(\widetilde{\delta})^{-2}\widetilde{r}^{-2}_m$, the solution is
well-defined on the whole parabolic region
$P(\bar{x},\bar{t},\theta_1\widetilde{r}_m,-\theta_2\widetilde{r}^2_m)$
(i.e., this parabolic region is unaffected by surgery). Thus by
combining (7.4.6) and (7.4.8), we know that the Li-Yau-Perelman
distance from $(x_0,t_0)$ to each point of the ball
$B_{\bar{t}-\theta_2\widetilde{r}^2_m}(\bar{x},
\theta_1\widetilde{r}_m)$ is uniformly bounded by some universal
constant. Let us define Perelman's reduced volume of the ball
$B_{\bar{t}-\theta_2\widetilde{r}^2_m}(\bar{x},
\theta_1\widetilde{r}_m)$, by
\begin{align*}
&\widetilde{V}_{t_0-\bar{t}+\theta_2\widetilde{r}^2_m}
(B_{\bar{t}-\theta_2\widetilde{r}^2_m}(\bar{x},\theta_1\widetilde{r}_m))\\
&=\int_{B_{\bar{t}-\theta_2\widetilde{r}^2_m}(\bar{x},
\theta_1\widetilde{r}_m)}(4\pi(t_0-\bar{t}+\theta_2\widetilde{r}^2_m))^
{-\frac{3}{2}} \\
&\qquad\cdot\exp(-l(q,t_0-\bar{t}+\theta_2\widetilde{r}^2_m))
dV_{\bar{t}-\theta_2\widetilde{r}^2_m}(q),
\end{align*}
where $l(q,\tau)$ is the Li-Yau-Perelman distance from
$(x_0,t_0)$. Hence by the $\kappa_m$-noncollapsing assumption on
the time interval $[(m-1)\varepsilon^2,m\varepsilon^2)$, we
conclude that Perelman's reduced volume of the ball
$B_{\bar{t}-\theta_2\widetilde{r}^2_m}(\bar{x},\theta_1\widetilde{r}_m)$
is bounded from below by a positive constant depending only on
$\kappa_m$ and $\widetilde{r}_m$.

Finally we want to apply a local version of the monotonicity of
Perelman's reduced volume to get a lower bound estimate for the
volume of the ball $B_{t_0}(x_0,r_0)$.

We have seen that the Li-Yau-Perelman distance from $(x_0,t_0)$ to
each point of the ball
$B_{\bar{t}-\theta_2\widetilde{r}^2_m}(\bar{x},\theta_1\widetilde{r}_m)$
is uniformly bounded by some universal constant. Now we can choose
a sufficiently small (universal) positive constant
$\bar{\varepsilon}_0$ such that when $0<\varepsilon \leq
\bar{\varepsilon}_0$, by (7.4.5), all the points in the ball
$B_{\bar{t}-\theta_2\widetilde{r}^2_m}(\bar{x},\theta_1\widetilde{r}_m)$
can be connected to $(x_0,t_0)$ by shortest
$\mathcal{L}$-geodesics, and all of these $\mathcal{L}$-geodesics
are admissible (i.e., they stay in the region unaffected by
surgery). The union of all shortest $\mathcal{L}$-geodesics from
$(x_0,t_0)$ to the ball
$B_{\bar{t}-\theta_2\widetilde{r}^2_m}(\bar{x},
\theta_1\widetilde{r}_m)$ defined by
$CB_{\bar{t}-\theta_2\widetilde{r}^2_m}(\bar{x},
\theta_1\widetilde{r}_m)=\{(x,t)\ |\ (x,t)$ lies in a shortest
$\mathcal{L}$-geodesic from $(x_0,t_0)$ to a point in
$B_{\bar{t}-\theta_2\widetilde{r}^2_m}(\bar{x},
\theta_1\widetilde{r}_m)\}$, forms a cone-like subset in
space-time with the vertex $(x_0,t_0)$. Denote $B(t)$ by the
intersection of the cone-like subset
$CB_{\bar{t}-\theta_2\widetilde{r}^2_m}(\bar{x},\theta_1\widetilde{r}_m)$
with the time-slice at $t$. Perelman's reduced volume of the
subset $B(t)$ is given by
$$
\widetilde{V}_{t_0-t}(B(t))
=\int_{B(t)}(4\pi(t_0-t))^{-\frac{3}{2}}exp(-l(q,t_0-t))dV_t(q).
$$
Since the cone-like subset
$CB_{\bar{t}-\theta_2\widetilde{r}^2_m}(\bar{x},\theta_1\widetilde{r}_m)$
lies entirely in the region unaffected by surgery, we can apply
Perelman's Jacobian comparison theorem (Theorem 3.2.7) to conclude
that
\begin{align}
\widetilde{V}_{t_0-t}(B(t))& \geq
\widetilde{V}_{t_0-\bar{t}+\theta_2\widetilde{r}^2_m}
(B_{\bar{t}-\theta_2\widetilde{r}^2_m}(\bar{x},
\theta_1\widetilde{r}_m))\\
&  \geq c(\kappa_m,\widetilde{r}_m), \nonumber
\end{align}  %\eqno (7.4.9)
for all $t\in[\bar{t}-\theta_2\widetilde{r}^2_m,t_0]$, where
$c(\kappa_m,\widetilde{r}_m)$ is some positive constant depending
only on $\kappa_m$ and $\widetilde{r}_m$.

Set $\xi=r^{-1}_0\Vol_{t_0}(B_{t_0}(x_0,r_0))^{\frac{1}{3}}$. Our
purpose is to give a positive lower bound for $\xi$. Without loss
of generality, we may assume $\xi<\frac{1}{4}$, thus $0<\xi
r^2_0<t_0-\bar{t}+\theta_2\widetilde{r}^2_m$. Denote by
$\widetilde{B}(t_0-\xi r^2_0)$ the subset of the time-slice
$\{t=t_0-\xi r^2_0\}$ of which every point can be connected to
$(x_0,t_0)$ by an admissible shortest $\mathcal{L}$-geodesic.
Clearly, $B(t_0-\xi r^2_0)\subset\widetilde{B}(t_0-\xi r^2_0)$. We
now argue as in the proof of Theorem 3.3.2 to bound Perelman's
reduced volume of $\widetilde{B}(t_0-\xi r^2_0)$ from above.

Since $r_0\geq \frac{1}{2\eta}r$ and
$\tilde{\delta}=\tilde{\delta}(r,\varepsilon,\widetilde{r}_m)$
sufficiently small, the whole region $P(x_0,t_0,r_0,$ $-r^2_0)$ is
unaffected by surgery. Then by exactly the same argument as in
deriving (3.3.5), we see that there exists a universal positive
constant $\xi_0$ such that when $0<\xi\leq\xi_0$, there holds \begin{equation}
\mathcal{L} \exp_{\{|\upsilon|\leq
\frac{1}{4}\xi^{-\frac{1}{2}}\}}(\xi r^2_0)\subset
B_{t_0}(x_0,r_0). %\eqno (7.4.10)
\end{equation} Perelman's reduced volume of $\widetilde{B}(t_0-\xi r^2_0)$ is
given by
\begin{align}
& \widetilde{V}_{\xi r^2_0}(\widetilde{B}(t_0-\xi r^2_0))\\
&  = \int_{\widetilde{B}(t_0-\xi r^2_0)}(4\pi\xi
r^2_0)^{-\frac{3}{2}}\exp(-l(q,\xi r^2_0))dV_{t_0-\xi
r^2_0}(q)\nonumber\\
&  = \int_{\widetilde{B}(t_0-\xi r^2_0)\cap \mathcal{L}
\exp_{\{|\upsilon|\leq \frac{1}{4}\xi^{-\frac{1}{2}}\}}(\xi
r^2_0)}(4\pi\xi r^2_0)^{-\frac{3}{2}}\exp(-l(q,\xi
r^2_0))dV_{t_0-\xi r^2_0}(q)\nonumber\\
& \quad +\int_{\widetilde{B}(t_0-\xi r^2_0)\setminus
\mathcal{L}\exp_{\{|\upsilon|\leq
\frac{1}{4}\xi^{-\frac{1}{2}}\}}(\xi r^2_0)}(4\pi\xi
r^2_0)^{-\frac{3}{2}}\exp(-l(q,\xi r^2_0))dV_{t_0-\xi
r^2_0}(q).\nonumber
\end{align}  %\eqno (7.4.11)
The first term on the RHS of (7.4.11) can be estimated by
\begin{align}
&\int_{\widetilde{B}(t_0-\xi r^2_0)\cap \mathcal{L}
\exp_{\{|\upsilon| \leq\frac{1}{4}\xi^{-\frac{1}{2}}\}}(\xi
r^2_0)}(4\pi\xi
r^2_0)^{-\frac{3}{2}}\exp(-l(q,\xi r^2_0))dV_{t_0-\xi r^2_0}(q)\\
&\leq e^{C\xi}(4\pi)^{-\frac{3}{2}}\cdot\xi^{\frac{3}{2}} \nonumber
\end{align}  %\eqno(7.4.12)
for some universal constant $C$, as in deriving (3.3.7). While as
in deriving (3.3.8), the second term on the RHS of (7.4.11) can be
estimated by
\begin{align}
&  \int_{\widetilde{B}(t_0-\xi r^2_0)\setminus\mathcal{L}
\exp_{\{|\upsilon|\leq\frac{1}{4}\xi^{-\frac{1}{2}}\}}{(\xi
r^2_0)}}(4\pi \xi r^2_0)^{-\frac{3}{2}}\exp(-l(q,\xi
r^2_0))dV_{t_0-\xi r^2_0}(q)\\
&  \leq \int_{\{|\upsilon|>\frac{1}{4}\xi^{-\frac{1}{2}}\}}
(4\pi\tau)^{-\frac{3}{2}}\exp(-l(\tau))\mathcal{J}
(\tau)|_{\tau=0}d\upsilon\nonumber\\
& =(4\pi)^{-\frac{3}{2}}\int_{\{|\upsilon|>\frac{1}{4}
\xi^{-\frac{1}{2}}\}}\exp(-|\upsilon|^2)d\upsilon,\nonumber
\end{align}  % \eqno(7.4.13)
where we have used Perelman's Jacobian comparison theorem (Theorem
3.2.7) in the first inequality. Hence the combination of (7.4.9),
(7.4.11), (7.4.12) and (7.4.13) bounds $\xi$ from below by a
positive constant depending only on $\kappa_m$ and
$\widetilde{r}_m$. Therefore we have completed the proof of the
lemma.
%$$\eqno \#$$ \vskip 0.3cm
\mbox{ \ \ }\end{proof}

We are ready to prove the proposition.

%\vskip0.1cm \noindent{\bf Proof of Proposition 7.4.1}:
\medskip
{\bf\em Proof of Proposition} {\bf 7.4.1.} \ We now follow
Perelman \cite{P2} to prove the proposition by induction: having
constructed our sequences for $1\leq j\leq m$, we make one more
step, defining $\widetilde{r}_{m+1}$, $\kappa_{m+1}$,
$\widetilde{\delta}_{m+1}$, and redefining
$\widetilde{\delta}_m=\widetilde{\delta}_{m+1}$. In view of the
previous lemma, we only need to define $\widetilde{r}_{m+1}$ and
$\widetilde{\delta}_{m+1}$.

In Theorem 7.1.1 we have obtained the canonical neighborhood
structure for smooth solutions. When adapting the arguments in the
proof of Theorem 7.1.1 to the present surgical solutions, we will
encounter the new difficulty of how to take a limit for the
surgically modified solutions. The idea to overcome the difficulty
consists of two parts. The first part, due to Perelman \cite{P2},
is to choose $\widetilde{\delta}_{m}$ and
$\widetilde{\delta}_{m+1}$ small enough to push the surgical
regions to infinity in space. (This is the reason why we need to
redefine $\widetilde{\delta}_{m} = \widetilde{\delta}_{m+1}$.) The
second part (see Assertions 1-3 proved in step 2 below), due to
the authors and Bing-Long Chen, is to show that solutions are
smooth on some small, but uniform, time intervals (on compact
subsets) so that we can apply Hamilton's compactness theorem,
since we only have curvature bounds; otherwise Shi's interior
derivative estimate may not be applicable.

We now argue by contradiction. Suppose for sequence of positive
numbers $r^{\alpha}$ and $\widetilde{\delta}^{\alpha\beta}$,
satisfying $r^{\alpha}\rightarrow0$ as $\alpha\rightarrow\infty$ and
$\widetilde{\delta}^{\alpha\beta}\leq\frac{1}{\alpha \cdot
\beta}(\rightarrow0)$, there exist sequences of solutions
$g^{\alpha\beta}_{ij}$ to the Ricci flow with surgery, where each of
them has only a finite number of cutoff surgeries and has a compact
orientable normalized three-manifold as initial data, so that the
following two assertions hold:
\begin{itemize}
\item[(i)] each $\delta$-cutoff at a time
$t\in[(m-1)\varepsilon^2,(m+1)\varepsilon^2]$ satisfies $\delta
\leq \widetilde{\delta}^{\alpha\beta}$; and \item[(ii)] the
solutions satisfy the statement of the proposition on
$[0,m\varepsilon^2]$, but violate the canonical neighborhood
assumption (with accuracy $\varepsilon$) with $r=r^{\alpha}$ on
$[m\varepsilon^2,(m+1)\varepsilon^2]$.
\end{itemize}

For each solution $g^{\alpha\beta}_{ij}$, we choose $\bar{t}$
(depending on $\alpha$ and $\beta$) to be the nearly first time
for which the canonical neighborhood assumption (with accuracy
$\varepsilon$) is violated. More precisely, we choose
$\bar{t}\in[m\varepsilon^2,(m+1)\varepsilon^2]$ so that the
canonical neighborhood assumption with $r=r^{\alpha}$ and with
accuracy parameter $\varepsilon$ is violated at some
$(\bar{x},\bar{t})$, however the canonical neighborhood assumption
with accuracy parameter $2\varepsilon$ holds on
$t\in[m\varepsilon^2,\bar{t}]$. After passing to subsequences, we
may assume each $\widetilde{\delta}^{\alpha\beta}$ is less than
the $\widetilde{\delta}$ in Lemma 7.4.2 with $r=r^{\alpha}$ when
$\alpha$ is fixed. Then by Lemma 7.4.2 we have uniform
$\kappa$-noncollapsing on all scales less than $\varepsilon$ on
$[0,\bar{t}]$ with some $\kappa>0$ independent of $\alpha,\beta$.

Slightly abusing notation, we will often drop the indices $\alpha$
and $\beta$.

Let $\widetilde{g}^{\alpha\beta}_{ij}$ be the rescaled solutions
around $(\bar{x},\bar{t})$ with factors $R(\bar{x},\bar{t})(\geq
r^{-2}\rightarrow+\infty)$ and shift the times $\bar{t}$ to zero.
We hope to take a limit of the rescaled solutions for subsequences
of $\alpha,\beta\rightarrow\infty$ and show the limit is an
orientable ancient $\kappa$-solution, which will give the desired
contradiction. We divide our arguments into the following six
steps.

\medskip
{\it Step} 1. \  Let $(y,\hat{t})$ be a point on the rescaled
solution $\widetilde{g}^{\alpha\beta}_{ij}$ with
$\widetilde{R}(y,\hat{t})\leq A$ (for some $A\geq 1$) and
$\hat{t}\in [-(\bar{t} -
(m-1)\varepsilon^{2})R(\bar{x},\bar{t}),0]$. Then we have estimate
\begin{equation}
\widetilde{R}(x,t)\leq10A %\eqno (7.4.14)
\end{equation} for those $(x,t)$ in the parabolic neighborhood
$P(y,\hat{t},\frac{1}{2}\eta^{-1}A^{-\frac{1}{2}},
-\frac{1}{8}\eta^{-1}A^{-1})$ $\triangleq\{(x',t')\ |\ x'\in
\widetilde{B}_{t'}(y,\frac{1}{2}\eta^{-1}A^{-\frac{1}{2}}),
t'\in[\hat{t}-\frac{1}{8}\eta^{-1}A^{-1},\hat{t}]\}$, for which
the rescaled solution is defined.

Indeed, as in the first step of the proof of Theorem 7.1.1, this
follows directly from the gradient estimates (7.3.4) in the
canonical neighborhood assumption with parameter $2\varepsilon$.

\medskip
{\it Step} 2. \ In this step, we will prove three time extension
results.

\medskip
{\bf Assertion 1.} \ For arbitrarily fixed $\alpha$,
$0<A<+\infty$,  $1 \leq C<+\infty$ and $0 \leq B <
\frac{1}{2}\varepsilon^2(r^{\alpha})^{-2}-
\frac{1}{8}\eta^{-1}C^{-1}$, there is a
$\beta_0=\beta_0(\varepsilon, A, B, C)$ (independent of $\alpha$)
such that if $\beta\geq\beta_0$ and the rescaled solution
$\widetilde{g}^{\alpha\beta}_{ij}$ on the ball
$\widetilde{B}_{0}(\bar{x},A)$ is defined on a time interval
$[-b,0]$ with $0 \leq b \leq B$ and the scalar curvature satisfies
$$\widetilde{R}(x,t)\leq C , \ \ \mbox{ on }
\widetilde{B}_{0}(\bar{x},A) \times [-b,0],$$ then the rescaled
solution $\widetilde{g}^{\alpha\beta}_{ij}$ on the ball
$\widetilde{B}_{0}(\bar{x},A)$ is also defined on the extended
time interval $ [-b-\frac{1}{8}\eta^{-1}C^{-1},0]$.

\medskip
Before giving the proof, we make a simple \textbf{observation}:
once a space point in the Ricci flow with surgery is removed by
surgery at some time, then it never appears for later time; if a
space point at some time $t$ cannot be defined before the time $t$
, then either the point lies in a gluing cap of the surgery at
time $t$ or the time $t$ is the initial time of the Ricci flow.

\medskip
{\bf \em Proof of Assertion} {\bf 1.} \ Firstly we claim that
there exists $\beta_0=\beta_0( \varepsilon, A, B, C)$ such that
when $\beta\geq\beta_0$, the rescaled solution
$\widetilde{g}^{\alpha\beta}_{ij}$ on the ball
$\widetilde{B}_{0}(\bar{x},A)$ can be defined before the time $-b$
(i.e., there are no surgeries interfering in
$\widetilde{B}_{0}(\bar{x},A)\times [-b-\epsilon',-b]$ for some
$\epsilon'>0$).

We argue by contradiction. Suppose not, then there is some point
$\tilde{x}\in \widetilde{B}_{0}(\bar{x},A)$ such that the rescaled
solution $\widetilde{g}^{\alpha\beta}_{ij}$ at $\tilde{x}$ cannot
be defined before the time $-b$. By the above observation, there
is a surgery at the time $-b$ such that the point $\tilde{x}$ lies
in the instant gluing cap.

Let\; $\tilde{h}$\; $(=R(\bar{x},\,\bar{t})^{\frac{1}{2}}h$)\;
be\; the\; cut-off\; radius at the\; time\; $-b$\; for the
rescaled solution.  Clearly, there is a universal constant $D$
such that
${D}^{-1}\tilde{h}\leq\widetilde{R}(\tilde{x},-b)^{-\frac{1}{2}}\leq
D\tilde{h}$.

By Lemma 7.3.4 and looking at the rescaled solution at the time
$-b$, the gluing cap and the adjacent $\delta$-neck, of radius
$\tilde{h}$, constitute a
${(\widetilde{\delta}^{\alpha\beta})}^{\frac{1}{2}}$-cap
$\mathcal{K}$. For any fixed small positive constant
$\delta^{\prime}$ (much smaller than $\varepsilon$), we see that
$$
\widetilde{B}_{(-b)}(\tilde{x},{(\delta^{\prime})}^{-1}
\widetilde{R}(\tilde{x}, -b)^{-\frac{1}{2}})\subset\mathcal{K}
$$
when $\beta$ large enough.  We first verify the following

\medskip
{\bf Claim 1.} \ For any small constants $0<\tilde{\theta}<1$,
$\delta'>0$, there exists a $\beta(\delta',\varepsilon,
\tilde{\theta})> 0$ such that when $\beta \geq
\beta(\delta',\varepsilon, \tilde{\theta})$, we have
\begin{itemize}
\item[(i)] the rescaled solution
$\widetilde{g}^{\alpha\beta}_{ij}$ over
$\widetilde{B}_{(-b)}(\tilde{x},{(\delta^{\prime})}^{-1}
\tilde{h})$  is defined on the time interval
$[-b,0]\cap[-b,-b+(1-\tilde{\theta})\tilde{h}^{2}]$; \item[(ii)]
the ball
$\widetilde{B}_{(-b)}(\tilde{x},{(\delta^{\prime})}^{-1}\tilde{h})$
in the ${(\widetilde{\delta}^{\alpha\beta})}^{\frac{1}{2}}$-cap
$\mathcal{K}$ evolved by the Ricci flow  on the time interval
$[-b,0]\cap[-b,-b+(1-\tilde{\theta})\tilde{h}^{2}]$ is, after
scaling with factor $\tilde{h}^{-2}$, ${\delta}^{\prime}$-close
(in the $C^{[\delta'^{-1}]}$ topology) to the corresponding subset
of the standard solution.
\end{itemize}

\medskip
This claim essentially follows from Lemma 7.3.6. Indeed, suppose
there is a surgery at some time $\tilde{\tilde{t}}\in
[-b,0]\cap(-b,-b+(1-\tilde{\theta})\tilde{h}^{2}]$ which removes
some point  $\tilde{\tilde{x}}\in \widetilde
B_{(-b)}(\tilde{x},{(\delta^{\prime})}^{-1}\tilde{h})$. We assume
$\tilde{\tilde{t}}\in (-b,0]$ is the first time with that
property.

Then by Lemma 7.3.6, there is a
$\bar{\delta}=\bar{\delta}(\delta',\varepsilon,\tilde{\theta})$
such that if $\widetilde{\delta}^{\alpha\beta}<\bar{\delta}$, then
the ball
$\widetilde{B}_{(-b)}(\tilde{x},{(\delta^{\prime})}^{-1}\tilde{h})$
in the ${(\widetilde{\delta}^{\alpha\beta})}^{\frac{1}{2}}$-cap
$\mathcal{K}$ evolved by the Ricci flow  on the time interval
$[-b,\tilde{\tilde{t}})$ is, after scaling with factor
$\tilde{h}^{-2}$, ${\delta}^{\prime}$-close to the corresponding
subset of the standard solution. Note that the metrics for times
in $[-b,\tilde{\tilde{t}})$ on $\widetilde
B_{(-b)}(\tilde{x},{(\delta^{\prime})}^{-1}\tilde{h})$  are
equivalent. By Lemma 7.3.6, the solution on $\widetilde
B_{(-b)}(\tilde{x},{(\delta^{\prime})}^{-1}\tilde{h})$ keeps
looking like a cap for $t\in [-b,\tilde{\tilde{t}})$. On the other
hand, by the definition, the surgery is always done along the
middle two-sphere of a $\delta$-neck with $\delta <
\widetilde{\delta}^{\alpha\beta}$. Then for $\beta$ large, all the
points in $\widetilde
B_{(-b)}(\tilde{x},{(\delta^{\prime})}^{-1}\tilde{h})$ are removed
(as a part of a capped horn) at the time $\tilde{\tilde{t}}$. But
$\tilde{x}$ (near the tip of the cap) exists past the time
$\tilde{\tilde{t}}$. This is a contradiction. Hence we have proved
that
$\widetilde{B}_{(-b)}(\tilde{x},{(\delta^{\prime}})^{-1}\tilde{h})$
is defined on the time interval
$[-b,0]\cap[-b,-b+(1-\tilde{\theta})\tilde{h}^{2}]$.

The $\delta'$-closeness of the solution on $\widetilde{B}_{(-b)}
(\tilde{x},{(\delta^{\prime})}^{-1}h)\times ([-b,0]\cap[-b,
-b+(1-\tilde{\theta})\tilde{h}^{2}])$ with the corresponding
subset of the standard solution follows from Lemma 7.3.6. Then we
have proved Claim 1.

\smallskip
We next verify the following

\medskip
{\bf Claim 2.} \ There is $\tilde{\theta}=\tilde{\theta}(CB)$,
$0<\tilde{\theta}<1$, such that
$b\leq(1-\tilde{\theta})\tilde{h}^{2}$ when $\beta$ large.

\medskip
Note from Proposition 7.3.3, there is a universal constant $D'>0$
such that the standard solution satisfies the following curvature
estimate
$$
R(y,s) \geq \frac{2D'}{1-s}.
$$
We choose $\tilde{\theta}= {D'}/{2(D'+CB)}$. Then for $\beta$
large enough, the rescaled solution satisfies \begin{equation}
\widetilde{R}(x,t)\geq
\frac{D'}{1-(t+b)\tilde{h}^{-2}}\tilde{h}^{-2} %\eqno (7.4.15)
\end{equation} on
$\widetilde{B}_{(-b)}(\tilde{x},{(\delta^{\prime})}^{-1}\tilde{h})\times
([-b,0]\cap[-b,-b+(1-\tilde{\theta})\tilde{h}^{2}])$.

Suppose $b\geq (1-\tilde{\theta})\tilde{h}^{2}$. Then by combining
with the assumption $\widetilde{R}(\tilde{x},t)\leq {C}$ for
$t=(1-\tilde{\theta})\tilde{h}^{2}-b$, we have
$$
C\geq \frac{D'}{1-(t+b)\tilde{h}^{-2}}\tilde{h}^{-2},
$$
and then
$$
1\geq(1-\tilde{\theta})\left(1+\frac{D'}{CB}\right).
$$
This is a contradiction. Hence we have proved Claim 2.

The combination of the above two claims shows that there is a
positive constant $0 <\tilde{\theta}=\tilde{\theta}(CB)<1$ such
that for any small $\delta'>0$, there is a positive
$\beta(\delta',\varepsilon, \tilde{\theta})$ such that when $\beta
\geq \beta(\delta',\varepsilon, \tilde{\theta})$, we have $b\leq
(1-\tilde{\theta})\tilde{h}^{2}$ and the rescaled solution in the
ball
$\widetilde{B}_{(-b)}(\tilde{x},{(\delta^{\prime}})^{-1}\tilde{h})$
on the time interval $[-b,0]$ is, after scaling with factor
$\tilde{h}^{-2}$, ${\delta}^{\prime}$-close ( in the
$C^{[(\delta')^{-1}]}$ topology) to the corresponding subset of
the standard solution.

By (7.4.15) and the assumption $\widetilde{R}\leq {C}$ on
$\widetilde{B}_{0}(\bar{x},A) \times [-b,0],$ we know that the
cut-off radius $\tilde{h}$ at the time $-b$ for the rescaled
solution satisfies
$$\tilde{h} \geq \sqrt{\frac{D'}{C}}.$$

Let $\delta'>0$ be much smaller than $\varepsilon$ and
$\min\{A^{-1},A\}$. Since $\tilde{d}_0(\tilde{x},\bar{x})\leq A$,
it follows that there is constant $C(\tilde{\theta})$ depending
only on $\tilde{\theta}$ such that
$\tilde{d}_{(-b)}(\tilde{x},\bar{x})\leq C(\tilde{\theta})A \ll
(\delta')^{-1}\tilde{h}$. We now apply Lemma 7.3.5 with the
accuracy parameter ${\varepsilon}/{2}$. Let $C({\varepsilon}/{2})$
be the positive constant in Lemma 7.3.5. Without loss of
generality, we may assume the positive constant $C_1(\varepsilon)$
in the canonical neighborhood assumption is larger than
$4C({\varepsilon}/{2})$.  When $\delta'(>0)$ is much smaller than
$\varepsilon$ and $\min\{A^{-1},A\}$, the point $\bar{x}$ at the
time $\bar{t}$ has a neighborhood which is either a
$\frac{3}{4}\varepsilon$-cap or a $\frac{3}{4}\varepsilon$-neck.

Since the canonical neighborhood assumption with accuracy
parameter $\varepsilon$ is violated at $(\bar{x},\bar{t})$, the
neighborhood of the point $\bar{x}$ at the new time zero for the
rescaled solution must be a $\frac{3}{4}\varepsilon$-neck. By
Lemma 7.3.5 (b), we know the neighborhood is the slice at the time
zero of the parabolic neighborhood
$$
P(\bar{x},0,\frac{4}{3}\varepsilon^{-1}\widetilde{R}
(\bar{x},0)^{-\frac{1}{2}},
-\min\{\widetilde{R}(\bar{x},0)^{-1},b\})
$$
(with $\widetilde{R}(\bar{x},0)=1$) which is
$\frac{3}{4}\varepsilon$-close (in the
$C^{[\frac{4}{3}\varepsilon^{-1}]}$ topology) to the corresponding
subset of the evolving standard cylinder $\mathbb{S}^2 \times
\mathbb{R}$ over the time interval $[-\min \{b,1\},0]$ with scalar
curvature $1$ at the time zero. If $b \geq 1$, the
$\frac{3}{4}\varepsilon$-neck is strong, which is a contradiction.
While if $b<1$, the $\frac{3}{4}\varepsilon$-neck at time $-b$ is
contained in the union of the gluing cap and the adjacent
$\delta$-neck where the $\delta$-cutoff surgery took place. Since
$\varepsilon$ is small (say $\varepsilon< 1/100$), it is clear
that the point $\bar{x}$ at time $-b$ is the center of an
$\varepsilon$-neck which is entirely contained in the adjacent
$\delta$-neck. By the proof of Lemma 7.3.2, the adjacent
$\delta$-neck approximates an ancient $\kappa$-solution. This
implies the point $\bar{x}$ at the time $\bar{t}$ has a strong
$\varepsilon$-neck, which is also a contradiction.

Hence we have proved that there exists $\beta_0=\beta_0(
\varepsilon, A, B, C)$ such that when $\beta\geq\beta_0$, the
rescaled solution on the ball $\widetilde{B}_{0}(\bar{x},A)$ can
be defined before the time $-b$.

Let $[t_{A}^{\alpha\beta},0]\supset[-b,0]$ be the largest time
interval so that the rescaled solution
$\widetilde{g}^{\alpha\beta}_{ij}$ can be defined on
$\widetilde{B}_{0}(\bar{x},A)\times[t_{A}^{\alpha\beta},0]$. We
finally claim that $t_{A}^{\alpha\beta}\le
-b-\frac{1}{8}\eta^{-1}C^{-1}$ for $\beta$ large enough.

Indeed, suppose not, by the gradient estimates as in Step 1, we
have the curvature estimate
$$
\widetilde{R}(x,t)\leq {10C}
$$
on $\widetilde{B}_{0}(\bar{x},A)\times [t_{A}^{\alpha\beta},-b]$.
Hence we have the curvature estimate
$$
\widetilde{R}(x,t)\leq {10C}
$$
on $\widetilde{B}_{0}(\bar{x},A)\times [t_{A}^{\alpha\beta},0]$.
By the above argument there is a $\beta_0=\beta_0( \varepsilon, A,
B + \frac{1}{8}\eta^{-1}C^{-1}, 10C)$ such that for
$\beta\geq\beta_0$, the solution in the ball
$\widetilde{B}_{0}(\bar{x},A)$ can be defined before the time
$t_{A}^{\alpha\beta}$. This is a contradiction.

Therefore we have proved Assertion 1.

\medskip
{\bf Assertion 2.} \ For arbitrarily fixed $\alpha$,
$0<A<+\infty$, $1 \leq C<+\infty$ and $0 < B
<\frac{1}{2}\varepsilon^2 (r^{\alpha})^{-2}-
\frac{1}{50}\eta^{-1}$, there is a $\beta_0=\beta_0(\varepsilon,
A,B,C)$ (independent of $\alpha$) such that if $\beta\geq\beta_0$
and the rescaled solution $\widetilde{g}^{\alpha\beta}_{ij}$ on
the ball $\widetilde{B}_{0}(\bar{x},A)$ is defined on a time
interval $[-b+\epsilon',0]$ with $0 < b \leq B$ and $0<\epsilon'
<\frac{1}{50}\eta^{-1}$ and the scalar curvature satisfies
$$
\widetilde{R}(x,t)\leq {C} \ \ \mbox{ on } \
\widetilde{B}_{0}(\bar{x},A)\times [-b+\epsilon',0],
$$
 and there is a point $y\in \widetilde{B}_{0}(\bar{x},A)$ such
that $\widetilde{R}(y,-b+\epsilon')\leq \frac{3}{2}$, then the
rescaled solution $\widetilde{g}^{\alpha\beta}_{ij}$ at $y$ is
also defined on the extended time interval
$[-b-\frac{1}{50}\eta^{-1},0]$ and satisfies the estimate
$$
\widetilde{R}(y,t)\leq 15
$$
for $t\in[-b-\frac{1}{50}\eta^{-1},-b+\epsilon']$.

\medskip
{\bf \em Proof of Assertion} {\bf 2.} \ We imitate the proof of
Assertion 1. If the rescaled solution
$\widetilde{g}^{\alpha\beta}_{ij}$ at $y$ cannot be defined for
some time in $[-b-\frac{1}{50}\eta^{-1},-b+\epsilon')$, then there
is a surgery at some time $\tilde{\tilde{t}} \in
[-b-\frac{1}{50}\eta^{-1},-b+\epsilon']$ such that $y$ lies in the
instant gluing cap.  Let $\tilde{h}$
$(=R(\bar{x},\bar{t})^{\frac{1}{2}}h$) be the cutoff radius at the
time $\tilde{\tilde{t}}$ for the rescaled solution. Clearly, there
is a universal constant $D>1$ such that
${D}^{-1}\tilde{h}\leq\widetilde{R}(y,\tilde{\tilde{t}})^{-\frac{1}{2}}
\leq D\tilde{h}$. By the gradient estimates as in Step 1, the
cutoff radius satisfies
$$
\tilde{h} \geq D^{-1}15^{-\frac{1}{2}}.
$$

As in Claim 1 (i) in the proof of Assertion 1, for any small
constants $0<\tilde{\theta}<\frac{1}{2}$, $\delta'>0$, there
exists a $\beta(\delta',\varepsilon, \tilde{\theta})> 0$ such that
for $\beta \geq \beta(\delta',\varepsilon, \tilde{\theta})$, there
is no surgery interfering in
$\widetilde{B}_{\tilde{\tilde{t}}}(y,(\delta')^{-1}\tilde{h})\times
([\tilde{\tilde{t}},(1-\tilde{\theta})\tilde{h}^{2}
+\tilde{\tilde{t}}]\cap(\tilde{\tilde{t}},0])$. Without loss of
generality, we may assume that the universal constant $\eta$ is
much larger than $D$. Then we have
$(1-\tilde{\theta})\tilde{h}^{2}
+\tilde{\tilde{t}}>-b+\frac{1}{50}\eta^{-1}$. As in Claim 2 in the
proof of Assertion 1, we can use the curvature bound assumption to
choose $\tilde{\theta}=\tilde{\theta}(B,C)$ such that
$(1-\tilde{\theta})\tilde{h}^{2} +\tilde{\tilde{t}}\geq 0$;
otherwise
$$
C\geq\frac{D'}{\tilde{\theta}\tilde{h}^{2}}
$$
for some universal constant $D'>1$, and
$$
|\tilde{\tilde{t}}+b|\leq \frac{1}{50}\eta^{-1},
$$
which implies
$$
1\geq(1-\tilde{\theta})\left(1+\frac{D'}{C\left(B+\frac{1}{50}\eta^{-1}\right)}\right).
$$
This is a contradiction if we choose
$\tilde{\theta}={D'}/{2(D'+C(B+\frac{1}{50}\eta^{-1}))}$.

So there is a positive constant $0
<\tilde{\theta}=\tilde{\theta}(B,C)<1$ such that for any
$\delta'>0$, there is a positive $\beta(\delta',\varepsilon,
\tilde{\theta})$ such that when $\beta \geq
\beta(\delta',\varepsilon, \tilde{\theta})$, we have
$-\tilde{\tilde{t}}\leq (1-\tilde{\theta})\tilde{h}^{2}$  and the
solution in the ball $\widetilde{B}_{\tilde{\tilde{t}}}(\tilde{x},
{(\delta^{\prime})}^{-1}\tilde{h})$ on the time interval
$[\tilde{\tilde{t}},0]$ is, after scaling with factor
$\tilde{h}^{-2}$, ${\delta}^{\prime}$-close (in the
$C^{[\delta'^{-1}]}$ topology) to the corresponding subset of the
standard solution.

Then exactly as in the proof of Assertion 1, by using the
canonical neighborhood structure of the standard solution in Lemma
7.3.5, this gives the desired contradiction with the hypothesis
that the canonical neighborhood assumption with accuracy parameter
$\varepsilon$ is violated at $(\bar{x},\bar{t})$, for $\beta$
sufficiently large.

The curvature estimate at the point $y$ follows from Step 1.
Therefore the proof of Assertion 2 is complete.

Note that the standard solution satisfies $R(x_1,t)\leq D''
R(x_2,t)$ for any $t\in [0,\frac{1}{2}]$ and any two points
$x_1,x_2$, where $D''\geq 1$ is a universal constant.

\medskip
{\bf Assertion 3.} \ For arbitrarily fixed $\alpha$,
$0<A<+\infty$, $1 \leq C<+\infty$, there is a
$\beta_0=\beta_0(\varepsilon, AC^{\frac{1}{2}})$ such that if any
point $(y_0,t_0)$  with  $0 \leq -t_0 <\frac{1}{2}\varepsilon^2
(r^{\alpha})^{-2}- \frac{1}{8}\eta^{-1}C^{-1}$ of the rescaled
solution $\widetilde{g}^{\alpha\beta}_{ij}$ for $\beta\geq\beta_0$
satisfies $\widetilde{R}(y_0,t_0)\leq C$ , then either the
rescaled solution at $y_0$ can be defined at least on
$[t_0-\frac{1}{16}\eta^{-1}C^{-1},t_0]$ and the rescaled scalar
curvature satisfies
$$
\widetilde{R}(y_0,t)\leq 10 {C} \; \mbox{ for } t\in
\Big[t_0-\frac{1}{16}\eta^{-1}C^{-1},t_0\Big],
$$
or we have
$$
\widetilde{R}(x_1,t_0)\leq 2D'' \widetilde{R}(x_2,t_0)
$$
for any two points $x_1, x_2\in \widetilde{B}_{t_0}(y_0,A)$, where
$D''$ is the above universal constant.

\medskip
{\bf \em Proof of Assertion} {\bf 3.} \ Suppose the rescaled
solution $\widetilde{g}^{\alpha\beta}_{ij}$ at $y_0$ cannot be
defined for some $t\in [t_0-\frac{1}{16}\eta^{-1}C^{-1},t_0)$;
then there is a surgery at some time $\tilde{t} \in
[t_0-\frac{1}{16}\eta^{-1}C^{-1},t_0]$ such that $y_0 $ lies in
the instant gluing cap. Let $\tilde{h}$
$(=R(\bar{x},\bar{t})^{\frac{1}{2}}h$) be the cutoff radius at the
time $\tilde{t}$ for the rescaled solution
$\widetilde{g}^{\alpha\beta}_{ij}$.  By the gradient estimates as
in Step 1, the cutoff radius satisfies
$$
\tilde{h} \geq D^{-1}10^{-\frac{1}{2}}C^{-\frac{1}{2}},
$$
where $D$ is the universal constant in the proof of the Assertion
1. Since we assume $\eta$ is suitably larger than $D$ as before,
we have $\frac{1}{2}\tilde{h}^{2} +\tilde{t}>t_0$. As in Claim 1
(ii) in the proof of Assertion 1, for arbitrarily small
$\delta'>0$, we know that for $\beta$ large enough the rescaled
solution on
$\widetilde{B}_{\tilde{t}}(y_0,{(\delta^{\prime})}^{-1}\tilde{h})\times
[\tilde{t},t_0]$ is, after scaling with factor $\tilde{h}^{-2}$,
${\delta}^{\prime}$-close (in the $C^{[(\delta')^{-1}]}$ topology)
to the corresponding subset of the standard solution. Since
$(\delta')^{-1}\tilde{h}\gg A$ for $\beta$ large enough, Assertion
3 follows from the curvature estimate of standard solution in the
time interval $ [0,\frac{1}{2}]$.

\medskip
{\it Step} 3. \ For any subsequence $(\alpha_k,\beta_k)$ of
$(\alpha,\beta)$ with $r^{\alpha_k}\rightarrow 0$ and
$\delta^{\alpha_k\beta_k}\rightarrow 0$ as $k\rightarrow \infty$,
we next argue as in the second step of the proof of Theorem 7.1.1
to show that the curvatures of the rescaled solutions
$\tilde{g}^{\alpha_k\beta_k}_{ij}$ at the new times zero (after
shifting) stay uniformly bounded at bounded distances from
$\bar{x}$ for all sufficiently large $k$. More precisely, we will
prove the following assertion:

\medskip
{\bf Assertion 4.} \ Given any subsequence of the rescaled
solutions $\tilde{g}^{\alpha_k\beta_k}_{ij}$ with
$r^{\alpha_k}\rightarrow 0$ and
$\delta^{\alpha_k\beta_k}\rightarrow 0$ as $k\rightarrow \infty$,
then for any $L>0$, there are constants $C(L)>0$ and $k(L)$ such
that the rescaled solutions $\tilde{g}^{\alpha_k\beta_k}_{ij}$
satisfy
\begin{itemize}
\item[(i)] $\tilde{R}(x,0)\leq C(L)$ for all points $x$ with
$\tilde{d}_{0}(x,\bar{x})\leq L$ and all $k \geq 1$; \item[(ii)]
the rescaled solutions over the ball $\tilde{B}_{0}(\bar{x},L)$
are defined at least on the time interval
$[-\frac{1}{16}\eta^{-1}C(L)^{-1},0]$ for all $k\geq k(L)$.
\end{itemize}

\medskip
{\bf \em Proof of Assertion} {\bf 4.} \ For each $\rho>0$, set
\begin{multline*}
M(\rho)=\sup\Big\{\tilde{R}(x,0)\ |\ k \geq 1\; \mbox{ and }\;
\tilde{d}_0(x,\bar{x})\leq\rho \\
\mbox{in the rescaled solutions } \;
\tilde{g}^{\alpha_k\beta_k}_{ij}\Big\}
\end{multline*}
and
$$
\rho_0=\sup\{\rho > 0\ |\ M(\rho)<+\infty\}.
$$
Note that the estimate (7.4.14) implies that $\rho_0>0$. For (i),
it suffices to prove $\rho_0=+\infty$.

We argue by contradiction. Suppose $\rho_0<+\infty$. Then there is
a sequence of points $y$ in the rescaled solutions
$\tilde{g}^{\alpha_k\beta_k}_{ij}$ with
$\tilde{d}_0(\bar{x},y)\rightarrow\rho_0<+\infty$ and
$\tilde{R}(y,0)\rightarrow +\infty$. Denote by $\gamma$ a
minimizing geodesic segment from $\bar{x}$ to $y$ and denote by
$\tilde{B}_0(\bar{x},\rho_0)$ the open geodesic ball centered at
$\bar{x}$ of radius $\rho_0$ on the rescaled solution
$\tilde{g}^{\alpha_k\beta_k}_{ij}$.

First, we claim that for any $0<\rho<\rho_0$ with $\rho$ near
$\rho_0$, the rescaled solutions on the balls
$\tilde{B}_{0}(\bar{x},\rho)$ are defined on the time interval
$[-\frac{1}{16}\eta^{-1}M(\rho)^{-1},0]$ for all large $k$.
Indeed, this follows from Assertion 3 or Assertion 1. For the
later purpose in Step 6, we now present an argument by using
Assertion 3.  If the claim is not true, then there is a surgery at
some time $\tilde{t} \in [-\frac{1}{16}\eta^{-1}M(\rho)^{-1},0]$
such that some point $\tilde{y}\in \tilde{B}_{0}(\bar{x},\rho) $
lies in the instant gluing cap.  We can choose sufficiently small
$\delta'>0$ such that $2\rho_0<(\delta')^{-\frac{1}{2}}\tilde{h}$,
where $\tilde{h} \geq
D^{-1}20^{-\frac{1}{2}}M(\rho)^{-\frac{1}{2}}$ is the cutoff
radius of the rescaled solutions at $\tilde{t}$. By applying
Assertion 3 with $(\tilde{y},0)=(y_0,t_0)$, we see that there is a
$k(\rho_0,M(\rho))>0$ such that when $k\geq k(\rho_0,M(\rho))$,
$$
\widetilde{R}(x,0)\leq 2D''
$$
for all $x\in \widetilde{B}_{0}(\bar{x},\rho)$. This is a
contradiction as $\rho\rightarrow\rho_0$.

Since for each fixed $0<\rho<\rho_0$ with $\rho$ near $\rho_0$,
the rescaled solutions are defined on
$\tilde{B}_{0}(\bar{x},\rho)\times
[-\frac{1}{16}\eta^{-1}M(\rho)^{-1},0]$ for all large $k$, by Step
1 and Shi's derivative estimate, we know that the covariant
derivatives and higher order derivatives of the curvatures on
$\tilde{B}_{0}(\bar{x},\rho -
\frac{(\rho_0-\rho)}{2})\times[-\frac{1}{32}\eta^{-1}M(\rho)^{-1},0]$
are also uniformly bounded.

By the uniform $\kappa$-noncollapsing property and Hamilton's
compactness theorem (Theorem 4.1.5), after passing to a
subsequence, we can assume that the marked sequence
$(\tilde{B}_0(\bar{x},\rho_0),\widetilde{g}^{\alpha_k\beta_k}_{ij},
\bar{x})$ converges in the $C^{\infty}_{loc}$ topology to a marked
(noncomplete) manifold
($B_{\infty},\widetilde{g}^{\infty}_{ij},\bar{x})$ and the
geodesic segments $\gamma$ converge to a geodesic segment (missing
an endpoint) $\gamma_{\infty}\subset B_{\infty}$ emanating from
$\bar{x}$.

Clearly, the limit has nonnegative sectional curvature by the
pinching assumption. Consider a tubular neighborhood along
$\gamma_{\infty}$ defined by
$$
V=\bigcup_{q_0\in\gamma_{\infty}}B_{\infty}
(q_0,4\pi(\widetilde{R}_{\infty}(q_0))^{-\frac{1}{2}}),
$$
where $\widetilde{R}_{\infty}$ denotes the scalar curvature of the
limit and
$$
B_{\infty}(q_0,4\pi(\widetilde{R}_{\infty}(q_0))^{-\frac{1}{2}})
$$
is the ball centered at $q_0\in B_{\infty}$ with the radius
$4\pi(\widetilde{R}_{\infty}(q_0))^{-\frac{1}{2}}$. Let
$\bar{B}_{\infty}$ denote the completion of
$(B_{\infty},\widetilde{g}^{\infty}_{ij})$, and
$y_{\infty}\in\bar{B}_{\infty}$ the limit point of
$\gamma_{\infty}$. Exactly as in the second step of the proof of
Theorem 7.1.1, it follows from the canonical neighborhood
assumption with accuracy parameter $2\varepsilon$ that the
limiting metric $\widetilde{g}^{\infty}_{ij}$ is cylindrical at
any point $q_0\in\gamma_{\infty}$ which is sufficiently close to
$y_{\infty}$ and then the metric space $\bar{V}=V\cup
\{y_{\infty}\}$ by adding the point $y_{\infty}$ has nonnegative
curvature in the Alexandrov sense. Consequently we have a
three-dimensional non-flat tangent cone $C_{y_{\infty}}\bar{V}$ at
$y_{\infty}$ which is a metric cone with aperture $\leq
20\varepsilon$.

On the other hand, note that by the canonical neighborhood
assumption, the canonical $2\varepsilon$-neck neighborhoods are
strong. Thus at each point $q\in V$ near $y_{\infty}$, the
limiting metric $\widetilde{g}^{\infty}_{ij}$ actually exists on
the whole parabolic neighborhood
$$
V \bigcap
P\left(q,0,\frac{1}{3}\eta^{-1}(\widetilde{R}_{\infty}(q))^{-\frac{1}{2}},
-\frac{1}{10}\eta^{-1}(\widetilde{R}_{\infty}(q))^{-1}\right),
$$
and is a smooth solution of the Ricci flow there. Pick $z\in
C_{y_{\infty}}\bar{V}$ with distance one from the vertex
$y_{\infty}$ and it is nonflat around $z$. By definition the ball
$B(z,\frac{1}{2})\subset C_{y_{\infty}}\bar{V}$ is the
Gromov-Hausdorff convergent limit of the scalings of a sequence of
balls $B_{\infty}(z_{\ell},\sigma_{\ell})
(\subset(V,\widetilde{g}^{\infty}_{ij}))$ where
$\sigma_{\ell}\rightarrow0$. Since the estimate (7.4.14) survives
on $(V,\widetilde{g}^{\infty}_{ij})$ for all $A < +\infty$, and
the tangent cone is three-dimensional and nonflat around $z$, we
see that this convergence is actually in the $C^{\infty}_{\rm
loc}$ topology and over some ancient time interval. Since the
limiting $B_{\infty}(z,\frac{1}{2})(\subset
C_{y_{\infty}}\bar{V})$ is a piece of nonnegatively curved nonflat
metric cone, we get a contradiction with Hamilton's strong maximum
principle (Theorem 2.2.1) as before. So we have proved
$\rho_0=\infty$. This proves (i).

By the same proof of Assertion 1 in Step 2, we can further show
that for any $L$, the rescaled solutions on the balls
$\tilde{B}_{0}(\bar{x},L)$ are defined at least on the time
interval $[-\frac{1}{16}\eta^{-1}C(L)^{-1},0]$ for all
sufficiently large $k$. This proves (ii).

\medskip
{\it Step} 4. \ For any subsequence $(\alpha_k,\beta_k)$ of
$(\alpha,\beta)$ with $r^{\alpha_k}\rightarrow 0$ and
$\widetilde{\delta}^{\alpha_k\beta_k}\rightarrow 0$ as
$k\rightarrow \infty$, by Step 3, the $\kappa$-noncollapsing
property and Hamilton's compactness theorem, we can extract a
$C^{\infty}_{loc}$ convergent subsequence of
$\tilde{g}^{\alpha_k\beta_k}_{ij}$ over some space-time open
subsets containing the slice $\{t=0\}$. We now want to show
\textbf{any} such limit has bounded curvature at $t=0$. We prove
by contradiction. Suppose not, then there is a sequence of points
$z_{\ell}$ divergent to infinity in the limiting metric at time
zero with curvature divergent to infinity. Since the curvature at
$z_{\ell}$ is large (comparable to one), $z_{\ell}$ has a
canonical neighborhood which is a $2\varepsilon$-cap or strong
$2\varepsilon$-neck. Note that the boundary of $2\varepsilon$-cap
lies in some $2\varepsilon$-neck. So we get a sequence of
$2\varepsilon$-necks with radius going to zero. Note also that the
limit has nonnegative sectional curvature. Without loss of
generality, we may assume $2\varepsilon < \varepsilon_0$, where
$\varepsilon_0$ is the positive constant in Proposition 6.1.1.
Thus this arrives at a contradiction with Proposition 6.1.1.

\medskip
{\it Step} 5. \ In this step, we will choose some subsequence
$(\alpha_k,\beta_k)$ of $(\alpha,\beta)$ so that we can extract a
complete smooth limit of the rescaled solutions
$\widetilde{g}^{\alpha_k \beta_k}_{ij}$ to the Ricci flow with
surgery on a time interval $[-a,0]$ for some $a>0$.

Choose $\alpha_k,\beta_k\rightarrow\infty$ so that
$r^{\alpha_k}\rightarrow0$,
$\widetilde{\delta}^{\alpha_{k}\beta_{k}}\rightarrow 0$, and
Assertion 1, 2, 3 hold with $\alpha=\alpha_k, \beta=\beta_k$ for
all $A \in \{{p}/{q} \ |\ p, q=1, 2, \ldots, k\}$, and $B,C \in
\{1,2,\ldots,k\}$. By Step 3, we may assume the rescaled solutions
$\widetilde{g}^{\alpha_k \beta_k}_{ij}$ converge in the
$C^{\infty}_{\rm loc}$ topology at the time $t=0$. Since the
curvature of the limit at $t=0$ is bounded by Step 4, it follows
from Assertion 1 in Step 2 and the choice of the sequence
$(\alpha_k,\beta_k)$ that the limiting
$(M_{\infty},\widetilde{g}^{\infty}_{ij}(\cdot,t))$ is defined at
least on a backward time interval $[-a,0]$ for some positive
constant $a$ and is a smooth solution to the Ricci flow there.

\medskip
{\it Step} 6. \ We further want to extend the limit in Step 5
backwards in time to infinity to get an ancient $\kappa$-solution.
Let $\widetilde{g}^{\alpha_k \beta_k}_{ij}$ be the convergent
sequence obtained in the above Step 5.

Denote by
\begin{eqnarray*}
t_{\max} = \sup \Big\{\ t'&  |&  \mbox {we can take a
smooth limit on } (-t',0] \mbox{ (with bounded}  \\
& &   \mbox{curvature at each time slice) from a subsequence
of}\\
&   &   \mbox{the rescaled solutions
}\widetilde{g}^{\alpha_k\beta_k}_{ij} \Big\}.
\end{eqnarray*}
We first claim that there is a subsequence of the rescaled
solutions $\widetilde{g}^{\alpha_k\beta_k}_{ij}$ which converges
in the $C^{\infty}_{\rm loc}$ topology to a smooth limit
$(M_{\infty},\widetilde{g}^{\infty}_{ij}(\cdot,t))$ on the maximal
time interval $(-t_{\max},0]$.

Indeed, let $t_{\ell}$ be a sequence of positive numbers such that
$t_{\ell}\rightarrow t_{\max}$ and there exist smooth limits
$(M_{\infty},\widetilde{g}^{\infty}_{{\ell}}(\cdot,t))$ defined on
$(-t_{\ell},0]$. For each ${\ell}$, the limit has nonnegative
sectional curvature and has bounded curvature at each time slice.
Moreover by the gradient estimate in canonical neighborhood
assumption with accuracy parameter $2\varepsilon$, the limit has
bounded curvature on each subinterval
$[-b,0]\subset(-t_{\ell},0]$. Denote by $\widetilde{Q}$ the scalar
curvature upper bound of the limit at time zero ($\widetilde{Q}$
is independent of ${\ell}$). Then we can apply Li-Yau-Hamilton
inequality (Corollary 2.5.5) to get
$$
 \widetilde{R}^{\infty}_{{\ell}}(x,t)\leq
 \frac{t_{\ell}}{t+t_{\ell}}\widetilde{Q},
$$
where $\widetilde{R}^{\infty}_{{\ell}}(x,t)$ are the scalar
curvatures of the limits
$(M_{\infty},\widetilde{g}^{\infty}_{{\ell}}(\cdot,t))$. Hence by
the definition of convergence and the above curvature estimates,
we can find a subsequence of the rescaled solutions
$\widetilde{g}^{\alpha_k\beta_k}_{ij}$ which converges in the
$C^{\infty}_{loc}$ topology to a smooth limit
$(M_{\infty},\widetilde{g}^{\infty}_{ij}(\cdot,t))$ on the maximal
time interval $(-t_{\max},0]$.

We need to show $-t_{\max}=-\infty$. Suppose $-t_{\max}>-\infty$,
there are only the following two possibilities: either
\begin{itemize}
\item[(1)] The curvature of the limiting solution
$(M_{\infty},\widetilde{g}^{\infty}_{ij}(\cdot,t))$ becomes
unbounded as $t\searrow -t_{\max}$; or \item[(2)] For each small
constant $\theta>0$ and each large integer $k_0>0$, there is some
$k\geq k_0$ such that the rescaled solution
$\widetilde{g}^{\alpha_k\beta_k}_{ij}$ has a surgery time
$T_k\in[-t_{\max}-\theta,0]$ and a surgery point $x_k$ lying in a
gluing cap at the times $T_k$ so that $d^2_{T_k}(x_k,\bar{x})$ is
uniformly bounded from above by a constant independent of $\theta$
and $k_0$.
\end{itemize}

We next claim that the possibility (1) always occurs. Suppose not;
then the curvature of the limiting solution
$(M_{\infty},\widetilde{g}^{\infty}_{ij}(\cdot,t))$ is bounded on
$M_{\infty}\times(-t_{\max},0]$ by some positive constant
$\hat{C}$. In particular, for any $A>0$, there is a sufficiently
large integer $k_1>0$ such that any rescaled solution
$\widetilde{g}^{\alpha_k\beta_k}_{ij}$ with $k\geq k_1$ on the
geodesic ball $\widetilde{B}_{0}(\bar{x},A)$ is defined on the
time interval $[-t_{\max}+\frac{1}{50}\eta^{-1}{\hat{C}}^{-1},0]$
and its scalar curvature is bounded by $2\hat{C}$ there. (Here,
without loss of generality, we may assume that the upper bound
$\hat{C}$ is so large that
$-t_{\max}+\frac{1}{50}\eta^{-1}{\hat{C}}^{-1} < 0$.) By Assertion
1 in Step 2, for $k$ large enough, the rescaled solution
$\widetilde{g}^{\alpha_k \beta_k}_{ij}$ over
$\widetilde{B}_{0}(\bar{x},A)$ can be defined on the extended time
interval $[-t_{\max}-\frac{1}{50}\eta^{-1}{\hat{C}}^{-1},0]$ and
has the scalar curvature $\widetilde{R}\leq 10 \hat{C}$ on
$\widetilde{B}_{0}(\bar{x},A)\times
[-t_{\max}-\frac{1}{50}\eta^{-1}{\hat{C}}^{-1},0]$. So we can
extract a smooth limit from the sequence to get the limiting
solution which is defined on a larger time interval
$[-t_{\max}-\frac{1}{50}\eta^{-1}{\hat{C}}^{-1},0]$. This
contradicts the definition of the maximal time $-t_{\max}$.

It remains to exclude the possibility (1).

By using Li-Yau-Hamilton inequality (Corollary 2.5.5) again, we
have
$$
\widetilde{R}_{\infty}(x,t)\leq
\frac{t_{\max}}{t+t_{\max}}\widetilde{Q}.
$$
So we only need to control the curvature near $-t_{\max}$. Exactly
as in Step 4 in the proof of Theorem 7.1.1, it follows from
Li-Yau-Hamilton inequality that \begin{equation} \tilde{d}_0(x,y)\leq
\tilde{d}_t(x,y)\leq \tilde{d}_0(x,y)
+30t_{\max}\sqrt{\widetilde{Q}}
%\eqno (7.4.16)
\end{equation} for any $x,y\in M_{\infty}$ and $t\in(-t_{\max},0]$.

Since the infimum of the scalar curvature is nondecreasing in
time, we have some point $y_{\infty}\in M_{\infty}$  and some time
$-t_{\max} < t_{\infty} < -t_{\max}+\frac{1}{50}\eta^{-1}$ such
that $\widetilde{R}_{\infty}(y_{\infty},t_{\infty})<{5}/{4}$. By
(7.4.16), there is a constant $\widetilde{A}_0>0$ such that
$\tilde{d}_t(\bar{x},y_{\infty})\leq \widetilde{A}_0/2$ for all $t
\in (-t_{\max},0]$.

Now we come back to the rescaled solution
$\widetilde{g}^{\alpha_k\beta_k}_{ij}$. Clearly, for arbitrarily
given small $\epsilon' > 0$, when $k$ large enough, there is a
point $y_k$ in the underlying manifold of
$\widetilde{g}^{\alpha_k\beta_k}_{ij}$ at time $0$ satisfying the
following properties \begin{equation}
\widetilde{R}(y_k,t_{\infty})<\frac{3}{2},\qquad
\widetilde{d}_{t}(\bar{x},y_k)\leq \widetilde{A}_0 %\eqno (7.4.17)
\end{equation} for $t\in [-t_{\max} + \epsilon',0]$. By the definition of
convergence, we know that for any fixed $A_0 \geq
2\widetilde{A}_0$, for $k$ large enough, the rescaled solution
over $\widetilde{B}_{0}(\bar{x},A_0)$ is defined on the time
interval $[t_{\infty},0]$ and satisfies
$$
\widetilde{R}(x,t)\leq \frac{2t_{\max}}{t+t_{\max}}\widetilde{Q}
$$
on $\widetilde{B}_{0}(\bar{x},A_0)\times [t_{\infty},0]$. Then by
Assertion 2 of Step 2, we have proved that there is a sufficiently
large integer $\bar{k}_0$ such that when $k \geq \bar{k}_0$, the
rescaled solutions $\widetilde{g}^{\alpha_k\beta_k}_{ij}$ at $y_k$
can be defined on $[-t_{\max}-\frac{1}{50}\eta^{-1},0]$, and
satisfy
$$
\widetilde{R}(y_k,t)\leq 15
$$
for $t\in[-t_{\max}-\frac{1}{50}\eta^{-1},t_{\infty}]$.

We now prove a statement analogous to Assertion 4 (i) of Step 3.

\medskip
{\bf Assertion 5.} \ For the above rescaled solutions
$\widetilde{g}^{\alpha_k\beta_k}_{ij}$ and $\bar{k}_0$, we have
that for any $L>0$, there is a positive constant $\omega(L)$ such
that the rescaled solutions $\widetilde{g}^{\alpha_k\beta_k}_{ij}$
satisfy
$$
\widetilde{R}(x,t)\leq\omega(L)
$$
for all $(x,t)$ with $\tilde{d}_{t}(x,y_k)\leq L$ and $t\in
[-t_{\max}-\frac{1}{50}\eta^{-1},t_{\infty}]$, and for all $k \geq
\bar{k}_0$.

\medskip
{\bf \em Proof of Assertion} {\bf 5.} \ We slightly modify the
argument in the proof of Assertion 4 (i). Let

\begin{eqnarray*}
M(\rho)=\sup\Big\{\widetilde{R}(x,t) &  |&
\tilde{d}_t(x,y_k)\leq\rho\; \mbox{ and }\; t \in
[-t_{\max}-\frac{1}{50}\eta^{-1},t_{\infty}]\\
&   &   \mbox{in the rescaled solutions } \;
\widetilde{g}^{\alpha_k\beta_k}_{ij}, k\geq \bar{k}_0 \Big\}
\end{eqnarray*}
and
$$
\rho_0=\sup\{\rho > 0\ |\  M(\rho)<+\infty\}.
$$
Note that the estimate (7.4.14) implies that $\rho_0>0$. We only
need to show $\rho_0 = +\infty$.

We argue by contradiction. Suppose $\rho_0<+\infty$. Then, after
passing to a subsequence, there is a sequence $(\tilde{y}_k,t_k)$
in the rescaled solutions $\widetilde{g}^{\alpha_k\beta_k}_{ij}$
with $t_k \in [-t_{\max}-\frac{1}{50}\eta^{-1},t_{\infty}]$ and
$\tilde{d}_{t_k}(y_k,\tilde{y}_k)\rightarrow\rho_0<+\infty$ such
that $\widetilde{R}(\tilde{y}_k,t_k)\rightarrow +\infty$. Denote
by $\gamma_k$ a minimizing geodesic segment from $y_k$ to
$\tilde{y}_k$ at the time $t_k$ and denote by
$\widetilde{B}_{t_k}(y_k,\rho_0)$ the open geodesic ball centered
at $y_k$ of radius $\rho_0$ on the rescaled solution
$\widetilde{g}^{\alpha_k\beta_k}_{ij}(\cdot,t_k)$.

For any $0<\rho<\rho_0$ with $\rho$ near $\rho_0$, by applying
Assertion 3 as before, we get that  the rescaled solutions on the
balls $\widetilde{B}_{t_k}(y_k,\rho)$ are defined on the time
interval $[t_k-\frac{1}{16}\eta^{-1}M(\rho)^{-1},t_k]$ for all
large $k$. By Step 1 and Shi's derivative estimate, we further
know that the covariant derivatives and higher order derivatives
of the curvatures on
$\widetilde{B}_{t_k}(y_k,\rho-\frac{(\rho_0-\rho)}{2})
\times[t_k-\frac{1}{32}\eta^{-1}M(\rho)^{-1},t_k]$ are also
uniformly bounded. Then by the uniform $\kappa$-noncollapsing
property and Hamilton's compactness theorem (Theorem 4.1.5), after
passing to a subsequence, we can assume that the marked sequence
$(\tilde{B}_{t_k}(y_k,\rho_0),\widetilde{g}^{\alpha_k\beta_k}_{ij}
(\cdot,t_k),y_k)$ converges in the $C^{\infty}_{loc}$ topology to
a marked (noncomplete) manifold
($B_{\infty},\widetilde{g}^{\infty}_{ij},y_{\infty})$ and the
geodesic segments $\gamma_k$ converge to a geodesic segment
(missing an endpoint) $\gamma_{\infty}\subset B_{\infty}$
emanating from $y_{\infty}$.

Clearly, the limit also has nonnegative sectional curvature by the
pinching assumption. Then by repeating the same argument as in the
proof of Assertion 4 (i) in the rest, we derive a contradiction
with Hamilton's strong maximum principle. This proves Assertion 5.

We then apply the second estimate of (7.4.17) and Assertion 5 to
conclude that for any large constant $0<A<+\infty$, there is a
positive constant $C(A)$ such that for any small $\epsilon'>0$,
the rescaled solutions $\widetilde{g}^{\alpha_k\beta_k}_{ij}$
satisfy \begin{equation}
\widetilde{R}(x,t)\leq C(A), %\eqno (7.4.18)
\end{equation} for all $x \in \widetilde{B}_{0}(\bar{x},A)$ and $t\in
[-t_{\max} + \epsilon',0]$, and for all sufficiently large $k$.
Then by applying Assertion 1 in Step 2, we conclude that the
rescaled solutions $\widetilde{g}^{\alpha_k\beta_k}_{ij}$ on the
geodesic balls $\widetilde{B}_{0}(\bar{x},A)$ are also defined on
the extended time interval $ [-t_{\max} +
\epsilon'-\frac{1}{8}\eta^{-1}C(A)^{-1},0]$ for all sufficiently
large $k$. Furthermore, by the gradient estimates as in Step 1, we
have
$$
\widetilde{R}(x,t)\leq 10C(A),
$$
for $x \in \widetilde{B}_{0}(\bar{x},A)$ and $t\in [-t_{\max} +
\epsilon'-\frac{1}{8}\eta^{-1}C(A)^{-1},0]$. Since $\epsilon'>0$
is arbitrarily small and the positive constant $C(A)$ is
independent of $\epsilon'$, we conclude that the rescaled
solutions $\widetilde{g}^{\alpha_k\beta_k}_{ij}$ on
$\widetilde{B}_{0}(\bar{x},A)$ are defined on the extended time
interval $ [-t_{\max} -\frac{1}{16}\eta^{-1}C(A)^{-1},0]$ and
satisfy \begin{equation}
\widetilde{R}(x,t)\leq 10C(A), %\eqno (7.4.19)
\end{equation} for $x \in \widetilde{B}_{0}(\bar{x},A)$ and $t\in [-t_{\max}
-\frac{1}{16}\eta^{-1}C(A)^{-1},0]$, and for all sufficiently
large $k$.

Now, by taking convergent subsequences from the rescaled solutions
$\widetilde{g}^{\alpha_k\beta_k}_{ij}$, we see that the limit
solution is defined smoothly on a space-time open subset of
$M_{\infty}\times (-\infty,0]$ containing $M_{\infty}\times
[-t_{\max},0]$. By Step 4, we see that the limiting metric
$\widetilde{g}^{\infty}_{ij}(\cdot,-t_{\max})$ at time $-t_{\max}$
has bounded curvature. Then by combining with the canonical
neighborhood assumption of accuracy $2\varepsilon$, we conclude
that the curvature of the limit is uniformly bounded on the time
interval $[-t_{\max},0]$. So we have excluded the possibility (1).

Hence we have proved a subsequence of the rescaled solutions
converges to an orientable ancient $\kappa$-solution.

Finally by combining with the canonical neighborhood theorem
(Theorem 6.4.6), we see that $(\bar{x},\bar{t})$ has a canonical
neighborhood with parameter $\varepsilon$, which is a
contradiction. Therefore we have completed the proof of the
proposition.
%$$\eqno \#$$ \vskip 0.3cm
\endproof

Summing up, we have proved that for any $\varepsilon>0$, (without
loss of generality, we may assume $\varepsilon \leq
\bar{\varepsilon}_0$), there exist nonincreasing (continuous)
positive functions $\widetilde{\delta}(t)$ and $\widetilde{r}(t)$,
defined on $[0,+\infty)$ with
$$
\widetilde{\delta}(t) \leq \bar{\delta}(t) = \min
\left\{\frac{1}{{2e^2\log(1+t)}}, \delta_0\right\},
$$
such that for arbitrarily given (continuous) positive function
$\delta(t)$ with $\delta(t)<\widetilde{\delta}(t)$ on
$[0,+\infty)$, and arbitrarily given a compact orientable
normalized three-manifold as initial data, the Ricci flow with
surgery has a solution on $[0,T)$ obtained by evolving the Ricci
flow and by performing $\delta$-cutoff surgeries at a sequence of
times $0<t_1 < t_2 < \cdots < t_i < \cdots<T$, with ${\delta}(t_i)
\leq \delta \leq \widetilde{\delta}(t_i)$ at each time $t_i$, so
that the pinching assumption and the canonical neighborhood
assumption (with accuracy $\varepsilon$) with $r=\widetilde{r}(t)$
are satisfied.  (At this moment we still do not know whether the
surgery times $t_i$ are discrete.)

Since the $\delta$-cutoff surgeries occur at the points lying
deeply in the $\varepsilon$-horns, the minimum of the scalar
curvature $R_{min}(t)$ of the solution to the Ricci flow with
surgery at each time-slice is achieved in the region unaffected by
the surgeries. Thus we know from the evolution equation of the
scalar curvature that \begin{equation}
\frac{d}{dt}R_{\min}(t) \geq \frac{2}{3}R^2_{\min}(t).%\eqno (7.4.20)
\end{equation} In particular, the minimum of the scalar curvature
$R_{\min}(t)$ is nondecreasing in time. Also note that each
$\delta$-cutoff surgery decreases volume. Then the upper
derivative of the volume in time satisfies
\begin{align*}
\bar{\left(\frac{d}{dt}\right)}V(t) &\triangleq \lim\sup_{\triangle
t\rightarrow
0}\frac{V(t+\triangle t)-V(t)}{\triangle t}\\
&\leq-R_{\min}(0)V(t)
\end{align*}
which implies that
$$
V(t)\leq V(0)e^{-R_{\min}(0)t}.
$$

On the other hand, by Lemma 7.3.2 and the $\delta$-cutoff procedure
given in the previous section, we know that at each time $t_i$, each
$\delta$-cutoff surgery cuts down the volume at least at an amount
of $h^3(t_i)$ with $h(t_i)$ depending only on ${\delta}(t_i)$ and
$\widetilde{r}(t_i)$. Thus the surgery times $t_i$ cannot accumulate
in any finite interval. When the solution becomes extinct at some
finite time $T$, the solution at time near $T$ is entirely covered
by canonical neighborhoods and then the initial manifold is
diffeomorphic to a connected sum of a finite copies of
$\mathbb{S}^2\times \mathbb{S}^1$ and $\mathbb{S}^3/\Gamma$ (the
metric quotients of round three-sphere). So we have proved the
following long-time existence result which was given by Perelman in
\cite{P2}.

%\vskip 0.2cm \noindent{\bf Theorem 7.4.3} (Long-time existence
%theorem)
%CZ-SOURCE-LINE-20634
\begin{theorem}[long-time existence of Ricci flow with surgery]
  \label{thm:cz-ricci-flow-surgery-long-time-existence}
  \dcref{Cao--Zhu Theorem 7.4.3}
  \group{Canonical-Neighborhood Induction}
  \source{cao-zhu}

For any fixed constant $\varepsilon >0$, there exist nonincreasing
$($continuous$)$ positive functions $\widetilde{\delta}(t)$ and
$\widetilde{r}(t)$, defined on $[0,+\infty)$, such that for an
arbitrarily given $($continuous$)$ positive function $\delta(t)$
with $\delta(t) \leq \widetilde{\delta}(t)$ on $[0,+\infty)$, and
arbitrarily given a compact orientable normalized three-manifold
as initial data, the Ricci flow with surgery has a solution with
the following properties: either
\begin{itemize}
\item[(i)] it is defined on a finite interval $[0,T)$ and obtained
by evolving the Ricci flow and by performing a finite number of
cutoff surgeries, with each $\delta$-cutoff at a time $t \in
(0,T)$ having ${\delta}= \delta(t)$, so that the solution becomes
extinct at the finite time $T$, and the initial manifold is
diffeomorphic to a connected sum of a finite copies of
$\mathbb{S}^2\times \mathbb{S}^1$ and $\mathbb{S}^3/\Gamma$ $($the
metric quotients of round three-sphere$)$ ; or \item[(ii)] it is
defined on $[0,+\infty)$ and obtained by evolving the Ricci flow
and by performing at most countably many cutoff surgeries, with
each $\delta$-cutoff at a time $t \in [0,+\infty)$ having
${\delta}=\delta(t)$, so that the pinching assumption and the
canonical neighborhood assumption $($with accuracy $\varepsilon)$
with $r=\widetilde{r}(t)$ are satisfied, and there exist at most a
finite number of surgeries on every finite time interval.
%$$\eqno \#$$\vskip 0.3cm
\end{itemize}
\end{theorem}

In particular, if the initial manifold has positive scalar
curvature, say $R\geq a>0$, then by (7.4.20), the solution becomes
extinct at $T\leq \frac{3}{2}a$. Hence we have the following
topological description of compact three-manifolds with
nonnegative scalar curvature which improves the well-known work of
Schoen-Yau \cite{ScY79m}, \cite{ScY79}.

%\vskip 0.2cm \noindent{\bf Corollary 7.4.4} (Perelman \cite{P2})\
%CZ-SOURCE-LINE-20671
\begin{corollary}[finite extinction topological classification]
  \label{cor:cz-finite-extinction-topological-classification}
  \dcref{Cao--Zhu Corollary 7.4.4}
  \group{Canonical-Neighborhood Induction}
  \source{cao-zhu}

 Let $M$ be a compact orientable three-manifold with nonnegative
scalar curvature. Then either $M$ is flat or it is diffeomorphic
to a connected sum of a finite copies of $\mathbb{S}^2\times
\mathbb{S}^1$ and $\mathbb{S}^3/\Gamma$ $($the metric quotients of
the round three-sphere$).$
%$$\eqno \#$$ \vskip 0.3cm
\end{corollary}

The famous \textbf{Poincar\'{e} conjecture}\index{Poincar\'{e}
conjecture} states that every compact three-manifold with trivial
fundamental group is diffeomorphic to $\mathbb{S}^3$. Developing
tools to attack the conjecture formed the basis for much of the
works in three-dimensional topology over the last one hundred
years. Now we use the Ricci flow to discuss the Poincar$\acute{e}$
conjecture.

Let $M$ be a compact three-manifold with trivial fundamental
group. In particular, the three-manifold $M$ is orientable.
Arbitrarily given a Riemannian metric on $M$, by scaling we may
assume the metric is normalized. With this normalized metric as
initial data, we consider the solution to the Ricci flow with
surgery. If one can show the solution becomes extinct in finite
time, it will follow from Theorem 7.4.3 (i) that the
three-manifold $M$ is diffeomorphic to the three-sphere
$\mathbb{S}^3$. Such a finite extinction time result was first
proposed by Perelman in \cite{P3}. Recently, Colding-Minicozzi has
published a proof of it in \cite{CM}. So \textbf{the combination
of Theorem 7.4.3 (i) and the finite extinction result \cite{P3,
CM} gives a proof of the Poincar\'{e} conjecture}.

We also remark that the above long-time existence result of
Perelman has been extended to compact four-manifolds with positive
isotropic curvature by Chen and the second author in \cite{CZ05F}.
As a consequence it gave a complete proof of the following
classification theorem of compact four-manifolds, with no
essential incompressible space-form and with a metric of positive
isotropic curvature. The theorem was first proved by Hamilton in
(\cite{Ha97}), though it was later found that the proof contains
some gaps (see for example the comment of Perelman in Page 1, the
second paragraph, of \cite{P2}).

%\vskip 0.2cm \noindent{\bf Theorem 7.4.5} \emph{
%CZ-SOURCE-LINE-20714
\begin{theorem}[four-dimensional positive-isotropic-curvature classification]
  \label{thm:cz-positive-isotropic-curvature-four-manifold-classification}
  \dcref{Cao--Zhu Theorem 7.4.5}
  \group{Canonical-Neighborhood Induction}
  \source{cao-zhu}

A compact four-manifold with no essential incompressible
space-form and with a metric of positive isotropic curvature is
diffeomorphic to $\mathbb{S}^4$, or $\mathbb{RP}^4$, or
$\mathbb{S}^3\times \mathbb{S}^1$, or
$\mathbb{S}^3\widetilde{\times}\mathbb{S}^1$ $($the $\mathbb{Z}_2$
quotient of $\mathbb{S}^{3}\times \mathbb{S}^{1}$ where
$\mathbb{Z}_{2}$ flips $\mathbb{S}^{3}$ antipodally and rotates
$\mathbb{S}^{1}$ by $180^{0}),$ or a connected sum of them.
%$$\eqno \#$$ \vskip 1cm
\end{theorem}


\section{Curvature Estimates for Surgically Modified Solutions}

This section is a detailed exposition of section 6 of Perelman
\cite{P2}. Here we will generalize the curvature estimates for
smooth solutions in Section 7.2 to that of solutions with cutoff
surgeries. We first state and prove a version of Theorem 7.2.1.

%\vskip 0.2cm \noindent{\bf Theorem 7.5.1} (Perelman \cite{P2})
%CZ-SOURCE-LINE-20735
\begin{theorem}[curvature estimate near surgery at bounded times]
  \label{thm:cz-bounded-time-surgery-curvature-estimate}
  \dcref{Cao--Zhu Theorem 7.5.1}
  \group{Curvature Estimates after Surgery}
  \source{cao-zhu}
  \uses{thm:cz-ricci-flow-surgery-long-time-existence}

For any $\varepsilon > 0$ and  $1 \leq A<+\infty$, one can find
$\kappa=\kappa(A,\varepsilon)>0$,
$K_1=K_1(A,\varepsilon)<+\infty$, $K_2=K_2(A,\varepsilon)<+\infty$
and $\bar{r}=\bar{r}(A,\varepsilon)>0$ such that for any
$t_0<+\infty$ there exists $\bar{\delta}_A=\bar{\delta}_A(t_0)>0$
$($depending also on $\varepsilon),$ nonincreasing in $t_0$, with
the following property. Suppose we have a solution, constructed by
Theorem $7.4.3$ with the nonincreasing (continuous) positive
functions $\widetilde{\delta}(t)$ and $\widetilde{r}(t)$, to the
Ricci flow with $\delta$-cutoff surgeries on time interval $[0,T]$
and with a compact orientable normalized three-manifold as initial
data, where each $\delta$-cutoff at a time $t$ satisfies
$\delta=\delta(t) \leq \widetilde{\delta}(t)$ on $[0,T]$ and
$\delta=\delta(t) \leq \bar{\delta}_A$ on $[\frac{t_0}{2},t_0]$;
assume that the solution is defined on the whole parabolic
neighborhood $P(x_0,t_0,r_0,-r^2_0)\triangleq\{(x,t)\ |\ x\in
B_t(x_0,r_0),t\in[t_0-r^2_0,t_0]\}$, $2r^2_0<t_0$, and satisfies
$$
|Rm|\leq r^{-2}_0 \ \ on \ \ P(x_0,t_0,r_0,-r^2_0),
$$
$$
\text{and} \qquad \Vol_{t_0}(B_{t_0}(x_0,r_0))\geq A^{-1}r^3_0.
$$
Then
\begin{itemize}
\item[(i)] the solution is $\kappa$-noncollapsed on all scales
less than $r_0$ in the ball $B_{t_0}(x_0,Ar_0)$; \item[(ii)] every
point $x\in B_{t_0}(x_0,Ar_0)$ with $R(x,t_0)\geq K_1r^{-2}_0$ has
a canonical neighborhood $B$, with $ B_{t_0}(x,\sigma)\subset
B\subset B_{t_0}(x,2\sigma)$ for some
$0<\sigma<C_1(\varepsilon)R^{-\frac{1}{2}}(x,t_0), $ which is
either a strong $\varepsilon$-neck or an $\varepsilon$-cap;
\item[(iii)] if $r_0\leq \bar{r}\sqrt{t_0}$ then $R\leq
K_2r^{-2}_0$ in $B_{t_0}(x_0,Ar_0)$.
\end{itemize}
Here $C_1(\varepsilon)$ is the positive constant in the canonical
neighborhood assumption.

\end{theorem}

%\vskip 0.1cm \noindent{\bf Proof.}
\begin{proof}
  \uses{lem:cz-admissible-minimizer-reduced-length-bound}
  \uses{lem:cz-surgery-noncollapsing-canonical-induction-step, thm:cz-almost-euclidean-volume-curvature-control, prop:cz-canonical-neighborhood-induction-surgery}
Without loss of generality, we may assume $0<\varepsilon \leq
\bar{\varepsilon}_0$, where $\bar{\varepsilon}_0$ is the
sufficiently small (universal) positive constant in Lemma 7.4.2.

\medskip
(i) This is analog of no local collapsing theorem II (Theorem
3.4.2). In comparison with the no local collapsing theorem II,
this statement gives $\kappa$-noncollapsing property no matter how
big the time is and it also allows the solution to be modified by
surgery.

Let $\eta (\geq 10)$ be the universal constant in the definition
of the canonical neighborhood assumption. Recall that we had
removed every component which has positive sectional curvature in
our surgery procedure. By the same argument as in the first part
of the proof of Lemma 7.4.2, the canonical neighborhood assumption
of the solution implies the $\kappa$-noncollapsing on the scales
less than $\frac{1}{2\eta}\widetilde{r}(t_0)$ for some positive
constant $\kappa$ depending only on $C_1(\varepsilon)$ and
$C_2(\varepsilon)$ (in the definition of the canonical
neighborhood assumption). So we may assume
$\frac{1}{2\eta}\widetilde{r}(t_0)\leq r_0\leq
\sqrt{\frac{t_0}{2}}$, and study the scales $\rho$,
$\frac{1}{2\eta}\widetilde{r}(t_0)\leq \rho\leq r_0$. Let $x\in
B_{t_0}(x_0,Ar_0)$ and assume that the solution satisfies
$$
|Rm|\leq \rho^{-2}
$$
for those points in $P(x,t_0,\rho,-\rho^{2})\triangleq\{(y,t)\ |\
y\in B_t(x,\rho),t\in[t_0-\rho^2,t_0]\}$ for which the solution is
defined. We want to bound the ratio
$\Vol_{t_0}(B_{t_0}(x,\rho))/\rho^3$ from below.

Recall that a space-time curve is called admissible if it stays in
the region unaffected by surgery, and a space-time curve on the
boundary of the set of admissible curves is called a barely
admissible curve. Consider any barely admissible curve $\gamma$,
parametrized by $t\in[t_{\gamma},t_0]$, $t_0-r^2_0\leq
t_{\gamma}\leq t_0$, with $\gamma(t_0)=x$. The same proof for the
assertion (7.4.2) (in the proof of Lemma 7.4.2) shows that for
arbitrarily large $L>0$ (to be determined later), one can find a
sufficiently small
$\bar{\delta}(L,t_0,\widetilde{r}(t_0),\widetilde{r}
(\frac{t_0}{2}),\varepsilon)>0$ such that when each
$\delta$-cutoff in $[\frac{t_0}{2},t_0]$ satisfies $\delta \leq
\bar{\delta}(L,t_0,\widetilde{r}(t_0),\widetilde{r}
(\frac{t_0}{2}),\varepsilon)$, there holds \begin{equation}
\int^{t_0}_{t_{\gamma}}\sqrt{t_0-t}(R_+(\gamma(t),t)
+|\dot{\gamma}(t)|^2)dt\geq Lr_0. %\eqno (7.5.1)
\end{equation}

{}From now on, we assume that each $\delta$-cutoff of the solution
in the time interval $[\frac{t_0}{2},t_0]$ satisfies $\delta \leq
\bar{\delta}(L,t_0,\widetilde{r}(t_0),
\widetilde{r}(\frac{t_0}{2}),\varepsilon).$

Let us scale the solution, still denoted by $g_{ij}(\cdot,t)$, to
make $r_0=1$ and the time as $t_0=1$. By the maximum principle, it
is easy to see that the (rescaled) scalar curvature satisfies
$$
R \geq -\frac{3}{2t}
$$
on $(0,1]$. Let us consider the time interval $[\frac{1}{2},1]$
and define a function of the form
$$
h(y,t)=\phi(d_t(x_0,y)-A(2t-1))(\bar{L}(y,\tau)+2\sqrt{\tau})
$$
where $\tau=1-t$, $\phi$ is the function of one variable chosen in
the proof of Theorem 3.4.2 which is equal to one on
$(-\infty,\frac{1}{20})$, rapidly increasing to infinity on
$(\frac{1}{20},\frac{1}{10})$, and satisfies
$2\frac{(\phi')^2}{\phi}-\phi''\geq(2A+300)\phi'-C(A)\phi$ for
some constant $C(A)<+\infty$, and $\bar{L}$ is the function
defined by
\begin{align*}
\bar{L}(q,\tau) &=\inf\Big\{2\sqrt{\tau}\int^{\tau}_0\sqrt{s}
(R+|\dot{\gamma}|^2)ds\ |\ (\gamma(s),s),
s\in[0,\tau]\\
&\text{is a space-time curve with $\gamma(0)=x$ and
$\gamma(\tau)=q$}\Big\}.
\end{align*}
Note that
\begin{align}
\bar{L}(y,\tau)
& \geq 2\sqrt{\tau}\int^{\tau}_0\sqrt{s}Rds\\
& \geq -4\tau^2 \nonumber\\
&  >   -2\sqrt{\tau} \nonumber% \eqno (7.5.2)
\end{align}
since $R\geq -3$ and $0<\tau\leq \frac{1}{2}$. This says $h$ is
positive for $t\in[\frac{1}{2},1]$. Also note that \begin{equation}
\frac{\partial}{\partial \tau}\bar{L}+\triangle \bar{L}\leq 6
%\eqno (7.5.3)
\end{equation} as long as $\bar{L}$ is achieved by admissible curves. Then as
long as the shortest $\mathcal{L}$-geodesics from $(x_0,0)$ to
$(y,\tau)$ are admissible, there holds at $y$ and $t=1-\tau$,
\begin{align*}
\left(\frac{\partial}{\partial t}-\triangle\right)h &\geq
\left(\phi'\left[\left(\frac{\partial}{\partial t}
-\triangle\right)d_t-2A\right]-\phi''\right)\cdot(\bar{L}+2\sqrt{\tau})\\
&\quad-\left(6+\frac{1}{\sqrt{\tau}}\right)\phi-2\langle\nabla
\phi,\nabla\bar{L}\rangle.
\end{align*}

Firstly,\; we\; may\; assume\; the\; constant\; $L$\; in (7.5.1)\;
is\; not\; less\; than $2\exp(C(A)+100)$. We claim that Lemma
3.4.1(i) is applicable for $d=d_t(\cdot,x_0)$ at $y$ and $
t=1-\tau $ (with $\tau\in[0,\frac{1}{2}]$) whenever
$\bar{L}(y,\tau)$ is achieved by admissible curves and satisfies
the estimate
$$
\bar{L}(y,\tau) \leq 3\sqrt{\tau}\exp(C(A)+100).
$$
Indeed, since the solution is defined on the whole neighborhood
$P(x_0,t_0,r_0,$ $-r^2_0)$ with $r_0=1$ and $t_0=1$, the point
$x_0$ at the time $t=1-\tau$ lies on the region unaffected by
surgery. Note that $R \geq -3$ for $t \in [\frac{1}{2},1]$.  When
$\bar{L}(y,\tau)$ is achieved by admissible curves and satisfies
$\bar{L}(y,\tau) \leq 3\sqrt{\tau}\exp(C(A)+100)$, the estimate
(7.5.1) implies that the point $y$ at the time $t=1-\tau$ does not
lie in the collars of the gluing caps. Thus any minimal geodesic
(with respect to the metric $g_{ij}(\cdot,t)$ with $t=1-\tau$)
connecting $x_0$ and $y$ also lies in the region unaffected by
surgery; otherwise the geodesic is not minimal.  Then from the
proof of Lemma 3.4.1(i), we see that it is applicable.

Assuming the minimum of $h$ at a time, say $t=1-\tau$, is achieved
at a point, say $y$, and assuming $\bar{L}(y,\tau)$ is achieved by
admissible curves and satisfies $\bar{L}(y,\tau) \leq
3\sqrt{\tau}\exp(C(A)+100)$, we have
$$
(\bar{L}+2\sqrt{\tau})\nabla\phi=-\phi\nabla\bar{L},
$$
and then by the computations and estimates in the proof of Theorem
3.4.2,
\begin{align*}
&\left(\frac{\partial}{\partial t}-\triangle\right)h  \\
&  \geq \left(\phi'\left[\left(\frac{\partial}{\partial
t}-\triangle\right)d_t-2A\right]-\phi''
+2\frac{(\phi')^2}{\phi}\right)\cdot(\bar{L}+2\sqrt{\tau})
-\left(6+\frac{1}{\sqrt{\tau}}\right)\phi\\
& \geq -C(A)h-\left(6+\frac{1}{\sqrt{\tau}}\right)
\frac{h}{(2\sqrt{\tau}-4\tau^2)},
\end{align*}
at $y$ and $t=1-\tau$. Here we used (7.5.2) and Lemma 3.4.1(i).

As before, denoting by $h_{min}(\tau)=\min_{z} h(z,1-\tau)$, we
obtain
\begin{align}
\frac{d}{d\tau}\left(\log\left(\frac{h_{\min}(\tau)}{\sqrt{\tau}}\right)\right) &
\leq C(A)+\frac{6\sqrt{\tau}+1}{2\tau-4\tau^2\sqrt{\tau}}
-\frac{1}{2\tau}\\
& \leq C(A)+\frac{50}{\sqrt{\tau}}, \nonumber
\end{align} %\eqno (7.5.4)
as long as the associated shortest $\mathcal{L}$-geodesics are
admissible with $\bar{L} \leq 3\sqrt{\tau}\exp(C(A)+100)$. On the
other hand, by definition, we have \begin{equation} \lim_{\tau\rightarrow
0^+}\frac{h_{\min}(\tau)}{\sqrt{\tau}}\leq
\phi(d_1(x_0,x)-A)\cdot 2=2. %\eqno (7.5.5)
\end{equation} The combination of (7.5.4) and (7.5.5) gives the following
assertion:

\vskip 0.2cm \emph{Let $\tau\in[0,\frac{1}{2}]$. If for each
$s\in[0,\tau]$, $\inf\{\bar{L}(y,s)\ |\ \mbox{ } d_t(x_0,y)\leq
A(2t-1)+\frac{1}{10}\mbox{ with }s=1-t\}$ is achieved by
admissible curves, then we have
\begin{align}
&\inf\left\{\bar{L}(y,\tau)\ |\ \mbox{ } d_t(x_0,y)
\leq A(2t-1)+\frac{1}{10}\mbox{ with }\tau=1-t\right\}\\
&\leq 2\sqrt{\tau}\exp(C(A)+100). \nonumber  %\eqno (7.5.6)
\end{align}
}

\vskip 0.1cm Note again that $R \geq -3$ for $t \in
[\frac{1}{2},1]$. By combining with (7.5.1), we know that any
barely admissible curve $\gamma$, parametrized by $s\in[0,\tau]$,
$0\leq \tau \leq \frac{1}{2}$, with $\gamma(0)=x$, satisfies
$$
\int^{\tau}_0\sqrt{s}(R + |\dot{\gamma}|^2)ds \geq
\frac{7}{4}\exp(C(A)+100),
$$
by assuming $L\geq 2\exp(C(A)+100)$.

Since $|Rm|\leq \rho^{-2}$ on $P(x,t_0,\rho,-\rho^{2})$ with
$\rho\geq \frac{1}{2\eta}\widetilde{r}(t_0)$ (and $t_0=1$) and
$\bar{\delta}(L,t_0,\widetilde{r}(t_0),
\widetilde{r}(\frac{t_0}{2}),\varepsilon)>0$ is sufficiently
small, the parabolic neighborhood $P(x,1,\rho,-\rho^2)$ around the
point $(x,1)$ is contained in the region unaffected by the
surgery. Thus as $\tau=1-t$ is sufficiently close to zero,
$\frac{1}{2\sqrt{\tau}}\inf\bar{L}$ can be bounded from above by a
small positive constant and then the infimum
$\inf\{\bar{L}(y,\tau)\ |\ \mbox{ } d_t(x_0,y)\leq
A(2t-1)+\frac{1}{10}\mbox{ with }\tau=1-t\}$ is achieved by
admissible curves.

Hence we conclude that for each $\tau\in [0,\frac{1}{2}]$, any
minimizing curve $\gamma_{\tau}$ of $\inf\{\bar{L}(y,\tau)|\mbox{
} d_t(x_0,y)\leq A(2t-1)+\frac{1}{10}\mbox{ with }\tau=1-t\}$ is
admissible and satisfies
$$
\int^{\tau}_0\sqrt{s}(R+|\dot{\gamma}_{\tau}|^2)ds\leq
\exp(C(A)+100).
$$

Now we come back to the unrescaled solution. It then follows that
the Li-Yau-Perelman distance $l$ from $(x,t_0)$ satisfies the
following estimate \begin{equation} \min\left\{l\left(y,t_0-\frac{1}{2}r^2_0\right)\
\Big|\ y\in
B_{t_0-\frac{1}{2}r^2_0}\left(x_0,\frac{1}{10}r_0\right)\right\}
\leq\exp(C(A)+100), %\eqno (7.5.7)
\end{equation} by noting the (parabolic) scaling invariance of the
Li-Yau-Perelman distance.

By the assumption that $|Rm|\leq r^{-2}_0$ on
$P(x_0,t_0,r_0,-r^2_0)$, exactly as before, for any $q\in
B_{t_0-r^2_0}(x_0,r_0)$, we can choose a path $\gamma$
parametrized by $\tau\in[0,r^2_0]$ with $\gamma(0)=x$,
$\gamma(r^2_0)=q$, and $\gamma(\frac{1}{2}r^2_0)=y\in
B_{t_0-\frac{1}{2}r^2_0}(x_0,\frac{1}{10}r_0)$, where
$\gamma|_{[0,\frac{1}{2}r^2_0]}$ achieves the minimum
$\min\{l(y,t_0-\frac{1}{2}r^2_0)\ |\ y\in
B_{t_0-\frac{1}{2}r^2_0}(x_0,\frac{1}{10}r_0)\}$ and
$\gamma|_{[\frac{1}{2}r^2_0,r^2_0]}$ is a suitable curve
satisfying $\gamma|_{[\frac{1}{2}r^2_0,r^2_0]}(\tau) \in B_{t_0 -
\tau}(x_0,r_0)$, for each $\tau \in [\frac{1}{2}r^2_0,r^2_0]$, so
that the $\mathcal{L}$-length of $\gamma$ is uniformly bounded
from above by a positive constant (depending only on $A$)
multiplying $r_0$. This implies that the Li-Yau-Perelman distance
from $(x,t_0)$ to the ball $B_{t_0-r^2_0}(x_0,r_0)$ is uniformly
bounded by a positive constant $L(A)$ (depending only on $A$). Now
we can choose the constant $L$ in (7.5.1) by
$$
L=\max\{2L(A),2\exp(C(A)+100)\}.
$$
Thus every shortest $\mathcal{L}$-geodesic from $(x,t_0)$ to the
ball $B_{t_0-r^2_0}(x_0,r_0)$ is necessarily admissible. By
combining with the assumption that $\Vol_{t_0}(B_{t_0}(x_0,r_0))$
$\geq\; \;A^{-1}\,r^3_0$,\; we\; conclude\; that\; Perelman's\;
reduced\; volume\; of\; the\; ball $B_{t_0-r^2_0}(x_0,r_0)$
satisfies the estimate
\begin{align}
\widetilde{V}_{r^2_0}(B_{t_0-r^2_0}(x_0,r_0)) &  =
\int_{B_{t_0-r^2_0}(x_0,r_0)}
(4\pi r^2_0)^{-\frac{3}{2}}\exp(-l(q,r^2_0))dV_{t_0-r^2_0}(q)\\
&  \geq c(A) \nonumber
\end{align}  %\eqno (7.5.8)
for some positive constant $c(A)$ depending only on $A$.

We can now argue as in the last part of the proof of Lemma 7.4.2
to get a lower bound estimate for the volume of the ball
$B_{t_0}(x,\rho)$. The union of all shortest
$\mathcal{L}$-geodesics from $(x,t_0)$ to the ball
$B_{t_0-r^2_0}(x_0,r_0)$, defined by
\begin{multline*}
CB_{t_0-r^2_0}(x_0,r_0)=\{(y,t)\ |\ (y,t)
\mbox{ lies in a shortest $\mathcal{L}$-geodesic from}\\
(x,t_0)\mbox{ to a point in }B_{t_0-r^2_0}(x_0,r_0)\},
\end{multline*}
forms a cone-like subset in space-time with vertex $(x,t_0)$.
Denote by $B(t)$ the intersection of the cone-like subset
$CB_{t_0-r^2_0}(x_0,r_0)$ with the time-slice at $t$. Perelman's
reduced volume of the subset $B(t)$ is given by
$$
\widetilde{V}_{t_0-t}(B(t))=\int_{B(t)}(4\pi(t_0-t))^{-\frac{3}{2}}
\exp(-l(q,t_0-t))dV_t(q).
$$
Since the cone-like subset $CB_{t_0-r^2_0}(x_0,r_0)$ lies entirely
in the region unaffected by surgery, we can apply Perelman's
Jacobian comparison Theorem 3.2.7 and the estimate (7.5.8) to
conclude that
\begin{align}
\widetilde{V}_{t_0-t}(B(t))&  \geq
\widetilde{V}_{r^2_0}(B_{t_0-r^2_0}(x_0,r_0))\\
&  \geq c(A) \nonumber
\end{align}         %\eqno (7.5.9)
for all $t\in[t_0-r^2_0,t_0]$.

As before, denoting by
$\xi=\rho^{-1}\Vol_{t_0}(B_{t_0}(x,\rho))^{\frac{1}{3}}$, we only
need to get a positive lower bound for $\xi$. Of course we may
assume $\xi<1$. Consider $\widetilde{B}(t_0-\xi\rho^2)$, the
subset at the time-slice $\{t=t_0-\xi\rho^2\}$ where every point
can be connected to $(x,t_0)$ by an admissible shortest
$\mathcal{L}$-geodesic. Perelman's reduced volume of
$\widetilde{B}(t_0-\xi\rho^2)$ is given by
\begin{align}
& \widetilde{V}_{\xi\rho^2}(\widetilde{B}(t_0-\xi\rho^2))\\
&  = \int_{\widetilde{B}(t_0-\xi\rho^2)}
(4\pi\xi\rho^2)^{-\frac{3}{2}}
\exp(-l(q,\xi\rho^2))dV_{t_0-\xi\rho^2}(q) \nonumber\\
& =\int_{\widetilde{B}(t_0-\xi\rho^2)\cap\mathcal{L}
\exp_{\{|\upsilon|\leq\frac{1}{4}
\xi^{-\frac{1}{2}}\}}(\xi\rho^2)} (4\pi\xi\rho^2)^{-\frac{3}{2}}
\exp(-l(q,\xi\rho^2))dV_{t_0-\xi\rho^2}(q) \nonumber\\
&\quad  +\int_{\widetilde{B}(t_0-\xi\rho^2)
\setminus\mathcal{L}\exp_{\{|\upsilon|\leq\frac{1}{4}
\xi^{-\frac{1}{2}}\}}(\xi\rho^2)} (4\pi\xi\rho^2)^{-\frac{3}{2}}
\exp(-l(q,\xi\rho^2))dV_{t_0-\xi\rho^2}(q). \nonumber
\end{align}  %\eqno (7.5.10)
Note that the whole region $P(x,t_0,\rho,-\rho^2)$ is unaffected
by surgery because $\rho\geq \frac{1}{2\eta}\widetilde{r}(t_0)$
and $\bar{\delta}(L,t_0,\widetilde{r}(t_0),
\widetilde{r}(\frac{t_0}{2}),\varepsilon)>0$ is sufficiently
small. Then exactly as before, there is a universal positive
constant $\xi_0$ such that when $0<\xi\leq\xi_0$, there holds
$$
\mathcal{L}\exp_{\{|\upsilon|\leq\frac{1}{4}
\xi^{-\frac{1}{2}}\}}(\xi\rho^2)\subset B_{t_0}(x,\rho)
$$
and the first term on RHS of (7.5.10) can be estimated by
\begin{align}
&\int_{\widetilde{B}(t_0-\xi\rho^2)\cap\mathcal{L}
\exp_{\{|\upsilon|\leq\frac{1}{4}\xi^{-\frac{1}{2}}\}}
(\xi\rho^2)}(4\pi\xi\rho^2)^{-\frac{3}{2}}
\exp(-l(q,\xi\rho^2))dV_{t_0-\xi\rho^2}(q)\\
&\leq e^{C\xi}(4\pi)^{-\frac{3}{2}}\xi^{\frac{3}{2}}\nonumber  % \eqno(7.5.11)
\end{align}
for some universal constant $C$; while the second term on RHS of
(7.5.10) can be estimated by
\begin{align}
&\int_{\widetilde{B}(t_0-\xi\rho^2)\setminus\mathcal{L}
\exp_{\{|\upsilon|\leq\frac{1}{4}\xi^{-\frac{1}{2}}\}}
(\xi\rho^2)}(4\pi\xi\rho^2)^{-\frac{3}{2}}
\exp(-l(q,\xi\rho^2))dV_{t_0-\xi\rho^2}(q)\\
&\leq(4\pi)^{-\frac{3}{2}}\int_{\{|\upsilon|>\frac{1}{4}
\xi^{-\frac{1}{2}}\}}\exp(-|\upsilon|^2)d\upsilon. \nonumber
%\eqno (7.5.12)
\end{align}
Since $B(t_0-\xi\rho^2)\subset\widetilde{B}(t_0-\xi\rho^2)$, the
combination of (7.5.9)-(7.5.12) bounds $\xi$ from below by a
positive constant depending only on $A$. This proves the statement
(i).

\medskip
(ii) This is analogous to the claim in the proof of Theorem 7.2.1.
We argue by contradiction. Suppose that for some $A<+\infty$ and a
sequence $K^{\alpha}_1\rightarrow\infty$, there exists a sequence
$t^{\alpha}_0$ such that for any sequences
$\bar{\delta}^{\alpha\beta}>0$ with
$\bar{\delta}^{\alpha\beta}\rightarrow0$ for fixed $\alpha$, we
have sequences of solutions $g^{\alpha\beta}_{ij}$ to the Ricci
flow with surgery and sequences of points $x^{\alpha\beta}_0$, of
radii $r^{\alpha\beta}_0$, which satisfy the assumptions but
violate the statement (ii) at some $x^{\alpha\beta}\in
B_{t^{\alpha}_0}(x^{\alpha\beta}_0,Ar^{\alpha\beta}_0)$ with
$R(x^{\alpha\beta},t^{\alpha}_0)\geq
K^{\alpha}_1(r^{\alpha\beta}_0)^{-2}$. Slightly abusing notation,
we will often drop the indices $\alpha,\beta$ in the following
argument.

Exactly as in the proof of Theorem 7.2.1, we need to adjust the
point $(x,t_0)$. More precisely, we claim that there exists a
point $(\bar{x},\bar{t})\!\in\! B_{\bar{t}}(x_0,2Ar_0)$
$\times[t_0-\frac{r^2_0}{2},t_0]$ with $\bar{Q}\triangleq
R(\bar{x},\bar{t})\geq K_1r^{-2}_0$ such that the point
$(\bar{x},\bar{t})$ does not satisfy the canonical neighborhood
statement, but each point $(y,t)\in \bar{P}$ with $R(y,t)\geq
4\bar{Q}$ does, where $\bar{P}$ is the set of all $(x',t')$
satisfying $\bar{t}-\frac{1}{4}K_1\bar{Q}^{-1}\leq t'\leq
\bar{t}$, $d_{t'}(x_0,x')\leq
d_{\bar{t}}(x_0,\bar{x})+K^{\frac{1}{2}}_1\bar{Q}^{-\frac{1}{2}}$.
Indeed as before, the point $(\bar{x},\bar{t})$ is chosen by an
induction argument. We first choose $(x_1,t_1)=(x,t_0)$ which
satisfies $d_{t_1}(x_0,x_1)\leq Ar_0$ and $R(x_1,t_1)\geq
K_1r^{-2}_0$, but does not satisfy the canonical neighborhood
statement. Now if $(x_k,t_k)$ is already chosen and is not the
desired $(\bar{x},\bar{t})$, then some point $(x_{k+1},t_{k+1})$
satisfies $t_k-\frac{1}{4}K_1R(x_k,t_k)^{-1}\leq t_{k+1}\leq t_k$,
$d_{t_{k+1}}(x_0,x_{k+1})\leq
d_{t_k}(x_0,x_k)+K^{\frac{1}{2}}_1R(x_k,t_k)^{-\frac{1}{2}}$, and
$R(x_{k+1},t_{k+1})\geq 4R(x_k,t_k)$, but $(x_{k+1},t_{k+1})$ does
not satisfy the canonical neighborhood statement. Then we have
\begin{align*}
R(x_{k+1},t_{k+1})&\geq 4^kR(x_1,t_1)\geq 4^kK_1r^{-2}_0,\\
d_{t_{k+1}}(x_0,x_{k+1}) &\leq d_{t_1}(x_0,x_1)+K^{\frac{1}{2}}_1
\sum^k_{i=1}R(x_i,t_i)^{-\frac{1}{2}} \leq Ar_0+2r_0,
\end{align*}
and
$$
t_0\geq t_{k+1}\geq
t_0-\frac{1}{4}K_1\sum^k_{i=1}R(x_i,t_i)^{-1}\geq
t_0-\frac{1}{2}r^2_0.
$$
So the sequence must be finite and its last element is the desired
$(\bar{x},\bar{t})$.

Rescale the solutions along $(\bar{x},\bar{t})$ with factor
$R(\bar{x},\bar{t})(\geq K_1r^{-2}_0)$ and shift the times
$\bar{t}$ to zero. We will adapt both the proof of Proposition
7.4.1 and that of Theorem 7.2.1 to show that a sequence of the
rescaled solutions $\widetilde{g}_{ij}^{\alpha\beta}$ converges to
an ancient $\kappa$-solution, which will give the desired
contradiction. Since we only need to consider the scale of the
curvature less than $\widetilde{r}({\bar{t}})^{-2}$, the present
situation is much easier than that of Proposition 7.4.1.

Firstly as before, we need to get a local curvature estimate.

For each adjusted $(\bar x,\bar t)$, let $[t',\bar t]$ be the
maximal subinterval of $[\bar
t-\frac{1}{20}\eta^{-1}\bar{Q}^{-1},\bar t]$\; so\; that\; for\;
each\; sufficiently\; large $\alpha$ and then sufficiently large
$\beta$, the canonical neighborhood statement holds for any
$(y,t)$ in $P(\bar x,\bar t,$ $\frac{1}{10}K_1^{\frac{1}{2}}
\bar{Q}^{-\frac{1}{2}},t'-\bar t)=\{ ({x},{t}) \ |\ {x} \in
B_{{t}}(\bar
x,\frac{1}{10}K_1^{\frac{1}{2}}\bar{Q}^{-\frac{1}{2}}), {t} \in
[t',\bar{t}]\}$ with $R(y,t)\geq 4\bar{Q}$, where $\eta$ is the
universal positive constant in the definition of canonical
neighborhood assumption. We want to show \begin{equation}
t' = \bar t-\frac{1}{20}\eta^{-1}\bar{Q}^{-1}. %\eqno(7.5.13)
\end{equation}

Consider the scalar curvature $R$ at the point $\bar x$ over the
time interval $[t',\bar t]$. If there is a time $\tilde{t} \in
[t',\bar t]$ satisfying $R(\bar x,\tilde{t}) \geq 4\bar Q$, we let
$\tilde{t}$ be the first of such time from $\bar t$. Since the
chosen point $(\bar{x},\bar{t})$ does not satisfy the canonical
neighborhood statement, we know $R(\bar{x},\bar{t})\leq
\widetilde{r}(\bar{t})^{-2}$. Recall from our designed surgery
procedure that if there is a cutoff surgery at a point $x$ at a
time $t$, the scalar curvature at $(x,t)$ is at least
$(\bar{\delta}^{\alpha\beta})^{-2}\widetilde{r}({t})^{-2}$. Then
for each fixed $\alpha$, for $\beta$ large enough, the solution
$g^{\alpha\beta}_{ij}(\cdot,t)$ around the point $\bar x$ over the
time interval $[\tilde{t}-
\frac{1}{20}\eta^{-1}\bar{Q}^{-1},{\tilde{t}}]$ is well defined
and satisfies the following curvature estimate
$$
R(\bar x,t) \leq 8\bar Q,
$$
for $ t \in [\tilde{t}- \frac{1}{20}\eta^{-1}\bar{Q}^{-1},\bar t]$
(or $t \in [t',\bar t]$ if there is no such time $\tilde{t}$). By
the assumption that $t_0>2r^2_0$, we have
\begin{align*}
\bar{t}R(\bar{x},\bar{t})&  \geq
\frac{t_0}{2}R(\bar{x},\bar{t})\\
&  \geq r^2_0(K_1r^{-2}_0)\\
&  = K_1\rightarrow +\infty.
\end{align*}
Thus by using the pinching assumption and the gradient estimates
in the canonical neighborhood assumption, we further have
$$
|Rm(x,t)| \leq 30\bar{Q},
$$
for all $x \in B_t(\bar
x,\frac{1}{10}\eta^{-1}\bar{Q}^{-\frac{1}{2}})$ and $t \in
[\tilde{t}- \frac{1}{20}\eta^{-1}\bar{Q}^{-1},\bar t]$ (or $t \in
[t',\bar t]$) and all sufficiently large $\alpha$ and $\beta$.
Observe that Lemma 3.4.1 (ii) is applicable for $d_t(x_0,\bar{x})$
with $t \in [\tilde{t}- \frac{1}{20}\eta^{-1}\bar{Q}^{-1}, \bar
t]$ (or $t \in [t',\bar t]$) since any minimal geodesic, with
respect to the metric $g_{ij}(\cdot,t)$, connecting $x_0$ and
$\bar{x}$ lies in the region unaffected by surgery; otherwise the
geodesic is not minimal. After having obtained the above curvature
estimate, we can argue as deriving (7.2.2) and (7.2.3) in the
proof of Theorem 7.2.1 to conclude that any point $(x,t)$, with $
\tilde{t}- \frac{1}{20}\eta^{-1}\bar{Q}^{-1} \leq t \leq \bar t$
(or $t \in [t',\bar t]$) and $d_t(x,\bar x) \leq
\frac{1}{10}K_1^{\frac{1}{2}}\bar{Q}^{-\frac{1}{2}}$, satisfies
$$
d_t(x,x_{0}) \leq d_{\bar t}(\bar x,x_{0}) +
\frac{1}{2}K_1^{\frac{1}{2}}\bar{Q}^{-\frac{1}{2}},
$$
for all sufficiently large $\alpha$ and $\beta$. Then by combining
with the choice of the points $(\bar x,\bar t)$, we prove $t'=
\bar t-\frac{1}{20}\eta^{-1}\bar{Q}^{-1}$ (i.e., the canonical
neighborhood statement holds for any point $(y,t)$ in the
parabolic neighborhood $P(\bar x,\bar
t,\frac{1}{10}K_1^{\frac{1}{2}}\bar{Q}^{-\frac{1}{2}}$,$
-\frac{1}{20}\eta^{-1}\bar{Q}^{-1})$ with $R(y,t) \geq 4\bar Q$)
for all sufficiently large $\alpha$ and then sufficiently large
$\beta$.

Now it follows from the gradient estimates in the canonical
neighborhood assumption that the scalar curvatures of the rescaled
solutions $\widetilde{g}^{\alpha\beta}_{ij}$ satisfy
$$
\widetilde{R}(x,t)\leq 40
$$
for those $(x,t)\in P(\bar{x},0,\frac{1}{10}\eta^{-1},
-\frac{1}{20}\eta^{-1})\triangleq\{(x',t')\ |\
x'\in\widetilde{B}_{t'}(\bar{x},\frac{1}{10}\eta^{-1}),
t'\in[-\frac{1}{20}\eta^{-1},0]\}$, for which the rescaled
solution is defined. (Here $\widetilde{B}_{t'}$ denotes the
geodesic ball in the rescaled solution at time $t'$). Note again
that $R(\bar{x},\bar{t})\leq \widetilde{r}(\bar{t})^{-2}$ and
recall from our designed surgery procedure that if there is a
cutoff surgery at a point $x$ at a time $t$, the scalar curvature
at $(x,t)$ is at least
$(\bar{\delta}^{\alpha\beta})^{-2}\widetilde{r}({t})^{-2}$. Then
for each fixed sufficiently large $\alpha$, for $\beta$ large
enough, the rescaled solution $\widetilde{g}^{\alpha\beta}_{ij}$
is defined on the whole parabolic neighborhood
$P(\bar{x},0,\frac{1}{10}\eta^{-1},-\frac{1}{20}\eta^{-1})$. More
generally, for arbitrarily fixed $0<\widetilde{K}<+\infty$, there
is a positive integer $\alpha_0$ so that for each $\alpha \geq
\alpha_0$ we can find $\beta_0 >0$ (depending on $\alpha$) such
that if $\beta\geq \beta_0$ and $(y,0)$ is a point on the rescaled
solution $\widetilde{g}^{\alpha\beta}_{ij}$ with
$\widetilde{R}(y,0)\leq\widetilde{K}$ and
$\widetilde{d}_0(y,\bar{x}) \leq \widetilde{K}$, we have estimate
\begin{equation}
\widetilde{R}(x,t)\leq 40\widetilde{K} %\eqno (7.5.14)
\end{equation} for\; $(x,\,t)\;\in\; P(y,\,0,\,\frac{1}{10}\,\eta^{-1}\,
\widetilde{K}^{-\frac{1}{2}},\;
-\frac{1}{20}\,\eta^{-1}\widetilde{K}^{-1})\;\triangleq\;\{(x',\,t')\
|\ x'\;\in\; \widetilde{B}_{t'}(y,$ $\frac{1}{10}\eta^{-1}
\widetilde{K}^{-\frac{1}{2}}),
t'\in[-\frac{1}{20}\eta^{-1}\widetilde{K}^{-1},0]\}$. In
particular, the rescaled solution is defined on the whole
parabolic neighborhood
$P(y,0,\frac{1}{10}\eta^{-1}\widetilde{K}^{-\frac{1}{2}},
-\frac{1}{20}\eta^{-1}\widetilde{K}^{-1})$.

Next, we want to show the curvature of the rescaled solutions at
the new times zero (after shifting) stay uniformly bounded at
bounded distances from $\bar{x}$ for some subsequences of $\alpha$
and $\beta$. Let $\alpha_m,\beta_m\rightarrow+\infty$ be chosen so
that the estimate (7.5.14) holds with $\widetilde{K}=m$. For all
$\rho \geq 0$, set
$$
M(\rho)=\sup\{\widetilde{R}(x,0)\ |\ m \geq 1, d_0(x,\bar{x})\leq
\rho\mbox{ in the rescaled solutions }
\widetilde{g}^{\alpha_m\beta_m}_{ij}\}
$$
and
$$
\rho_0=\sup\{\rho\geq0\ |\ M(\rho)<+\infty\}.
$$
Clearly the estimate (7.5.14) yields $\rho_0>0$. As we consider
the unshifted time $\bar{t}$, by combining with the assumption
that $t_0>2r^2_0$, we have
\begin{align}
\bar{t}R(\bar{x},\bar{t})&  \geq
\frac{t_0}{2}R(\bar{x},\bar{t})\\
&  \geq r^2_0(K_1r^{-2}_0) \nonumber\\
&  = K_1\rightarrow +\infty. \nonumber
\end{align}  %\eqno (7.5.15)
It then follows from the pinching assumption that we only need to
show $\rho_0=+\infty$. As before, we argue by contradiction.
Suppose we have a sequence of points $y_m$ in the rescaled
solutions $\widetilde{g}^{\alpha_m\beta_m}_{ij}$ with
$\widetilde{d}_0(\bar{x},y_m)\rightarrow \rho_0<+\infty$ and
$\widetilde{R}(y_m,0)\rightarrow+\infty$. Denote by $\gamma_m$ a
minimizing geodesic segment from $\bar{x}$ to $y_m$ and denote by
$\widetilde{B}_0(\bar{x},\rho_0)$ the open geodesic balls centered
at $\bar{x}$ of radius $\rho_0$ of the rescaled solutions. By
applying the assertion in statement (i), we have uniform
$\kappa$-noncollapsing at the points $(\bar{x},\bar{t})$. By
combining with the local curvature estimate (7.5.14) and
Hamilton's compactness theorem, we can assume that, after passing
to a subsequence, the marked sequence
$(\widetilde{B}_0(\bar{x},\rho_0),
\widetilde{g}^{\alpha_m\beta_m}_{ij},\bar{x})$ converges in the
$C^{\infty}_{\rm loc}$ topology to a marked (noncomplete) manifold
($B_{\infty},\widetilde{g}^{\infty}_{ij},x_{\infty}$) and the
geodesic segments $\gamma_m$ converge to a geodesic segment
(missing an endpoint) $\gamma_{\infty}\subset B_{\infty}$
emanating from $x_{\infty}$. Moreover, by the pinching assumption
and the estimate (7.5.15), the limit has nonnegative sectional
curvature.

Then exactly as before, we consider the tubular neighborhood along
$\gamma_{\infty}$
$$
V=\bigcup_{q_0\in\gamma_{\infty}}B_{\infty}
(q_0,4\pi(\widetilde{R}_{\infty}(q_0))^{-\frac{1}{2}})
$$
and the completion $\bar{B}_{\infty}$ of
($B_{\infty},\widetilde{g}^{\infty}_{ij}$) with $y_{\infty}\in
\bar{B}_{\infty}$ the limit point of $\gamma_{\infty}$. As before,
by the choice of the points $(\bar{x},\bar{t})$, we know that the
limiting metric $\widetilde{g}^{\infty}_{ij}$ is cylindrical at
any point $q_0\in \gamma_{\infty}$ which is sufficiently close to
$y_{\infty}$. Then by the same reason as before the metric space
$\bar{V}=V\cup\{y_{\infty}\}$ has nonnegative curvature in
Alexandrov sense, and we have a three-dimensional nonflat tangent
cone $C_{y_{\infty}}\bar{V}$ at $y_{\infty}$. Pick $z\in
C_{y_{\infty}}\bar{V}$ with distance one from the vertex and it is
nonflat around $z$. By definition $B(z,\frac{1}{2}) (\subset
C_{y_{\infty}}\bar{V})$ is the Gromov-Hausdoff convergent limit of
the scalings of a sequence of balls $B_{\infty}(z_k,\sigma_k)
(\subset(V,\widetilde{g}^{\infty}_{ij}))$ with
$\sigma_k\rightarrow0$. Since the estimate (7.5.14) survives on
$(V,\widetilde{g}^{\infty}_{ij})$ for all $\widetilde{K}<
+\infty$, we know that this convergence is actually in the
$C^{\infty}_{\rm loc}$ topology and over some time interval. Since
the limit $B(z,\frac{1}{2}) (\subset C_{y_{\infty}}\bar{V})$ is a
piece of a nonnegatively curved nonflat metric cone, we get a
contradiction with Hamilton's strong maximum Principle (Theorem
2.2.1) as before. Hence we have proved that a subsequence of the
rescaled solution $\widetilde{g}^{\alpha_m\beta_m}_{ij}$ has
uniformly bounded curvatures at bounded distance from $\bar{x}$ at
the new times zero.

Further, by the uniform $\kappa$-noncollapsing at the points
$(\bar{x},\bar{t})$ and the estimate (7.5.14) again, we can take a
$C^{\infty}_{\rm loc}$ limit
($M_{\infty},\widetilde{g}^{\infty}_{ij},x_{\infty}$), defined on
a space-time subset which contains the time slice $\{t=0\}$ and is
relatively open in $M_{\infty}\times(-\infty,0]$, for the
subsequence of the rescaled solutions. The limit is a smooth
solution to the Ricci flow, and is complete at $t=0$, as well as
has nonnegative sectional curvature by the pinching assumption and
the estimate (7.5.15). Thus by repeating the same argument as in
the Step 4 of the proof Proposition 7.4.1, we conclude that the
curvature of the limit at $t=0$ is bounded.

Finally we try to extend the limit backwards in time to get an
ancient $\kappa$-solution. Since the curvature of the limit is
bounded at $t=0$, it follows from the estimate (7.5.14) that the
limit is a smooth solution to the Ricci flow defined at least on a
backward time interval $[-a,0]$ for some positive constant $a$.
Let $(t_{\infty},0]$ be the maximal time interval over which we
can extract a smooth limiting solution. It suffices to show
$t_{\infty}=-\infty$. If $t_{\infty}>-\infty$, there are only two
possibilities: either there exist surgeries in finite distance
around the time $t_{\infty}$ or the curvature of the limiting
solution becomes unbounded as $t\searrow t_{\infty}$.

Let $c>0$ be a positive constant much smaller than
$\frac{1}{100}\eta^{-1}$. Note again that the infimum of the
scalar curvature is nondecreasing in time. Then we can find some
point $y_{\infty} \in M_{\infty}$ and some time $t=t_{\infty} +
\theta$ with $0<\theta<\frac{c}{3}$ such that
$\widetilde{R}_{\infty}(y_{\infty},t_{\infty}+\theta) \leq 2$.

Consider the (unrescaled) scalar curvature $R$ of
${g}_{ij}^{\alpha_m \beta_m}(\cdot,t)$ at the point $\bar x$ over
the time interval $[\bar
t+(t_{\infty}+\frac{\theta}{2})\bar{Q}^{-1},\bar t]$. Since the
scalar curvature $R_{\infty}$ of the limit on $M_{\infty}\times
[t_{\infty}+\frac{\theta}{3},0]$ is uniformly bounded by some
positive constant $C$, we have the curvature estimate
$$
R(\bar x,t) \leq 2C\bar{Q}
$$
for all $t \in [\bar
t+(t_{\infty}+\frac{\theta}{2})\bar{Q}^{-1},\bar t]$ and all
sufficiently large $m$. For each fixed $m$ and $\alpha_m$, we may
require the chosen $\beta_m$ to satisfy
$$
(\bar{\delta}^{\alpha_m\beta_m})^{-2}\left(\widetilde{r}
\left(\frac{t^{\alpha_m}_0}{2}\right)\right)^{-2} \geq
m\widetilde{r}({t^{\alpha_m}_0})^{-2}\geq m\bar{Q}.
$$
When $m$ is large enough, we observe again that Lemma 3.4.1 (ii)
is applicable for $d_t(x_0,\bar{x})$ with $t \in [\bar
t+(t_{\infty}+\frac{\theta}{2})\bar{Q}^{-1},\bar t]$. Then by
repeating the argument as in the derivation of (7.2.1), (7.2.2)
and (7.2.3), we deduce that for all sufficiently large $m$, the
canonical neighborhood statement holds for any $(y,t)$ in the
parabolic neighborhood $P(\bar x,\bar
t,\frac{1}{10}K_1^{\frac{1}{2}}\bar{Q}^{-\frac{1}{2}},
(t_{\infty}+\frac{\theta}{2})\bar{Q}^{-1})$ with $R(y,t) \geq
4\bar{Q}$.

Let $(y_m,\bar t+(t_{\infty}+\theta_m)\bar{Q}^{-1})$ be a sequence
of associated points and times in the (unrescaled) solutions
$g_{ij}^{\alpha_m \beta_m}(\cdot,t)$ so that after rescaling, the
sequence converges to $(y_{\infty},t_{\infty}+\theta)$ in the
limit. Clearly $\frac{\theta}{2}\leq \theta_m \leq 2\theta$ for
all sufficiently large $m$. Then by the argument as in the
derivation of (7.5.13), we know that for all sufficiently large
$m$, the solution $g_{ij}^{\alpha_m \beta_m}(\cdot,t)$ at $y_m$ is
defined on the whole time interval $[\bar
t+(t_{\infty}+\theta_m-\frac{1}{20}\eta^{-1})\bar{Q}^{-1}, \bar
t+(t_{\infty}+\theta_m)\bar{Q}^{-1}]$ and satisfies the curvature
estimate
$$
R(y_m,t) \leq 8\bar{Q}
$$
there; moreover the canonical neighborhood statement holds for any
$(y,t)$ with $R(y,t) \geq 4\bar Q$ in the parabolic neighborhood
$P(y_m,\bar t,\frac{1}{10}K_1^{\frac{1}{2}}\bar{Q}^{-\frac{1}{2}},
(t_{\infty}-\frac{c}{3})\bar{Q}^{-1}).$

We now consider the rescaled sequence
$\widetilde{g}_{ij}^{\alpha_m\beta_m}(\cdot,t)$ with the marked
points replaced by $y_m$ and the times replaced by $s_m \in [\bar
t+(t_{\infty}-\frac{c}{4})\bar{Q}^{-1},
\bar{t}+(t_{\infty}+\frac{c}{4})\bar{Q}^{-1}]$. As before the
Li-Yau-Hamilton inequality implies the rescaling limit around
$(y_m,s_m)$ agrees with the original one. Then the arguments in
previous paragraphs imply the limit is well-defined and smooth on
a space-time open neighborhood of the maximal time slice
$\{t=t_{\infty}\}$. Particularly this excludes the possibility of
existing surgeries in finite distance around the time
$t_{\infty}$. Moreover, the limit at $t=t_{\infty}$ also has
bounded curvature. By using the gradient estimates in the
canonical neighborhood assumption on the parabolic neighborhood
$P(y_m,\bar t,\frac{1}{10}K_1^{\frac{1}{2}}\bar{Q}^{-\frac{1}{2}},
(t_{\infty}-\frac{c}{3})\bar{Q}^{-1})$, we see that the second
possibility is also impossible. Hence we have proved a subsequence
of the rescaled solutions converges to an ancient
$\kappa$-solution.

Therefore we have proved the canonical neighborhood statement
(ii).

\medskip
(iii) This is analogous to Theorem 7.2.1. We also argue by
contradiction. Suppose for some $A<+\infty$ and sequences of
positive numbers $K^{\alpha}_2\rightarrow+\infty$,
$\bar{r}^{\alpha}\rightarrow0$ there exists a sequence of times
$t^{\alpha}_0$ such that for any sequences
$\bar{\delta}^{\alpha\beta}>0$ with
$\bar{\delta}^{\alpha\beta}\rightarrow0$ for fixed $\alpha$, we
have sequences of solutions $g^{\alpha\beta}_{ij}$ to the Ricci
flow with surgery and sequences of points $x^{\alpha\beta}_0$, of
positive constants $r^{\alpha\beta}_0$ with $r^{\alpha\beta}_0\leq
\bar{r}^{\alpha}\sqrt{t^{\alpha}_0}$ which satisfy the
assumptions, but for all $\alpha,\beta$ there hold \begin{equation}
R(x^{\alpha\beta},t^{\alpha}_0)>K^{\alpha}_2(r^{\alpha\beta}_0)^{-2},\
\mbox{ for some } \ x^{\alpha\beta}\in
B_{t^{\alpha}_0}(x^{\alpha\beta}_0,Ar^{\alpha\beta}_0).
%\eqno(7.5.16)
\end{equation}

We may assume that $\bar{\delta}^{\alpha\beta}\leq
\bar{\delta}_{4A}(t^{\alpha}_0)$ for all $\alpha,\beta$, where
$\bar{\delta}_{4A}(t^{\alpha}_0)$ is chosen so that the statements
(i) and (ii) hold on
$B_{t^{\alpha}_0}(x^{\alpha\beta}_0,4Ar^{\alpha\beta}_0)$. Let
$\hat{g}^{\alpha\beta}_{ij}$ be the rescaled solutions of
$g^{\alpha\beta}_{ij}$ around the origins $x^{\alpha\beta}_0$ with
factor $(r^{\alpha\beta}_0)^{-2}$ and shift the times
$t^{\alpha}_0$ to zero. Then by applying the statement (ii), we
know that the regions, where the scalar curvature of the rescaled
solutions $\hat{g}^{\alpha\beta}_{ij}$ is at least
$K_1(=K_1(4A))$, are canonical neighborhood regions. Note that
canonical $\varepsilon$-neck neighborhoods are strong. Also note
that the pinching assumption and the assertion
$$
t^{\alpha}_0(r^{\alpha\beta}_0)^{-2} \geq
(\bar{r}^{\alpha})^{-2}\rightarrow+\infty, \mbox{ as
}\alpha\rightarrow+\infty,
$$
imply that any subsequent limit of the rescaled solutions
$\hat{g}^{\alpha\beta}_{ij}$ must have nonnegative sectional
curvature. Thus by the above argument in the proof of the
statement (ii) (or the argument in Step 2 of the proof of Theorem
7.1.1), we conclude that there exist subsequences
$\alpha=\alpha_m,\beta=\beta_m$ such that the curvatures of the
rescaled solutions $\hat{g}^{\alpha_m\beta_m}_{ij}$ stay uniformly
bounded at distances from the origins $x^{\alpha_m\beta_m}_0$ not
exceeding $2A$. This contradicts (7.5.16) for $m$ sufficiently
large. This proves the statements (iii).

Clearly for fixed $A$, after defining the $\bar{\delta}_A(t_0)$
for each $t_0$, one can adjust the $\bar{\delta}_A(t_0)$ so that
it is nonincreasing in $t_0$.

Therefore we have completed the proof of the theorem.
%$$\eqno \#$$\vskip 0.3cm
\end{proof}

{}From now on we redefine the function $\widetilde{\delta}(t)$  so
that it is also less than $\bar{\delta}_{2(t+1)}(2t)$ and then the
above theorem always holds for $A \in [1,2(t_0+1)]$. Particularly,
we still have
$$
\widetilde{\delta}(t)\leq \bar{\delta}(t) = \min
\{\frac{1}{{2e^2\log(1+t)}}, \delta_0\},
$$
which tends to zero as $t\rightarrow+\infty$. We may also require
that $\widetilde{r}(t)$ tends to zero as $t\rightarrow+\infty$.

The next result is a version of Theorem 7.2.2 for solutions with
surgery.

%\vskip 0.2cm\noindent{\bf Theorem 7.5.2} (Perelman \cite{P2})
%CZ-SOURCE-LINE-21551
\begin{theorem}[large-time curvature estimate for surgically modified flows]
  \label{thm:cz-large-time-surgery-curvature-estimate}
  \dcref{Cao--Zhu Theorem 7.5.2}
  \group{Curvature Estimates after Surgery}
  \source{cao-zhu}
  \uses{thm:cz-ricci-flow-surgery-long-time-existence}

For any $\varepsilon > 0$ and $w>0$, there exist
$\tau=\tau(w,\varepsilon)>0$, $K=K(w,\varepsilon)<+\infty$,
$\bar{r}=\bar{r}(w,\varepsilon)>0$,
$\theta=\theta(w,\varepsilon)>0$ and $T=T(w) < +\infty$ with the
following property. Suppose we have a solution, constructed by
Theorem $7.4.3$ with the nonincreasing $($continuous$)$ positive
functions $\widetilde{\delta}(t)$ and $\widetilde{r}(t)$, to the
Ricci flow with surgery on the time interval $[0,t_0]$ with a
compact orientable normalized three-manifold as initial data,
where each $\delta$-cutoff at a time $t \in [0,t_0]$ has
$\delta=\delta(t) \leq \min
\{\widetilde{\delta}(t),\widetilde{r}(2t)\}$. Let $r_0,t_0$
satisfy $\theta^{-1}h\leq r_0\leq \bar{r}\sqrt{t_0}$ and $t_0 \geq
T$, where $h$ is the maximal cutoff radius for surgeries in
$[\frac{t_0}{2},t_0]$ $($if there is no surgery in the time
interval $[\frac{t_0}{2},t_0]$, we take $h=0),$ and assume that
the solution on the ball $B_{t_0}(x_0,r_0)$ satisfies
$$
Rm(x,t_0) \geq -r^{-2}_0,\; \mbox{ on }\; B_{t_0}(x_0,r_0),
$$
$$
\mbox{and }\quad  \Vol_{t_0}(B_{t_0}(x_0,r_0)) \geq w r^3_0.
$$
Then the solution is well defined and satisfies
$$
R(x,t)<Kr^{-2}_0
$$
in the whole parabolic neighborhood
$$
P\left(x_0,t_0,\frac{r_0}{4},-\tau r^2_0\right) =\left\{(x,t)\ |\ x\in
B_t\left(x_0,\frac{r_0}{4}\right),t\in[t_0-\tau r^2_0,t_0]\right\}.
$$
\end{theorem}

%\vskip 0.1cm\noindent{\bf Proof.}
\begin{proof}
  \uses{thm:cz-bounded-time-surgery-curvature-estimate, thm:cz-local-volume-growth-curvature-estimate, thm:cz-lower-volume-ratio-curvature-control}
We are given that $Rm(x,t_0)\geq -r^{-2}_0$ for $x \in
B_{t_0}(x_0,r_0)$, and $\Vol_{t_0}(B_{t_0}(x_0,r_0))\geq w r^3_0$.
The same argument in the derivation of (7.2.7) and (7.2.8) (by
using the Alexandrov space theory) implies that there exists a
ball $B_{t_0}(x',r')\subset B_{t_0}(x_0,r_0)$ with \begin{equation}
\Vol_{t_0}(B_{t_0}(x',r'))\geq \frac{1}{2}\alpha_3(r')^3
%\eqno(7.5.17)
\end{equation} and with \begin{equation}
r'\geq c(w)r_0 %\eqno (7.5.18)
\end{equation} for some small positive constant $c(w)$ depending only on $w$,
where $\alpha_3$ is the volume of the unit ball in $\mathbb{R}^3$.

As in (7.1.2), we can rewrite the pinching assumption (7.3.3) as
$$
Rm \geq-[f^{-1}(R(1+t))/(R(1+t))]R,
$$
where
$$
y= f(x) = x(\log x - 3), \quad \mbox{ for }\; e^2 \leq x <
+\infty,
$$
is increasing and convex with range $-e^2 \leq y < +\infty$, and
its inverse function is also increasing and satisfies
$$
\lim_{y\rightarrow +\infty} f^{-1}(y)/y = 0.
$$

Note that $t_0 r^{-2}_0 \geq \bar{r}^{-2}$ by the hypotheses. We
may require $T(w) \geq 8c(w)^{-1}$. Then by applying Theorem 7.5.1
(iii) with $A=8c(w)^{-1}$ and combining with the pinching
assumption, we can reduce the proof of the theorem to the special
case $w=\frac{1}{2}\alpha_3$. In the following we simply assume
$w=\frac{1}{2}\alpha_3$.

Let us first consider the case $r_0 < \widetilde{r}(t_0)$. We
claim that $R(x,t_0)\leq C^2_0r^{-2}_0$ on
$B_{t_0}(x_0,\frac{r_0}{3})$, for some sufficiently large positive
constant $C_0$ depending only on $\varepsilon$. If not, then there
is a canonical neighborhood around $(x,t_0)$. Note that the type
(c) canonical neighborhood has already been ruled out by our
design of cutoff surgeries. Thus $(x,t_0)$ belongs to an
$\varepsilon$-neck or an $\varepsilon$-cap. This tells us that
there is a nearby point $y$, with $R(y,t_0)\geq C^{-1}_2R(x,t_0) >
C^{-1}_2C^2_0r^{-2}_0$ and $d_{t_0}(y,x)\leq
C_1R(x,t_0)^{-\frac{1}{2}}\leq C_1C^{-1}_0r_0$, which is the
center of the $\varepsilon$-neck
$B_{t_0}(y,\varepsilon^{-1}R(y,t_0)^{-\frac{1}{2}})$. (Here
$C_1,C_2$ are the positive constants in the definition of
canonical neighborhood assumption). Clearly, when we choose $C_0$
to be much larger than $C_1,C_2$ and $\varepsilon^{-1}$, the whole
$\varepsilon$-neck
$B_{t_0}(y,\varepsilon^{-1}R(y,t_0)^{-\frac{1}{2}})$ is contained
in $B_{t_0}(x_0,\frac{r_0}{2})$ and we have \begin{equation}
\frac{\Vol_{t_0}(B_{t_0}(y,\varepsilon^{-1}
R(y,t_0)^{-\frac{1}{2}}))}{(\varepsilon^{-1}
R(y,t_0)^{-\frac{1}{2}})^3}\leq
8\pi\varepsilon^2. %\eqno (7.5.19)
\end{equation} Without loss of generality, we may assume $\varepsilon>0$ is
very small. Since we have assumed that $Rm\geq -r^{-2}_0$ on
$B_{t_0}(x_0,r_0)$ and $\Vol_{t_0}(B_{t_0}(x_0,r_0))\geq
\frac{1}{2}\alpha_3r^3_0$, we then get a contradiction by applying
the standard Bishop-Gromov volume comparison. Thus we have the
desired curvature estimate $R(x,t_0)\leq C^2_0r^{-2}_0$ on
$B_{t_0}(x_0,\frac{r_0}{3})$.

Furthermore, by using the gradient estimates in the definition of
canonical neighborhood assumption, we can take
$K=10C^2_0,\tau=\frac{1}{100}\eta^{-1}C^{-2}_0$ and
$\theta=\frac{1}{5}C^{-1}_0$ in this case. And since
$r_0\geq\theta^{-1}h$, we have $R<10C^2_0r^{-2}_0\leq
\frac{1}{2}h^{-2}$ and the surgeries do not interfere in
$P(x_0,t_0,\frac{r_0}{4},-\tau r^2_0)$.

We now consider the remaining case $\widetilde{r}(t_0)\leq r_0\leq
\bar{r}\sqrt{t_0}$. Let us redefine
\begin{align*}
\tau&=\min\left\{\frac{\bar{\tau}_0}{2},\frac{1}{100}
\eta^{-1}C^{-2}_0 \right\},\\
K&=\max\left\{2\left(\bar{C}+\frac{2\bar{B}}{\bar{\tau}_0}\right),25C^2_0\right\},
\end{align*}
and
$$
\theta=\frac{1}{2}K^{-\frac{1}{2}}
$$
where $\bar{\tau}_0=\tau_0(w),\bar{B}=B(w)$ and $\bar{C}=C(w)$ are
the positive constants in Theorem 6.3.3(ii) with
$w=\frac{1}{2}\alpha_3$, and $C_0$ is the positive constant chosen
above. We will show there is a sufficiently small $\bar{r}>0$ such
that the conclusion of the theorem for $w=\frac{1}{2}\alpha_3$
holds for the chosen $\tau,K$ and $\theta$.

Argue by contradiction. Suppose not, then there exist a sequence
of $\bar{r}^{\alpha}\rightarrow0$, and a sequence of solutions
$g^{\alpha}_{ij}$ with points $(x^{\alpha}_0,t^{\alpha}_0)$ and
radii $r^{\alpha}_0$ such that the assumptions of the theorem do
hold with $\widetilde{r}(t^{\alpha}_0)\leq r^{\alpha}_0\leq
\bar{r}^{\alpha}\sqrt{t^{\alpha}_0}$ whereas the conclusion does
not. Similarly as in the proof of Theorem 7.2.2, we claim that we
may assume that for all sufficiently large $\alpha$, any other
point $(x^{\alpha},t^{\alpha})$ and radius $r^{\alpha}>0$ with
that property has either $t^{\alpha}> t^{\alpha}_0$ or
$t^{\alpha}= t^{\alpha}_0$ with $r^{\alpha}\geq r^{\alpha}_0$;
moreover $t^{\alpha}$ tends to $+\infty$ as $\alpha \rightarrow
+\infty$. Indeed, for fixed $\alpha$ and the solution
$g^{\alpha}_{ij}$, let $t^{\alpha}_{\min}$ be the infimum of all
possible times $t^{\alpha}$ with some point $x^{\alpha}$ and some
radius $r^{\alpha}$ having that property. Since each such
$t^{\alpha}$ satisfies $\bar{r}^{\alpha}\sqrt{t^{\alpha}}\geq
r^{\alpha}\geq \widetilde{r}(t^{\alpha})$, it follows that when
$\alpha$ is large, $t^{\alpha}_{\min}$ must be positive and very
large. Clearly for each fixed sufficiently large $\alpha$, by
passing to a limit, there exist some point $x^{\alpha}_{\min}$ and
some radius $r^{\alpha}_{\min}(\geq
\widetilde{r}(t^{\alpha}_{\min})>0)$ so that all assumptions of
the theorem still hold for $(x^{\alpha}_{\min},t^{\alpha}_{\min})$
and $r^{\alpha}_{\min}$, whereas the conclusion of the theorem
does not hold with $R\geq K(r^{\alpha}_{\min})^{-2}$ somewhere in
$P(x^{\alpha}_{\min},t^{\alpha}_{\min},\frac{1}{4}
r^{\alpha}_{\min},-\tau(r^{\alpha}_{\min})^2)$ for all
sufficiently large $\alpha$. Here we used the fact that if $R <
K(r^{\alpha}_{\min})^{-2}$ on
$P(x^{\alpha}_{\min},t^{\alpha}_{\min},\frac{1}{4}
r^{\alpha}_{\min},-\tau(r^{\alpha}_{\min})^2)$, then there is no
$\delta$-cutoff surgery there; otherwise there must be a point
there with scalar curvature at least
$\frac{1}{2}\delta^{-2}(\frac{t^{\alpha}_{\min}}{2})
(\widetilde{r}(\frac{t^{\alpha}_{\min}}{2}))^{-2}
\geq\frac{1}{2}(\widetilde{r}(t^{\alpha}_{\min}))^{-2}
(\widetilde{r}(\frac{t^{\alpha}_{\min}}{2}))^{-2} \gg K
(r^{\alpha}_{\min})^{-2}$ since
$\bar{r}^{\alpha}\sqrt{t^{\alpha}_{\min}}\geq
r^{\alpha}_{\min}\geq \widetilde{r}(t^{\alpha}_{\min})$ and
$\bar{r}^{\alpha} \rightarrow 0$, which is a contradiction.

After choosing the first time $t^{\alpha}_{\min}$, by passing to a
limit again, we can then choose $r^{\alpha}_{\min}$ to be the
smallest radius for all possible
$(x^{\alpha}_{\min},t^{\alpha}_{\min})$'s and
$r^{\alpha}_{\min}$'s with that property. Thus we have verified
the claim.

For simplicity, we will drop the index $\alpha$ in the following
arguments. By the assumption and the standard volume comparison,
we have
$$
\Vol_{t_0}\left(B_{t_0}\left(x_0,\frac{1}{2}r_0\right)\right)\geq \xi_0r^3_0
$$
for some universal positive $\xi_0$. As in deriving (7.2.7) and
(7.2.8), we can find a ball $B_{t_0}(x'_0,r'_0)\subset
B_{t_0}(x_0,\frac{r_0}{2})$ with
$$
\Vol_{t_0}(B_{t_0}(x'_0,r'_0)) \geq \frac{1}{2}\alpha_3(r'_0)^3\;
\mbox{ and }\; \frac{1}{2}r_0\geq r'_0\geq \xi'_0r_0
$$
for some universal positive constant $\xi'_0$. Then by what we had
proved in the previous case and by the choice of the first time
$t_0$ and the smallest radius $r_0$, we know that the solution is
defined in $P(x'_0,t_0,\frac{r'_0}{4},-\tau (r'_0)^2)$ with the
curvature bound $$R<K(r'_0)^{-2}\leq K(\xi'_0)^{-2}r^{-2}_0.$$
Since $\bar{r}\sqrt{t_0}\geq r_0\geq \widetilde{r}(t_0)$ and
$\bar{r}\rightarrow 0$ as $\alpha\rightarrow\infty$, we see that
$t_0\rightarrow +\infty$ and $t_0r^{-2}_0\rightarrow+\infty$ as
$\alpha\rightarrow+\infty$. Define $T(w) = 8c(w)^{-1} +
\bar{\xi}$, for some suitable large universal positive constant
$\bar{\xi}$. Then for $\alpha$ sufficiently large, we can apply
Theorem 7.5.1(iii) and the pinching assumption to conclude that
\begin{equation} R\leq K'r^{-2}_0,\quad \mbox{ on }\;
P\left(x_0,t_0,4r_0,-\frac{\tau}{2}(\xi'_0)^2r^2_0\right), %\eqno (7.5.20)
\end{equation} for some positive constant $K'$ depending only on $K$ and
$\xi'_0$.

Furthermore, by combining with the pinching assumption, we deduce
that when $\alpha$ sufficiently large,
\begin{align}
Rm &  \geq -[f^{-1}(R(1+t))/(R(1+t))]R\\
&  \geq -r^{-2}_0 \nonumber
\end{align}  %\eqno (7.5.21)
on $P(x_0,t_0,r_0,-\frac{\tau}{2}(\xi'_0)^2r^2_0)$.  So by
applying Theorem 6.3.3(ii) with $w=\frac{1}{2}\alpha_3$, we have
that when $\alpha$ sufficiently large, \begin{equation}
\Vol_t(B_t(x_0,r_0)) \geq \xi_1 r^3_0,  %\eqno (7.5.22)
\end{equation} for all $ t \in [t_0 -\frac{\tau}{2}(\xi'_0)^2r^2_0, t_0]$,
and \begin{equation} R\leq\left(\bar{C}+\frac{2\bar{B}}{\bar{\tau}_0}\right)r^{-2}_0\leq
\frac{1}{2}Kr^{-2}_0 %\eqno (7.5.23)
\end{equation} on $P(x_0,t_0,\frac{r_0}{4},-\frac{\tau}{2}(\xi'_0)^2r^2_0)$,
where $\xi_1$ is some universal positive constant.

Next we want to extend the estimate (7.5.23) backwards in time.
Denote by $t_1=t_0-\frac{\tau}{2}(\xi'_0)^2r^2_0$. The estimate
(7.5.22) gives
$$
\Vol_{t_1}(B_{t_1}(x_0,r_0))\geq \xi_1r^3_0.
$$
By the same argument in the derivation of (7.2.7) and (7.2.8)
again, we can find a ball $B_{t_1}(x_1,r_1)\subset
B_{t_1}(x_0,r_0)$ with
$$
\Vol_{t_1}(B_{t_1}(x_1,r_1)) \geq \frac{1}{2}\alpha_3r^3_1
$$
and with
$$
r_1\geq \xi'_1r_0
$$
for some universal positive constant $\xi'_1$. Then by what we had
proved in the previous case and by the lower bound (7.5.21) at
$t_1$ and the choice of the first time $t_0$, we know that the
solution is defined on $P(x_1,t_1,\frac{r_1}{4},-\tau r^2_1)$ with
the curvature bound $R<K r^{-2}_1$. By applying Theorem 7.5.1(iii)
and the pinching assumption again we get that for $\alpha$
sufficiently large,
$$
R\leq K''r^{-2}_0 \leqno(7.5.20)'
$$ on
$P(x_0,t_1,4r_0,-\frac{\tau}{2}(\xi'_1)^2r^2_0)$, for some
positive constant $K''$ depending only on $K$ and $\xi'_1$.
Moreover, by combining with the pinching assumption, we have
\begin{align*}
Rm&  \geq-[f^{-1}(R(1+t))/(R(1+t))]R \tag*{(7.5.21)$'$}\\
&  \geq -r^{-2}_0
\end{align*}  %\eqno (7.5.21)'
on $P(x_0,t_1,4r_0,-\frac{\tau}{2}(\xi'_1)^2r^2_0)$, for $\alpha$
sufficiently large. So by applying Theorem 6.3.3 (ii) with
$w=\frac{1}{2}\alpha_3$ again, we have that for $\alpha$
sufficiently large,
$$
\Vol_t(B_t(x_0,r_0)) \geq \xi_1 r^3_0,\leqno (7.5.22)'
$$
for all $ t \in [t_0 -\frac{\tau}{2}(\xi'_0)^2r^2_0
-\frac{\tau}{2}(\xi'_1)^2r^2_0, t_0]$, and
$$
R\leq \left(\bar{C}+\frac{2\bar{B}}{\bar{\tau}_0}\right)r^{-2}_0\leq
\frac{1}{2}Kr^{-2}_0 \leqno (7.5.23)'
$$
on $P(x_0,t_0,\frac{r_0}{4},-\frac{\tau}{2}(\xi'_0)^2
r^2_0-\frac{\tau}{2}(\xi'_1)^2r^2_0)$.

Note that the constants $\xi_0$, $\xi'_0$, $\xi_1$ and $\xi'_1$
are universal, independent of the time $t_1$ and the choice of the
ball $B_{t_1}(x_1,r_1)$. Then we can repeat the above procedure as
many times as we like, until we reach the time $t_0-\tau r^2_0$.
Hence we obtain the estimate
$$
R\leq(\bar{C}+\frac{2\bar{B}}{\bar{\tau}_0})r^{-2}_0\leq
\frac{1}{2}Kr^{-2}_0 \leqno (7.5.23)''
$$
on $P(x_0,t_0,\frac{r_0}{4},-\tau r^2_0)$, for sufficiently large
$\alpha$. This contradicts the choice of the point $(x_0,t_0)$ and
the radius $r_0$ which make $R\geq Kr^{-2}_0$ somewhere in
$P(x_0,t_0,\frac{r_0}{4},-\tau r^2_0)$.

Therefore we have completed the proof of the theorem.
%$$\eqno \#$$\vskip 0.3cm
\end{proof}

Consequently, we have the following result which is analog of
Corollary 7.2.4. This result is a weak version of a claim in the
section 7.3 of {Perelman \cite{P2}}.

%\vskip 0.2cm \noindent{\bf Corollary 7.5.3} \emph{
%CZ-SOURCE-LINE-21847
\begin{corollary}[large-time canonical neighborhoods under local volume control]
  \label{cor:cz-large-time-local-volume-canonical-neighborhoods}
  \dcref{Cao--Zhu Corollary 7.5.3}
  \group{Curvature Estimates after Surgery}
  \source{cao-zhu}

For any $\varepsilon > 0$ and $w>0$, there exist
$\bar{r}=\bar{r}(w,\varepsilon)>0$,
$\theta=\theta(w,\varepsilon)>0$ and $T=T(w)$ with the following
property. Suppose we have a solution, constructed by Theorem
$7.4.3$ with the positive functions $\widetilde{\delta}(t)$ and
$\widetilde{r}(t)$, to the Ricci flow with surgery with a compact
orientable normalized three-manifold as initial data, where each
$\delta$-cutoff at a time $t$ has $\delta=\delta(t)\leq
\min\{\widetilde{\delta}(t), \widetilde{r}(2t)\}$. If
$B_{t_0}(x_0,r_0)$ is a geodesic ball at time $t_0$, with
$\theta^{-1}h\leq r_0\leq \bar{r}\sqrt{t_0}$ and $t_0 \geq T$,
where $h$ is the maximal cutoff radii for surgeries in
$[\frac{t_0}{2},t_0]$ (if there is no surgery in the time interval
$[\frac{t_0}{2},t_0]$, we take $h=0$), and satisfies
$$
\min\{Rm(x,t_0)\ |\ x\in B_{t_0}(x_0,r_0)\}=-r^{-2}_0,
$$
then
$$
\Vol_{t_0}(B_{t_0}(x_0,r_0))<w r^3_0.
$$
\end{corollary}

%\noindent{\bf Proof.}
\begin{proof}
  \uses{thm:cz-large-time-surgery-curvature-estimate, thm:cz-bounded-time-surgery-curvature-estimate}
We argue by contradiction. Let $\theta=\theta(w,\varepsilon)$  and
$T=2T(w)$, where $\theta(w,\varepsilon)$ and $T(w)$ are the
positive constant in Theorem 7.5.2. Suppose for any $\bar{r}>0$
there is a solution and a geodesic ball $B_{t_0}(x_0,r_0)$
satisfying the assumptions of the corollary with $\theta^{-1}h\leq
r_0\leq \bar{r}\sqrt{t_0}$ and $t_0 \geq T$, and with
$$
\min\{Rm(x,t_0)\ |\ x\in B_{t_0}(x_0,r_0)\}=-r^{-2}_0,
$$
but
$$
\Vol_{t_0}(B_{t_0}(x_0,r_0))\geq w r^3_0.
$$
Without loss of generality, we may assume that $\bar{r}$ is less
than the corresponding constant in Theorem 7.5.2. We can then
apply Theorem 7.5.2 to get $$R(x,t)\leq Kr^{-2}_0 $$ whenever
$t\in[t_0-\tau r^2_0,t_0]$ and $d_t(x,x_0)\leq \frac{r_0}{4}$,
where $\tau$ and $K$ are the positive constants in Theorem 7.5.2.
Note that $t_0r^{-2}_0\geq \bar{r}^{-2}\rightarrow+\infty$ as
$\bar{r}\rightarrow 0$. By combining with the pinching assumption
we have
\begin{align*}
Rm&  \geq-[f^{-1}(R(1+t))/(R(1+t))]R\\
&  \geq -\frac{1}{2}r^{-2}_0
\end{align*}
in the region $P(x_0,t_0,\frac{r_0}{4},-\tau r^2_0)= \{(x,t)\ |\
x\in B_t(x_0,\frac{r_0}{4}),t\in[t_0-\tau r^2_0,t_0]\}$, provided
$\bar{r}>0$ is sufficiently small. Thus we get the estimate
$$|Rm|\leq K'r^{-2}_0
$$
in $P(x_0,t_0,\frac{r_0}{4},-\tau r^2_0)$, where $K'$ is a
positive constant depending only on $w$ and $\varepsilon$.

We can now apply Theorem 7.5.1 (iii) to conclude that
$$
R(x,t)\leq \widetilde{K}r^{-2}_0
$$
whenever $t\in[t_0-\frac{\tau}{2}r^2_0,t_0]$ and $d_t(x,x_0)\leq
r_0$, where $\widetilde{K}$ is a positive constant depending only
on $w$ and $\varepsilon$. By using the pinching assumption again
we further have
$$
Rm(x,t)\geq -\frac{1}{2}r^{-2}_0
$$
in the region $P(x_0,t_0,r_0,-\frac{\tau}{2}r^2_0)=\{(x,t)\ |\
x\in B_t(x_0,r_0),t\in[t_0-\frac{\tau}{2}r^2_0,t_0]\}$, as long as
$\bar{r}$ is sufficiently small. In particular, this would imply
$$
\min\{Rm(x,t_0)\ |\ x\in B_{t_0}(x_0,r_0)\}>-r^{-2}_0,
$$
which is a contradiction.  %$$\eqno \#$$\vskip 0.3cm
\end{proof}

%\noindent{\bf Remark 7.5.4} \ \
%CZ-SOURCE-LINE-21927
\begin{remark}[scope of the large-time canonical-neighborhood estimate]
  \label{rem:cz-large-time-canonical-neighborhood-scope}
  \dcref{Cao--Zhu Remark 7.5.4}
  \group{Curvature Estimates after Surgery}
  \source{cao-zhu}
  \uses{cor:cz-large-time-local-volume-canonical-neighborhoods}

In section 7.3 of \cite{P2}, Perelman claimed a stronger statement
than the above Corollary 7.5.3 that allows $r_0 < \theta^{-1}h$ in
the assumptions. Nevertheless, the above weaker statement is
sufficient to deduce the geometrization result.
\end{remark}


\section{Long Time Behavior}

In Section 5.3, we obtained the long time behavior for smooth
(compact) solutions to the three-dimensional Ricci flow with
bounded normalized curvature. The purpose of this section is to
adapt Hamilton's arguments there to solutions of the Ricci flow
with surgery and to drop the bounded normalized curvature
assumption as sketched by Perelman \cite{P2}.

Recall from Corollary 7.4.4 that we have completely understood the
topological structure of a compact, orientable three-manifold with
nonnegative scalar curvature. From now on we assume that our
initial manifold does not admit any metric with nonnegative scalar
curvature, and that \textbf{once we get a compact component with
nonnegative scalar curvature, it is immediately removed}.
Furthermore, if a solution to the Ricci flow with surgery becomes
extinct in a finite time, we have also obtained the topological
structure of the initial manifold. So in the following we only
consider those solutions to the Ricci flow with surgery which
exist for all times $t \geq 0$.

Let $g_{ij}(t)$, $0 \leq t < +\infty$, be a solution to the Ricci
flow with $\delta$-cutoff surgeries, constructed by Theorem 7.4.3
with normalized initial data. Let $0<t_1<t_2<\cdots<t_k<\cdots$ be
the surgery times, where each $\delta$-cutoff at a time $t_k$ has
$\delta = \delta(t_k) \leq \min
\{\widetilde{\delta}(t_k),\widetilde{r}(2t_k)\}$. On each time
interval $(t_{k-1},t_k)$ (denote by $t_0=0$), the scalar curvature
satisfies the evolution equation \begin{equation} \frac{\partial}{\partial
t}R=\Delta
R+2|\stackrel{\circ}{\Ric}|^2+\frac{2}{3}R^2 %\eqno (7.6.1)
\end{equation} where $\stackrel{\circ}{\Ric}$ is the trace-free part of
$\Ric.$ Then $R_{\min}(t)$, the minimum of the scalar curvature at
the time $t$, satisfies
$$
\frac{d}{dt}R_{\min}(t)\geq \frac{2}{3}R^2_{\min}(t)
$$
for $t\in(t_{k-1},t_k)$, for each $k = 1, 2, \ldots$. Since our
surgery procedure had removed all components with nonnegative
scalar curvature, the minimum $R_{\min}(t)$ is negative for all $t
\in [0,+\infty)$. Also recall that the cutoff surgeries were
performed only on $\delta$-necks. Thus the surgeries do not occur
at the parts where $R_{\min}(t)$ are achieved. So the differential
inequality
$$
\frac{d}{dt}R_{\min}(t)\geq \frac{2}{3}R^2_{\min}(t)
$$
holds for all $t\geq 0$, and then by normalization,
$R_{\min}(0)\geq -1$, we have \begin{equation} R_{\min}(t)\geq
-\frac{3}{2}\cdot\frac{1}{t+\frac{3}{2}}, \
\mbox{ for all } \ t\geq 0. %\eqno (7.6.2)
\end{equation} Meanwhile, on each time interval $(t_{k-1},t_k)$, the volume
satisfies the evolution equation
$$
\frac{d}{dt}V=-\int RdV
$$
and then by (7.6.2),
$$
\frac{d}{dt}V\leq \frac{3}{2}\cdot\frac{1}{(t+\frac{3}{2})}V.
$$
Since the cutoff surgeries do not increase volume, we thus have
\begin{equation} \frac{d}{dt}\log\left(V(t)\left(t+\frac{3}{2}\right)^{-\frac{3}{2}}\right)\leq 0
%\eqno(7.6.3)
\end{equation} for all $t\geq 0$. Equivalently, the function
$V(t)(t+\frac{3}{2})^{-\frac{3}{2}}$ is nonincreasing on
$[0,+\infty)$.

We can now use the monotonicity of the function
$V(t)(t+\frac{3}{2})^{-\frac{3}{2}}$ to extract the information of
the solution at large times. On each time interval
$(t_{k-1},t_k)$, we have
\begin{align*}
&\frac{d}{dt}\log\left(V(t)\left(t+\frac{3}{2}\right)^{-\frac{3}{2}}\right) \\
&=-\left(R_{\min}(t)+\frac{3}{2\left(t+\frac{3}{2}\right)}\right)
+\frac{1}{V}\int_M(R_{\min}(t)-R)dV.
\end{align*}
Then by noting that the cutoff surgeries do not increase volume,
we get
\begin{align}
\frac{V(t)}{(t+\frac{3}{2})^{\frac{3}{2}}}
&\leq\frac{V(0)}{(\frac{3}{2})^{\frac{3}{2}}}
\exp\Bigg\{-\int^t_0\left(R_{\min}(t)
+\frac{3}{2(t+\frac{3}{2})}\right)dt \\
&\qquad-\int^t_0\frac{1}{V}
\int_M(R-R_{\min}(t))dVdt\Bigg\} \nonumber %\eqno (7.6.4)
\end{align}
for all $t>0$. Now by this inequality and the equation (7.6.1), we
obtain the following consequence (cf. Lemma 7.1 of Hamilton
\cite{Ha99} and section 7.1 of Perelman \cite{P2}).

%\vskip 0.2cm\noindent{\bf Lemma 7.6.1} \emph{
%CZ-SOURCE-LINE-22026
\begin{lemma}[flatness of bounded-curvature parabolic blow-downs]
  \label{lem:cz-parabolic-blowdown-flatness}
  \dcref{Cao--Zhu Lemma 7.6.1}
  \group{Long-Time Thick--Thin Geometry}
  \source{cao-zhu}
  \uses{thm:cz-ricci-flow-surgery-long-time-existence}

Let $g_{ij}(t)$ be a solution to the Ricci flow with surgery,
constructed by Theorem $7.4.3$ with normalized initial data. If for
a fixed $0<r<1$ and a sequence of times
$t^{\alpha}\rightarrow\infty$, the rescalings of the solution on the
parabolic neighborhoods
$P(x^{\alpha},t^{\alpha},r\sqrt{t^{\alpha}},-r^2t^{\alpha})=\{(x,t)\
|\ x\in B_t(x^{\alpha},r\sqrt{t^{\alpha}}),
t\in[t^{\alpha}-r^2t^{\alpha},t^{\alpha}]\}$, with factor
$(t^{\alpha})^{-1}$ and shifting the times $t^{\alpha}$ to $1$,
converge in the $C^{\infty}$ topology to some smooth limiting
solution, defined in an abstract parabolic neighborhood
$P(\bar{x},1,r,-r^2)$, then this limiting solution has constant
sectional curvature $-1/4t$ at any time $t\in [1-r^2,1]$.
%$$ \eqno \#$$ \vskip 0.3cm
\end{lemma}

In the previous section we obtained several curvature estimates
for the solutions to the Ricci flow with surgery. Now we combine
the curvature estimates with the above lemma to derive the
following asymptotic result for the curvature.

%\vskip 0.2cm \noindent{\bf Lemma 7.6.2} (Perelman \cite{P2})\
%CZ-SOURCE-LINE-22049
\begin{lemma}[large-time local curvature alternatives]
  \label{lem:cz-large-time-local-curvature-alternatives}
  \dcref{Cao--Zhu Lemma 7.6.2}
  \group{Long-Time Thick--Thin Geometry}
  \source{cao-zhu}
  \uses{thm:cz-ricci-flow-surgery-long-time-existence}

For any $\varepsilon > 0$, let $g_{ij}(t)$, $0 \leq t < +\infty$,
be a solution to the Ricci flow with surgery, constructed by
Theorem $7.4.3$ with normalized initial data.
\begin{itemize}
\item[(i)] Given $w>0$, $r>0$, $\xi>0$, one can find
$T=T(w,r,\xi,\varepsilon)<+\infty$ such that if the geodesic ball
$B_{t_0}(x_0,r\sqrt{t_0})$ at some time $t_0\geq T$ has volume at
least $w r^3t^{\frac{3}{2}}_0$ and the sectional curvature at
least $-r^{-2}t^{-1}_0$, then the curvature at $x_0$ at time
$t=t_0$ satisfies \begin{equation}
|2tR_{ij}+g_{ij}|<\xi. %\eqno (7.6.5)
\end{equation} \item[(ii)] Given in addition $1 \leq A<\infty$ and allowing
$T$ to depend on $A$, we can ensure $(7.6.5)$ for all points in
$B_{t_0}(x_0,Ar\sqrt{t_0})$. \item[(iii)] The same is true for all
points in the forward parabolic neighborhood
$P(x_0,t_0,Ar\sqrt{t_0},Ar^2t_0)\triangleq\{(x,t)\ |\ x\in
B_t(x_0,Ar\sqrt{t_0}),t\in[t_0,t_0+Ar^2t_0]\}.$
\end{itemize}
\end{lemma}

%\vskip 0.1cm \noindent{\bf Proof.}
\begin{proof}
  \uses{thm:cz-large-time-surgery-curvature-estimate, lem:cz-parabolic-blowdown-flatness, thm:cz-bounded-time-surgery-curvature-estimate}
(i) By the assumptions and the standard volume comparison, we have
$$
\Vol_{t_0}(B_{t_0}(x_0,\rho))\geq  cw \rho^3
$$
for all $0<\rho\leq r\sqrt{t_0}$, where $c$ is a universal
positive constant. Let $\bar{r}=\bar{r}(cw,\varepsilon)$ be the
positive constant in Theorem 7.5.2 and set
$r_0=\min\{r,\bar{r}\}$. On $B_{t_0}(x_0,r_0\sqrt{t_0})(\subset
B_{t_0}(x_0,r\sqrt{t_0}))$, we have
\begin{gather}
Rm\geq -(r_0\sqrt{t_0})^{-2} \\
\mbox{and }\qquad \Vol_{t_0}(B_{t_0}(x_0,r_0\sqrt{t_0}))\geq
cw(r_0\sqrt{t_0})^3.\nonumber
%\eqno (7.6.6)
\end{gather}
Obviously, there holds $\theta^{-1}h\leq r_0\sqrt{t_0}\leq
\bar{r}\sqrt{t_0}$ when $t_0$ is large enough, where
$\theta=\theta(cw,\varepsilon)$ is the positive constant in
Theorem 7.5.2 and $h$ is the maximal cutoff radius for surgeries
in $[\frac{t_0}{2},t_0]$ (if there is no surgery in the time
interval $[\frac{t_0}{2},t_0]$, we take $h=0$). Then it follows
from Theorem 7.5.2 that the solution is defined and satisfies
$$
R<K(r_0\sqrt{t_0})^{-2}
$$
on whole parabolic neighborhood
$P(x_0,t_0,\frac{r_0\sqrt{t_0}}{4},-\tau(r_0\sqrt{t_0})^2)$. Here
$\tau=\tau(cw,\varepsilon)$ and $K=K(cw,\varepsilon)$ are the
positive constants in Theorem 7.5.2. By combining with the
pinching assumption we have
\begin{align*}
Rm&  \geq-[f^{-1}(R(1+t))/(R(1+t))]R\\
&  \geq -{\rm const.}\,K(r_0\sqrt{t_0})^{-2}
\end{align*}
in the region $P(x_0,t_0, \frac{r_0\sqrt{t_0}}{4},
-\tau(r_0\sqrt{t_0})^2)$. Thus we get the estimate \begin{equation}
|Rm|\leq K'(r_0\sqrt{t_0})^{-2} %\eqno (7.6.7)
\end{equation} on
$P(x_0,t_0,\frac{r_0\sqrt{t_0}}{4},-\tau(r_0\sqrt{t_0})^2)$, for
some positive constant $K'=K'(w,\varepsilon)$ depending only on
$w$ and $\varepsilon$.

The curvature estimate (7.6.7) and the volume estimate (7.6.6)
ensure that as $t_0\rightarrow+\infty$ we can take smooth
(subsequent) limits for the rescalings of the solution with factor
$(t_0)^{-1}$ on parabolic neighborhoods
$P(x_0,t_0,\frac{r_0\sqrt{t_0}}{4},$ $-\tau(r_0\sqrt{t_0})^2)$.
Then by applying Lemma 7.6.1, we can find
$T=T(w,r,\xi,\varepsilon)<+\infty$ such that when $t_0\geq T$,
there holds \begin{equation}
|2tR_{ij}+g_{ij}|<\xi, %\eqno(7.6.8)
\end{equation} on
$P(x_0,t_0,\frac{r_0\sqrt{t_0}}{4},-\tau(r_0\sqrt{t_0})^2)$, in
particular,
$$
|2tR_{ij}+g_{ij}|(x_0,t_0)<\xi.
$$
This proves the assertion (i).

\medskip
(ii) In view of the above argument, to get the estimate (7.6.5)
for all points in $B_{t_0}(x_0,Ar\sqrt{t_0})$, the key point is to
get a upper bound for the scalar curvature on the parabolic
neighborhood $P(x_0,t_0,Ar\sqrt{t_0},-\tau(r_0\sqrt{t_0})^2)$.
After having the estimates (7.6.6) and (7.6.7), one would like to
use Theorem 7.5.1(iii) to obtain the desired scalar curvature
estimate. Unfortunately it does not work since our $r_0$ may be
much larger than the constant $\bar{r}(A,\varepsilon)$ there when
$A$ is very large. In the following we will use Theorem 7.5.1(ii)
to overcome the difficulty.

Given $1 \leq A<+\infty$, based on (7.6.6) and (7.6.7), we can use
Theorem 7.5.1(ii) to find a positive constant
$K_1=K_1(w,r,A,\varepsilon)$ such that each point in
$B_{t_0}(x_0,2Ar\sqrt{t_0})$ with its scalar curvature at least
$K_1(r\sqrt{t_0})^{-2}$ has a canonical neighborhood. We claim
that there exists $T=T(w,r,A,\varepsilon)<+\infty$ so that when
$t_0\geq T$, we have \begin{equation} R<K_1(r\sqrt{t_0})^{-2},\; \mbox{ on } \;
B_{t_0}(x_0,2Ar\sqrt{t_0}). %\eqno (7.6.9)
\end{equation}

Argue by contradiction. Suppose not; then there exist a sequence
of times $t^{\alpha}_0\rightarrow+\infty$ and sequences of points
$x^{\alpha}_0$, $x^{\alpha}$ with $x^{\alpha}\in
B_{t^{\alpha}_0}(x^{\alpha}_0,2Ar\sqrt{t^{\alpha}_0})$ and
$R(x^{\alpha},t^{\alpha}_0)=K_1(r\sqrt{t^{\alpha}_0})^{-2}$. Since
there exist canonical neighborhoods ($\varepsilon$-necks or
$\varepsilon$-caps) around the points $(x^{\alpha},t^{\alpha}_0)$,
there exist positive constants $c_1$, $C_2$ depending only on
$\varepsilon$ such that
$$
\Vol_{t^{\alpha}_0}(B_{t^{\alpha}_0}(x^{\alpha},
K^{-\frac{1}{2}}_1(r\sqrt{t^{\alpha}_0})))\geq
c_1(K^{-\frac{1}{2}}_1(r\sqrt{t^{\alpha}_0}))^3
$$
and
$$
C^{-1}_2K_1(r\sqrt{t^{\alpha}_0})^{-2} \leq R(x,t^{\alpha}_0) \leq
C_2K_1(r\sqrt{t^{\alpha}_0})^{-2},
$$
on $B_{t^{\alpha}_0} (x^{\alpha}, K^{-\frac{1}{2}}_1
(r\sqrt{t^{\alpha}_0}))$, for all $\alpha$. By combining with the
pinching assumption we have
\begin{align*}
Rm&  \geq-[f^{-1}(R(1+t))/(R(1+t))]R\\
&  \geq -{\rm const.}\, C_2K_1(r\sqrt{t^{\alpha}_0})^{-2},
\end{align*}
on
$B_{t^{\alpha}_0}(x^{\alpha},K^{-\frac{1}{2}}_1(r\sqrt{t^{\alpha}_0}))$,
for all $\alpha$. It then follows from the assertion (i) we just
proved that
$$
\lim_{\alpha\rightarrow+\infty}|2tR_{ij}+g_{ij}|
(x^{\alpha},t^{\alpha}_0)=0.
$$
In particular, we have
$$
t^{\alpha}_0R(x^{\alpha},t^{\alpha}_0)<-1
$$
for $\alpha$ sufficiently large. This contradicts our assumption
that $R(x^{\alpha},t^{\alpha}_0)=K_1(r\sqrt{t^{\alpha}_0})^{-2}$.
So we have proved assertion (7.6.9).

Now by combining (7.6.9) with the pinching assumption as before,
we have \begin{equation}
Rm\geq  -K_2(r\sqrt{t_0})^{-2} %\eqno (7.6.10)
\end{equation} on $B_{t_0}(x_0,2Ar\sqrt{t_0})$, where
$K_2=K_2(w,r,A,\varepsilon)$ is some positive constant depending
only on $w$, $r$, $A$ and $\varepsilon$. Thus by (7.6.9) and
(7.6.10) we have \begin{equation} |Rm|\leq K'_1(r\sqrt{t_0})^{-2},\
 \ \mbox{on}\ \ B_{t_0}(x_0,2Ar\sqrt{t_0}), %\eqno (7.6.11)
\end{equation} for some positive constant $K'_1=K'_1(w,r,A,\varepsilon)$
depending only on $w$, $r$, $A$ and $\varepsilon$. This gives us
the curvature estimate on $B_{t_0}(x_0,2Ar\sqrt{t_0})$ for all
$t_0 \geq T(w,r,A,\varepsilon)$.

{}From the arguments in proving the above assertion (i), we have
the estimates (7.6.6) and (7.6.7) and the solution is well-defined
on the whole parabolic neighborhood
$P(x_0,t_0,\frac{r_0\sqrt{t_0}}{4},-\tau(r_0\sqrt{t_0})^2)$ for
all $t_0 \geq T(w,r,A,\varepsilon)$. Clearly we may assume that
$(K'_1)^{-\frac{1}{2}}r < \min \{\frac{r_0}{4}, \sqrt{\tau}r_0
\}$. Thus by combining with the curvature estimate (7.6.11), we
can apply Theorem 7.5.1(i) to get the following volume control \begin{equation}
\Vol_{t_0}(B_{t_0}(x,(K'_1)^{-\frac{1}{2}}r\sqrt{t_0}))
\geq \kappa((K'_1)^{-\frac{1}{2}}r\sqrt{t_0})^3 %\eqno (7.6.12)
\end{equation} for any $x\in B_{t_0}(x_0,Ar\sqrt{t_0})$, where
$\kappa=\kappa(w,r,A,\varepsilon)$ is some positive constant
depending only on $w$, $r$, $A$ and $\varepsilon$. So by using the
assertion (i), we see that for $t_0\geq T$ with
$T=T(w,r,\xi,A,\varepsilon)$ large enough, the curvature estimate
(7.6.5) holds for all points in $B_{t_0}(x_0,Ar\sqrt{t_0})$.

\medskip
(iii) We next want to extend the curvature estimate (7.6.5) to all
points in the forward parabolic neighborhood
$P(x_0,t_0,Ar\sqrt{t_0},Ar^2t_0)$. Consider the time interval
$[t_0,t_0+Ar^2t_0]$ in the parabolic neighborhood. In assertion
(ii), we have obtained the desired estimate (7.6.5) at $t=t_0$.
Suppose estimate (7.6.5) holds on a maximal time interval
$[t_0,t')$ with $t'\leq t_0+Ar^2t_0$. This says that we have \begin{equation}
|2tR_{ij}+g_{ij}|<\xi %\eqno (7.6.13)
\end{equation} on $P(x_0,t_0,Ar\sqrt{t_0},t'-t_0)\triangleq\{(x,t)\ |\ x\in
B_t(x_0,Ar\sqrt{t_0}),t\in[t_0,t')\}$ so that either there exists
a surgery in the ball $B_{t'}(x_0,Ar\sqrt{t_0})$ at $t=t'$, or
there holds $|2tR_{ij}+g_{ij}|=\xi$ somewhere in
$B_{t'}(x_0,Ar\sqrt{t_0})$ at $t=t'$. Since the Ricci curvature is
near $-\frac{1}{2t'}$ in the geodesic ball, the surgeries cannot
occur there. Thus we only need to consider the latter possibility.

Recall that the evolution of the length of a curve $\gamma$ and
the volume of a domain $\Omega$ are given by
\begin{align*}
\frac{d}{dt}L_t(\gamma)
&=-\int_{\gamma}\Ric(\dot{\gamma},\dot{\gamma})ds_t \\
\mbox{and } \quad \frac{d}{dt}Vol_t(\Omega) &=-\int_{\Omega}RdV_t.
\end{align*}
By substituting the curvature estimate (7.6.13) into the above two
evolution equations and using the volume lower bound (7.6.6), it
is not hard to see \begin{equation} \Vol_{t'}(B_{t'}(x_0,\sqrt{t'}))\geq
 \kappa'(t')^{\frac{3}{2}} % \eqno (7.6.14)
\end{equation} for some positive constant
$\kappa'=\kappa'(w,r,\xi,A,\varepsilon)$ depending only on $w$,
$r$, $\xi$, $A$ and $\varepsilon$. Then by the above assertion
(ii), the combination of the curvature estimate (7.6.13) and the
volume lower bound (7.6.14) implies that the curvature estimate
(7.6.5) still holds for all points in $B_{t'}(x_0,Ar\sqrt{t_0})$
provided $T=T(w,r,\xi,A,\varepsilon)$ is chosen large enough. This
is a contradiction.  Therefore we have proved assertion (iii).
%$$\eqno \#$$ \vskip 0.3cm
\end{proof}

We now state and prove the following important \textbf{Thick-thin
decomposition theorem}\index{thick-thin decomposition theorem}. A
more general version (without the restriction on $\varepsilon$)
was implicitly claimed by Perelman in \cite{P1} and \cite{P2}.

%\vskip 0.2cm\noindent{\bf Theorem 7.6.3}
%CZ-SOURCE-LINE-22270
\begin{theorem}[thick--thin decomposition for Ricci flow with surgery]
  \label{thm:cz-surgery-thick-thin-decomposition}
  \dcref{Cao--Zhu Theorem 7.6.3}
  \group{Long-Time Thick--Thin Geometry}
  \source{cao-zhu}
  \uses{thm:cz-ricci-flow-surgery-long-time-existence}

For any $w>0$ and $0 < \varepsilon \leq \frac{1}{2}w$, there
exists a positive constant $\rho = \rho(w,\varepsilon) \leq 1$
with the following property. Suppose $g_{ij}(t)$ $(t \in
[0,+\infty))$ is a solution, constructed by Theorem $7.4.3$ with
the nonincreasing $($continuous$)$ positive functions
$\widetilde{\delta}(t)$ and $\widetilde{r}(t)$, to the Ricci flow
with surgery and with a compact orientable normalized
three-manifold as initial data, where each $\delta$-cutoff at a
time $t$ has $\delta=\delta(t) \leq \min
\{\widetilde{\delta}(t),\widetilde{r}(2t)\}$. Then for any
arbitrarily fixed $\xi >0$, for $t$ large enough, the manifold
$M_t$ at time $t$ admits a decomposition $M_t=M_{\rm
thin}(w,t)\cup M_{\rm thick}(w,t)$ with the following properties:
\begin{itemize}
\item[(a)] For every $x\in M_{\rm thin}(w,t)$, there exists some
$r=r(x,t)>0$, with $0<r\sqrt{t}<\rho\sqrt{t}$, such that
$$
Rm\geq -(r\sqrt{t})^{-2} \ \ on \ \ B_t(x,r\sqrt{t}), \ \
\mbox{and }
$$
$$
\Vol_t(B_t(x,r\sqrt{t}))<w(r\sqrt{t})^3.
$$
\item[(b)] For every $x\in M_{\rm thick}(w,t)$, we have
$$
|2tR_{ij}+g_{ij}|<\xi\ \ on \ \ B_t(x,\rho\sqrt{t}), \ \ \mbox{and
}
$$
$$
\Vol_t(B_t(x,\rho\sqrt{t}))\geq \frac{1}{10}w(\rho\sqrt{t})^3.
$$
\end{itemize}
Moreover, if we take any sequence of points $x^{\alpha}\in M_{\rm
thick}(w,t^{\alpha}),t^{\alpha}\rightarrow+\infty$, then the
scalings of $g_{ij}(t^{\alpha})$ around $x^{\alpha}$ with factor
$(t^{\alpha})^{-1}$ converge smoothly, along a subsequence of
$\alpha\rightarrow+\infty$, to a complete hyperbolic manifold of
finite volume with constant sectional curvature $-\frac{1}{4}$.
\end{theorem}

%\vskip 0.1cm\noindent{\bf Proof.}
\begin{proof}
  \uses{cor:cz-large-time-local-volume-canonical-neighborhoods, lem:cz-large-time-local-curvature-alternatives, thm:cz-large-time-surgery-curvature-estimate}
Let $\bar{r}=\bar{r}(w,\varepsilon),\theta=\theta(w,\varepsilon)$
and $h$ be the positive constants obtained in Corollary 7.5.3. We
may assume $\rho\leq\bar{r} \leq e^{-3}$. For any point $x\in
M_t$, there are two cases: either
$$
{\rm (i)}\ \ \ \mbox{ } \min\{Rm\ |\ B_t(x,\rho\sqrt{t})\}\geq
-(\rho\sqrt{t})^{-2},
$$
or
$$
{\rm (ii)}\ \ \ \mbox{    } \min\{Rm\ |\
B_t(x,\rho\sqrt{t})\}<-(\rho\sqrt{t})^{-2}.$$

Let us first consider Case (i). If
$\Vol_t(B_t(x,\rho\sqrt{t}))<\frac{1}{10}w(\rho\sqrt{t})^3$, then
we can choose $r$ slightly less than $\rho$ so that
$$
Rm\geq -(\rho\sqrt{t})^{-2}\geq -(r\sqrt{t})^{-2}
$$
on $B_t(x,r\sqrt{t})(\subset B_t(x,\rho\sqrt{t}))$, and
$$
\Vol_t(B_t(x,r\sqrt{t}))<\frac{1}{10}w(\rho\sqrt{t})^3<w(r\sqrt{t})^3;
$$
thus $x\in M_{thin}(w,t)$. If $\Vol_t(B_t(x,\rho\sqrt{t}))\geq
\frac{1}{10}w(\rho\sqrt{t})^3$, we can apply Lemma 7.6.2(ii) to
conclude that for $t$ large enough,
$$
|2tR_{ij}+g_{ij}|<\xi \quad \mbox{ on }\; B_t(x,\rho\sqrt{t});
$$
thus $x\in M_{thick}(w,t)$.

Next we consider Case (ii). By continuity, there exists
$0<r=r(x,t)<\rho$ such that \begin{equation}
\min\{Rm\ |\ B_t(x,r\sqrt{t})\}=-(r\sqrt{t})^{-2}. %\eqno (7.6.15)
\end{equation} If $\theta^{-1}h\leq r\sqrt{t}\mbox{  }(\leq
\bar{r}\sqrt{t})$, we can apply Corollary 7.5.3 to conclude
$$
\Vol_t(B_t(x,r\sqrt{t}))<w(r\sqrt{t})^3;
$$
thus $x\in M_{thin}(w,t)$.

We now consider the difficult subcase $\mbox{}r\sqrt{t}
<\theta^{-1}h$. By the pinching assumption, we have
\begin{align*}
R &  \geq (r\sqrt{t})^{-2}(\log[(r\sqrt{t})^{-2}(1+t)] -3)\\
&  \geq (\log r^{-2}-3)(r\sqrt{t})^{-2}\\
&  \geq 2(r\sqrt{t})^{-2}\\
&  \geq 2\theta^{2}h^{-2}
\end{align*}
somewhere in $B_t(x,r\sqrt{t})$. Since $h$ is the maximal cutoff
radius for surgeries in $[\frac{t}{2},t]$, by the design of the
$\delta$-cutoff surgery, we have
\begin{align*}
 h &  \leq \sup \left\{{\delta}^2(s)\widetilde{r}(s)\ |\ s\in
\left[\frac{t}{2},t\right]\right\}\\
& \leq\widetilde{\delta}\left(\frac{t}{2}\right)
\widetilde{r}(t)\widetilde{r}\left(\frac{t}{2}\right).
\end{align*}
Note also $\widetilde{\delta}(\frac{t}{2})\rightarrow0$ as
$t\rightarrow+\infty$. Thus from the canonical neighborhood
assumption, we see that for $t$ large enough, there exists a point
in the ball $B_t(x,r\sqrt{t})$ which has a canonical neighborhood.

We claim that for $t$ sufficiently large, the point $x$ satisfies
$$
R(x,t)\geq \frac{1}{2}(r\sqrt{t})^{-2},
$$
and then the above argument shows that the point $x$ also has a
canonical $\varepsilon$-neck or $\varepsilon$-cap neighborhood.
Otherwise, by continuity, we can choose a point $x^*\in
B_t(x,r\sqrt{t})$ with $R(x^*,t)=\frac{1}{2}(r\sqrt{t})^{-2}$.
Clearly the new point $x^*$ has a canonical neighborhood $B^*$ by
the above argument. In particular, there holds
$$
C^{-1}_2(\varepsilon)R \leq  \frac{1}{2}(r\sqrt{t})^{-2} \leq
C_2(\varepsilon)R
$$
on the canonical neighborhood $B^*$. By the definition of
canonical neighborhood assumption, we have
$$
B_t(x^*,\sigma^*)\subset B^* \subset B_t(x^*,2\sigma^*)
$$
for some $\sigma^* \in
(0,C_1(\varepsilon)R^{-\frac{1}{2}}(x^*,t))$. Clearly, without
loss of generality, we may assume (in the definition of canonical
neighborhood assumption) that $\sigma^* >
2R^{-\frac{1}{2}}(x^*,t).$ Then
\begin{align*}
R(1+t) &  \geq\frac{1}{2}C^{-1}_2(\varepsilon)r^{-2}\\
&  \geq \frac{1}{2}C^{-1}_2(\varepsilon)\rho^{-2}
\end{align*}
on $B_t(x^*,2r\sqrt{t}).$ Thus when we choose
$\rho=\rho(w,\varepsilon)$ small enough, it follows from the
pinching assumption that
\begin{align*}
Rm&  \geq-[f^{-1}(R(1+t))/(R(1+t))]R\\
&  \geq -\frac{1}{2}(r\sqrt{t})^{-2},
\end{align*}
on $B_t(x^*,2r\sqrt{t}).$ This is a contradiction with (7.6.15).

We have seen that $tR(x,t)\geq \frac{1}{2}r^{-2} (\geq
\frac{1}{2}\rho^{-2})$. Since $r^{-2}
> \theta^2 h^{-2} t$ in this subcase, we conclude that for
arbitrarily given $A<+\infty$, \begin{equation}
tR(x,t)>A^2\rho^{-2} %\eqno (7.6.16)
\end{equation} as long as $t$ is large enough.

Let $B$, with $B_t(x,\sigma)\subset B \subset B_t(x,2\sigma)$, be
the canonical $\varepsilon$-neck or $\varepsilon$-cap neighborhood
of $(x,t)$. By the definition of the canonical neighborhood
assumption, we have
$$
0 < \sigma < C_1(\varepsilon)R^{-\frac{1}{2}}(x,t),
$$
$$
C^{-1}_2(\varepsilon)R \leq  R(x,t) \leq C_2(\varepsilon)R, \
\mbox{ on } \ B,
$$
and \begin{equation} \Vol_t(B) \leq \varepsilon \sigma^3 \leq
\frac{1}{2}w \sigma^3. %\eqno (7.6.17)
\end{equation} Choose $0 < A < C_1(\varepsilon)$ so that $\sigma =
AR^{-\frac{1}{2}}(x,t)$. For sufficiently large $t$, since
\begin{align*}
R(1+t) &  \geq C^{-1}_2(\varepsilon)(tR(x,t))\\
&  \geq \frac{1}{2}C^{-1}_2(\varepsilon)\rho^{-2},
\end{align*}
on $B$, we can require $\rho=\rho(w,\varepsilon)$ to be smaller
still, and use the pinching assumption to conclude
\begin{align}
Rm&  \geq-[f^{-1}(R(1+t))/(R(1+t))]R\\
&  \geq -(AR^{-\frac{1}{2}}(x,t))^{-2} \nonumber\\
&  = -\sigma^{-2}, \nonumber
\end{align} %\eqno (7.6.18)
on $B$. For sufficiently large $t$, we adjust
\begin{align}
r&  = \sigma (\sqrt{t})^{-1}\\
&  = (AR^{-\frac{1}{2}}(x,t))(\sqrt{t})^{-1} \nonumber\\
&  < \rho, \nonumber
\end{align} %\eqno (7.6.19)
by (7.6.16). Then the combination of (7.6.17), (7.6.18) and
(7.6.19) implies that $ x\in M_{thin}(w,t).$

The last statement in (b) follows directly from Lemma 7.6.2. (Here
we also used Bishop-Gromov volume comparison, Theorem 7.5.2 and
Hamilton's compactness theorem to take a subsequent limit.)

Therefore we have completed the proof of the theorem.
%$$ \eqno \#$$ \vskip 0.3cm
\end{proof}

To state the long-time behavior of a solution to the Ricci flow
with surgery, we first recall some basic terminology in
three-dimensional topology. A three-manifold $M$ is called
\textbf{irreducible}\index{irreducible} if every smooth two-sphere
embedded in $M$ bounds a three-ball in $M$. If we have a solution
$(M_t,g_{ij}(t))$ obtained by Theorem 7.4.3 with a compact,
orientable and irreducible three-manifold $(M,g_{ij})$ as initial
data, then at each time $t>0$, by the cutoff surgery procedure,
the solution manifold $M_t$ consists of a finite number of
components where one of the components, called the
\textbf{essential component} \index{essential component} and
denoted by $M_t^{(1)}$, is diffeomorphic to the initial manifold
$M$ while the rest are diffeomorphic to the three-sphere
$\mathbb{S}^3$.

The main result of this section is the following generalization of
Theorem 5.3.4. A more general version of the result (without the
restriction on $\varepsilon$) was implicitly claimed by Perelman
in \cite{P2}.

%\vskip 0.2cm\noindent{\bf Theorem 7.6.4 }
%CZ-SOURCE-LINE-22484
\begin{theorem}[long-time geometry of Ricci flow with surgery]
  \label{thm:cz-surgery-long-time-geometry}
  \dcref{Cao--Zhu Theorem 7.6.4}
  \group{Long-Time Thick--Thin Geometry}
  \source{cao-zhu}

Let $w\!>0$ and $0<\!\varepsilon\leq\!\frac{1}{2}w$ be any small
positive constants and let $(M_t,g_{ij}(t)),$ $0<t<+\infty,$ be a
solution to the Ricci flow with surgery, constructed by Theorem
$7.4.3$ with the nonincreasing $($continuous$)$ positive functions
$\widetilde{\delta}(t)$ and $\widetilde{r}(t)$ and with a compact,
orientable, irreducible and normalized three-manifold $M$ as
initial data, where each $\delta$-cutoff at a time $t$ has
$\delta=\delta(t) \leq \min
\{\widetilde{\delta}(t),\widetilde{r}(2t)\}$. Then one of the
following holds: either
\begin{itemize}
\item[(i)] for all sufficiently large $t$, we have $M_t = M_{\rm
thin}(w,t)$; or \item[(ii)] there exists a sequence of times
$t^{\alpha}\rightarrow+\infty$ such that the scalings of
$g_{ij}(t^{\alpha})$ on the essential component
$M_{t^\alpha}^{(1)}$, with factor $(t^{\alpha})^{-1}$, converge in
the $C^{\infty}$ topology to a hyperbolic metric on the initial
compact manifold $M$ with constant sectional curvature
$-\frac{1}{4}$; or \item[(iii)] we can find a finite collection of
complete noncompact hyperbolic three-manifolds
$\mathcal{H}_1,\ldots,\mathcal{H}_m$, with finite volume, and
compact subsets $K_1,\ldots,K_m$ of
$\mathcal{H}_1,\ldots,\mathcal{H}_m$ respectively obtained by
truncating each cusp of the hyperbolic manifolds along constant
mean curvature torus of small area, and for all $t$ beyond some
time $T<+\infty$ we can find diffeomorphisms $\varphi_l,1\leq
l\leq m$, of $K_l$ into $M_t$ so that as long as $t$ is
sufficiently large, the metric $t^{-1}\varphi^*_l(t)g_{ij}(t)$ is
as close to the hyperbolic metric as we like on the compact sets
$K_1,\ldots,K_m$; moreover, the complement $M_t\backslash
(\varphi_1(K_1)\cup\cdots\cup\varphi_m(K_m))$ is contained in the
thin part $M_{\rm thin}(w,t)$, and the boundary tori of each $K_l$
are incompressible in the sense that each $\varphi_l$ injects
$\pi_1(\partial K_l)$ into $\pi_1(M_t)$.
\end{itemize}
\end{theorem}

%\vskip 0.1cm\noindent{\bf Proof.}
\begin{proof}
  \uses{lem:cz-finite-volume-hyperbolic-embedding-stability, lem:cz-concave-boundary-metric-perturbations, thm:cz-noncollapsing-nonsingular-flow-structure}
The proof of the theorem follows, with some modifications, the
same argument of Hamilton \cite{Ha99} as in the proof of Theorem
5.3.4.

Clearly we may assume that the thick part $M_{\rm thick}(w,t)$ is
not empty for a sequence $t^{\alpha}\rightarrow+\infty$, since
otherwise we have case (i). If we take a sequence of points
$x^{\alpha}\in M_{\rm thick}(w,t^{\alpha})$, then by Theorem
7.6.3(b) the scalings of $g_{ij}(t^{\alpha})$ around $x^{\alpha}$
with factor $(t^{\alpha})^{-1}$ converge smoothly, along a
subsequence of $\alpha\rightarrow+\infty$, to a complete
hyperbolic manifold of finite volume with constant sectional
curvature $-\frac{1}{4}$. The limits may be different for
different choices of $(x^{\alpha},t^{\alpha})$. If a limit is
compact, we have case (ii). Thus we assume that all limits are
noncompact.

Consider all the possible hyperbolic limits of the solution, and
among them choose one such complete noncompact hyperbolic
three-manifold $\mathcal{H}$ with the least possible number of
cusps. Denote by $h_{ij}$ the hyperbolic metric of $\mathcal{H}$.
For all small $a>0$ we can truncate each cusp of $\mathcal{H}$
along a constant mean curvature torus of area $a$ which is
uniquely determined; we denote the remainder by $\mathcal{H}_a$.
Fix $a>0$ so small that Lemma 5.3.7 is applicable for the compact
set $\mathcal{K}=\mathcal{H}_a$. Pick an integer $l_0$
sufficiently large and an $\epsilon_0$ sufficiently small to
guarantee from Lemma 5.3.8 that the identity map $Id$ is the only
harmonic map $F$ from $\mathcal{H}_a$ to itself with taking
$\partial \mathcal{H}_a$ to itself, with the normal derivative of
$F$ at the boundary of the domain normal to the boundary of the
target, and with $d_{C^{l_0}(\mathcal{H}_a)}(F,Id) < \epsilon_0$.
Then choose a positive integer $q_0$ and a small number
$\delta_0>0$ from Lemma 5.3.7 such that if $\widetilde{F}$ is a
diffeomorphism of $\mathcal{H}_a$ into another complete noncompact
hyperbolic three-manifold
$(\widetilde{\mathcal{H}},\widetilde{h}_{ij})$ with no fewer cusps
(than $\mathcal{H}$), of finite volume and satisfying
$$
\|\widetilde{F}^*\widetilde{h}_{ij}-h_{ij}\|_{C^{q_0}
(\mathcal{H}_a)}\leq\delta_0,
$$
then there exists an isometry $I$ of $\mathcal{H}$ to
$\widetilde{\mathcal{H}}$ such that
$$
d_{C^{l_0}(\mathcal{H}_a)}(\widetilde{F},I)<\epsilon_0.
$$
By Lemma 5.3.8 we further require $q_0$ and $\delta_0$ to
guarantee the existence of a harmonic diffeomorphism from
$(\mathcal{H}_a,\widetilde{g}_{ij})$ to $(\mathcal{H}_a,h_{ij})$
for any metric $\widetilde{g}_{ij}$ on $\mathcal{H}_a$ with
$\|\widetilde{g}_{ij}-h_{ij}\|_{C^{q_0}(\mathcal{H}_a)}\leq
\delta_0$.

Let $x^{\alpha}\in M_{\rm
thick}(w,t^{\alpha}),t^{\alpha}\rightarrow+\infty$, be a sequence
of points such that the scalings of $g_{ij}(t^{\alpha})$ around
$x^{\alpha}$ with factor $(t^{\alpha})^{-1}$ converge to $h_{ij}$.
Then there exist a marked point $x^{\infty}\in \mathcal{H}_a$ and
a sequence of diffeomorphisms $F_{\alpha}$ from $\mathcal{H}_a$
into $M_{t^{\alpha}}$ such that
$F_{\alpha}(x^{\infty})=x^{\alpha}$ and
$$
\|(t^{\alpha})^{-1}F^*_{\alpha}g_{ij}(t^{\alpha})
-h_{ij}\|_{C^m(\mathcal{H}_a)}\rightarrow 0
$$
as $\alpha\rightarrow\infty$ for all positive integers $m$. By
applying Lemma 5.3.8 and the implicit function theorem, we can
change $F_{\alpha}$ by an amount which goes to zero as
$\alpha\rightarrow\infty$ so as to make $F_{\alpha}$ a harmonic
diffeomorphism taking $\partial \mathcal{H}_a$ to a constant mean
curvature hypersurface $F_{\alpha}(\partial \mathcal{H}_a)$ of
$(M_{t^{\alpha}},(t^{\alpha})^{-1}g_{ij}(t^{\alpha}))$ with the
area $a$ and satisfying the free boundary condition that the
normal derivative of $F_{\alpha}$ at the boundary of the domain is
normal to the boundary of the target; and by combining with Lemma
7.6.2 (iii), we can smoothly continue each harmonic diffeomorphism
$F_{\alpha}$ forward in time a little to a family of harmonic
diffeomorphisms $F_{\alpha}(t)$ from $\mathcal{H}_a$ into $M_t$
with the metric $t^{-1}g_{ij}(t)$, with
$F_{\alpha}(t^{\alpha})=F_{\alpha}$ and with the time $t$ slightly
larger than $t^{\alpha}$, where $F_{\alpha}(t)$ takes $\partial
\mathcal{H}_a$ into a constant mean curvature hypersurface of
$(M_t,t^{-1}g_{ij}(t))$ with the area $a$ and also satisfies the
free boundary condition. Moreover, since the surgeries do not take
place at the points where the scalar curvature is negative, by the
same argument as in Theorem 5.3.4 for an arbitrarily given
positive integer $q\geq q_0$, positive number $\delta<\delta_0$,
and sufficiently large $\alpha$, we can ensure the extension
$F_{\alpha}(t)$ satisfies
$\|t^{-1}F^*_{\alpha}(t)g_{ij}(t)-h_{ij}\|_{C^q(\mathcal{H}_a)}\leq
\delta$ on a maximal time interval $t^{\alpha}\leq t\leq
\omega^{\alpha}$ (or $t^{\alpha}\leq t< \omega^{\alpha}$ when
$\omega^{\alpha}=+\infty$), and with
$\|(\omega^{\alpha})^{-1}F^*_{\alpha}
(\omega^{\alpha})g_{ij}(\omega^{\alpha})-h_{ij}\|_{C^q(\mathcal{H}_a)}=
\delta$, when $\omega^{\alpha}<+\infty$. Here we have implicitly
used the fact that $F_{\alpha}(\omega^{\alpha})(\partial
\mathcal{H}_a)$ is still strictly concave to ensure the map
$F_{\alpha}(\omega^{\alpha})$ is diffeomorphic.

We further claim that there must be some $\alpha$ such that
$\omega^{\alpha}=+\infty$; in other words, at least one hyperbolic
piece persists. Indeed, suppose that for each large enough
$\alpha$ we can only continue the family $F_{\alpha}(t)$ on a
finite interval $t^{\alpha}\leq t\leq \omega^{\alpha}<+\infty$
with
$$
\|(\omega^{\alpha})^{-1}F^*_{\alpha}(\omega^{\alpha})g_{ij}
(\omega^{\alpha})-h_{ij}\|_{C^q(\mathcal{H}_a)}=\delta.
$$
Consider the new sequence of manifolds
($M_{\omega^{\alpha}},g_{ij}(\omega^{\alpha})$). Clearly by Lemma
7.6.1, the scalings of $g_{ij}(\omega^{\alpha})$ around the new
origins $F_{\alpha}(\omega^{\alpha})(x^{\infty})$ with factor
$(\omega^{\alpha})^{-1}$ converge smoothly (by passing to a
subsequence) to a complete noncompact hyperbolic three-manifold
$\widetilde{\mathcal{H}}$ with the metric $\widetilde{h}_{ij}$ and
the origin $\widetilde{x}^{\infty}$ and with finite volume. By the
choice of the old limit $\mathcal{H}$, the new limit
$\widetilde{\mathcal{H}}$ has at least as many cusps as
$\mathcal{H}$. By the definition of convergence, we can find a
sequence of compact subsets $\widetilde{U}_{\alpha}$ exhausting
$\widetilde{\mathcal{H}}$ and containing $\widetilde{x}^{\infty}$,
and a sequence of diffeomorphisms $\widetilde{F}_{\alpha}$ of
neighborhood of $\widetilde{U}_{\alpha}$ into
$M_{\omega^{\alpha}}$ with
$\widetilde{F}_{\alpha}(\widetilde{x}^{\infty})
=F_{\alpha}(\omega^{\alpha})(x^{\infty})$ such that for each
compact subset $\widetilde{U}$ of $\widetilde{\mathcal{H}}$ and
each integer $m$,
$$
\|(\omega^{\alpha})^{-1}\widetilde{F}^*_{\alpha}(g_{ij}
(\omega^{\alpha}))-\widetilde{h}_{ij}
\|_{C^m(\widetilde{U})}\rightarrow 0
$$
as $\alpha\rightarrow+\infty$. Thus for sufficiently large
$\alpha$, we have the map
$$
G_{\alpha}=\widetilde{F}^{-1}_{\alpha}\circ
F_{\alpha}(\omega^{\alpha}):\ \ \mathcal{H}_a\rightarrow
\widetilde{\mathcal{H}}
$$
such that
$$
\|G^*_{\alpha}\widetilde{h}_{ij}-h_{ij}
\|_{C^q(\mathcal{H}_a)}<\widetilde{\delta}
$$
for any fixed $\widetilde{\delta}>\delta$. Then a subsequence of
$G_{\alpha}$ converges at least in the $C^{q-1}(\mathcal{H}_a)$
topology to a map $G_{\infty}$ of $\mathcal{H}_a$ into
$\widetilde{\mathcal{H}}$ which is a harmonic map from
$\mathcal{H}_a$ into $\widetilde{\mathcal{H}}$ and takes $\partial
\mathcal{H}_a$ to a constant mean curvature hypersurface
$G_{\infty}(\partial \mathcal{H}_a)$ of
$(\widetilde{\mathcal{H}},\widetilde{h}_{ij})$ with the area $a$,
as well as satisfies the free boundary condition. Clearly,
$G_{\infty}$ is at least a local diffeomorphism. Since
$G_{\infty}$ is the limit of diffeomorphisms, the only possibility
of overlap is at the boundary. Note that $G_{\infty}(\partial
\mathcal{H}_a)$ is still strictly concave. So $G_{\infty}$ is
still a diffeomorphism. Moreover by using the standard regularity
result of elliptic partial differential equations (see for example
\cite{GT}), we also have \begin{equation}
\|G^*_{\infty}\widetilde{h}_{ij}-h_{ij}
\|_{C^q(\mathcal{H}_a)}=\delta. %\eqno(7.6.16)      WRONG
\end{equation} Now by Lemma 5.3.7 we deduce that there exists an isometry $I$
of $\mathcal{H}$ to $\widetilde{\mathcal{H}}$ with
$$
d_{C^{l_0}(\mathcal{H}_a)}(G_{\infty},I)<\epsilon_0.
$$
Thus $I^{-1}\circ G_{\infty} $ is a harmonic diffeomorphism of
$\mathcal{H}_a$ to itself which satisfies the free boundary
condition and
$$
d_{C^{l_0}(\mathcal{H}_a)}(I^{-1}\circ G_{\infty},Id)<\epsilon_0.
$$
However the uniqueness in Lemma 5.3.8 concludes that $I^{-1}\circ
G_{\infty}=Id$ which contradicts (7.6.20). So we have shown that
at least one hyperbolic piece persists and the metric
$t^{-1}F^*_{\alpha}(t)g_{ij}(t)$, for $\omega^{\alpha}\leq
t<\infty$, is as close to the hyperbolic metric $h_{ij}$ as we
like.

We can continue to form other persistent hyperbolic pieces in the
same way as long as there is a sequence of points $y^{\beta}\in
M_{\rm thick}(w,t^{\beta}),t^{\beta}\rightarrow+\infty$, lying
outside the chosen pieces. Note that
$V(t)(t+\frac{3}{2})^{-\frac{3}{2}}$ is nonincreasing on
$[0,+\infty)$. Therefore by combining with Margulis lemma (see for
example \cite{Gro79} or \cite{KM}), we have proved that there
exists a finite collection of complete noncompact hyperbolic
three-manifolds $\mathcal{H}_1,\ldots,\mathcal{H}_m$ with finite
volume, a small number $a>0$ and a time $T<+\infty$ such that for
all $t$ beyond $T$ we can find diffeomorphisms $\varphi_l(t)$ of
$(\mathcal{H}_l)_a$ into $M_t$, $1\leq l\leq m$, so that as long
as $t$ is sufficiently large, the metric
$t^{-1}\varphi^*_l(t)g_{ij}(t)$ is as close to the hyperbolic
metrics as we like and the complement $M_t\backslash
(\varphi_1(t)((\mathcal{H}_1)_a)\cup\cdots\cup
\varphi_m(t)((\mathcal{H}_m)_a))$ is contained in the thin part
$M_{\rm thin}(w,t)$.

It remains to show the boundary tori of any persistent hyperbolic
piece are incompressible. Let $B$ be a small positive number and
assume the above positive number $a$ is much smaller than $B$. Let
$M_a(t)=\varphi_l(t)((\mathcal{H}_l)_a)$ ($1\leq l\leq m$) be such
a persistent hyperbolic piece of the manifold $M_t$ truncated by
boundary tori of area $at$ with constant mean curvature, and
denote by $M^c_a(t)=M_t\setminus \stackrel{\circ}{M}_a(t)$ the
part of $M_t$ exterior to $M_a(t)$. Thus there exists a family of
subsets $M_B(t)\subset M_a(t)$ which is a persistent hyperbolic
piece of the manifold $M_t$ truncated by boundary tori of area
$Bt$ with constant mean curvature. We also denote by
$M^c_B(t)=M_t\setminus \stackrel{\circ}{M}_B(t)$. By Van Kampen's
Theorem, if $\pi_1(\partial M_B(t))$ injects into
$\pi_1(M_B^c(t))$ then it injects into $\pi_1(M_t)$ also. Thus we
only need to show $\pi_1(\partial M_B(t))$ injects into
$\pi_1(M_B^c(t))$.

As before we will use a contradiction argument to show
$\pi_1(\partial M_B(t))$ injects into $\pi_1(M^c_B(t))$. Let $T$
be a torus in $\partial M_B(t)$ and suppose $\pi_1(T)$ does not
inject into $\pi_1(M^c_B(t))$. By Dehn's Lemma we know that the
kernel is a cyclic subgroup of $\pi_1(T)$ generated by a primitive
element. Consider the normalized metric
$\widetilde{g}_{ij}(t)=t^{-1}g_{ij}(t)$ on $M_t$. Then by the work
of Meeks-Yau \cite{MY} or Meeks-Simon-Yau \cite{MSY}, we know that
among all disks in $M^c_B(t)$ whose boundary curve lies in $T$ and
generates the kernel of $\pi_1(T)$, there is a smooth embedded
disk normal to the boundary which has the least possible area
(with respect to the normalized metric $\widetilde{g}_{ij}(t)$).
Denote by $D$ the minimal disk and
$\widetilde{A}=\widetilde{A}(t)$ its area. We will show that
$\widetilde{A}(t)$ decreases at a certain rate which will arrive
at a contradiction.

We first consider the case that there exist no surgeries at the
time $t$. Exactly as in Part III of the proof of Theorem 5.3.4,
the change of the area $\widetilde{A}(t)$ comes from the change in
the metric and the change in the boundary. For the change in the
metric, we choose an orthonormal frame $X,Y,Z$ at a point $x$ in
the disk $D$ so that $X$ and $Y$ are tangent to the disk $D$ while
$Z$ is normal. Since the normalized metric $\widetilde{g}_{ij}$
evolves by
$$
\frac{\partial}{\partial t}\widetilde{g}_{ij}
=-t^{-1}(\widetilde{g}_{ij}+2\widetilde{R}_{ij}),
$$
the (normalized) area element $d\widetilde{\sigma}$ of the disk
$D$ around $x$ satisfies
$$\frac{\partial}{\partial t}d\widetilde{\sigma}
=-t^{-1}(1+\widetilde{\Ric}(X,X)+\widetilde{\Ric}(Y,Y))
d\widetilde{\sigma}.
$$
For the change in the boundary, we notice that the tensor
$\widetilde{g}_{ij}+2\widetilde{R}_{ij}$ is very small for the
persistent hyperbolic piece. Then by using the Gauss-Bonnet
theorem as before, we obtain the rate of change of the area \begin{equation}
\frac{d\widetilde{A}}{dt}\leq
-\int_D\left(\frac{1}{t}+\frac{\widetilde{R}}{2t}\right)
d\widetilde{\sigma}+\frac{1}{t}\int_{\partial
D}\widetilde{k}d\widetilde{s}-\frac{2\pi}{t}+
o\left(\frac{1}{t}\right)\widetilde{L}, %\eqno (7.6.17)-->(7.6.21) WRONG
\end{equation} where $\widetilde{k}$ is the geodesic curvature of the
boundary and $\widetilde{L}$ is the length of the boundary curve
$\partial D$ (with respect to the normalized metric
$\widetilde{g}_{ij}(t)$).  Since $\widetilde{R}\geq
-{3t}/{2(t+\frac{3}{2})}$ for all $t\geq 0$ by (7.6.2), the first
term on the RHS of (7.6.21) is bounded above by
$$
-\int_D\left(\frac{1}{t}+\frac{\widetilde{R}}{2t}\right)d\widetilde{\sigma}\leq
-\frac{1}{t}\left(\frac{1}{4}- o(1)\right)\widetilde{A};
$$
while the second term on the RHS of (7.6.21) can be estimated
exactly as before by
$$
\frac{1}{t}\int_{\partial D}\widetilde{k}d\widetilde{s}\leq
\frac{1}{t}\left(\frac{1}{4}+o(1)\right)\widetilde{L}.
$$
Thus we obtain \begin{equation} \frac{d\widetilde{A}}{dt}\leq
\frac{1}{t}\left[\left(\frac{1}{4}+o(1)\right)
\widetilde{L}-\left(\frac{1}{4}-o(1)\right)\widetilde{A}-2\pi\right].
%\eqno (7.6.18)                     WRONG
\end{equation}

Next we show that these arguments also work for the case that
there exist surgeries at the time $t$. To this end, we only need
to check that the embedded minimal disk $D$ lies in the region
which is unaffected by surgery. Our surgeries for the irreducible
three-manifold took place on $\delta$-necks in
$\varepsilon$-horns, where the scalar curvatures are at least
$\delta^{-2} (\widetilde{r}(t))^{-1}$, and the components with
nonnegative scalar curvature have been removed. So the hyperbolic
piece is not affected by the surgeries. In particular, the
boundary $\partial D$ is unaffected by the surgeries. Thus if
surgeries occur on the minimal disk, the minimal disk has to pass
through a long thin neck before it reaches the surgery regions.
Look at the intersections of the embedding minimal disk with a
generic center two-sphere $\mathbb{S}^2$ of the long thin neck;
these are circles. Since the two-sphere $\mathbb{S}^2$ is simply
connected, we can replace the components of the minimal disk $D$
outside the center two-sphere $\mathbb{S}^2$ by some corresponding
components on the center two-sphere $\mathbb{S}^2$ to form a new
disk which also has $\partial D$ as its boundary. Since the metric
on the long thin neck is nearly a product metric, we could choose
the generic center two-sphere $\mathbb{S}^2$ properly so that the
area of the new disk is strictly less than the area of the
original disk $D$. This contradiction proves the minimal disk lies
entirely in the region unaffected by surgery.

Since $a$ is much smaller than $B$, the region within a long
distance from $\partial M_B(t)$ into $M^c_B(t)$ will look nearly
like a hyperbolic cusplike collar and is unaffected by the
surgeries. So we can repeat the arguments in the last part of the
proof of Theorem 5.3.4 to bound the length $\widetilde{L}$ by the
area $\widetilde{A}$ and to conclude
$$\frac{d\widetilde{A}}{dt}\leq -\frac{\pi}{t}$$ for all
sufficiently large times $t$, which is impossible because the RHS
is not integrable. This proves that the boundary tori of any
persistent hyperbolic piece are incompressible.

Therefore we have proved the theorem.
%$$\eqno \#$$ \vskip 1cm
\end{proof}


\section{Geometrization of Three-manifolds}

In the late 70's and early 80's, Thurston \cite{Th82, Th86, Th88}
proved a number of remarkable results on the existence of
geometric structures on a class of three-manifolds:
\textbf{Haken}\index{Haken} manifolds (i.e. each of them contains
an incompressible surface of genus $\geq$ 1). These results
motivated him to formulate a profound conjecture which roughly
says every compact three-manifold admits a canonical decomposition
into domains, each of which has a canonical geometric structure.
To give a detailed description of the conjecture, we recall some
terminology as follows.

An $n$-dimensional complete Riemannian manifold $(M,g)$ is called
a \textbf{homogeneous manifold}\index{homogeneous manifold} if its
group of isometries acts transitively on the manifold. This means
that the homogeneous manifold looks the same metrically at
everypoint. For example, the round $n$-sphere $\mathbb{S}^n$, the
Euclidean space $\mathbb{R}^n$ and the standard hyperbolic space
$\mathbb{H}^n$ are homogeneous manifolds.  A Riemannian manifold
is said to \textbf{be modeled}\index{be modeled} on a given
homogeneous manifold $(M,g)$ if every point of the manifold has a
neighborhood isometric to an open set of $(M,g)$. And an
$n$-dimensional Riemannian manifold is called a \textbf{locally
homogeneous manifold}\index{homogeneous manifold!locally} if it is
complete and is modeled on a homogeneous manifold. By a theorem of
Singer \cite{Sin}, the universal cover of a locally homogeneous
manifold (with the pull-back metric) is a homogeneous manifold.

In dimension three, every locally homogeneous manifold with finite
volume is modeled on one of the following eight homogeneous
manifolds (see for example Theorem 3.8.4 of \cite{Th97}):
\begin{itemize}
\item[(1)] $\mathbb{S}^3$, the round three-sphere; \item[(2)]
$\mathbb{R}^3$, the Euclidean space ; \item[(3)] $\mathbb{H}^3$,
the standard hyperbolic space; \item[(4)] $\mathbb{S}^2 \times
\mathbb{R}$; \item[(5)] $\mathbb{H}^2 \times \mathbb{R}$;
\item[(6)] $Nil$, the three-dimensional nilpotent Heisenberg group
(consisting of upper triangular $3\times 3$ matrices with diagonal
entries 1); \item[(7)] $\widetilde{PSL}(2,\mathbb{R})$, the
universal cover of the unit sphere bundle of $\mathbb{H}^2$;
\item[(8)] $Sol$, the three-dimensional solvable Lie group.
\end{itemize}

\medskip
A three-manifold $M$ is called \textbf{prime}\index{prime} if it
is not diffeomorphic to $\mathbb{S}^3$ and if every (topological)
$\mathbb{S}^2 \subset M$, which separates $M$ into two pieces, has
the property that one of the two pieces is diffeomorphic to a
three-ball. Recall that a three-manifold is irreducible if every
embedded two-sphere bounds a three-ball in the manifold. Clearly
an irreducible three-manifold is either prime or is diffeomorphic
to $\mathbb{S}^3$. Conversely, an orientable prime three-manifold
is either irreducible or is diffeomorphic to $\mathbb{S}^2 \times
\mathbb{S}^1$ (see for example \cite{Hempel}). One of the first
results in three-manifold topology is the following prime
decomposition theorem obtained by Kneser \cite{Kn} in 1929 (see
also Theorem 3.15 of \cite{Hempel}).

\medskip
{\bf Prime Decomposition Theorem. }\index{prime! decomposition
theorem} \emph{ Every compact orientable three-manifold admits a
decomposition as a finite connected sum of orientable prime
three-manifolds. }

\medskip
In \cite{Mil62}, Milnor showed that the factors involved in the
above Prime Decomposition are unique.  Based on the prime
decomposition, the question about topology of compact orientable
three-manifolds is reduced to the question about prime
three-manifolds. Thurston's Geometrization Conjecture is about
prime three-manifolds.

\medskip
{\bf Thurston's Geometrization Conjecture. }\index{Thurston's
geometrization conjecture} Let $M$ be a compact, orientable and
prime three-manifold. Then there is an embedding of a finite
number of disjoint unions, possibly empty, of incompressible
two-tori $\coprod_{i} T_i^2 \subset M$ such that every component
of the complement admits a locally homogeneous Riemannian metric
of finite volume.

\medskip
We remark that the existence of a torus decomposition, also called
JSJ-decomposition, was already obtained by Jaco-Shalen \cite{Jaco}
and Johannsen \cite{Joha}. The JSJ-decomposition states that any
compact, orientable, and prime three-manifold has a finite
collection, possibly empty, of disjoint incompressible embedding
two-tori $\{T_i^2\}$ which separate the manifold into a finite
collection of compact three-manifolds (with toral boundary), each
of which is either a graph manifold or is \textbf{atoroidal}
\index{atoroidal} in the sense that any subgroup of its
fundamental group isomorphic to $\mathbb{Z}\times\mathbb{Z}$ is
conjugate into the fundamental group of some component of its
boundary. A compact three-manifold $X$, possibly with boundary, is
called a \textbf{graph manifold}\index{graph manifold} if there is
a finite collection of disjoint embedded tori $T_i \subset X$ such
that each component of $X \setminus \bigcup T_i$ is an
$\mathbb{S}^1$ bundle over a surface. Thus the point of the
conjecture is that the components should all be geometric.

The geometrization conjecture for a general compact orientable
3-man\-ifold is the statement that each of its prime factors
satisfies the above conjecture. We say a compact orientable
three-manifold is \textbf{geometrizable}\index{geometrizable} if
it satisfies the geometric conjecture.

We also remark, as is well-known, that the Poincar\'{e} conjecture
can be deduced from Thur\-ston's geometrization conjecture.
Indeed, suppose that we have a compact simply connected
three-manifold that satisfies the conclusion of the geometrization
conjecture. If it were not diffeomorphic to the three-sphere
$\mathbb{S}^3$, there would be a prime factor in the prime
decomposition of the manifold. Since the prime factor still has
vanishing fundamental group, the (torus) decomposition of the
prime factor in the geometrization conjecture must be trivial.
Thus the prime factor is a compact homogeneous manifold model.
>From the list of above eight models, we see that the only compact
three-dimensional model is $\mathbb{S}^3$. This is a
contradiction. Consequently, the compact simply connected
three-manifold is diffeomorphic to $\mathbb{S}^3$.

Now, based on the long-time behavior result (Theorem 7.6.4) we
apply the Ricci flow to discuss Thurston's geometrization
conjecture. Let $M$ be a compact, orientable and prime
three-manifold. Since a prime orientable three-manifold is either
irreducible or is diffeomorphic to $\mathbb{S}^2 \times
\mathbb{S}^1$, we may thus assume the manifold $M$ is irreducible
also. Arbitrarily given a (normalized) Riemannian metric for the
manifold $M$, we use it as initial data to evolve the metric by
the Ricci flow with surgery. From Theorem 7.4.3, we know that the
Ricci flow with surgery has a long-time solution on a maximal time
interval $[0,T)$ which satisfies the a priori assumptions and has
a finite number of surgeries on each finite time interval.
Furthermore, from the long-time behavior theorem (Theorem 7.6.4),
we have well-understood geometric structures on the thick part.
Whereas, to understand the thin part, Perelman announced the
following result in \cite{P2}.

\medskip
{\bf Perelman's Claim}\index{Perelman's claim} (cf. Theorem 7.4 of
\cite{P2}){\bf .} \ Suppose $(M^{\alpha},g^{\alpha}_{ij})$ is a
sequence of compact orientable three-manifolds, closed or with
convex boundary, and $w^{\alpha}\rightarrow  0$. Assume that
\begin{itemize}
\item[(1)] for each point $x \in M^{\alpha}$ there exists a radius
$\rho = \rho^{\alpha}(x), 0<\rho<1,$ not exceeding the diameter of
the manifold, such that the ball $B(x,\rho)$ in the metric
$g^{\alpha}_{ij}$ has volume at most $w^{\alpha}\rho^3$ and
sectional curvatures at least $-\rho^{-2}$; \item[(2)] each
component of the boundary of $M^{\alpha}$ has diameter at most
$w^{\alpha}$, and has a (topologically trivial) collar of length
one, where the sectional curvatures are between $-1/4 - \epsilon$
and $-1/4 + \epsilon$.
\end{itemize}

Then $M^{\alpha}$ for sufficiently large $\alpha$ are
diffeomorphic to graph manifolds.

\medskip
The topology of graph manifolds is well understood; in particular,
every graph manifold is geometrizable (see \cite{W}).

The proof of Perelman's Claim promised in \cite{P2} is still not
available in literature. Nevertheless, recently in \cite{ShY},
Shioya and Yamaguchi provided a proof of Perelman's Claim for the
special case when all the manifolds $(M^{\alpha},g^{\alpha}_{ij})$
are closed. That is, they proved the following weaker assertion.

\medskip
{\bf Weaker Assertion} (Theorem 8.1 of Shioya-Yamaguchi
\cite{ShY}){\bf .} \index{Weaker Assertion} \ Suppose
$(M^{\alpha},g^{\alpha}_{ij})$ is a sequence of compact orientable
three-manifolds without boundary, and $w^{\alpha}\rightarrow 0$.
Assume that for each point $x \in M^{\alpha}$ there exists a
radius $\rho = \rho^{\alpha}(x),$ not exceeding the diameter of
the manifold, such that the ball $B(x,\rho)$ in the metric
$g^{\alpha}_{ij}$ has volume at most $w^{\alpha}\rho^3$ and
sectional curvatures at least $-\rho^{-2}$. Then $M^{\alpha}$ for
sufficiently large $\alpha$ are diffeomorphic to graph manifolds.

\medskip
Based on the the long-time behavior theorem (Theorem 7.6.4) and
using the above Weaker Assertion, we can now give a proof for
Thurston's geometrization conjecture \index{Thurston's
geometrization conjecture}. We remark that if we assume the above
Perelman's Claim, then we does not need to use Thurston's theorem
for Haken manifolds in the proof of Theorem 7.7.1.

%\vskip 0.2cm\noindent{\bf Theorem 7.7.1}
%CZ-SOURCE-LINE-23041
\begin{theorem}[geometrization of compact orientable three-manifolds]
  \label{thm:cz-thurston-geometrization}
  \dcref{Cao--Zhu Theorem 7.7.1}
  \group{Geometrization}
  \source{cao-zhu}

Thurston's geometrization conjecture is true.
\end{theorem}

%\vskip 0.1cm\noindent \textbf{Proof.}
\begin{proof}
  \uses{thm:cz-shioya-yamaguchi-closed-collapsing-graph-manifolds}
  \uses{thm:cz-ricci-flow-surgery-long-time-existence, thm:cz-surgery-long-time-geometry, thm:cz-surgery-thick-thin-decomposition, thm:cz-large-time-surgery-curvature-estimate, cor:cz-localized-volume-sectional-curvature-control, thm:cz-almost-euclidean-volume-curvature-control, thm:cz-lower-volume-ratio-curvature-control}
Let $M$ be a compact, orientable, and prime three-manifold
(without boundary). Without loss of generality, we may assume that
the manifold $M$ is irreducible also.

Recall that the theorem of Thurston (see for example Theorem A and
Theorem B in the third section of \cite{Morgan}, see also
\cite{Mc} and \cite{Ot})) says that any compact, orientable, and
irreducible Haken three-manifold (with or without boundary) is
geometrizable. Thus in the following, we may assume that the
compact three-manifold $M$ (without boundary) is atoroidal, and
then the fundamental group $\pi_1(M)$ contains no noncyclic,
abelian subgroup.

Arbitrarily given a (normalized) Riemannian metric on the manifold
$M$, we use it as initial data for the Ricci flow. Arbitrarily
take a sequence of small positive constants $w^{\alpha}\rightarrow
0$ as $\alpha \rightarrow +\infty$. For each fixed $\alpha$, we
set $\varepsilon={w^{\alpha}}/{2}>0$. Then by Theorem 7.4.3, the
Ricci flow with surgery has a long-time solution
$(M^{\alpha}_t,g^{\alpha}_{ij}(t))$ on a maximal time interval
$[0,T^{\alpha})$, which satisfies the a priori assumptions (with
the accuracy parameter $\varepsilon={w^{\alpha}}/{2}$) and has a
finite number of surgeries on each finite time interval. Since the
initial manifold is irreducible, by the surgery procedure, we know
that for each $\alpha$ and each $t>0$ the solution manifold
$M^{\alpha}_t$ consists of a finite number of components where the
essential component $(M^{\alpha}_t)^{(1)}$ is diffeomorphic to the
initial manifold $M$ and the others are diffeomorphic to the
three-sphere $\mathbb{S}^3$.

If for some $\alpha=\alpha_0$ the maximal time $T^{\alpha_0}$ is
finite, then the solution $(M^{\alpha_0}_t,g^{\alpha_0}_{ij}(t))$
becomes extinct at $T^{\alpha_0}$ and the (irreducible) initial
manifold $M$ is diffeomorphic to $\mathbb{S}^3/\Gamma$ (the metric
quotients of round three-sphere); in particular, the manifold $M$
is geometrizable. Thus we may assume that the maximal time
$T^{\alpha} = +\infty$ for all $\alpha$.

We now apply the long-time behavior theorem (Theorem 7.6.4). If
there is some $\alpha$ such that case (ii) of Theorem 7.6.4
occurs, then for some sufficiently large time $t$, the essential
component $(M^{\alpha}_t)^{(1)}$ of the solution manifold
$M^{\alpha}_t$ is diffeomorphic to a compact hyperbolic space, so
the initial manifold $M$ is geometrizable. Whereas if there is
some sufficiently large $\alpha$ such that case (iii) of Theorem
7.6.4  occurs, then it follows that for all sufficiently large
$t$, there is an embedding of a (nonempty) finite number of
disjoint unions of incompressible two-tori $\coprod_{i} T_i^2 $ in
the essential component $(M^{\alpha}_t)^{(1)}$ of $M^{\alpha}_t$.
This is a contradiction since we have assumed the initial manifold
$M$ is atoroidal.

It remains to deal with the situation that there is a sequence of
positive $\alpha_k \rightarrow +\infty$ such that the solutions
$(M^{\alpha_k}_t,g^{\alpha_k}_{ij}(t))$ always satisfy case (i) of
Theorem 7.6.4. That is, for each $\alpha_k$, $M^{\alpha_k}_t =
M_{\rm thin}(w^{\alpha_k},t)$ when the time $t$ is sufficiently
large. By the Thick-thin decomposition theorem (Theorem 7.6.3),
there is a positive constant, $0<\rho(w^{\alpha_k}) \leq 1$, such
that as long as $t$ is sufficiently large, for every $x\in
M^{\alpha_k}_t = M_{\rm thin}(w^{\alpha_k},t)$, we have some
$r=r(x,t)$, with $0<r\sqrt{t}<\rho(w^{\alpha_k})\sqrt{t}$, such
that \begin{equation} Rm\geq -(r\sqrt{t})^{-2} \quad \mbox{ on } \;
B_t(x,r\sqrt{t}),
%\eqno (7.7.1)
\end{equation} and \begin{equation}
\Vol_t(B_t(x,r\sqrt{t}))<w^{\alpha_k}(r\sqrt{t})^3. %\eqno (7.7.2)
\end{equation}

Clearly we only need to consider the essential component
$(M^{\alpha_k}_{t})^{(1)}$.  We divide the discussion into the
following two cases:

\smallskip
(1) there is a positive constant $1<C<+\infty$ such that for each
$\alpha_k$ there is a sufficiently large time $t_k>0$ such that
\begin{equation} r(x,t_k)\sqrt{t_k} < C \cdot {\rm
diam}\,\left((M^{\alpha_k}_{t_k})^{(1)}\right)
%\eqno (7.7.3)
\end{equation} for all $x\in (M^{\alpha_k}_{t_k})^{(1)} \subset M_{\rm
thin}(w^{\alpha_k},{t_k})$;

\smallskip
(2) there are a subsequence $\alpha_k$ (still denoted by
$\alpha_k$), and sequences of positive constants $C_k \rightarrow
+\infty$ and times $T_k < +\infty$ such that for each $t \geq
T_k$, we have \begin{equation} r(x(t),t)\sqrt{t} \geq C_k \cdot {\rm diam}\,
\left((M^{\alpha_k}_{t})^{(1)}\right)
%\eqno (7.7.4)
\end{equation} for some $x(t) \in (M^{\alpha_k}_{t})^{(1)}$, $k=1,2,\ldots.$
Here we denote by ${\rm diam}\,((M^{\alpha}_t)^{(1)})$ the
diameter of the essential component $(M^{\alpha}_t)^{(1)}$ with
the metric $g^{\alpha}_{ij}(t)$ at the time $t$.

Let us first consider case (1). For each point $x \in
(M^{\alpha_k}_{t_k})^{(1)} \subset M_{\rm thin}(w^{\alpha_k},$
${t_k})$, we denote by $\rho_k(x) = C^{-1}r(x,t_k)\sqrt{t_k}$.
Then by (7.7.1), (7.7.2) and (7.7.3), we have
$$
\rho_k(x) <{\rm diam}\,\left((M^{\alpha_k}_{t_k})^{(1)}\right),
$$
$$
\Vol_{t_k}(B_{t_k}(x,\rho_k(x))) \leq
\Vol_{t_k}(B_{t_k}(x,r(x,t_k)\sqrt{t_k}))<C^3w^{\alpha_k}(\rho_k(x))^3,
$$
and
$$
Rm\geq -(r(x,t_k)\sqrt{t_k})^{-2} \geq -(\rho_k(x))^{-2}
$$
on  $B_{t_k}(x,\rho_k(x)).$ Then it follows from the above Weaker
Assertion that $(M^{\alpha_k}_{t_k})^{(1)}$, for sufficiently
large $k$, are diffeomorphic to graph manifolds. This implies that
the (irreducible) initial manifold $M$ is diffeomorphic to a graph
manifold. So the manifold $M$ is geometrizable in case (1).

We next consider case (2). Clearly, for each $\alpha_k$ and the
chosen $T_k$, we may assume that the estimates (7.7.1) and (7.7.2)
hold for all $t \geq T_k$ and $x\in (M^{\alpha_k}_t)^{(1)}$. The
combination of (7.7.1) and (7.7.4) gives \begin{equation} Rm \geq
-C^{-2}_k({\rm diam}\,((M^{\alpha_k}_{t})^{(1)}))^{-2}
\quad \mbox{ on } \; (M^{\alpha_k}_{t})^{(1)}, %\eqno (7.7.5)
\end{equation} for all $t \geq T_k$. If there are a subsequence $\alpha_k$
(still denoted by $\alpha_k$) and a sequence of times $t_k \in
(T_k,+\infty)$ such that \begin{equation}
\Vol_{t_k}((M^{\alpha_k}_{t_k})^{(1)}) <
w'_k({\rm diam}\,((M^{\alpha_k}_{t_k})^{(1)}))^3 %\eqno (7.7.6)
\end{equation} for some sequence $w'_k \rightarrow 0$, then it follows from
the Weaker Assertion that $(M^{\alpha_k}_{t_k})^{(1)}$, for
sufficiently large $k$, are diffeomorphic to graph manifolds which
implies the initial manifold $M$ is geometrizable. Thus we may
assume that there is a positive constant $w'$ such that \begin{equation}
\Vol_t((M^{\alpha_k}_t)^{(1)})
\geq w'({\rm diam}\,((M^{\alpha_k}_t)^{(1)}))^3 %\eqno (7.7.7)
\end{equation} for each $k$ and all $t \geq T_k$.

In view of the estimates (7.7.5) and (7.7.7), we now want to use
Theorem 7.5.2 to get a uniform upper bound  for the curvatures of
the essential components
$((M^{\alpha_k}_t)^{(1)},g^{\alpha_k}_{ij}(t))$ with sufficiently
large time $t$. Note that the estimate in Theorem 7.5.2 depends on
the parameter $\varepsilon$ and our $\varepsilon$'s depend on
$w^{\alpha_k}$ with $0<\varepsilon = {w^{\alpha_k}}/{2}$; so it
does not work in the present situation. Fortunately we notice that
the curvature estimate for smooth solutions in Corollary 7.2.3 is
independent of $\varepsilon$. In the following we try to use
Corollary 7.2.3 to obtain the desired curvature estimate.

We first claim that for each $k$, there is a sufficiently large
$T'_k \in (T_k,+\infty)$ such that the solution, when restricted
to the essential component $((M^{\alpha_k}_t)^{(1)},$
$g^{\alpha_k}_{ij}(t))$, has no surgery for all $t\geq T'_k$.
Indeed, for each fixed $k$, if there is a $\delta(t)$-cutoff
surgery at a sufficiently large time $t$, then the manifold
$((M^{\alpha_k}_t)^{(1)},g^{\alpha_k}_{ij}(t))$ would contain a
$\delta(t)$-neck $B_t(y,{\delta(t)}^{-1}R(y,t)^{-\frac{1}{2}})$
for some $y \in (M^{\alpha_k}_t)^{(1)}$ with the volume ratio \begin{equation}
\frac{\Vol_t(B_t(y,{\delta(t)}^{-1}
R(y,t)^{-\frac{1}{2}}))}{({\delta(t)}^{-1}R(y,t)^{-\frac{1}{2}})^3}
\leq 8\pi{\delta(t)}^2. %\eqno (7.7.8)
\end{equation} On the other hand, by (7.7.5) and (7.7.7), the standard
Bishop-Gromov volume comparison implies that
$$
\frac{\Vol_t(B_t(y,{\delta(t)}^{-1}
R(y,t)^{-\frac{1}{2}}))}{({\delta(t)}^{-1}R(y,t)^{-\frac{1}{2}})^3}
\geq c(w')
$$
for some positive constant $c(w')$ depending only on $w'$. Since
$\delta(t)$ is very small when $t$ is large, this arrives at a
contradiction with (7.7.8). So for each $k$, the essential
component $((M^{\alpha_k}_t)^{(1)},g^{\alpha_k}_{ij}(t))$ has no
surgery for all sufficiently large $t$.

For each $k$, we consider any fixed time $\tilde{t}_k > 3 T'_k$.
Let us scale the solution $g^{\alpha_k}_{ij}(t)$ on the essential
component $(M^{\alpha_k}_t)^{(1)}$ by
$$
\tilde{g}^{\alpha_k}_{ij}(\cdot,s) =
(\tilde{t}_{k})^{-1}g^{\alpha_k}_{ij}(\cdot,\tilde{t}_{k}s).
$$
Note that $(M^{\alpha_k}_t)^{(1)}$ is diffeomorphic to $M$ for all
$t$. By the above claim, we see that the rescaled solution
$(M,\tilde{g}^{\alpha_k}_{ij}(\cdot,s))$ is a smooth solution to
the Ricci flow on the time interval $s \in [\frac{1}{2},1]$. Set
$$
\tilde{r}_k = \left(\sqrt{\tilde{t}_k}\right)^{-1} {\rm
diam}\,\left((M^{\alpha_k}_{\tilde{t}_k})^{(1)}\right).
$$
Then by (7.7.4), (7.7.5) and (7.7.7), we have
$$
\tilde{r}_k \leq C^{-1}_k \rightarrow 0, \quad \mbox{ as }\; k
\rightarrow +\infty,
$$
$$
\widetilde{Rm} \geq -C^{-2}_k(\tilde{r}_k)^{-2}, \quad \mbox{on }
\; B_1(x(\tilde{t}_k),\tilde{r}_k),
$$
and
$$
\Vol_1(B_1(x(\tilde{t}_k),\tilde{r}_k)) \geq w'(\tilde{r}_k)^3,
$$
where $\widetilde{Rm}$ is the rescaled curvature, $x(\tilde{t}_k)$
is the point given by (7.7.4) and
$B_1(x(\tilde{t}_k),\tilde{r}_k)$ is the geodesic ball of rescaled
solution at the time $s=1$. Moreover, the closure of
$B_1(x(\tilde{t}_k),\tilde{r}_k)$ is the whole manifold
$(M,\tilde{g}^{\alpha_k}_{ij}(\cdot,1)).$

Note that in Theorem 7.2.1, Theorem 7.2.2 and Corollary 7.2.3, the
condition about normalized initial metrics is just to ensure that
the solutions satisfy the Hamilton-Ivey pinching estimate. Since
our solutions $(M^{\alpha_k}_t,g^{\alpha_k}_{ij}(t))$ have already
satisfied the pinching assumption, we can then apply Corollary
7.2.3 to conclude
$$
|\widetilde{Rm}(x,s)| \leq K(w')(\tilde{r}_k)^{-2},
$$
whenever $s \in [1- \tau(w')(\tilde{r}_k)^{2}, 1]$, $x \in
(M,\tilde{g}^{\alpha_k}_{ij}(\cdot,s))$ and $k$ is sufficiently
large. Here $K(w')$ and $\tau(w')$ are positive constants
depending only on $w'$. Equivalently, we have the curvature
estimates \begin{equation} |Rm(\cdot,t)| \leq K(w') ({\rm
diam}\,((M^{\alpha_k}_{\tilde{t}_k})^{(1)}))^{-2},
\quad \mbox{ on }\;  M, %\eqno (7.7.9)
\end{equation} whenever $t \in [\tilde{t}_k - \tau(w')({\rm
diam}\,((M^{\alpha_k}_{\tilde{t}_k})^{(1)}))^{2},\tilde{t}_k]$ and
$k$ is sufficiently large.

\smallskip\noindent
For each $k$, let us scale
$((M^{\alpha_k}_t)^{(1)},g^{\alpha_k}_{ij}(t))$ with the factor
$({\rm diam}((M^{\alpha_k}_{\tilde{t}_k})^{(1)}))^{-2}$ and shift
the time $\tilde{t}_k$ to the new time zero. By the curvature
estimate (7.7.9) and Hamilton's compactness theorem (Theorem
4.1.5), we can take a subsequential limit (in the $C^{\infty}$
topology) and get a smooth solution to the Ricci flow on $M \times
(-\tau(w'),0]$.  Moreover, by (7.7.5), the limit has nonnegative
sectional curvature on $M \times (-\tau(w'),0]$. Recall that we
have removed all compact components with nonnegative scalar
curvature. By combining this with the strong maximum principle, we
conclude that the limit is a flat metric. Hence in case (2), $M$
is diffeomorphic to a flat manifold and then it is also
geometrizable.

Therefore we have completed the proof of the theorem.
%$$\eqno \#$$ \vskip 0.3cm
\end{proof}


%%%\backmatter
\newpage

\section{Named intermediate formalization obligations}

\begin{lemma}[local ancient-model alternative without a small-radius hypothesis]
  \label{lem:cz-local-ancient-model-alternative}
  \dcref{Claim in the proof of Theorem 7.2.1}
  \group{Curvature Estimates for Smooth Solutions}
  \source{cao-zhu}
  \uses{prop:cz-neck-radius-lower-bound-at-infinity}
For the above fixed $\varepsilon>0$, one can find
$K=K(A,\varepsilon)<+\infty$ such that if we have a
three-dimensional complete orientable solution with normalized
initial metric and satisfying
$$
|Rm|(x,t)\le r_0^{-2}\quad \mbox{for }\; 0\le t\le
r_0^2,\;d_0(x,x_0)\le r_0,
$$
and
$$
\Vol_0(B_0(x_0,r_0))\ge A^{-1}r_0^3
$$
for some $x_0 \in M$ and some $r_0 > 0$, then for any point $x\in
M$ with $d_{r_0^2}(x,x_0)<3Ar_0$, either $$R(x,r_0^2)< Kr_0^{-2}$$
or the subset $\{(y,t)\ |\ d_{r_0^2}^2(y,x)\leq \varepsilon^{-2}
R(x,r_0^2)^{-1},\;r_0^2-\varepsilon^{-2} R(x,r_0^2)^{-1}\le t\le
r_0^2\}$ around the point $(x,r_0^2)$ is $\varepsilon$-close to
the corresponding subset of an orientable ancient
$\kappa$-solution.
\medskip

\end{lemma}
\begin{proof}
arbitrary in proving the above claim. Note that the assumption on
the normalization of the initial metric is just to ensure the
pinching estimate. By scaling, we may assume $r_0=1$. The proof of
the claim is essentially adapted from that of Theorem 7.1.1. But
we will meet the difficulties of adjusting points and verifying a
local curvature estimate.

Suppose that the claim is not true. Then there exist a sequence of
solutions $(M_k,g_k(\cdot,t))$ to the Ricci flow satisfying the
assumptions of the claim with the origins $x_{0_k}$, and a
sequence of positive numbers $K_k\rightarrow\infty$, times $t_k=1$
and points $x_k\in M_k$ with $d_{t_k}(x_k,x_{0_k})<3A$ such that
$Q_k=R_k(x_k,t_k)\ge K_k$ and the solution in $\{(x,t)\ |\
t_k-C(\varepsilon)Q_k^{-1}\le t\le t_k,\;d_{t_k}^2(x,x_k)\le
C(\varepsilon)Q_k^{-1}\}$ is not, after scaling by the factor
$Q_k$, $\varepsilon$-close to the corresponding subset of any
orientable ancient $\kappa$-solution, where $R_k$ denotes the
scalar curvature of $(M_k,g_k(\cdot,t))$ and $C(\varepsilon)(\geq
\varepsilon^{-2})$ is the constant defined in the proof of Theorem
7.1.1. As before we need to first adjust the point $(x_k,t_k)$
with $t_k\ge\frac{1}{2}$ and $d_{t_k}(x_k,x_{0_k})<4A$ so that
$Q_k=R_k(x_k,t_k)\ge K_k$ and the conclusion of the claim fails at
$(x_k,t_k)$, but holds for any $(x,t)$ satisfying $R_k(x,t)\ge
2Q_k,$ $t_k-H_kQ_k^{-1}\le t\le t_k$ and
$d_t(x,x_{0_k})<d_{t_k}(x_k,x_{0_k})
+H_k^{\frac{1}{2}}Q_k^{-\frac{1}{2}}$, where
$H_k=\frac{1}{4}K_k\rightarrow\infty,$ as $k\rightarrow+\infty$.

Indeed, by starting with $(x_{k_1},t_{k_1})=(x_k,1)$ we can choose
$(x_{k_2},t_{k_2})\in M_k\times (0,1]$ with
$t_{k_1}-H_kR_k(x_{k_1},t_{k_1})^{-1}\le t_{k_2}\le t_{k_1}$, and
$d_{t_{k_2}}(x_{k_2},x_{0_k})<d_{t_{k_1}}(x_{k_1},x_{0_k})
+H_k^\frac{1}{2}R_k(x_{k_1},t_{k_1})^{-\frac{1}{2}}$ such that
$R_k(x_{k_2},t_{k_2})\ge 2R_k(x_{k_1},t_{k_1})$ and the conclusion
of the claim fails at $(x_{k_2},t_{k_2})$; otherwise we have the
desired point. Repeating this process, we can choose points
$(x_{k_i},t_{k_i})$, $i=2, \ldots, j$, such that
$$
R_k(x_{k_i},t_{k_i})\ge 2R_k(x_{k_{i-1}},t_{k_{i-1}}),
$$
$$
t_{k_{i-1}}-H_kR_k(x_{k_{i-1}},t_{k_{i-1}})^{-1} \le t_{k_{i}}\le
t_{k_{i-1}},
$$
$$
d_{t_{k_{i}}}(x_{k_{i}},x_{0_{k}})
<d_{t_{k_{i-1}}}(x_{k_{i-1}},x_{0_k})+H_k^\frac{1}{2}
R_k(x_{k_{i-1}},t_{k_{i-1}})^{-\frac{1}{2}},
$$
and the conclusion of the claim fails at the points
$(x_{k_i},t_{k_i})$, $i=2, \ldots, j$. These inequalities imply
$$
R_k(x_{k_{j}},t_{k_{j}})\ge 2^{j-1}R_k(x_{k_1},t_{k_1})\ge
2^{j-1}K_k,
$$
$$
1\ge t_{k_{j}}\ge t_{k_{1}}
-H_k\sum_{i=0}^{j-2}\frac{1}{2^i}R_k(x_{k_1},t_{k_1})^{-1}
\ge\frac{1}{2},
$$
and
$$
d_{t_{k_{j}}}(x_{k_j},x_{0_k}) <d_{t_{k_1}}(x_{k_1},x_{0_k})
+H_k^\frac{1}{2}\sum_{i=0}^{j-2}\frac{1}{(\sqrt{2})^i}
R_k(x_{k_1},t_{k_1})^{-\frac{1}{2}}<4A.
$$
Since the solutions are smooth, this process must terminate after
a finite number of steps to give the desired point, still denoted
by $(x_k,t_k).$

For each adjusted $(x_k,t_k)$, let $[t',t_k]$ be the maximal
subinterval of $[t_k-\frac{1}{2}\varepsilon^{-2}Q_k^{-1},t_k]$ so
that the conclusion of the claim with $K=2Q_k$ holds on
\begin{align*}
& P\left(x_k,t_k,\frac{1}{10}H_k^{\frac{1}{2}}Q_k^{-\frac{1}{2}},t'-t_k\right)\\
&=\left\{ ({x},{t}) \ |\ {x} \in
B_{{t}}\left(x_k,\frac{1}{10}H_k^{\frac{1}{2}}Q_k^{-\frac{1}{2}}\right),
{t} \in [t',t_k]\right\}
\end{align*}
for all sufficiently large $k$. We now want to show $t' =
t_k-\frac{1}{2}\varepsilon^{-2}Q_k^{-1}$.

Consider the scalar curvature $R_k$ at the point $x_k$ over the
time interval $[t',t_k]$. If there is a time $\tilde{t} \in
[t',t_k]$ satisfying $R_k(x_k,\tilde{t}) \geq 2Q_k$, we let
$\tilde{t}$ be the first of such time from $t_k$. Then the
solution $(M_k,g_k(\cdot,t))$ around the point $x_k$ over the time
interval $[\tilde{t}-
\frac{1}{2}\varepsilon^{-2}Q_k^{-1},{\tilde{t}}]$ is
$\varepsilon$-close to some orientable ancient $\kappa$-solution.
Note from the Li-Yau-Hamilton inequality that the scalar curvature
of any ancient $\kappa$-solution is pointwise nondecreasing in
time. Consequently, we have the following curvature estimate
$$
R_k(x_k,t) \leq 2(1+\varepsilon)Q_k
$$
for $ t \in [\tilde{t}- \frac{1}{2}\varepsilon^{-2}Q_k^{-1},t_k]$
(or $t \in [t',t_k]$ if there is no such time $\tilde{t}$).  By
combining with the elliptic type estimate for ancient
$\kappa$-solutions (Theorem 6.4.3) and the Hamilton-Ivey pinching
estimate, we further have \begin{equation}
|Rm(x,t)| \leq 5\omega(1)Q_k  %\eqno (7.2.1)
\end{equation} for all $x \in B_t(x_k,(3Q_k)^{-\frac{1}{2}})$ and $t \in
[\tilde{t}- \frac{1}{2}\varepsilon^{-2}Q_k^{-1},t_k]$ (or $t \in
[t',t_k]$) and all sufficiently large $k$, where $\omega$ is the
positive function in Theorem 6.4.3.

For any point $(x,t)$ with $ \tilde{t}-
\frac{1}{2}\varepsilon^{-2}Q_k^{-1} \leq t \leq t_k$ (or $t \in
[t',t_k]$) and $d_t(x,x_k) \leq
\frac{1}{10}H_k^{\frac{1}{2}}Q_k^{-\frac{1}{2}}$, we divide the
discussion into two cases.

\medskip
{\it Case} (1): $d_t(x_k,x_{0_k}) \leq
\frac{3}{10}H_k^{\frac{1}{2}}Q_k^{-\frac{1}{2}}.$
\begin{align}
d_t(x,x_{0_k})&  \leq
d_t(x,x_k) + d_t(x_k,x_{0_k})\\
&  \leq\frac{1}{10}H_k^{\frac{1}{2}}Q_k^{-\frac{1}{2}}
+ \frac{3}{10}H_k^{\frac{1}{2}}Q_k^{-\frac{1}{2}}\nonumber\\
&  \leq \frac{1}{2}H_k^{\frac{1}{2}}Q_k^{-\frac{1}{2}}.\nonumber
\end{align} %\eqno (7.2.2)

\medskip
{\it Case} (2): $d_t(x_k,x_{0_k}) >
\frac{3}{10}H_k^{\frac{1}{2}}Q_k^{-\frac{1}{2}}.$

\smallskip
{}From the curvature bound $(7.2.1)$ and the assumption, we apply
Lemma 3.4.1(ii) with $r_0 = Q_k^{-\frac{1}{2}}$ to get
$$
\frac{d}{dt}(d_t(x_k,x_{0_k})) \geq -20(\omega(1) +
1)Q_k^{\frac{1}{2}},
$$
and then for $k$ large enough,
\begin{align*}
d_t(x_k,x_{0_k})&  \leq d_{\hat{t}}(x_k,x_{0_k})
+ 20(\omega(1)+1)\varepsilon^{-2}Q_k^{-\frac{1}{2}}\\
&  \leq d_{\hat{t}}(x_k,x_{0_k}) +
\frac{1}{10}H_k^{\frac{1}{2}}Q_k^{-\frac{1}{2}},
\end{align*}
where $\hat{t} \in (t,t_k]$ satisfies the property that
$d_s(x_k,x_{0_k}) \geq
\frac{3}{10}H_k^{\frac{1}{2}}Q_k^{-\frac{1}{2}}$ whenever $s \in
[t,\hat{t}]$. So we have
\begin{align}
d_t(x,x_{0_k})&  \leq d_t(x,x_k) + d_t(x_k,x_{0_k})\\
&  \leq\frac{1}{10}H_k^{\frac{1}{2}}Q_k^{-\frac{1}{2}} +
d_{\hat{t}}(x_k,x_{0_k})
+ \frac{1}{10}H_k^{\frac{1}{2}}Q_k^{-\frac{1}{2}} \nonumber\\
&  \leq d_{t_k}(x_k,x_{0_k}) +
\frac{1}{2}H_k^{\frac{1}{2}}Q_k^{-\frac{1}{2}},\nonumber
\end{align}   %\eqno (7.2.3)
for all sufficiently large $k$. Then the combination of (7.2.2),
(7.2.3) and the choice of the points $(x_k,t_k)$ implies $t'=
t_k-\frac{1}{2}\varepsilon^{-2}Q_k^{-1}$ for all sufficiently
large $k$. (Here we also used the maximality of the subinterval
$[t',t_k]$ in the case that there is no time in $ [t',t_k]$ with
$R_k(x_k,\cdot) \geq 2Q_k$.)

Now we rescale the solutions $(M_k,g_k(\cdot,t))$ into
$(M_k,\tilde{g}_k(\cdot,t))$ around the points $x_k$ by the
factors $Q_k=R_k(x_k,t_k)$ and shift the times $t_k$ to the new
times zero. Then the same arguments from Step 1 to Step 3 in the
proof of Theorem 7.1.1 prove that a subsequence of the rescaled
solutions $(M_k,\tilde{g}_k(\cdot,t))$ converges in the
$C^\infty_{\rm loc}$ topology to a limiting (complete) solution
$(M_{\infty},\tilde{g}_{\infty}(\cdot,t))$, which is defined on a
backward time interval $[-a,0]$ for some $a>0$. (The only
modification is in Lemma 7.1.2 of Step 1 by further requiring
$t_k- \frac{1}{4}\varepsilon^{-2}Q_k^{-1} \leq \bar{t} \leq t_k$).

We next study how to adapt the argument of Step 4 in the proof of
Theorem 7.1.1. As before, we have a maximal time interval
$(t_{\infty},0]$ for which we can take a smooth limit
$(M_{\infty},\tilde{g}_{\infty}(\cdot,t),x_{\infty})$ from a
subsequence of the rescaled solutions
$(M_k,\tilde{g}_k(\cdot,t),x_k)$. We want to show
$t_{\infty}=-\infty$.

Suppose not; then $t_{\infty} > -\infty.$ Let $c>0$ be a positive
constant much smaller than $\frac{1}{10}\varepsilon^{-2}$. Note
that the infimum of the scalar curvature is nondecreasing in time.
Then we can find some point $y_{\infty} \in M_{\infty}$ and some
time $t=t_{\infty} + \theta$ with $0<\theta<\frac{c}{3}$ such that
$\tilde{R}_{\infty}(y_{\infty},t_{\infty}+\theta) \leq
\frac{3}{2}$.

Consider the (unrescaled) scalar curvature $R_k$ of
$(M_k,{g}_k(\cdot,t))$ at the point $x_k$ over the time interval
$[t_k+(t_{\infty}+\frac{\theta}{2})Q_k^{-1},t_k]$. Since the
scalar curvature $\tilde{R}_{\infty}$ of the limit on
$M_{\infty}\times [t_{\infty}+\frac{\theta}{3},0]$ is uniformly
bounded by some positive constant $C$, we have the curvature
estimate $$R_k(x_k,t) \leq 2CQ_k$$ for all $t \in
[t_k+(t_{\infty}+\frac{\theta}{2})Q_k^{-1},t_k]$ and all
sufficiently large $k$. Then by repeating the same arguments as in
deriving (7.2.1), (7.2.2) and (7.2.3), we deduce that the
conclusion of the claim with $K=2Q_k$ holds on the parabolic
neighborhood
$P(x_k,t_k,\frac{1}{10}H_k^{\frac{1}{2}}Q_k^{-\frac{1}{2}},
(t_{\infty}+\frac{\theta}{2})Q_k^{-1})$ for all sufficiently large
$k$.

Let $(y_k,t_k+(t_{\infty}+\theta_k)Q_k^{-1})$ be a sequence of
associated points and times in the (unrescaled) solutions
$(M_k,g_k(\cdot,t))$ so that after rescaling, the sequence
converges to the $(y_{\infty},t_{\infty}+\theta)$ in the limit.
Clearly $\frac{\theta}{2}\leq \theta_k \leq 2\theta$ for all
sufficiently large $k$. Then, by considering the scalar curvature
$R_k$ at the point $y_k$ over the time interval
$[t_k+(t_{\infty}-\frac{c}{3})Q_k^{-1},t_k+(t_{\infty}+\theta_k)Q_k^{-1}]$,
the above argument (as in deriving the similar estimates
(7.2.1)-(7.2.3)) implies that the conclusion of the claim with
$K=2Q_k$ holds on the parabolic neighborhood
$P(y_k,t_k,\frac{1}{10}H_k^{\frac{1}{2}}Q_k^{-\frac{1}{2}},
(t_{\infty}-\frac{c}{3})Q_k^{-1})$ for all sufficiently large $k$.
In particular, we have the curvature estimate
$$
R_k(y_k,t) \leq 4(1+\varepsilon)Q_k
$$
for $t \in [t_k+(t_{\infty}-\frac{c}{3})Q_k^{-1},
t_k+(t_{\infty}+\theta_k)Q_k^{-1}]$ for all sufficiently large
$k$.

We now consider the rescaled sequence $(M_k,\tilde{g}_k(\cdot,t))$
with the marked points replaced by $y_k$ and the times replaced by
$s_k \in [t_k+(t_{\infty}-\frac{c}{4})Q_k^{-1},
t_k+(t_{\infty}+\frac{c}{4})Q_k^{-1}]$. By applying the same
arguments from Step 1 to Step 3 in the proof of Theorem 7.1.1 and
the Li-Yau-Hamilton inequality as in Step 4 of Theorem 7.1.1, we
conclude that there is some small constant $a'>0$ such that the
original limit $(M_{\infty},\tilde{g}_{\infty}(\cdot,t))$ is
actually well defined on $M_{\infty} \times [t_{\infty}-a',0]$
with uniformly bounded curvature. This is a contradiction.
Therefore we have checked the claim.
\end{proof}

\begin{lemma}[spatial scalar-curvature comparison for the standard solution]
  \label{lem:cz-standard-solution-spatial-curvature-comparison}
  \dcref{Claim 1 in the proof of Proposition 7.3.3}
  \group{Ricci Flow with Surgery}
  \source{cao-zhu}
There is a positive function
$\omega:[0,\infty)\rightarrow[0,\infty)$ depending only on the
initial metric and $\kappa$ such that
$$
  R(x,t)\leq R(y,t)\omega(R(y,t)d_{t}^{2}(x,y))
$$
for all $x,y\in \mathbb{R}^{n}$, $t\in[0,T)$.
\end{lemma}
\begin{proof}
The proof is similar to that of Theorem 6.4.3. Notice that the
initial metric has nonnegative curvature operator and its scalar
curvature satisfies the bounds \begin{equation}
C^{-1}\leqslant R(x)\leqslant C %\eqno (7.3.5)
\end{equation} for some positive constant $C$. By the maximum principle, we
know $T\geq\frac{1}{2nC}$ and $R(x,t)\leq 2C$ for
$t\in[0,\frac{1}{4nC}]$. The assertion is clearly true for
$t\in[0,\frac{1}{4nC}]$.

Now fix $(y,t_{0})\in \mathbb{R}^{n}\times[0,T)$ with
$t_{0}\geq\frac{1}{4nC}$. Let $z$ be the closest point to $y$ with
the property $R(z,t_{0})d^{2}_{t_{0}}(z,y)=1$ (at time $t_{0}$).
Draw a shortest geodesic from $y$ to $z$ and choose a point
$\tilde{z}$ on the geodesic satisfying $d_{t_{0}}(z,\tilde{z})
=\frac{1}{4}R(z,t_{0})^{-\frac{1}{2}}$, then we have
$$
 R(x,t_{0})\leq\frac{1}{(\frac{1}{2}R(z,t_{0})^{-\frac{1}{2}})^{2}}
    \qquad \mbox{on}\ \
 B_{t_{0}}\left(\tilde{z},\frac{1}{4}R(z,t_{0}\right)^{-\frac{1}{2}}).
$$

Note that $R(x,t)\geqslant C^{-1}$ everywhere by the evolution
equation of the scalar curvature. Then by the Li-Yau-Hamilton
inequality (Corollary 2.5.5), for all $(x,t)\in
B_{t_{0}}(\tilde{z},\frac{1}{8nC}R(z,t_{0})^{-\frac{1}{2}})
\times[t_{0}-(\frac{1}{8nC}R(z,t_{0})^{-\frac{1}{2}})^{2},t_{0}]
$, we have
\begin{equation*}
\begin{split}
R(x,t)&  \leq \left(\frac{t_{0}}{t_{0}-\left(\frac{1}{8n\sqrt{C}}\right)^{2}}\right)
\frac{1}{\left(\frac{1}{2}R(z,t_{0})^{-\frac{1}{2}}\right)^{2}}\\
&  \leq \left[\frac{1}{8nC}R(z,t_{0})^{-\frac{1}{2}}\right]^{-2}.
\end{split}
\end{equation*}
Combining this with the $\kappa$-noncollapsing property, we have
$$
\Vol
\left(B_{t_{0}}\left(\tilde{z},\frac{1}{8nC}R(z,t_{0})^{-\frac{1}{2}}\right)\right)
\geq \kappa \left(\frac{1}{8nC}R(z,t_{0})^{-\frac{1}{2}}\right)^{n}
$$
and then
$$
\Vol \left(B_{t_{0}}\left(z,8R(z,t_{0})^{-\frac{1}{2}}\right)\right) \geq \kappa
\left(\frac{1}{64nC}\right)^{n}\left(8R(z,t_{0})^{-\frac{1}{2}}\right)^{n}.
$$
So by Theorem 6.3.3 (ii), we have
$$
R(x,t_{0})\leq C(\kappa)R(z,t_{0})\quad \text{for all }\; x\in
B_{t_{0}}\left(z,4R(z,t_{0})^{-\frac{1}{2}}\right).
$$
Here and in the following we denote by $C(\kappa)$ various
positive constants depending only on $\kappa,n$ and the initial
metric.

Now by the Li-Yau-Hamilton inequality (Corollary 2.5.5) and local
gradient estimate of Shi (Theorem 1.4.2), we obtain
\begin{align*}
R(x,t)\leq C(\kappa)R(z,t_{0})\; \text{ and }\;
\left|\frac{\partial}{\partial t}R\right|(x,t) \leq
C(\kappa)(R(z,t_{0}))^{2}
\end{align*}
for all $ (x,t)\in
B_{t_{0}}(z,2R(z,t_{0})^{-\frac{1}{2}}))\times[t_{0}-
(\frac{1}{8nC}R(z,t_{0})^{-\frac{1}{2}})^{2},t_{0}]$. Therefore by
combining with the Harnack estimate (Corollary 2.5.7), we obtain
\begin{align*}
R(y,t_{0})&  \geq C(\kappa)^{-1}R(z,t_{0}-C(\kappa)^{-1}R(z,t_{0})^{-1})\\
&  \geq C(\kappa)^{-2} R(z,t_{0})
\end{align*}

Consequently, we have showed that there is a constant $C(\kappa)$
such that
$$
\Vol \left(B_{t_{0}}\left(y,R(y,t_{0})^{-\frac{1}{2}}\right)\right) \geq
C(\kappa)^{-1}\left(R(y,t_{0})^{-\frac{1}{2}}\right)^{n}
$$
and
$$
R(x,t_{0})\leq C(\kappa)R(y,t_{0}) \quad \text{for all }\; x\in
B_{t_{0}}\left(y,R(y,t_{0})^{-\frac{1}{2}}\right).
$$
In general, for any $r\geq R(y,t_{0})^{-\frac{1}{2}}$, we have
$$
\Vol (B_{t_{0}}(y,r)) \geq
C(\kappa)^{-1}(r^{2}R(y,t_{0}))^{-\frac{n}{2}}r^{n}.
$$
By\; applying\; Theorem\; 6.3.3(ii)\; again,\; there\; exists\;
a\; positive\; constant $\omega(r^{2}R(y,t_{0}))$ depending only
on the constant $r^{2}R(y,t_{0})$ and $\kappa$ such that
$$
R(x,t_{0})\leq R(y,t_{0})\omega(r^{2}R(y,t_{0}))\quad \mbox{for
all }\;  x\in B_{t_{0}}\left(y,\frac{1}{4}r\right).
$$
This proves the desired Claim 1.
\end{proof}

\begin{lemma}[standard-solution curvature blow-up at maximal time]
  \label{lem:cz-standard-solution-maximal-time-curvature-blowup}
  \dcref{Claim 2 in the proof of Proposition 7.3.3}
  \group{Ricci Flow with Surgery}
  \source{cao-zhu}
  \uses{lem:cz-standard-solution-spatial-curvature-comparison}
Suppose $T<\frac{n-1}{2}$. Fix a point $x_{0}\in
\mathbb{R}^{n} $, then there is a $\delta>0$, such that for any
$x\in M$ with $d_{0}(x,x_{0})\geq \delta^{-1}$, we have
$$
R(x,t)\leq 2C+\frac{n-1}{\frac{n-1}{2}-t}\quad \text{ for all }\;
t\in[0,T),
$$
where $C$ is the constant in (7.3.5).
\end{lemma}
\begin{proof}
In view of Claim 1, if Claim 2 holds, then
\begin{align*}
\sup_{M^{n}\times[0,T)}R(y,t) &  \leq
\omega\left(\delta^{-2}\left(2C+\frac{n-1}{\frac{n-1}{2}-T}\right)\right)
\left(2C+\frac{n-1}{\frac{n-1}{2}-T}\right)\\
&  < \infty
\end{align*}
which will contradict the definition of $T$.

To show Claim 2, we argue by contradiction. Suppose for each
$\delta>0$, there is a point $(x_{\delta},t_{\delta})$ with
$0<t_{\delta}<T $ such that
$$
R(x_{\delta},t_{\delta})>2C+\frac{n-1}{\frac{n-1}{2}-t_{\delta}}\;
\mbox{ and }\; d_{0}(x_{\delta},x_{0})\geq \delta^{-1}.
$$
Let
$$
\bar{t}_{\delta}=\sup\left\{t\; \Big|\; \sup_{M^{n} \setminus
B_{0}(x_{0},\delta^{-1})}
R(y,t)<2C+\frac{n-1}{\frac{n-1}{2}-t}\right\}.
$$
Since $\lim\limits_{d_{0}(y,x_{0})\rightarrow\infty}R(y,t)
=\frac{\frac{n-1}{2}}{\frac{n-1}{2}-t}$ and
$\sup_{M\times[0,\frac{1}{4nC}]}R(y,t)\leq 2C$, we know
$\frac{1}{4nC}\leq \bar{t}_{\delta}\leq t_{\delta}$ and there is a
$\bar{x}_{\delta}$ such that
$d_{0}(x_{0},\bar{x}_{\delta})\geq\delta^{-1}$ and
$R(\bar{x}_{\delta},\bar{t}_{\delta})
=2C+\frac{n-1}{\frac{n-1}{2}-\bar{t}_{\delta}}.$ By Claim 1 and
Hamilton's compactness theorem (Theorem 4.1.5), for $\delta\
\rightarrow 0$ and after taking a subsequence, the metrics
$g_{ij}(x,t)$ on $B_{0}(\bar{x}_{\delta},\frac{\delta^{-1}}{2})$
over the time interval $[0,\bar{t}_{\delta}]$ will converge to a
solution $\tilde{g}_{ij}$ on $\mathbb{R}\times \mathbb{S}^{n-1}$
with the standard metric of scalar curvature 1 as initial data
over the time interval $[0,\bar{t}_{\infty}]$, and its scalar
curvature satisfies
\begin{align*}
\tilde{R}(\bar{x}_{\infty},\bar{t}_{\infty})
&  = 2C+\frac{n-1}{\frac{n-1}{2}-\bar{t}_{\infty}},\\
\tilde{R}(x,t)&  \leqslant
2C+\frac{n-1}{\frac{n-1}{2}-\bar{t}_{\infty}},
 \quad \mbox{for all }\;  t\in [0,\bar{t}_{\infty}],
\end{align*}
where $(\bar{x}_{\infty},\bar{t}_{\infty})$ is the limit of
$(\bar{x}_{\delta},\bar{t}_{\delta})$. On the other hand, by the
uniqueness theorem (Theorem 1.2.4) again, we know
$$
\tilde{R}(\bar{x}_{\infty},\bar{t}_{\infty})
=\frac{\frac{n-1}{2}}{\frac{n-1}{2}-\bar{t}_{\infty}}
$$
which is a contradiction.  Hence we have proved Claim 2 and then
\end{proof}

\begin{lemma}[reduced-length barrier for curves entering a surgery cap]
  \label{lem:cz-surgery-cap-reduced-length-barrier}
  \dcref{Assertion (7.4.2)}
  \group{Canonical-Neighborhood Induction}
  \source{cao-zhu}
For any $L<\infty$ one can find
$\bar{\delta}=\bar{\delta}(L,r,\widetilde{r}_m,\varepsilon)>0$
with the following property. Suppose that we have a curve
$\gamma$, parametrized by $t\in[T_0,t_0]$, $(m-1)\varepsilon^2\leq
T_0<t_0$, such that $\gamma(t_0)=x_0$, $T_0$ is a surgery time,
and $\gamma(T_0)$ lies in the gluing cap. Suppose also each
$\delta$-cutoff surgery at a time in
$[(m-1)\varepsilon^2,\bar{T}]$ has $\delta \leq \bar{\delta}$.
Then we have an estimate \begin{equation}
\int^{t_0}_{T_0}\sqrt{t_0-t}(R_+(\gamma(t),t)
+|\dot{\gamma}(t)|^2)dt\geq L  %\eqno(7.4.2)
\end{equation} where $R_+=\max\{R,0\}$.

\medskip
Since $r_0\geq \frac{1}{2\eta}r$ and $|Rm|\leq r^{-2}_0$ on
\end{lemma}
\begin{proof}
$\bar{\delta}>0$, depending on $r$ and $\widetilde{r}_m$, to be so
small that $\gamma(T_0)$ does not lie in the region
$P(x_0,t_0,r_0,-r^2_0)$. Let $\Delta t$ be the maximal number such
that $\gamma|_{[t_0-\Delta t,t_0]}\subset P(x_0,t_0,r_0,-\Delta
t)$ (i.e., $t_0-\Delta t$ is the first time when $\gamma$ escapes
the parabolic region $P(x_0,t_0,r_0,-r^2_0)).$ Obviously we only
need to consider the case:
$$
\int^{t_0}_{t_0 -\Delta t}\sqrt{t_0 -
t}(R_+(\gamma(t),t)+|\dot{\gamma}(t)|^2)dt<L.
$$

We observe that $\Delta t$ can be bounded from below in terms of
$L$ and $r_0$. Indeed, if $\Delta t \geq r_0^2$, there is nothing
to prove. Thus we may assume $\Delta t<r_0^2$. By the curvature
bound $|Rm|\leq r^{-2}_0$ on $P(x_0,t_0,r_0,-r^2_0)$ and the Ricci
flow equation we see
$$
\int^{t_0}_{t_0-\Delta t}|\dot{\gamma}(t)|dt\geq cr_0
$$
for some universal positive constant $c$. On the other hand, by
the Cauchy-Schwarz inequality, we have
\begin{align*}
\int^{t_0}_{t_0-\Delta t}|\dot{\gamma}(t)|dt & \leq
\left(\int^{t_0}_{t_0-\Delta t}\sqrt{t_0-t}(R_++|
\dot{\gamma}|^2)dt\right)^{\frac{1}{2}} \cdot\left(\int^{t_0}_{t_0-\Delta
t}\frac{1}{\sqrt{t_0-t}}dt\right)^{\frac{1}{2}}\\
&  \leq (2L)^{\frac{1}{2}}(\Delta t)^{\frac{1}{4}},
\end{align*}
which implies \begin{equation}
(\Delta t)^{\frac{1}{2}}\geq\frac{c^2r^2_0}{2L}. %\eqno (7.4.3)
\end{equation} Thus
\begin{align*}
\int^{t_0}_{T_0}\sqrt{t_0-t}(R_++|\dot{\gamma}|^2)dt& \geq
\int^{t_0-\Delta
t}_{T_0}\sqrt{t_0-t}(R_++|\dot{\gamma}|^2)dt\\
&  \geq (\Delta t)^{\frac{1}{2}}\int^{t_0-\Delta
t}_{T_0}(R_++|\dot{\gamma}|^2)dt\\
&  \geq \left(\min \left\{\frac{c^2r^2_0}{2L},r_0 \right\}\right)
\int^{t_0-\Delta t}_{T_0}(R_++|\dot{\gamma}|^2)dt,
\end{align*}
while by Corollary 7.3.7, we can find
$\bar{\delta}=\bar{\delta}(L,r,\widetilde{r}_m,\varepsilon)>0$ so
small that
$$
\int^{t_0-\Delta t}_{T_0}(R_++|\dot{\gamma}|^2)dt \geq
L\left(\min\left\{\frac{c^2r^2_0}{2L},r_0 \right\}\right)^{-1}.
$$
Then we have proved the desired assertion (7.4.2).
\end{proof}

\begin{lemma}[backward extension on a bounded-curvature rescaled ball]
  \label{lem:cz-rescaled-ball-backward-extension}
  \dcref{Assertion 1 in Step 2 of Lemma 7.4.2}
  \group{Canonical-Neighborhood Induction}
  \source{cao-zhu}
  \uses{lem:cz-surgery-stability-near-standard-cap}
For arbitrarily fixed $\alpha$,
$0<A<+\infty$,  $1 \leq C<+\infty$ and $0 \leq B <
\frac{1}{2}\varepsilon^2(r^{\alpha})^{-2}-
\frac{1}{8}\eta^{-1}C^{-1}$, there is a
$\beta_0=\beta_0(\varepsilon, A, B, C)$ (independent of $\alpha$)
such that if $\beta\geq\beta_0$ and the rescaled solution
$\widetilde{g}^{\alpha\beta}_{ij}$ on the ball
$\widetilde{B}_{0}(\bar{x},A)$ is defined on a time interval
$[-b,0]$ with $0 \leq b \leq B$ and the scalar curvature satisfies
$$\widetilde{R}(x,t)\leq C , \ \ \mbox{ on }
\widetilde{B}_{0}(\bar{x},A) \times [-b,0],$$ then the rescaled
solution $\widetilde{g}^{\alpha\beta}_{ij}$ on the ball
$\widetilde{B}_{0}(\bar{x},A)$ is also defined on the extended
time interval $ [-b-\frac{1}{8}\eta^{-1}C^{-1},0]$.

\medskip
Before giving the proof, we make a simple \textbf{observation}:
once a space point in the Ricci flow with surgery is removed by
surgery at some time, then it never appears for later time; if a
space point at some time $t$ cannot be defined before the time $t$
, then either the point lies in a gluing cap of the surgery at
time $t$ or the time $t$ is the initial time of the Ricci flow.
\end{lemma}
\begin{proof}
Firstly we claim that
there exists $\beta_0=\beta_0( \varepsilon, A, B, C)$ such that
when $\beta\geq\beta_0$, the rescaled solution
$\widetilde{g}^{\alpha\beta}_{ij}$ on the ball
$\widetilde{B}_{0}(\bar{x},A)$ can be defined before the time $-b$
(i.e., there are no surgeries interfering in
$\widetilde{B}_{0}(\bar{x},A)\times [-b-\epsilon',-b]$ for some
$\epsilon'>0$).

We argue by contradiction. Suppose not, then there is some point
$\tilde{x}\in \widetilde{B}_{0}(\bar{x},A)$ such that the rescaled
solution $\widetilde{g}^{\alpha\beta}_{ij}$ at $\tilde{x}$ cannot
be defined before the time $-b$. By the above observation, there
is a surgery at the time $-b$ such that the point $\tilde{x}$ lies
in the instant gluing cap.

Let\; $\tilde{h}$\; $(=R(\bar{x},\,\bar{t})^{\frac{1}{2}}h$)\;
be\; the\; cut-off\; radius at the\; time\; $-b$\; for the
rescaled solution.  Clearly, there is a universal constant $D$
such that
${D}^{-1}\tilde{h}\leq\widetilde{R}(\tilde{x},-b)^{-\frac{1}{2}}\leq
D\tilde{h}$.

By Lemma 7.3.4 and looking at the rescaled solution at the time
$-b$, the gluing cap and the adjacent $\delta$-neck, of radius
$\tilde{h}$, constitute a
${(\widetilde{\delta}^{\alpha\beta})}^{\frac{1}{2}}$-cap
$\mathcal{K}$. For any fixed small positive constant
$\delta^{\prime}$ (much smaller than $\varepsilon$), we see that
$$
\widetilde{B}_{(-b)}(\tilde{x},{(\delta^{\prime})}^{-1}
\widetilde{R}(\tilde{x}, -b)^{-\frac{1}{2}})\subset\mathcal{K}
$$
when $\beta$ large enough.  We first verify the following

\medskip
{\bf Claim 1.} For any small constants $0<\tilde{\theta}<1$,
$\delta'>0$, there exists a $\beta(\delta',\varepsilon,
\tilde{\theta})> 0$ such that when $\beta \geq
\beta(\delta',\varepsilon, \tilde{\theta})$, we have
\begin{itemize}
\item[(i)] the rescaled solution
$\widetilde{g}^{\alpha\beta}_{ij}$ over
$\widetilde{B}_{(-b)}(\tilde{x},{(\delta^{\prime})}^{-1}
\tilde{h})$  is defined on the time interval
$[-b,0]\cap[-b,-b+(1-\tilde{\theta})\tilde{h}^{2}]$; \item[(ii)]
the ball
$\widetilde{B}_{(-b)}(\tilde{x},{(\delta^{\prime})}^{-1}\tilde{h})$
in the ${(\widetilde{\delta}^{\alpha\beta})}^{\frac{1}{2}}$-cap
$\mathcal{K}$ evolved by the Ricci flow  on the time interval
$[-b,0]\cap[-b,-b+(1-\tilde{\theta})\tilde{h}^{2}]$ is, after
scaling with factor $\tilde{h}^{-2}$, ${\delta}^{\prime}$-close
(in the $C^{[\delta'^{-1}]}$ topology) to the corresponding subset
of the standard solution.
\end{itemize}

\medskip
This claim essentially follows from Lemma 7.3.6. Indeed, suppose
there is a surgery at some time $\tilde{\tilde{t}}\in
[-b,0]\cap(-b,-b+(1-\tilde{\theta})\tilde{h}^{2}]$ which removes
some point  $\tilde{\tilde{x}}\in \widetilde
B_{(-b)}(\tilde{x},{(\delta^{\prime})}^{-1}\tilde{h})$. We assume
$\tilde{\tilde{t}}\in (-b,0]$ is the first time with that
property.

Then by Lemma 7.3.6, there is a
$\bar{\delta}=\bar{\delta}(\delta',\varepsilon,\tilde{\theta})$
such that if $\widetilde{\delta}^{\alpha\beta}<\bar{\delta}$, then
the ball
$\widetilde{B}_{(-b)}(\tilde{x},{(\delta^{\prime})}^{-1}\tilde{h})$
in the ${(\widetilde{\delta}^{\alpha\beta})}^{\frac{1}{2}}$-cap
$\mathcal{K}$ evolved by the Ricci flow  on the time interval
$[-b,\tilde{\tilde{t}})$ is, after scaling with factor
$\tilde{h}^{-2}$, ${\delta}^{\prime}$-close to the corresponding
subset of the standard solution. Note that the metrics for times
in $[-b,\tilde{\tilde{t}})$ on $\widetilde
B_{(-b)}(\tilde{x},{(\delta^{\prime})}^{-1}\tilde{h})$  are
equivalent. By Lemma 7.3.6, the solution on $\widetilde
B_{(-b)}(\tilde{x},{(\delta^{\prime})}^{-1}\tilde{h})$ keeps
looking like a cap for $t\in [-b,\tilde{\tilde{t}})$. On the other
hand, by the definition, the surgery is always done along the
middle two-sphere of a $\delta$-neck with $\delta <
\widetilde{\delta}^{\alpha\beta}$. Then for $\beta$ large, all the
points in $\widetilde
B_{(-b)}(\tilde{x},{(\delta^{\prime})}^{-1}\tilde{h})$ are removed
(as a part of a capped horn) at the time $\tilde{\tilde{t}}$. But
$\tilde{x}$ (near the tip of the cap) exists past the time
$\tilde{\tilde{t}}$. This is a contradiction. Hence we have proved
that
$\widetilde{B}_{(-b)}(\tilde{x},{(\delta^{\prime}})^{-1}\tilde{h})$
is defined on the time interval
$[-b,0]\cap[-b,-b+(1-\tilde{\theta})\tilde{h}^{2}]$.

The $\delta'$-closeness of the solution on $\widetilde{B}_{(-b)}
(\tilde{x},{(\delta^{\prime})}^{-1}h)\times ([-b,0]\cap[-b,
-b+(1-\tilde{\theta})\tilde{h}^{2}])$ with the corresponding
subset of the standard solution follows from Lemma 7.3.6. Then we
have proved Claim 1.

\smallskip
We next verify the following

\medskip
{\bf Claim 2.} \ There is $\tilde{\theta}=\tilde{\theta}(CB)$,
$0<\tilde{\theta}<1$, such that
$b\leq(1-\tilde{\theta})\tilde{h}^{2}$ when $\beta$ large.

\medskip
Note from Proposition 7.3.3, there is a universal constant $D'>0$
such that the standard solution satisfies the following curvature
estimate
$$
R(y,s) \geq \frac{2D'}{1-s}.
$$
We choose $\tilde{\theta}= {D'}/{2(D'+CB)}$. Then for $\beta$
large enough, the rescaled solution satisfies \begin{equation}
\widetilde{R}(x,t)\geq
\frac{D'}{1-(t+b)\tilde{h}^{-2}}\tilde{h}^{-2} %\eqno (7.4.15)
\end{equation} on
$\widetilde{B}_{(-b)}(\tilde{x},{(\delta^{\prime})}^{-1}\tilde{h})\times
([-b,0]\cap[-b,-b+(1-\tilde{\theta})\tilde{h}^{2}])$.

Suppose $b\geq (1-\tilde{\theta})\tilde{h}^{2}$. Then by combining
with the assumption $\widetilde{R}(\tilde{x},t)\leq {C}$ for
$t=(1-\tilde{\theta})\tilde{h}^{2}-b$, we have
$$
C\geq \frac{D'}{1-(t+b)\tilde{h}^{-2}}\tilde{h}^{-2},
$$
and then
$$
1\geq(1-\tilde{\theta})\left(1+\frac{D'}{CB}\right).
$$
This is a contradiction. Hence we have proved Claim 2.

The combination of the above two claims shows that there is a
positive constant $0 <\tilde{\theta}=\tilde{\theta}(CB)<1$ such
that for any small $\delta'>0$, there is a positive
$\beta(\delta',\varepsilon, \tilde{\theta})$ such that when $\beta
\geq \beta(\delta',\varepsilon, \tilde{\theta})$, we have $b\leq
(1-\tilde{\theta})\tilde{h}^{2}$ and the rescaled solution in the
ball
$\widetilde{B}_{(-b)}(\tilde{x},{(\delta^{\prime}})^{-1}\tilde{h})$
on the time interval $[-b,0]$ is, after scaling with factor
$\tilde{h}^{-2}$, ${\delta}^{\prime}$-close ( in the
$C^{[(\delta')^{-1}]}$ topology) to the corresponding subset of
the standard solution.

By (7.4.15) and the assumption $\widetilde{R}\leq {C}$ on
$\widetilde{B}_{0}(\bar{x},A) \times [-b,0],$ we know that the
cut-off radius $\tilde{h}$ at the time $-b$ for the rescaled
solution satisfies
$$\tilde{h} \geq \sqrt{\frac{D'}{C}}.$$

Let $\delta'>0$ be much smaller than $\varepsilon$ and
$\min\{A^{-1},A\}$. Since $\tilde{d}_0(\tilde{x},\bar{x})\leq A$,
it follows that there is constant $C(\tilde{\theta})$ depending
only on $\tilde{\theta}$ such that
$\tilde{d}_{(-b)}(\tilde{x},\bar{x})\leq C(\tilde{\theta})A \ll
(\delta')^{-1}\tilde{h}$. We now apply Lemma 7.3.5 with the
accuracy parameter ${\varepsilon}/{2}$. Let $C({\varepsilon}/{2})$
be the positive constant in Lemma 7.3.5. Without loss of
generality, we may assume the positive constant $C_1(\varepsilon)$
in the canonical neighborhood assumption is larger than
$4C({\varepsilon}/{2})$.  When $\delta'(>0)$ is much smaller than
$\varepsilon$ and $\min\{A^{-1},A\}$, the point $\bar{x}$ at the
time $\bar{t}$ has a neighborhood which is either a
$\frac{3}{4}\varepsilon$-cap or a $\frac{3}{4}\varepsilon$-neck.

Since the canonical neighborhood assumption with accuracy
parameter $\varepsilon$ is violated at $(\bar{x},\bar{t})$, the
neighborhood of the point $\bar{x}$ at the new time zero for the
rescaled solution must be a $\frac{3}{4}\varepsilon$-neck. By
Lemma 7.3.5 (b), we know the neighborhood is the slice at the time
zero of the parabolic neighborhood
$$
P(\bar{x},0,\frac{4}{3}\varepsilon^{-1}\widetilde{R}
(\bar{x},0)^{-\frac{1}{2}},
-\min\{\widetilde{R}(\bar{x},0)^{-1},b\})
$$
(with $\widetilde{R}(\bar{x},0)=1$) which is
$\frac{3}{4}\varepsilon$-close (in the
$C^{[\frac{4}{3}\varepsilon^{-1}]}$ topology) to the corresponding
subset of the evolving standard cylinder $\mathbb{S}^2 \times
\mathbb{R}$ over the time interval $[-\min \{b,1\},0]$ with scalar
curvature $1$ at the time zero. If $b \geq 1$, the
$\frac{3}{4}\varepsilon$-neck is strong, which is a contradiction.
While if $b<1$, the $\frac{3}{4}\varepsilon$-neck at time $-b$ is
contained in the union of the gluing cap and the adjacent
$\delta$-neck where the $\delta$-cutoff surgery took place. Since
$\varepsilon$ is small (say $\varepsilon< 1/100$), it is clear
that the point $\bar{x}$ at time $-b$ is the center of an
$\varepsilon$-neck which is entirely contained in the adjacent
$\delta$-neck. By the proof of Lemma 7.3.2, the adjacent
$\delta$-neck approximates an ancient $\kappa$-solution. This
implies the point $\bar{x}$ at the time $\bar{t}$ has a strong
$\varepsilon$-neck, which is also a contradiction.

Hence we have proved that there exists $\beta_0=\beta_0(
\varepsilon, A, B, C)$ such that when $\beta\geq\beta_0$, the
rescaled solution on the ball $\widetilde{B}_{0}(\bar{x},A)$ can
be defined before the time $-b$.

Let $[t_{A}^{\alpha\beta},0]\supset[-b,0]$ be the largest time
interval so that the rescaled solution
$\widetilde{g}^{\alpha\beta}_{ij}$ can be defined on
$\widetilde{B}_{0}(\bar{x},A)\times[t_{A}^{\alpha\beta},0]$. We
finally claim that $t_{A}^{\alpha\beta}\le
-b-\frac{1}{8}\eta^{-1}C^{-1}$ for $\beta$ large enough.

Indeed, suppose not, by the gradient estimates as in Step 1, we
have the curvature estimate
$$
\widetilde{R}(x,t)\leq {10C}
$$
on $\widetilde{B}_{0}(\bar{x},A)\times [t_{A}^{\alpha\beta},-b]$.
Hence we have the curvature estimate
$$
\widetilde{R}(x,t)\leq {10C}
$$
on $\widetilde{B}_{0}(\bar{x},A)\times [t_{A}^{\alpha\beta},0]$.
By the above argument there is a $\beta_0=\beta_0( \varepsilon, A,
B + \frac{1}{8}\eta^{-1}C^{-1}, 10C)$ such that for
$\beta\geq\beta_0$, the solution in the ball
$\widetilde{B}_{0}(\bar{x},A)$ can be defined before the time
$t_{A}^{\alpha\beta}$. This is a contradiction.

Therefore we have proved Assertion 1.
\end{proof}

\begin{lemma}[backward pointwise extension from a low-curvature point]
  \label{lem:cz-low-curvature-point-backward-extension}
  \dcref{Assertion 2 in Step 2 of Lemma 7.4.2}
  \group{Canonical-Neighborhood Induction}
  \source{cao-zhu}
  \uses{lem:cz-rescaled-ball-backward-extension}
For arbitrarily fixed $\alpha$,
$0<A<+\infty$, $1 \leq C<+\infty$ and $0 < B
<\frac{1}{2}\varepsilon^2 (r^{\alpha})^{-2}-
\frac{1}{50}\eta^{-1}$, there is a $\beta_0=\beta_0(\varepsilon,
A,B,C)$ (independent of $\alpha$) such that if $\beta\geq\beta_0$
and the rescaled solution $\widetilde{g}^{\alpha\beta}_{ij}$ on
the ball $\widetilde{B}_{0}(\bar{x},A)$ is defined on a time
interval $[-b+\epsilon',0]$ with $0 < b \leq B$ and $0<\epsilon'
<\frac{1}{50}\eta^{-1}$ and the scalar curvature satisfies
$$
\widetilde{R}(x,t)\leq {C} \ \ \mbox{ on } \
\widetilde{B}_{0}(\bar{x},A)\times [-b+\epsilon',0],
$$
 and there is a point $y\in \widetilde{B}_{0}(\bar{x},A)$ such
that $\widetilde{R}(y,-b+\epsilon')\leq \frac{3}{2}$, then the
rescaled solution $\widetilde{g}^{\alpha\beta}_{ij}$ at $y$ is
also defined on the extended time interval
$[-b-\frac{1}{50}\eta^{-1},0]$ and satisfies the estimate
$$
\widetilde{R}(y,t)\leq 15
$$
for $t\in[-b-\frac{1}{50}\eta^{-1},-b+\epsilon']$.

\end{lemma}
\begin{proof}
We imitate the proof of
Assertion 1. If the rescaled solution
$\widetilde{g}^{\alpha\beta}_{ij}$ at $y$ cannot be defined for
some time in $[-b-\frac{1}{50}\eta^{-1},-b+\epsilon')$, then there
is a surgery at some time $\tilde{\tilde{t}} \in
[-b-\frac{1}{50}\eta^{-1},-b+\epsilon']$ such that $y$ lies in the
instant gluing cap.  Let $\tilde{h}$
$(=R(\bar{x},\bar{t})^{\frac{1}{2}}h$) be the cutoff radius at the
time $\tilde{\tilde{t}}$ for the rescaled solution. Clearly, there
is a universal constant $D>1$ such that
${D}^{-1}\tilde{h}\leq\widetilde{R}(y,\tilde{\tilde{t}})^{-\frac{1}{2}}
\leq D\tilde{h}$. By the gradient estimates as in Step 1, the
cutoff radius satisfies
$$
\tilde{h} \geq D^{-1}15^{-\frac{1}{2}}.
$$

As in Claim 1 (i) in the proof of Assertion 1, for any small
constants $0<\tilde{\theta}<\frac{1}{2}$, $\delta'>0$, there
exists a $\beta(\delta',\varepsilon, \tilde{\theta})> 0$ such that
for $\beta \geq \beta(\delta',\varepsilon, \tilde{\theta})$, there
is no surgery interfering in
$\widetilde{B}_{\tilde{\tilde{t}}}(y,(\delta')^{-1}\tilde{h})\times
([\tilde{\tilde{t}},(1-\tilde{\theta})\tilde{h}^{2}
+\tilde{\tilde{t}}]\cap(\tilde{\tilde{t}},0])$. Without loss of
generality, we may assume that the universal constant $\eta$ is
much larger than $D$. Then we have
$(1-\tilde{\theta})\tilde{h}^{2}
+\tilde{\tilde{t}}>-b+\frac{1}{50}\eta^{-1}$. As in Claim 2 in the
proof of Assertion 1, we can use the curvature bound assumption to
choose $\tilde{\theta}=\tilde{\theta}(B,C)$ such that
$(1-\tilde{\theta})\tilde{h}^{2} +\tilde{\tilde{t}}\geq 0$;
otherwise
$$
C\geq\frac{D'}{\tilde{\theta}\tilde{h}^{2}}
$$
for some universal constant $D'>1$, and
$$
|\tilde{\tilde{t}}+b|\leq \frac{1}{50}\eta^{-1},
$$
which implies
$$
1\geq(1-\tilde{\theta})\left(1+\frac{D'}{C\left(B+\frac{1}{50}\eta^{-1}\right)}\right).
$$
This is a contradiction if we choose
$\tilde{\theta}={D'}/{2(D'+C(B+\frac{1}{50}\eta^{-1}))}$.

So there is a positive constant $0
<\tilde{\theta}=\tilde{\theta}(B,C)<1$ such that for any
$\delta'>0$, there is a positive $\beta(\delta',\varepsilon,
\tilde{\theta})$ such that when $\beta \geq
\beta(\delta',\varepsilon, \tilde{\theta})$, we have
$-\tilde{\tilde{t}}\leq (1-\tilde{\theta})\tilde{h}^{2}$  and the
solution in the ball $\widetilde{B}_{\tilde{\tilde{t}}}(\tilde{x},
{(\delta^{\prime})}^{-1}\tilde{h})$ on the time interval
$[\tilde{\tilde{t}},0]$ is, after scaling with factor
$\tilde{h}^{-2}$, ${\delta}^{\prime}$-close (in the
$C^{[\delta'^{-1}]}$ topology) to the corresponding subset of the
standard solution.

Then exactly as in the proof of Assertion 1, by using the
canonical neighborhood structure of the standard solution in Lemma
7.3.5, this gives the desired contradiction with the hypothesis
that the canonical neighborhood assumption with accuracy parameter
$\varepsilon$ is violated at $(\bar{x},\bar{t})$, for $\beta$
sufficiently large.

The curvature estimate at the point $y$ follows from Step 1.
Therefore the proof of Assertion 2 is complete.
\end{proof}

\begin{lemma}[short backward control or capwise scalar comparability]
  \label{lem:cz-backward-control-or-capwise-comparability}
  \dcref{Assertion 3 in Step 2 of Lemma 7.4.2}
  \group{Canonical-Neighborhood Induction}
  \source{cao-zhu}
  \uses{lem:cz-rescaled-ball-backward-extension}
For arbitrarily fixed $\alpha$,
$0<A<+\infty$, $1 \leq C<+\infty$, there is a
$\beta_0=\beta_0(\varepsilon, AC^{\frac{1}{2}})$ such that if any
point $(y_0,t_0)$  with  $0 \leq -t_0 <\frac{1}{2}\varepsilon^2
(r^{\alpha})^{-2}- \frac{1}{8}\eta^{-1}C^{-1}$ of the rescaled
solution $\widetilde{g}^{\alpha\beta}_{ij}$ for $\beta\geq\beta_0$
satisfies $\widetilde{R}(y_0,t_0)\leq C$ , then either the
rescaled solution at $y_0$ can be defined at least on
$[t_0-\frac{1}{16}\eta^{-1}C^{-1},t_0]$ and the rescaled scalar
curvature satisfies
$$
\widetilde{R}(y_0,t)\leq 10 {C} \; \mbox{ for } t\in
\Big[t_0-\frac{1}{16}\eta^{-1}C^{-1},t_0\Big],
$$
or we have
$$
\widetilde{R}(x_1,t_0)\leq 2D'' \widetilde{R}(x_2,t_0)
$$
for any two points $x_1, x_2\in \widetilde{B}_{t_0}(y_0,A)$, where
$D''$ is the above universal constant.

\end{lemma}
\begin{proof}
Suppose the rescaled
solution $\widetilde{g}^{\alpha\beta}_{ij}$ at $y_0$ cannot be
defined for some $t\in [t_0-\frac{1}{16}\eta^{-1}C^{-1},t_0)$;
then there is a surgery at some time $\tilde{t} \in
[t_0-\frac{1}{16}\eta^{-1}C^{-1},t_0]$ such that $y_0 $ lies in
the instant gluing cap. Let $\tilde{h}$
$(=R(\bar{x},\bar{t})^{\frac{1}{2}}h$) be the cutoff radius at the
time $\tilde{t}$ for the rescaled solution
$\widetilde{g}^{\alpha\beta}_{ij}$.  By the gradient estimates as
in Step 1, the cutoff radius satisfies
$$
\tilde{h} \geq D^{-1}10^{-\frac{1}{2}}C^{-\frac{1}{2}},
$$
where $D$ is the universal constant in the proof of the Assertion
1. Since we assume $\eta$ is suitably larger than $D$ as before,
we have $\frac{1}{2}\tilde{h}^{2} +\tilde{t}>t_0$. As in Claim 1
(ii) in the proof of Assertion 1, for arbitrarily small
$\delta'>0$, we know that for $\beta$ large enough the rescaled
solution on
$\widetilde{B}_{\tilde{t}}(y_0,{(\delta^{\prime})}^{-1}\tilde{h})\times
[\tilde{t},t_0]$ is, after scaling with factor $\tilde{h}^{-2}$,
${\delta}^{\prime}$-close (in the $C^{[(\delta')^{-1}]}$ topology)
to the corresponding subset of the standard solution. Since
$(\delta')^{-1}\tilde{h}\gg A$ for $\beta$ large enough, Assertion
3 follows from the curvature estimate of standard solution in the
time interval $ [0,\frac{1}{2}]$.

\end{proof}

\begin{lemma}[uniform curvature and backward existence on bounded rescaled balls]
  \label{lem:cz-rescaled-balls-uniform-curvature-backward-existence}
  \dcref{Assertion 4 in Step 3 of Lemma 7.4.2}
  \group{Canonical-Neighborhood Induction}
  \source{cao-zhu}
  \uses{lem:cz-backward-control-or-capwise-comparability}
Given any subsequence of the rescaled
solutions $\tilde{g}^{\alpha_k\beta_k}_{ij}$ with
$r^{\alpha_k}\rightarrow 0$ and
$\delta^{\alpha_k\beta_k}\rightarrow 0$ as $k\rightarrow \infty$,
then for any $L>0$, there are constants $C(L)>0$ and $k(L)$ such
that the rescaled solutions $\tilde{g}^{\alpha_k\beta_k}_{ij}$
satisfy
\begin{itemize}
\item[(i)] $\tilde{R}(x,0)\leq C(L)$ for all points $x$ with
$\tilde{d}_{0}(x,\bar{x})\leq L$ and all $k \geq 1$; \item[(ii)]
the rescaled solutions over the ball $\tilde{B}_{0}(\bar{x},L)$
are defined at least on the time interval
$[-\frac{1}{16}\eta^{-1}C(L)^{-1},0]$ for all $k\geq k(L)$.
\end{itemize}

\end{lemma}
\begin{proof}
For each $\rho>0$, set
\begin{multline*}
M(\rho)=\sup\Big\{\tilde{R}(x,0)\ |\ k \geq 1\; \mbox{ and }\;
\tilde{d}_0(x,\bar{x})\leq\rho \\
\mbox{in the rescaled solutions } \;
\tilde{g}^{\alpha_k\beta_k}_{ij}\Big\}
\end{multline*}
and
$$
\rho_0=\sup\{\rho > 0\ |\ M(\rho)<+\infty\}.
$$
Note that the estimate (7.4.14) implies that $\rho_0>0$. For (i),
it suffices to prove $\rho_0=+\infty$.

We argue by contradiction. Suppose $\rho_0<+\infty$. Then there is
a sequence of points $y$ in the rescaled solutions
$\tilde{g}^{\alpha_k\beta_k}_{ij}$ with
$\tilde{d}_0(\bar{x},y)\rightarrow\rho_0<+\infty$ and
$\tilde{R}(y,0)\rightarrow +\infty$. Denote by $\gamma$ a
minimizing geodesic segment from $\bar{x}$ to $y$ and denote by
$\tilde{B}_0(\bar{x},\rho_0)$ the open geodesic ball centered at
$\bar{x}$ of radius $\rho_0$ on the rescaled solution
$\tilde{g}^{\alpha_k\beta_k}_{ij}$.

First, we claim that for any $0<\rho<\rho_0$ with $\rho$ near
$\rho_0$, the rescaled solutions on the balls
$\tilde{B}_{0}(\bar{x},\rho)$ are defined on the time interval
$[-\frac{1}{16}\eta^{-1}M(\rho)^{-1},0]$ for all large $k$.
Indeed, this follows from Assertion 3 or Assertion 1. For the
later purpose in Step 6, we now present an argument by using
Assertion 3.  If the claim is not true, then there is a surgery at
some time $\tilde{t} \in [-\frac{1}{16}\eta^{-1}M(\rho)^{-1},0]$
such that some point $\tilde{y}\in \tilde{B}_{0}(\bar{x},\rho) $
lies in the instant gluing cap.  We can choose sufficiently small
$\delta'>0$ such that $2\rho_0<(\delta')^{-\frac{1}{2}}\tilde{h}$,
where $\tilde{h} \geq
D^{-1}20^{-\frac{1}{2}}M(\rho)^{-\frac{1}{2}}$ is the cutoff
radius of the rescaled solutions at $\tilde{t}$. By applying
Assertion 3 with $(\tilde{y},0)=(y_0,t_0)$, we see that there is a
$k(\rho_0,M(\rho))>0$ such that when $k\geq k(\rho_0,M(\rho))$,
$$
\widetilde{R}(x,0)\leq 2D''
$$
for all $x\in \widetilde{B}_{0}(\bar{x},\rho)$. This is a
contradiction as $\rho\rightarrow\rho_0$.

Since for each fixed $0<\rho<\rho_0$ with $\rho$ near $\rho_0$,
the rescaled solutions are defined on
$\tilde{B}_{0}(\bar{x},\rho)\times
[-\frac{1}{16}\eta^{-1}M(\rho)^{-1},0]$ for all large $k$, by Step
1 and Shi's derivative estimate, we know that the covariant
derivatives and higher order derivatives of the curvatures on
$\tilde{B}_{0}(\bar{x},\rho -
\frac{(\rho_0-\rho)}{2})\times[-\frac{1}{32}\eta^{-1}M(\rho)^{-1},0]$
are also uniformly bounded.

By the uniform $\kappa$-noncollapsing property and Hamilton's
compactness theorem (Theorem 4.1.5), after passing to a
subsequence, we can assume that the marked sequence
$(\tilde{B}_0(\bar{x},\rho_0),\widetilde{g}^{\alpha_k\beta_k}_{ij},
\bar{x})$ converges in the $C^{\infty}_{loc}$ topology to a marked
(noncomplete) manifold
($B_{\infty},\widetilde{g}^{\infty}_{ij},\bar{x})$ and the
geodesic segments $\gamma$ converge to a geodesic segment (missing
an endpoint) $\gamma_{\infty}\subset B_{\infty}$ emanating from
$\bar{x}$.

Clearly, the limit has nonnegative sectional curvature by the
pinching assumption. Consider a tubular neighborhood along
$\gamma_{\infty}$ defined by
$$
V=\bigcup_{q_0\in\gamma_{\infty}}B_{\infty}
(q_0,4\pi(\widetilde{R}_{\infty}(q_0))^{-\frac{1}{2}}),
$$
where $\widetilde{R}_{\infty}$ denotes the scalar curvature of the
limit and
$$
B_{\infty}(q_0,4\pi(\widetilde{R}_{\infty}(q_0))^{-\frac{1}{2}})
$$
is the ball centered at $q_0\in B_{\infty}$ with the radius
$4\pi(\widetilde{R}_{\infty}(q_0))^{-\frac{1}{2}}$. Let
$\bar{B}_{\infty}$ denote the completion of
$(B_{\infty},\widetilde{g}^{\infty}_{ij})$, and
$y_{\infty}\in\bar{B}_{\infty}$ the limit point of
$\gamma_{\infty}$. Exactly as in the second step of the proof of
Theorem 7.1.1, it follows from the canonical neighborhood
assumption with accuracy parameter $2\varepsilon$ that the
limiting metric $\widetilde{g}^{\infty}_{ij}$ is cylindrical at
any point $q_0\in\gamma_{\infty}$ which is sufficiently close to
$y_{\infty}$ and then the metric space $\bar{V}=V\cup
\{y_{\infty}\}$ by adding the point $y_{\infty}$ has nonnegative
curvature in the Alexandrov sense. Consequently we have a
three-dimensional non-flat tangent cone $C_{y_{\infty}}\bar{V}$ at
$y_{\infty}$ which is a metric cone with aperture $\leq
20\varepsilon$.

On the other hand, note that by the canonical neighborhood
assumption, the canonical $2\varepsilon$-neck neighborhoods are
strong. Thus at each point $q\in V$ near $y_{\infty}$, the
limiting metric $\widetilde{g}^{\infty}_{ij}$ actually exists on
the whole parabolic neighborhood
$$
V \bigcap
P\left(q,0,\frac{1}{3}\eta^{-1}(\widetilde{R}_{\infty}(q))^{-\frac{1}{2}},
-\frac{1}{10}\eta^{-1}(\widetilde{R}_{\infty}(q))^{-1}\right),
$$
and is a smooth solution of the Ricci flow there. Pick $z\in
C_{y_{\infty}}\bar{V}$ with distance one from the vertex
$y_{\infty}$ and it is nonflat around $z$. By definition the ball
$B(z,\frac{1}{2})\subset C_{y_{\infty}}\bar{V}$ is the
Gromov-Hausdorff convergent limit of the scalings of a sequence of
balls $B_{\infty}(z_{\ell},\sigma_{\ell})
(\subset(V,\widetilde{g}^{\infty}_{ij}))$ where
$\sigma_{\ell}\rightarrow0$. Since the estimate (7.4.14) survives
on $(V,\widetilde{g}^{\infty}_{ij})$ for all $A < +\infty$, and
the tangent cone is three-dimensional and nonflat around $z$, we
see that this convergence is actually in the $C^{\infty}_{\rm
loc}$ topology and over some ancient time interval. Since the
limiting $B_{\infty}(z,\frac{1}{2})(\subset
C_{y_{\infty}}\bar{V})$ is a piece of nonnegatively curved nonflat
metric cone, we get a contradiction with Hamilton's strong maximum
principle (Theorem 2.2.1) as before. So we have proved
$\rho_0=\infty$. This proves (i).

By the same proof of Assertion 1 in Step 2, we can further show
that for any $L$, the rescaled solutions on the balls
$\tilde{B}_{0}(\bar{x},L)$ are defined at least on the time
interval $[-\frac{1}{16}\eta^{-1}C(L)^{-1},0]$ for all
sufficiently large $k$. This proves (ii).

\end{proof}

\begin{lemma}[uniform space-time curvature around extended limit points]
  \label{lem:cz-extended-limit-uniform-spacetime-curvature}
  \dcref{Assertion 5 in Step 6 of Lemma 7.4.2}
  \group{Canonical-Neighborhood Induction}
  \source{cao-zhu}
  \uses{lem:cz-backward-control-or-capwise-comparability, lem:cz-rescaled-balls-uniform-curvature-backward-existence}
For the above rescaled solutions
$\widetilde{g}^{\alpha_k\beta_k}_{ij}$ and $\bar{k}_0$, we have
that for any $L>0$, there is a positive constant $\omega(L)$ such
that the rescaled solutions $\widetilde{g}^{\alpha_k\beta_k}_{ij}$
satisfy
$$
\widetilde{R}(x,t)\leq\omega(L)
$$
for all $(x,t)$ with $\tilde{d}_{t}(x,y_k)\leq L$ and $t\in
[-t_{\max}-\frac{1}{50}\eta^{-1},t_{\infty}]$, and for all $k \geq
\bar{k}_0$.

\end{lemma}
\begin{proof}
We slightly modify the
argument in the proof of Assertion 4 (i). Let

\begin{eqnarray*}
M(\rho)=\sup\Big\{\widetilde{R}(x,t) &  |&
\tilde{d}_t(x,y_k)\leq\rho\; \mbox{ and }\; t \in
[-t_{\max}-\frac{1}{50}\eta^{-1},t_{\infty}]\\
&   &   \mbox{in the rescaled solutions } \;
\widetilde{g}^{\alpha_k\beta_k}_{ij}, k\geq \bar{k}_0 \Big\}
\end{eqnarray*}
and
$$
\rho_0=\sup\{\rho > 0\ |\  M(\rho)<+\infty\}.
$$
Note that the estimate (7.4.14) implies that $\rho_0>0$. We only
need to show $\rho_0 = +\infty$.

We argue by contradiction. Suppose $\rho_0<+\infty$. Then, after
passing to a subsequence, there is a sequence $(\tilde{y}_k,t_k)$
in the rescaled solutions $\widetilde{g}^{\alpha_k\beta_k}_{ij}$
with $t_k \in [-t_{\max}-\frac{1}{50}\eta^{-1},t_{\infty}]$ and
$\tilde{d}_{t_k}(y_k,\tilde{y}_k)\rightarrow\rho_0<+\infty$ such
that $\widetilde{R}(\tilde{y}_k,t_k)\rightarrow +\infty$. Denote
by $\gamma_k$ a minimizing geodesic segment from $y_k$ to
$\tilde{y}_k$ at the time $t_k$ and denote by
$\widetilde{B}_{t_k}(y_k,\rho_0)$ the open geodesic ball centered
at $y_k$ of radius $\rho_0$ on the rescaled solution
$\widetilde{g}^{\alpha_k\beta_k}_{ij}(\cdot,t_k)$.

For any $0<\rho<\rho_0$ with $\rho$ near $\rho_0$, by applying
Assertion 3 as before, we get that  the rescaled solutions on the
balls $\widetilde{B}_{t_k}(y_k,\rho)$ are defined on the time
interval $[t_k-\frac{1}{16}\eta^{-1}M(\rho)^{-1},t_k]$ for all
large $k$. By Step 1 and Shi's derivative estimate, we further
know that the covariant derivatives and higher order derivatives
of the curvatures on
$\widetilde{B}_{t_k}(y_k,\rho-\frac{(\rho_0-\rho)}{2})
\times[t_k-\frac{1}{32}\eta^{-1}M(\rho)^{-1},t_k]$ are also
uniformly bounded. Then by the uniform $\kappa$-noncollapsing
property and Hamilton's compactness theorem (Theorem 4.1.5), after
passing to a subsequence, we can assume that the marked sequence
$(\tilde{B}_{t_k}(y_k,\rho_0),\widetilde{g}^{\alpha_k\beta_k}_{ij}
(\cdot,t_k),y_k)$ converges in the $C^{\infty}_{loc}$ topology to
a marked (noncomplete) manifold
($B_{\infty},\widetilde{g}^{\infty}_{ij},y_{\infty})$ and the
geodesic segments $\gamma_k$ converge to a geodesic segment
(missing an endpoint) $\gamma_{\infty}\subset B_{\infty}$
emanating from $y_{\infty}$.

Clearly, the limit also has nonnegative sectional curvature by the
pinching assumption. Then by repeating the same argument as in the
proof of Assertion 4 (i) in the rest, we derive a contradiction
with Hamilton's strong maximum principle. This proves Assertion 5.
\end{proof}

\begin{lemma}[reduced-length bound for admissible minimizers]
  \label{lem:cz-admissible-minimizer-reduced-length-bound}
  \dcref{Assertion (7.5.6)}
  \group{Curvature Estimates after Surgery}
  \source{cao-zhu}
  \uses{lem:cz-surgery-cap-reduced-length-barrier}
Let $\tau\in[0,\frac{1}{2}]$. If for each
$s\in[0,\tau]$, $\inf\{\bar{L}(y,s)\ |\ \mbox{ } d_t(x_0,y)\leq
A(2t-1)+\frac{1}{10}\mbox{ with }s=1-t\}$ is achieved by
admissible curves, then we have
\begin{align}
&\inf\left\{\bar{L}(y,\tau)\ |\ \mbox{ } d_t(x_0,y)
\leq A(2t-1)+\frac{1}{10}\mbox{ with }\tau=1-t\right\}\\
&\leq 2\sqrt{\tau}\exp(C(A)+100). \nonumber  %\eqno (7.5.6)
\end{align}
\end{lemma}
\begin{proof}
\begin{align}
\frac{d}{d\tau}\left(\log\left(\frac{h_{\min}(\tau)}{\sqrt{\tau}}\right)\right) &
\leq C(A)+\frac{6\sqrt{\tau}+1}{2\tau-4\tau^2\sqrt{\tau}}
-\frac{1}{2\tau}\\
& \leq C(A)+\frac{50}{\sqrt{\tau}}, \nonumber
\end{align} %\eqno (7.5.4)
as long as the associated shortest $\mathcal{L}$-geodesics are
admissible with $\bar{L} \leq 3\sqrt{\tau}\exp(C(A)+100)$. On the
other hand, by definition, we have \begin{equation} \lim_{\tau\rightarrow
0^+}\frac{h_{\min}(\tau)}{\sqrt{\tau}}\leq
\phi(d_1(x_0,x)-A)\cdot 2=2. %\eqno (7.5.5)
\end{equation} Integrating (7.5.4) from zero to $	au$ and using the initial bound (7.5.5) gives the asserted estimate.
\end{proof}

\begin{proposition}[Perelman collapsing criterion for graph manifolds]
  \label{prop:cz-perelman-collapsing-graph-manifold-criterion}
  \dcref{Perelman's Claim (P2, Theorem 7.4)}
  \group{Geometrization}
  \source{cao-zhu}
Suppose $(M^{\alpha},g^{\alpha}_{ij})$ is a
sequence of compact orientable three-manifolds, closed or with
convex boundary, and $w^{\alpha}\rightarrow  0$. Assume that
\begin{itemize}
\item[(1)] for each point $x \in M^{\alpha}$ there exists a radius
$\rho = \rho^{\alpha}(x), 0<\rho<1,$ not exceeding the diameter of
the manifold, such that the ball $B(x,\rho)$ in the metric
$g^{\alpha}_{ij}$ has volume at most $w^{\alpha}\rho^3$ and
sectional curvatures at least $-\rho^{-2}$; \item[(2)] each
component of the boundary of $M^{\alpha}$ has diameter at most
$w^{\alpha}$, and has a (topologically trivial) collar of length
one, where the sectional curvatures are between $-1/4 - \epsilon$
and $-1/4 + \epsilon$.
\end{itemize}

Then $M^{\alpha}$ for sufficiently large $\alpha$ are
diffeomorphic to graph manifolds.
\end{proposition}

\begin{theorem}[closed collapsing manifolds are graph manifolds]
  \label{thm:cz-shioya-yamaguchi-closed-collapsing-graph-manifolds}
  \dcref{Shioya--Yamaguchi Theorem 8.1}
  \group{Geometrization}
  \source{cao-zhu}
Suppose
$(M^{\alpha},g^{\alpha}_{ij})$ is a sequence of compact orientable
three-manifolds without boundary, and $w^{\alpha}\rightarrow 0$.
Assume that for each point $x \in M^{\alpha}$ there exists a
radius $\rho = \rho^{\alpha}(x),$ not exceeding the diameter of
the manifold, such that the ball $B(x,\rho)$ in the metric
$g^{\alpha}_{ij}$ has volume at most $w^{\alpha}\rho^3$ and
sectional curvatures at least $-\rho^{-2}$. Then $M^{\alpha}$ for
sufficiently large $\alpha$ are diffeomorphic to graph manifolds.
\end{theorem}
